% Engineering Decision Notebook (EDN), following the course instructions
% (references/Instruction_EDN_MLOps_v2026.md). Entries are written in
% reports/edn.md and transferred by `make edn` (tools/edn_latex.py), rerun in
% the same pull request as any change to it, which writes one file per entry
% into edn/ and the list of them into edn/entries.tex.
\documentclass[11pt, a4paper]{article}
% Shared preamble for report.tex and edn.tex.

% --- Build flags ------------------------------------------------------------
% Set from the command line by the Makefile (latexmk -usepretex), so a plain
% compile in an editor gives the defaults: initial report, guidance visible.
%   \ReportFinal       final report (Milestones 1-6) instead of initial (1-3)
%   \ReportSubmission  hide guidance boxes and example content
\newif\iffinalreport
\ifdefined\ReportFinal \finalreporttrue \fi
\newif\ifguidance
\guidancetrue
\ifdefined\ReportSubmission \guidancefalse \fi

% --- Packages ---------------------------------------------------------------
\usepackage[T1]{fontenc}
\usepackage[utf8]{inputenc}
\usepackage{lmodern}
\usepackage[english]{babel}
\usepackage{microtype}
% `--` inside \texttt stays two hyphens instead of the en dash the T1 ligature makes,
% so command-line options such as `ruff format --check` print as typed.
\DisableLigatures[-]{encoding=T1,family=tt*}
\usepackage[a4paper, margin=2.5cm, headheight=14pt]{geometry}
\usepackage{graphicx}
\graphicspath{{figures/}{../figures/}}
\usepackage[table]{xcolor}
\usepackage{booktabs}
\usepackage{tabularx}
\usepackage{array}
\usepackage{amssymb}
\usepackage{enumitem}
\usepackage{etoolbox}
\usepackage{fancyhdr}
\usepackage{listings}
\usepackage[most]{tcolorbox}
\usepackage{csquotes}
\usepackage[backend=biber, style=ieee]{biblatex}
\addbibresource{references.bib}
% Let long URLs in the bibliography break anywhere instead of overflowing.
\setcounter{biburlnumpenalty}{9000}
\setcounter{biburllcpenalty}{7000}
\setcounter{biburlucpenalty}{8000}
\usepackage[hidelinks]{hyperref}
\usepackage[capitalise, noabbrev]{cleveref}

% Team and project facts shown on the cover pages. Edit here, nowhere else.
% Replace every \tbd{...} before a submission.

\newcommand{\TeamName}{Recommenditos}
\newcommand{\ProjectTitle}{Car price estimation ML system}
\newcommand{\CourseName}{Machine Learning Systems in Production (MLOps)}
\newcommand{\CourseProgram}{Master in Data Science, Universitat Polit\`ecnica de Catalunya (UPC)}
\newcommand{\CourseTerm}{2026/27 Q1}

% Links (course rules: GitHub repo plus any other resources we use).
% Shown without the https:// prefix so they fit on one line.
\newcommand{\RepoLink}{\href{https://github.com/mlops-2627q1-mds-upc/MLOPS_Recommenditos}{\nolinkurl{github.com/mlops-2627q1-mds-upc/MLOPS_Recommenditos}}}
\newcommand{\BoardLink}{\href{https://github.com/orgs/mlops-2627q1-mds-upc/projects/1}{\nolinkurl{github.com/orgs/mlops-2627q1-mds-upc/projects/1}}}
\newcommand{\DiscordLink}{Discord (team channel, invite on request)}
\newcommand{\DagsHubLink}{\href{https://dagshub.com/recommenditos/MLOPS_Recommenditos}{\nolinkurl{dagshub.com/recommenditos/MLOPS_Recommenditos}}}
\newcommand{\SlidesLink}{\tbd{presentation slides URL}}
\newcommand{\ApiLink}{\tbd{deployed API URL}}

% One address that reaches the whole team (course rule). We have no shared
% mailing list, so this is the Scrum Master, who forwards to the team.
\newcommand{\TeamMail}[1]{\href{mailto:#1}{\nolinkurl{#1}}}
\newcommand{\TeamContact}{\TeamMail{lukas.haeussler@estudiantat.upc.edu} (Product Owner and Scrum Master)}

% Team members, used on the cover page and in the CATME table.
% Short names are what appears in the CATME table headers.
\newcommand{\MemberA}{Lukas}
\newcommand{\MemberB}{Kevin}
\newcommand{\MemberC}{Urszula}
\newcommand{\MemberD}{Michal}
\newcommand{\MemberE}{Mark}

% Rows for the cover page table: Surname & Name & UPC e-mail & GitHub account.
% Addresses exactly as recorded in the repository README. The last two are on
% estudiant.upc.edu while the other three are on estudiantat.upc.edu, so each
% of the two has to confirm their own address before the submission.
\newcommand{\TeamMemberRows}{%
  H\"au{\ss}ler & Lukas & \TeamMail{lukas.haeussler@estudiantat.upc.edu} & \texttt{lukas2510} \\
  Adameit & Kevin & \TeamMail{kevin.adameit@estudiantat.upc.edu} & \texttt{kadameit} \\
  Sawczuk & Urszula & \TeamMail{urszula.wanda.sawczuk@estudiantat.upc.edu} & \texttt{ulasawczuk} \\
  Dudek & Michal & \TeamMail{michal.feliks.dudek@estudiant.upc.edu} \tbd{Michal confirms this domain} & \texttt{michudud04} \\
  Atzberger & Mark & \TeamMail{mark.welf.atzberger@estudiant.upc.edu} \tbd{Mark confirms this domain} & \texttt{W11W11W11} \\
}


\hypersetup{
  pdftitle={MLOps\_\TeamName},
  pdfauthor={Team \TeamName},
}

\setlist{itemsep=2pt, topsep=4pt}
\setlength{\parindent}{0pt}
\setlength{\parskip}{0.5em}

\lstset{
  basicstyle=\ttfamily\small,
  breaklines=true,
  frame=single,
  columns=fullflexible,
}

\pagestyle{fancy}
\fancyhf{}
\renewcommand{\headrulewidth}{0.4pt}
\fancyfoot[C]{\thepage}

% --- Writing helpers --------------------------------------------------------

% \guidance{...}: what a section should contain, taken from the course
% template, the lab rubric and the demo repo. Hidden in submission builds.
\newtcolorbox{guidancebox}{
  breakable, colback=blue!4, colframe=blue!45!black, boxrule=0.5pt,
  fonttitle=\small\bfseries, title={Guidance (hidden in submission build)},
  fontupper=\small, left=4pt, right=4pt, top=2pt, bottom=2pt,
}
\newcommand{\guidance}[1]{\ifguidance\begin{guidancebox}#1\end{guidancebox}\fi}

% \tbd{...}: a placeholder that still has to be filled in. Always rendered in
% red, and counted by `make report-submit` so nothing slips into a submission.
\newcommand{\tbd}[1]{%
  \GenericWarning{}{Report warning: Unfilled placeholder: #1}%
  \textcolor{red}{\textbf{[TBD: #1]}}%
}

% Course repositories and documents we refer to.
\newcommand{\demorepo}{\url{https://github.com/mlops-2627q1-mds-upc/MLOps-2627q1-demos}}


\fancyhead[L]{\small MLOps\_\TeamName}
\fancyhead[R]{\small Engineering Decision Notebook}

% --- \ednentry{key=value, ...} ---------------------------------------------
% Keys mirror the fields of the course instructions. Required: id, title,
% milestone, participants, decision, rationale, ai. Brace any value that
% contains a comma or an equals sign. Values may span several paragraphs.
% `note` is ours: it says that a later entry amended this one, or that its
% evidence was corrected, and is printed first so it is read first. Each entry
% is labelled with its id, so \hyperref[EDN-14]{...} links to it.
\newtcolorbox{ednbox}[1]{
  breakable, enhanced, colback=white, colframe=black!60, boxrule=0.5pt,
  fonttitle=\bfseries, title={#1}, left=6pt, right=6pt,
}

\ExplSyntaxOn
\clist_const:Nn \c__edn_fields_clist
  { id, title, note, date, milestone, topic, participants, decision,
    rationale, ai, response, assessment, aievidence, evidence }
\clist_const:Nn \c__edn_required_clist
  { id, title, milestone, participants, decision, rationale, ai }
\clist_map_inline:Nn \c__edn_fields_clist
  {
    \tl_new:c { l__edn_#1_tl }
    \keys_define:nn { edn } { #1 .tl_set:c = { l__edn_#1_tl } }
  }

\cs_new_protected:Npn \__edn_field:nn #1#2
  {
    \tl_if_empty:cF { l__edn_#2_tl }
      { \item [#1] \tl_use:c { l__edn_#2_tl } }
  }
% The evidence fields are mostly links, paths and IDs, which justified text can
% only space out into loose lines, so they are set ragged-right; the group and
% the \par keep that from reaching the next field.
\cs_new_protected:Npn \__edn_ragged_field:nn #1#2
  {
    \tl_if_empty:cF { l__edn_#2_tl }
      { \item [#1] { \raggedright \tl_use:c { l__edn_#2_tl } \par } }
  }

\NewDocumentCommand \ednentry { +m }
  {
    \group_begin:
    \clist_map_inline:Nn \c__edn_fields_clist { \tl_clear:c { l__edn_##1_tl } }
    \keys_set:nn { edn } {#1}
    \clist_map_inline:Nn \c__edn_required_clist
      {
        \tl_if_empty:cT { l__edn_##1_tl }
          {
            \msg_error:nnVn { edn } { missing-field } \l__edn_id_tl {##1}
          }
      }
    \phantomsection
    \exp_args:Nx \label { \l__edn_id_tl }
    \addcontentsline { toc } { section } { \l__edn_id_tl :~ \l__edn_title_tl }
    \begin { ednbox } { \l__edn_id_tl :~ \l__edn_title_tl }
      \begin { description }
        [ leftmargin = 4.3cm, labelwidth = 4.1cm, labelsep = 0.2cm,
          align = parleft, font = \small\bfseries, itemsep = 3pt ]
        \__edn_field:nn { Note } { note }
        \__edn_field:nn { Date } { date }
        \__edn_field:nn { Milestone } { milestone }
        \__edn_field:nn { Activity~/~Topic } { topic }
        \__edn_field:nn { Participants } { participants }
        \__edn_field:nn { Decision } { decision }
        \__edn_field:nn { Rationale } { rationale }
        \__edn_field:nn { AI~involvement } { ai }
        \__edn_field:nn { Response~to~AI } { response }
        \__edn_field:nn { Assessment~of~the~AI~contribution } { assessment }
        \__edn_ragged_field:nn { AI~interaction~evidence } { aievidence }
        \__edn_ragged_field:nn { Other~evidence } { evidence }
      \end { description }
    \end { ednbox }
    \group_end:
  }

\msg_new:nnn { edn } { missing-field }
  { Entry~`#1'~is~missing~the~required~field~`#2'. }
\cs_generate_variant:Nn \msg_error:nnnn { nnVn }
\ExplSyntaxOff

% Quote a prompt or response inside the aievidence field.
\newcommand{\ednprompt}[1]{\textit{Prompt:} \enquote{#1}\par}
\newcommand{\ednresponse}[1]{\textit{Response:} \enquote{#1}\par}

\begin{document}

\begin{center}
  {\small \CourseName\ -- \CourseTerm\par}
  \vspace{0.5em}
  {\LARGE\bfseries MLOps\_\TeamName\par}
  \vspace{0.3em}
  {\Large Engineering Decision Notebook\par}
  \vspace{0.3em}
  {\small \RepoLink\par}
\end{center}

This notebook records a selected set of engineering decisions that meaningfully shaped our ML system, and how AI contributed to them.
It does not list every decision or every use of AI, nor the use of tools the course prescribes.

\tableofcontents
\bigskip

\guidance{
  Create an entry when a decision could reasonably have been made differently and the choice has meaningful consequences for the project.
  Entries are written in \texttt{reports/edn.md} and transferred here by \texttt{make edn}, rerun in the same pull request as any change to that file, which writes one file per entry into \texttt{edn/} (alternatives with pros and cons go into the rationale) and the \texttt{\textbackslash input} list into \texttt{edn/entries.tex}; never edit those files by hand.
  IDs are stable (EDN-01, EDN-02, \dots) so the report can refer to them; never renumber.
  Link issues, PRs and commits as other evidence.
  The example below is Example~3 of the course instructions and only shows in draft builds.
}

\ifguidance
  \ednentry{
    id = {EDN-00},
    title = {Example: test that the API rejects missing input},
    milestone = {M3: Quality Assurance},
    topic = {Testing Strategy},
    participants = {All team members},
    decision = {We decided to add a test to verify that the API rejects requests with missing input data.},
    rationale = {Missing input data is a likely source of invalid requests. We therefore decided to explicitly test this behavior before deployment.},
    ai = {No AI involvement.},
    evidence = {Link to the corresponding test.},
  }
\fi

% Entries, oldest first, generated from reports/edn.md by `make edn`.
% Generated by tools/edn_latex.py from reports/edn.md.
% Do not edit: change reports/edn.md and run `make edn`.
% One \input per entry, oldest first.
% Generated by tools/edn_latex.py from reports/edn.md.
% Do not edit: change reports/edn.md and run `make edn`.
\ednentry{
  id = {EDN-01},
  title = {Main dataset choice: AutoScout24 Car Listings},
  date = {2026-09-22},
  milestone = {M1: Project Inception},
  topic = {Dataset Selection},
  participants = {All team members},
  decision = {Use the AutoScout24 Car Listings Dataset (2025 snapshot, Zenodo DOI \texttt{10.\allowbreak{}5281/\allowbreak{}zenodo.\allowbreak{}17643343}) as the main training dataset for the used-car price component.},
  rationale = {\emph{Alternatives considered:}
\begin{itemize}[leftmargin=*, topsep=2pt]
\item \textbf{Option A (chosen): AutoScout24 Car Listings Dataset.} 118,382 listings, 75 columns, 8 countries, MIT license.
\newline Pros: real scraped listings, large enough for gradient boosting, rich feature set (equipment, history flags, free-text description), verifiable provenance and license.
\newline Cons: skewed toward premium brands (BMW, Porsche, Mercedes-Benz, Audi are 83\,\% of rows), mass-market brands nearly absent, single scrape with no listing date so no temporal split.
\item \textbf{Option B: smaller generic Kaggle used-car datasets.} Various single-market, pre-cleaned listing datasets evaluated per issue \#3's scope.
\newline Pros: often smaller and pre-cleaned, lower profiling effort.
\newline Cons: typically single-country, far fewer rows and features than AutoScout24, weaker or undocumented provenance/license.
\end{itemize}
\par
\emph{Why:} AutoScout24 is the only candidate large and feature-rich enough to support the planned gradient-boosting model with equipment/history features and SHAP explainability, and it has a clear, citable license. The brand skew is accepted as a stated scope limit (see project brief \S{}2) rather than a reason to switch datasets.},
  ai = {Information seeking, Alternative assessment},
  response = {Accepted},
  assessment = {AI profiled the raw dataset (row/column counts, country and brand distribution, missingness, known data issues) and summarized the trade-off against smaller generic alternatives; the team reviewed those facts and confirmed the choice in today's sprint planning.},
  aievidence = {Dataset profiling and trade-off summary in \texttt{docs/\allowbreak{}docs/\allowbreak{}project-\allowbreak{}brief.\allowbreak{}md} \S{}3.1 and \S{}2 (2026-09-22), confirmed in today's sprint planning session.},
  evidence = {\href{https://github.com/mlops-2627q1-mds-upc/MLOPS_Recommenditos/pull/13}{PR \#13}, \href{https://github.com/mlops-2627q1-mds-upc/MLOPS_Recommenditos/issues/3}{issue \#3}.},
}

% Generated by tools/edn_latex.py from reports/edn.md.
% Do not edit: change reports/edn.md and run `make edn`.
\ednentry{
  id = {EDN-02},
  title = {Model family: gradient boosting (LightGBM as main model)},
  date = {2026-09-22},
  milestone = {M1: Project Inception},
  topic = {Modelling Approach},
  participants = {All team members},
  decision = {Use gradient boosting as the model family, with LightGBM as the main model (basic features, then all features) and CatBoost as a challenger for high-cardinality categoricals.},
  rationale = {\emph{Alternatives considered:}
\begin{itemize}[leftmargin=*, topsep=2pt]
\item \textbf{Option A (chosen): Gradient boosting (LightGBM main, CatBoost challenger).}
\newline Pros: best-in-class on medium tabular data, native categoricals and missing values, trains in minutes on CPU, small artefacts, exact SHAP explanations.
\newline Cons: more tuning surface than a linear baseline; still needs a simpler baseline (median, Ridge) to prove its value.
\item \textbf{Option B: Linear / Ridge regression only.}
\newline Pros: simple, fast, highly interpretable depreciation baseline.
\newline Cons: cannot capture non-linear interactions (brand x age x mileage) well, lower expected accuracy.
\item \textbf{Option C: kNN, random forest, tabular NNs, hierarchical Bayes, LLM zero-shot.}
\newline Pros: each offers some property (kNN gives comparables for free, NNs could use text embeddings, LLM needs no training).
\newline Cons: kNN and random forest scale poorly or lag gradient boosting on this data size; tabular NNs need more data/tuning and lose native categorical handling; hierarchical Bayes is slow to fit at this scale; LLM zero-shot has no calibrated uncertainty and conflicts with the CPU-only, small-image constraint.
\end{itemize}
\par
\emph{Why:} The course grades MLOps practice over raw accuracy, so the team wants a model that trains fast on CPU, handles categoricals/missing values natively, keeps artefacts small, and supports exact SHAP explanations for UC1/UC2 -- gradient boosting (LightGBM) is the best fit, with a median and Ridge baseline kept in the experiment ladder to prove it earns the added complexity.},
  ai = {Information seeking, Alternative assessment},
  response = {Accepted},
  assessment = {AI laid out the alternatives and the CPU/size/explainability constraints behind the ``why gradient boosting'' rationale in the project brief; the team reviewed and confirmed the choice, including keeping CatBoost as a challenger rather than the main model, in today's sprint planning.},
  aievidence = {Modelling plan and ``Why gradient boosting'' write-up in \texttt{docs/\allowbreak{}docs/\allowbreak{}project-\allowbreak{}brief.\allowbreak{}md} \S{}4 (2026-09-22), confirmed in today's sprint planning session.},
  evidence = {\href{https://github.com/mlops-2627q1-mds-upc/MLOPS_Recommenditos/pull/13}{PR \#13}.},
}

% Generated by tools/edn_latex.py from reports/edn.md.
% Do not edit: change reports/edn.md and run `make edn`.
\ednentry{
  id = {EDN-03},
  title = {New-market drift scenario: hold out AutoScout24 Spain instead of using DataMarket},
  date = {2026-09-22},
  milestone = {M1: Project Inception (data plan; shapes M6: Monitoring)},
  topic = {Dataset Selection, Monitoring},
  participants = {lukas2510, after supervisor feedback (S. del Rey)},
  decision = {Hold out all AutoScout24 \texttt{ES} listings (8,015, 6.8\,\%) from training, validation and calibration, and use them as the ``new market'' drift scenario for M6. The DataMarket Spanish second-hand car sample is dropped from the pipeline.},
  rationale = {\emph{Alternatives considered:}
\begin{itemize}[leftmargin=*, topsep=2pt]
\item \textbf{Option A (chosen): Hold out one AutoScout24 country (\texttt{ES}).}
\newline Pros: same schema, scrape date and portal as the training data, so detected drift comes from the market alone; every held-out listing has a price, so we can show real MAE and interval-coverage degradation, not only input drift; the same prices serve as delayed labels for the feedback loop; no mapping stage and no license issue; \texttt{ES} costs only 6.8\,\% of training rows (holding out \texttt{IT} would cost 20\,\%).
\newline Cons: less training data; the served model does not cover Spain until retrained; \texttt{country} becomes a category unseen in training, which the API must handle.
\item \textbf{Option B: Keep DataMarket as an optional stress test, restricted to makes both datasets share.}
\newline Pros: real external data, from a different portal; answers the supervisor's question with an actual experiment.
\newline Cons: market, time (2021 vs 2025), portal and brand mix still shift together, so a drift alarm cannot be attributed; needs a schema-mapping stage for \textasciitilde{}10 overlapping fields with Spanish labels and free-text colour; no license file and a commercial sample, so it cannot live in a public DVC remote.
\item \textbf{Option C: Drop DataMarket without a replacement.}
\newline Pros: simplest.
\newline Cons: leaves M6 without a realistic drift scenario.
\end{itemize}
\par
\emph{Why:} The supervisor asked us to check whether DataMarket is feasible as a new-market set given the feature mismatch. Our checks found that Spain is already in AutoScout24 (8,015 \texttt{ES} listings), 36.5\,\% of DataMarket rows are makes the model never sees and the shared mass-market makes have under 400 training rows each, and several factors shift at once (e.g. BMW median price 16,000 EUR in DataMarket vs 24,819 EUR in AutoScout24 \texttt{ES}). A drift alarm on DataMarket would therefore not be explainable. The M6 rubric asks for ``timely, actionable signals about data changes, model performance degradation'', and holding out \texttt{ES} gives a drift with a single known cause plus ground-truth prices to show the performance drop, at the cost of one filter in the split stage.},
  ai = {Alternative assessment, Recommendation},
  response = {Accepted},
  assessment = {The three options came from our own feasibility check. AI (Claude Code) explained the trade-offs against the M6 rubric and recommended Option A with \texttt{ES} as the held-out country (smaller training loss than \texttt{IT}), keeping \texttt{country} as a feature with unseen countries treated as unknown, and an optional synthetic drift scenario with a controlled cause. We accepted it because it answers the supervisor's concern with evidence, keeps the pipeline to one data source, and lets M6 show both input drift and performance degradation.},
  aievidence = {Claude Code session, 2026-09-22: prompt ``explain this to me again and give me a recommendation [...] what is feasible for our project and still in the scope of the course'' on project brief \S{}3.2; the response compared options A-C against the M6 rubric and recommended A with \texttt{ES}.},
  evidence = {\href{https://github.com/mlops-2627q1-mds-upc/MLOPS_Recommenditos/pull/13}{PR \#13}, project brief \S{}3.2 (\texttt{docs/\allowbreak{}docs/\allowbreak{}project-\allowbreak{}brief.\allowbreak{}md}).},
}

% Generated by tools/edn_latex.py from reports/edn.md.
% Do not edit: change reports/edn.md and run `make edn`.
\ednentry{
  id = {EDN-04},
  title = {Model scope: used passenger cars only},
  date = {2026-09-22},
  milestone = {M1: Project Inception},
  topic = {Problem Specification},
  participants = {@lukas2510},
  decision = {The model only covers used passenger cars: \texttt{offer\_\allowbreak{}type =\allowbreak{} U}, not pre-registered, \texttt{vehicle\_\allowbreak{}type =\allowbreak{} Car}. New cars, pre-registered cars and transporters are filtered out in the pipeline and rejected by the API.},
  rationale = {\emph{Alternatives considered:}
\begin{itemize}[leftmargin=*, topsep=2pt]
\item \textbf{Option A (chosen): used cars only.} Drops about 5\,\% of the cleaned listings (4,252 new, 3,702 pre-registered and 456 transporter listings in the raw data).
\newline Pros: matches the product (``used-car price''); new cars follow list prices, not a second-hand market, so one model does not have to learn two price mechanisms; a clear, testable scope statement.
\newline Cons: slightly less data; the API has to reject new cars.
\item \textbf{Option B: keep everything, with the flags as features.}
\newline Pros: more data, broader API.
\newline Cons: blurs the product and mixes two price mechanisms in one model.
\end{itemize}
\par
\emph{Why:} An exploratory LightGBM run showed almost no accuracy difference between the two options (MdAPE 6.8\,\% vs 6.9\,\%), so the decision rests on product clarity: a used-car price component should be scoped to used cars.},
  ai = {Information seeking, Alternative generation, Alternative assessment, Recommendation},
  response = {Accepted},
  assessment = {AI profiled the data (counts of new, pre-registered and transporter listings), ran the comparison on both scopes and recommended option A. Lukas checked that the numbers came from the real dataset and accepted the recommendation because the accuracy argument was neutral and the product argument decisive.},
  aievidence = {Claude Code session on 2026-09-22 while working on issue \#2: prompt ``Which listings are in scope for the model (car type)?'' with options A and B and the exploratory results; Lukas chose option A.},
  evidence = {\href{https://github.com/mlops-2627q1-mds-upc/MLOPS_Recommenditos/blob/main/docs/docs/problem-spec.md}{Problem specification}, section 2; \href{https://github.com/mlops-2627q1-mds-upc/MLOPS_Recommenditos/issues/2}{issue \#2}, \href{https://github.com/mlops-2627q1-mds-upc/MLOPS_Recommenditos/pull/17}{PR \#17}.},
}

% Generated by tools/edn_latex.py from reports/edn.md.
% Do not edit: change reports/edn.md and run `make edn`.
\ednentry{
  id = {EDN-05},
  title = {Supported makes: minimum listing support per make},
  date = {2026-09-22},
  milestone = {M1: Project Inception},
  topic = {Problem Specification},
  participants = {@lukas2510},
  decision = {Only makes with at least 300 listings in the cleaned used-car data (counted after the \texttt{ES} holdout, before the split) are supported; the pipeline computes the list and the API rejects other makes. With the current data this gives 11 makes covering 98.6\,\% of the used listings after the \texttt{ES} holdout (EDN-03).},
  rationale = {\emph{Alternatives considered:}
\begin{itemize}[leftmargin=*, topsep=2pt]
\item \textbf{Option A (chosen): minimum support threshold, reject the rest.}
\newline Pros: honest and testable scope; no silent low-quality predictions; easy to state in the model card.
\newline Cons: smaller product (e.g. Ford, Renault, Opel, Kia are not supported).
\item \textbf{Option B: accept all 25 makes and flag rare ones as low confidence.}
\newline Pros: broader API.
\newline Cons: weak predictions for makes with a few dozen listings; the warning logic needs its own tests and thresholds anyway.
\item \textbf{Option C: accept all makes without any rule.}
\newline Pros: simplest.
\newline Cons: no clear scope statement; hidden weaknesses for rare makes.
\end{itemize}
\par
\emph{Why:} The dataset is skewed toward premium brands (EDN-01 accepts this as a scope limit). A threshold turns that limit into an explicit, testable rule. 300 keeps 11 makes and 98.5\,\% of the data (98.6\,\% after the \texttt{ES} holdout); 200 would add three makes (0.7\,\% of the data), 500 would drop Aston Martin and Volkswagen.},
  ai = {Information seeking, Alternative generation, Alternative assessment, Recommendation},
  response = {Accepted},
  assessment = {AI counted the listings per make after scoping and deduplication, compared thresholds of 200, 300 and 500, and recommended option A. Lukas accepted it because it states the dataset's brand skew openly instead of hiding it behind a warning.},
  aievidence = {Claude Code session on 2026-09-22 while working on issue \#2: prompt ``How do we handle the 25 makes, many of which have very few listings?'' with options A to C; Lukas chose option A.},
  evidence = {\href{https://github.com/mlops-2627q1-mds-upc/MLOPS_Recommenditos/blob/main/docs/docs/problem-spec.md}{Problem specification}, section 2; \href{https://github.com/mlops-2627q1-mds-upc/MLOPS_Recommenditos/issues/2}{issue \#2}, \href{https://github.com/mlops-2627q1-mds-upc/MLOPS_Recommenditos/pull/17}{PR \#17}.},
}

% Generated by tools/edn_latex.py from reports/edn.md.
% Do not edit: change reports/edn.md and run `make edn`.
\ednentry{
  id = {EDN-06},
  title = {Success criteria for the price model},
  date = {2026-09-22},
  milestone = {M1: Project Inception},
  topic = {Performance Criteria},
  participants = {@lukas2510},
  decision = {A model is deployable when it meets SC-01 to SC-05 on the grouped test set: MdAPE \ensuremath{\leq} 9\,\%; at least 85\,\% of predictions within \ensuremath{\pm}20\,\%; MdAPE at least 30\,\% lower than the median baseline B0; MdAPE \ensuremath{\leq} 15\,\% in every input-based segment with at least 500 test rows; empirical coverage of nominal 90\,\% intervals between 88\,\% and 92\,\% for full inputs and for a partial-input scenario.},
  rationale = {\emph{Alternatives considered:}
\begin{itemize}[leftmargin=*, topsep=2pt]
\item \textbf{Option A (chosen): absolute targets plus a baseline comparison, per segment.}
\newline Pros: says whether the product is useful to users and whether the model earns its complexity; segment criterion catches models that are good on average but bad for a group; grounded in an exploratory run.
\newline Cons: targets depend on this dataset and must be revisited if the scope changes.
\item \textbf{Option B: baseline comparison only.}
\newline Pros: independent of the data.
\newline Cons: says nothing about usefulness; a weak baseline makes it easy to pass.
\item \textbf{Option C: MAPE-based targets, as in the original plan.}
\newline Pros: common and familiar.
\newline Cons: dominated by cheap cars and outliers (the same model has 9.7\,\% MAPE but 6.8\,\% MdAPE).
\end{itemize}
\par
\emph{Why:} An exploratory run on the final scope (102,695 listings, 80/20 split grouped by seller) gave B0 11.6\,\% MdAPE / 71.4\,\% within \ensuremath{\pm}20\,\% and an untuned LightGBM 6.8\,\% / 90.7\,\%. The targets keep a margin to the LightGBM values so they hold across splits, while ruling out a model only marginally better than B0. Price buckets are excluded from the segment criterion because they condition on the target. The run already shows one gap: cars older than 20 years reach 16.7\,\% MdAPE, which is recorded as a known risk rather than hidden by loosening SC-04. Re-checked after the \texttt{ES} holdout (EDN-03) on 96,831 listings: B0 11.9\,\% / 70.9\,\%, LightGBM 6.7\,\% / 90.6\,\%, cars older than 20 years 15.3\,\%; the criteria stay unchanged.},
  ai = {Information seeking, Alternative generation, Alternative assessment, Recommendation},
  response = {Accepted},
  assessment = {AI ran the exploratory baseline and LightGBM models, proposed the criteria with concrete thresholds and recommended option A. Lukas accepted it because the thresholds are tied to measured values. After the decision, AI's per-segment check found that the proposed SC-04 is not yet met for cars older than 20 years; the criterion was kept as agreed and the gap documented.},
  aievidence = {Claude Code session on 2026-09-22 while working on issue \#2: prompt ``Which kind of success criteria?'' with options A to C and the exploratory results; Lukas chose option A.},
  evidence = {\href{https://github.com/mlops-2627q1-mds-upc/MLOPS_Recommenditos/blob/main/docs/docs/problem-spec.md}{Problem specification}, sections 6 to 8; \href{https://github.com/mlops-2627q1-mds-upc/MLOPS_Recommenditos/issues/2}{issue \#2}, \href{https://github.com/mlops-2627q1-mds-upc/MLOPS_Recommenditos/pull/17}{PR \#17}.},
}

% Generated by tools/edn_latex.py from reports/edn.md.
% Do not edit: change reports/edn.md and run `make edn`.
\ednentry{
  id = {EDN-07},
  title = {Raw data hosting: import from Zenodo, never push to our own remote},
  note = {\textbf{Amended by EDN-25} (2026-09-26, PR \#27): the raw file is tracked with \texttt{dvc add} and pushed to our DagsHub remote instead.
This entry is kept as the record of what was decided on 2026-09-22 and of the reasoning EDN-25 revised.},
  date = {2026-09-22},
  milestone = {M2: Reproducibility},
  topic = {Data Versioning, Privacy},
  participants = {@lukas2510},
  decision = {The raw AutoScout24 CSV is pulled with \texttt{dvc import-\allowbreak{}url} from its Zenodo DOI, pinned to version 1.0.0, and is never \texttt{dvc push}ed to our DagsHub remote. Only the PII-stripped \texttt{data/\allowbreak{}interim} output onward is tracked and pushed like a normal pipeline artefact. Enforced by \texttt{push: false} on the file's output in \texttt{data/\allowbreak{}raw/\allowbreak{}cars.\allowbreak{}csv.\allowbreak{}dvc}, which makes \texttt{dvc push} skip it, and a test that the field stays there.},
  rationale = {\emph{Alternatives considered:}
\begin{itemize}[leftmargin=*, topsep=2pt]
\item \textbf{Option A (chosen): \texttt{dvc import-\allowbreak{}url} from Zenodo, raw file never pushed.}
\newline Pros: the raw file carries PII (\texttt{street}, \texttt{zip}, exact coordinates, \texttt{seller\_\allowbreak{}company\_\allowbreak{}name} for private sellers); the Zenodo DOI is already immutable and versioned, so re-fetching it with \texttt{dvc update} gives the same reproducibility as a cached copy without a second, PII-bearing copy on our own remote (likely public); matches the data-minimization principle for sensitive raw data.
\newline Cons: fetching the raw data needs live access to Zenodo; reproducibility depends on Zenodo keeping this exact file available; \texttt{dvc push} uploads imported files like any other output, so the ``never pushed'' half needs \texttt{push: false} on the output, not just the command choice, to actually hold; because the file is never on the remote, a fresh clone gets it with \texttt{dvc update} instead of \texttt{dvc pull}, and \texttt{dvc repro} fails until it has.
\item \textbf{Option B: \texttt{dvc add} the raw file like any other pipeline input and push it to DagsHub.}
\newline Pros: consistent with how every other artefact in the project is tracked; does not depend on Zenodo staying reachable.
\newline Cons: re-hosts a PII-bearing file on our own remote, which is likely public; the second copy adds no reproducibility over the immutable Zenodo DOI, only exposure.
\end{itemize}
\par
\emph{Why:} Option A was already the direction noted (but left \texttt{[open]}) in the project brief's data-issues section. It was surfaced as an undecided requirement-vs-brief contradiction during the review of PR \#19 (``requirements.md'' stated the raw data is never re-hosted as a settled \texttt{NFR-\allowbreak{}08}, while the brief still called the same question open, and the existing \texttt{data-\allowbreak{}versioning.\allowbreak{}md} example showed the opposite pattern, \texttt{dvc add data/\allowbreak{}raw/\allowbreak{}cars.\allowbreak{}csv}). Data minimization for a file containing indirect identifiers of private sellers outweighs the convenience of a uniform tracking pattern; the Zenodo DOI's own immutability and versioning already give the reproducibility that pushing a copy would otherwise provide.},
  ai = {Information seeking, Alternative assessment, Recommendation},
  response = {Accepted},
  assessment = {While reviewing PR \#19 for SOTA fit and course requirements, AI noticed that \texttt{NFR-\allowbreak{}08}'s ``never re-hosted'' wording silently settled the brief's open \texttt{dvc import-\allowbreak{}url} question without going through the EDN process, and that the project's own \texttt{data-\allowbreak{}versioning.\allowbreak{}md} example contradicted it. Asked for a recommendation, AI laid out options A and B with pros and cons, recommended A on privacy and data-minimization grounds, and flagged that \texttt{dvc import-\allowbreak{}url} alone does not prevent a re-push and needs a CI check to be enforceable. Lukas reviewed the reasoning and confirmed option A. On 2026-09-23, asked whether NFR-08 was overkill, AI proposed \texttt{push: false} on the output instead of that CI check against the DagsHub remote, and verified it in a scratch repository with DVC 3.67.1: without the field \texttt{dvc push} uploads the imported file, with it the remote stays empty, also after a \texttt{dvc update} from a changed source. The same test showed that in a fresh clone \texttt{dvc repro} does not re-fetch an imported file and a plain \texttt{dvc pull} exits with an error, while \texttt{dvc update} does fetch it, which corrected the instructions in \texttt{data-\allowbreak{}versioning.\allowbreak{}md}. Lukas accepted the change of mechanism; the decision itself stayed the same.},
  aievidence = {Claude Code session on 2026-09-22 while reviewing PR \#19: AI flagged the NFR-08/brief contradiction, explained it in detail on request, then gave a recommendation and cited data-minimization and Zenodo DOI immutability as the SOTA rationale for option A; Lukas replied ``yes please do so''. Claude Code session on 2026-09-23: Lukas asked (translated from German) ``NFR-08 privacy, does that still need adjusting? how would you adjust it so it is no longer excessive'', AI proposed \texttt{push: false} with the scratch-repository evidence; Lukas replied ``yes please do that''.},
  evidence = {\href{https://github.com/mlops-2627q1-mds-upc/MLOPS_Recommenditos/blob/main/docs/docs/requirements.md}{requirements} NFR-08; \href{https://github.com/mlops-2627q1-mds-upc/MLOPS_Recommenditos/blob/main/docs/docs/specification.md}{specification} NFR-08 and ``Decisions pending confirmation''; \href{https://github.com/mlops-2627q1-mds-upc/MLOPS_Recommenditos/blob/main/docs/docs/data-versioning.md}{Data versioning}; \href{https://github.com/mlops-2627q1-mds-upc/MLOPS_Recommenditos/blob/main/docs/docs/project-brief.md}{project brief} section 3.3; \href{https://github.com/mlops-2627q1-mds-upc/MLOPS_Recommenditos/pull/19}{PR \#19}.},
}

% Generated by tools/edn_latex.py from reports/edn.md.
% Do not edit: change reports/edn.md and run `make edn`.
\ednentry{
  id = {EDN-08},
  title = {Model loading: bake into the API image via \texorpdfstring{\texttt{dvc pull}}{dvc pull} at CI build time, not the MLflow registry at runtime},
  date = {2026-09-22},
  milestone = {M4: Model Deployment},
  topic = {ML System Design},
  participants = {@lukas2510},
  decision = {The API image is built by CI with the current production model already baked in, fetched with \texttt{dvc pull} pinned to whatever \texttt{dvc.\allowbreak{}lock} on \texttt{main} points to. Promoting a model is a normal merge to \texttt{main} that updates that pointer. The API never calls the MLflow registry at build or run time; MLflow remains the experiment-tracking and audit record of which run was chosen as production.},
  rationale = {\emph{Alternatives considered:}
\begin{itemize}[leftmargin=*, topsep=2pt]
\item \textbf{Option A: pull from the MLflow model registry at API startup.} As originally stated in requirements.md FR-12.
\newline Pros: matches the target architecture's ``MLflow tracking + model registry'' component as originally drafted; a model can be promoted without a redeploy.
\newline Cons: startup now depends on DagsHub being reachable, which works against NFR-05 (ready within 30 s) and NFR-12 (a self-contained stack); a DagsHub outage right before a presentation could break the whole demo.
\item \textbf{Option B: CD resolves the MLflow registry alias and bakes the model into the image on promotion.}
\newline Pros: avoids the runtime dependency of option A while still using the registry as the deployment-time source of truth.
\newline Cons: needs a new trigger connecting an MLflow promotion event to a GitHub Actions build (webhook or poller), which is infrastructure the course does not demo yet and is disproportionate for a 5-person, one-semester, single-VM project; nothing watches it if it silently stops firing.
\item \textbf{Option C (chosen): bake into the image via \texttt{dvc pull} at CI build time; promotion is a merge to \texttt{main}.}
\newline Pros: matches the course demo's own deployment recipe (\texttt{dvc pull} the model before running \texttt{uvicorn}), just containerized; reuses GitHub Flow, already the project's branching model, as the promotion mechanism with no new trigger to build; leans on DVC, the single heaviest-weighted practice in the course rubric (15/100), for the thing that most needs to be reproducible; avoids the runtime dependency of option A.
\newline Cons: MLflow's model registry is no longer part of the deployment path, only tracking/audit, which is a smaller role than the target architecture originally sketched.
\end{itemize}
\par
\emph{Why:} Given the course's own demo pattern, the rubric's weight on DVC over MLflow, and the team's size and one-semester timeline, option C gives the same startup-independence as option B without inventing new promotion infrastructure. MLflow keeps its full credit for the experiment-tracking practice; it is just not on the deployment-time critical path.},
  ai = {Information seeking, Alternative generation, Alternative assessment, Recommendation},
  response = {Accepted},
  assessment = {While reviewing PR \#19, AI first flagged that FR-12's runtime registry pull (option A) worked against NFR-05 and NFR-12 and recommended option B. Asked to weigh that against the project's actual scope and tools, AI checked \texttt{references/\allowbreak{}course-\allowbreak{}demos.\allowbreak{}md}, found the demo's own recipe pulls the model via DVC rather than MLflow, cross-referenced the rubric's point weights, and revised its recommendation to option C, naming the trade-off (MLflow registry no longer in the deployment path) explicitly. Lukas reviewed the reasoning and confirmed option C.},
  aievidence = {Claude Code session on 2026-09-22 while reviewing PR \#19: AI recommended option B, Lukas asked ``considering our project scope and the tools we will use is this the best option for us?'', AI re-derived the recommendation from \texttt{references/\allowbreak{}course-\allowbreak{}demos.\allowbreak{}md} and the rubric weights and proposed option C instead; Lukas replied ``yes please do that''.},
  evidence = {\href{https://github.com/mlops-2627q1-mds-upc/MLOPS_Recommenditos/blob/main/docs/docs/specification.md}{Specification} FR-12; \href{https://github.com/mlops-2627q1-mds-upc/MLOPS_Recommenditos/blob/main/docs/docs/project-brief.md}{project brief} section 5 (target architecture); \href{https://github.com/mlops-2627q1-mds-upc/MLOPS_Recommenditos/blob/main/references/course-demos.md}{references/course-demos.md} ``API and deployment (M4)''; \href{https://github.com/mlops-2627q1-mds-upc/MLOPS_Recommenditos/pull/19}{PR \#19}.},
}

% Generated by tools/edn_latex.py from reports/edn.md.
% Do not edit: change reports/edn.md and run `make edn`.
\ednentry{
  id = {EDN-09},
  title = {Bump to Python 3.12, to use real SHAP instead of a workaround},
  date = {2026-09-22},
  milestone = {M2: Reproducibility},
  topic = {Tooling, Explainability},
  participants = {@lukas2510},
  decision = {\texttt{requires-\allowbreak{}python} is \texttt{\textasciitilde{}=\allowbreak{}3.\allowbreak{}12.\allowbreak{}0} (was \texttt{\textasciitilde{}=\allowbreak{}3.\allowbreak{}11.\allowbreak{}0}). \texttt{pyproject.\allowbreak{}toml}, \texttt{uv.\allowbreak{}lock}, and the Python version in \texttt{.\allowbreak{}github/\allowbreak{}workflows/\allowbreak{}ci.\allowbreak{}yml} are updated, and \texttt{shap} becomes a usable dependency for offline analysis. Which SHAP implementation FR-08's explanation endpoint uses is decided separately in EDN-11, which amends this entry: the API computes contributions from the booster's own export, and \texttt{shap} stays out of the API image.},
  rationale = {\emph{Alternatives considered:}
\begin{itemize}[leftmargin=*, topsep=2pt]
\item \textbf{Option A: use LightGBM's native \texttt{predict(pred\_\allowbreak{}contrib=\allowbreak{}True)} instead of the \texttt{shap} package.} Gives the same exact TreeSHAP values without touching \texttt{shap} or its Python 3.12 requirement, and avoids \texttt{numba}/\texttt{llvmlite} in the API image.
\newline Pros: no project-wide version bump needed this early in the semester; smaller API image.
\newline Cons: reimplements what \texttt{shap} already provides; still needs the real \texttt{shap} package for any offline/notebook analysis (global importance, summary plots), so the Python constraint would resurface there anyway; more code to maintain for a workaround to a constraint that turned out to be removable.
\item \textbf{Option B (chosen): bump to Python 3.12 project-wide.} Verified empirically: \texttt{uv lock -\allowbreak{}-\allowbreak{}python 3.\allowbreak{}12} and \texttt{uv sync -\allowbreak{}-\allowbreak{}python 3.\allowbreak{}12} resolved and installed cleanly with the full planned M2-M5 dependency set added (\texttt{shap}, \texttt{mlflow>=\allowbreak{}2.\allowbreak{}22,<3}, \texttt{lightgbm}, \texttt{catboost}, \texttt{mapie}, \texttt{fastapi}, \texttt{uvicorn}, \texttt{pydantic}, \texttt{great-\allowbreak{}expectations}, \texttt{dvc}, \texttt{pytest-\allowbreak{}cov}, \texttt{codecarbon}, \texttt{httpx}), and \texttt{shap.\allowbreak{}TreeExplainer} on a trained LightGBM model produced correct SHAP values under 3.12 in a functional test.
\newline Pros: uses \texttt{shap} directly with no workaround; every other planned dependency (checked against PyPI's \texttt{requires\_\allowbreak{}python} metadata for the exact pinned versions, not just ``latest'') only has a lower bound compatible with 3.12, so nothing else in the stack is put at risk; no Dockerfile exists yet to update, only three lines in \texttt{ci.\allowbreak{}yml}.
\newline Cons: a project-wide version bump this early, touching \texttt{pyproject.\allowbreak{}toml}, \texttt{uv.\allowbreak{}lock} and CI, for the sake of one requirement.
\end{itemize}
\par
\emph{Why:} The constraint was checked empirically rather than assumed from documentation: a scratch resolve of the full planned dependency set under Python 3.12 succeeded with no conflicts, and \texttt{shap.\allowbreak{}TreeExplainer} functionally worked. With no blocker found, option B removes the constraint at its source instead of working around it, and the change is small (no Dockerfile yet, three \texttt{ci.\allowbreak{}yml} lines, \texttt{pyproject.\allowbreak{}toml}, \texttt{uv.\allowbreak{}lock}).},
  ai = {Information seeking, Alternative generation, Alternative assessment, Recommendation},
  response = {Accepted},
  assessment = {While reviewing PR \#19, AI proposed option A as a way around the Python 3.11/SHAP 0.52 conflict documented in the brief. Asked to check whether the version could be bumped instead, AI verified this empirically (a scratch \texttt{uv lock}/\texttt{uv sync} under 3.12 with the full planned dependency set, plus a functional \texttt{TreeExplainer} test) rather than relying on PyPI metadata alone, found no blocker, and presented both options with the empirical evidence. Lukas reviewed it and chose to bump.},
  aievidence = {Claude Code session on 2026-09-22 while reviewing PR \#19: AI first recommended \texttt{pred\_\allowbreak{}contrib} as a workaround, Lukas asked ``please check if we can bump our python version,'' AI ran an empirical \texttt{uv lock}/\texttt{uv sync}/import/functional check under Python 3.12 with the full planned stack, found no blocker, and reported the evidence; Lukas replied ``ok please then bump the version in the PR an adjust the requriemnt.''},
  evidence = {\href{https://github.com/mlops-2627q1-mds-upc/MLOPS_Recommenditos/blob/main/docs/docs/specification.md}{Specification} FR-08; \href{https://github.com/mlops-2627q1-mds-upc/MLOPS_Recommenditos/blob/main/docs/docs/project-brief.md}{project brief} section 6 (tooling constraints); \texttt{pyproject.\allowbreak{}toml}, \texttt{.\allowbreak{}github/\allowbreak{}workflows/\allowbreak{}ci.\allowbreak{}yml}; \href{https://github.com/mlops-2627q1-mds-upc/MLOPS_Recommenditos/pull/19}{PR \#19}.},
}

% Generated by tools/edn_latex.py from reports/edn.md.
% Do not edit: change reports/edn.md and run `make edn`.
\ednentry{
  id = {EDN-10},
  title = {Availability target replaced by recovery time plus a presentation-window commitment},
  date = {2026-09-22},
  milestone = {M6: Monitoring},
  topic = {Resource Monitoring, ML System Design},
  participants = {@lukas2510},
  decision = {NFR-05 no longer commits to ``99\,\% uptime in the weeks before each presentation.'' It now commits to: (1) automatic recovery, the API reaches a healthy \texttt{/\allowbreak{}health} within 30 s of a crash or VM reboot, verified by a chaos test; (2) zero unplanned downtime during each live presentation and its warm-up, checked manually right before presenting; (3) outside those windows, uptime on the single, unattended, no-SLA VM is monitored and reported via Better Uptime and Grafana, not contractually targeted. Better Uptime is added to the architecture, the M6 tool list and the reference links, none of which named it before.},
  rationale = {\emph{Alternatives considered:}
\begin{itemize}[leftmargin=*, topsep=2pt]
\item \textbf{Option A: keep 99\,\% uptime over ``the weeks before each presentation.''}
\newline Pros: a single, simple number; easy to write into the report.
\newline Cons: nothing backs it: one VM, no redundancy, no SLA from FIB Virtech, no on-call team; the only listed check (Prometheus \texttt{up}) can't see the host itself going down; the measurement window is unbounded and unverifiable after the fact, so the number could not actually be defended if a teacher asked how it was checked.
\item \textbf{Option B: lower the aggregate target (e.g. 95\,\%) instead of dropping it.}
\newline Pros: keeps a single reportable number; slightly more honest than 99\,\%.
\newline Cons: still as unfalsifiable as 99\,\% without redundancy or an on-call team; still hides that the requirement was actually about two specific presentation days, not a general SLA the team can neither guarantee nor is expected to for a course project.
\item \textbf{Option C (chosen): split into recovery time, a narrow presentation-window commitment, and a monitored-not-targeted general period.}
\newline Pros: the recovery-time claim is testable today with a chaos test, not only after weeks of passive observation; the presentation-window claim is the one that actually matters for grading and is narrow enough for the team to actively watch; the monitored-general-period framing is honest about the single-VM, no-redundancy, no-SLA deployment while still giving the M6 monitoring practice something to show (Better Uptime + Grafana), which is what the rubric actually grades, not a specific SLA number.
\newline Cons: three clauses instead of one number, slightly longer to state.
\end{itemize}
\par
\emph{Why:} M6's Resource Monitoring criterion asks whether the monitoring anticipates real-world challenges and gives actionable insight, not whether a specific SLA percentage was hit. A number the team has no redundancy or on-call capacity to guarantee is a liability if questioned in the presentation, not evidence of the practice. Splitting the requirement matches what is actually controllable (recovery time, presentation-day behaviour) from what can only be observed and reported (general uptime on a single, unattended, no-SLA VM), and it adds Better Uptime, a tool the course's own M6 schedule already names but the docs previously omitted.},
  ai = {Information seeking, Alternative generation, Alternative assessment, Recommendation},
  response = {Accepted},
  assessment = {During the PR \#19 review, AI had already flagged that Prometheus's \texttt{up} metric cannot observe the host going down and that Better Uptime was missing from the docs despite being a named M6 course tool. Asked directly what availability target would actually be feasible given the single-VM, no-redundancy deployment, AI proposed the three-part split (recovery time, presentation window, monitored general period) with rationale grounded in the M6 rubric wording and the deployment's real constraints, and named option B as a weaker alternative. Lukas reviewed it and confirmed the split.},
  aievidence = {Claude Code session on 2026-09-22 while reviewing PR \#19: Lukas said ``i dont like that [99\% uptime]. what NFR that goes in the uptime direction is feasible with our system and constraints?''; AI proposed the three-part split and the addition of Better Uptime; Lukas replied ``ok do that.''},
  evidence = {\href{https://github.com/mlops-2627q1-mds-upc/MLOPS_Recommenditos/blob/main/docs/docs/requirements.md}{requirements} NFR-05; \href{https://github.com/mlops-2627q1-mds-upc/MLOPS_Recommenditos/blob/main/docs/docs/specification.md}{specification} NFR-05; \href{https://github.com/mlops-2627q1-mds-upc/MLOPS_Recommenditos/blob/main/docs/docs/project-brief.md}{project brief} section 5 (target architecture) and section 7 (milestones); \href{https://github.com/mlops-2627q1-mds-upc/MLOPS_Recommenditos/blob/main/references/MLOps-lab.md}{MLOps-lab.md} session 11; \href{https://github.com/mlops-2627q1-mds-upc/MLOPS_Recommenditos/pull/19}{PR \#19}.},
}

% Generated by tools/edn_latex.py from reports/edn.md.
% Do not edit: change reports/edn.md and run `make edn`.
\ednentry{
  id = {EDN-11},
  title = {Keep \texorpdfstring{\texttt{shap}}{shap} out of the API image; serving uses the booster's native SHAP export},
  date = {2026-09-22},
  milestone = {M4: Model Deployment},
  topic = {ML System Design, Explainability},
  participants = {@lukas2510},
  decision = {Amends EDN-09's serving-path clause. FR-08's per-request explanation is computed from the trained booster's own SHAP export (LightGBM \texttt{predict(pred\_\allowbreak{}contrib=\allowbreak{}True)}, or CatBoost \texttt{get\_\allowbreak{}feature\_\allowbreak{}importance(type=\allowbreak{}\textquotedbl{}ShapValues\textquotedbl{})} if CatBoost is ever the deployed model), not the \texttt{shap} package. \texttt{shap} stays a training/notebook-only dependency, used for offline global analysis (summary plots, dataset-level importance), and is never installed in the API image. The Python 3.12 bump from EDN-09 is unaffected and still stands, since it is what lets \texttt{shap} run at all for that offline use.},
  rationale = {\emph{Alternatives considered:}
\begin{itemize}[leftmargin=*, topsep=2pt]
\item \textbf{Option A: keep \texttt{shap.\allowbreak{}TreeExplainer} in the serving path, as EDN-09 originally set FR-08 to.}
\newline Pros: one implementation for both the offline/notebook analysis and the serving path, nothing to keep in sync.
\newline Cons: \texttt{shap} unconditionally imports \texttt{numba} and \texttt{llvmlite} at module load, even to use only \texttt{TreeExplainer}, confirmed by deleting them and watching a bare \texttt{import shap} fail; a real Docker build of the serving stack (\texttt{shap}, \texttt{lightgbm}, \texttt{mapie}, \texttt{fastapi}, \texttt{uvicorn}, \texttt{pydantic}, \texttt{numpy}, \texttt{pandas}) measured 702 MB, leaving only about 30\,\% of NFR-04's 1 GB budget for the baked-in model artefact, app code and everything else still to be added.
\item \textbf{Option B (chosen): compute serving-time explanations from the booster's native SHAP export; keep \texttt{shap} for offline analysis only.}
\newline Pros: the same Docker build without \texttt{shap} measured 486 MB, about 51\,\% headroom instead of 30\,\%; the native export is verified bit-identical to \texttt{shap.\allowbreak{}TreeExplainer} (max absolute difference 0.0 across all features and the base value in a direct comparison); the pattern generalises to CatBoost, the brief's stated challenger model (EDN-02), via its own native SHAP support, so it does not need revisiting if the challenger wins; \texttt{shap} is still fully available where it is actually needed, offline global analysis, unconstrained by image size.
\newline Cons: two call sites compute the same values (the booster's native export in the API, \texttt{shap.\allowbreak{}TreeExplainer} in training/notebook code) instead of one; a training-time test is needed to keep them provably in sync (added to FR-08's ``Verified by'').
\end{itemize}
\par
\emph{Why:} The image-size argument is independent of the Python-version problem EDN-09 solved, and EDN-09 folded both into one decision without weighing the size evidence, which was not yet measured at the time. Measured directly (Docker builds, a deletion test proving \texttt{numba}/\texttt{llvmlite} are unconditional runtime imports, and a numeric comparison proving the native export and \texttt{shap.\allowbreak{}TreeExplainer} produce identical values), option B costs nothing in correctness and gives back real budget headroom under NFR-04, while also removing a dependency choice tied to a single model family.},
  ai = {Information seeking, Alternative assessment, Recommendation},
  response = {Accepted},
  assessment = {AI had proposed the native-export approach earlier in the PR \#19 review, before EDN-09 was decided. Asked whether it was still relevant after the Python 3.12 bump, AI separated the two motivations (version conflict, now resolved; image size, still open), then built real Docker images for both options and a deletion test to confirm \texttt{numba}/\texttt{llvmlite} are unconditional, rather than assuming it from documentation. Asked whether this was the best option or if there were better ones, AI additionally ruled out a leaner Docker build (does not work, proven by the deletion test) and dropping per-instance SHAP entirely (would weaken FR-08 and lose the course's stated M3 credit for using SHAP), and pointed out CatBoost's equivalent native support. Lukas reviewed the evidence and confirmed option B.},
  aievidence = {Claude Code session on 2026-09-22 while reviewing PR \#19: AI's original \texttt{pred\_\allowbreak{}contrib} suggestion predated the Python 3.12 bump; Lukas asked ``is this still relevant?'', AI re-derived it with concrete measurements; Lukas asked ``is this what we would want? or are there better options?'', AI built real Docker images (702 MB vs. 486 MB) and ruled out the alternatives; Lukas replied ``ok apply it.''},
  evidence = {\href{https://github.com/mlops-2627q1-mds-upc/MLOPS_Recommenditos/blob/main/docs/docs/specification.md}{Specification} FR-08; \href{https://github.com/mlops-2627q1-mds-upc/MLOPS_Recommenditos/blob/main/docs/docs/project-brief.md}{project brief} section 6 (tooling constraints); \hyperref[EDN-09]{EDN-09}; \href{https://github.com/mlops-2627q1-mds-upc/MLOPS_Recommenditos/pull/19}{PR \#19}.},
}

% Generated by tools/edn_latex.py from reports/edn.md.
% Do not edit: change reports/edn.md and run `make edn`.
\ednentry{
  id = {EDN-12},
  title = {Retraining and promotion are human-triggered, not automated},
  date = {2026-09-22},
  milestone = {M6: Monitoring},
  topic = {ML System Design, Model Performance},
  participants = {@lukas2510},
  decision = {FR-15 formalises the retrain-after-drift loop that EDN-03 was designed to demonstrate: when FR-14's drift job flags input drift on the \texttt{ES} replay (per NFR-11), a team member retrains with \texttt{ES} included and opens the promotion PR by hand. No automation watches the drift job and triggers this on its own. The candidate still must pass NFR-01's gate before the PR is merged; merging is the promotion (EDN-08), and rollback is redeploying the previous image tag (EDN-08).},
  rationale = {\emph{Alternatives considered:}
\begin{itemize}[leftmargin=*, topsep=2pt]
\item \textbf{Option A: fully automated retraining and promotion.} A watcher on the drift job's output triggers \texttt{dvc repro} and opens the promotion PR automatically when drift is confirmed.
\newline Pros: reads as more ``extended/innovative'' for the M6 System Design criterion; no manual step to forget.
\newline Cons: needs new infrastructure (something connecting Alibi Detect's output to a CI trigger) built and debugged in the last two weeks before the Dec 9 presentation, exactly when there is the least slack to recover if it breaks; higher risk for a 5-person team with no prior automation of this kind in the project.
\item \textbf{Option B (chosen): human-triggered retraining and promotion.} The drift job surfaces the signal (Grafana); a team member runs the retrain and opens the PR once NFR-01's gate passes.
\newline Pros: still fully closes the loop EDN-03 was built to demonstrate, and still satisfies the M6 rubric's ``enabling confident investigation and response'', a human response is a legitimate response, not a lesser one; far less new infrastructure, and safer to demo live; matches the team's size and the amount of time left before M6.
\newline Cons: less visually impressive than full automation; a person has to actually notice the signal and act on it.
\end{itemize}
\par
\emph{Why:} The \texttt{ES} holdout exists specifically to demonstrate this loop closing (EDN-03), so leaving it as an unformalised sentence in the brief risked it quietly not getting built once M6 crunch hits. Given the timeline, automating the trigger is a real risk for a low grading upside over the human-triggered version, which already gets full credit for ``closing the loop'' and ``responding to a signal'' without the extra failure surface this late in the semester.},
  ai = {Information seeking, Alternative generation, Alternative assessment, Recommendation},
  response = {Accepted},
  assessment = {While reviewing PR \#19, AI noted that the retrain-after-drift plan in the brief had never become a testable requirement, unlike everything else on the requirements page, and that EDN-08 had already defined the promotion and rollback mechanism without it being referenced by any FR. Asked whether this and a related NFR-11 tightening were still relevant and worth doing, AI confirmed both were unaddressed, proposed the NFR-11 fix directly (no real alternative to weigh), and for retraining, named the automation-level choice as a genuine decision, laid out both options with the project's timeline as the deciding factor, and recommended the human-triggered option. Lukas confirmed it.},
  aievidence = {Claude Code session on 2026-09-22 while reviewing PR \#19: AI's original review had flagged both the trivially-passable NFR-11 and the missing retraining requirement; Lukas asked whether they were still relevant and whether they were overkill, AI confirmed relevance for both and recommended the human-triggered option for retraining with reasoning tied to the M6 timeline; Lukas replied ``yes.''},
  evidence = {\href{https://github.com/mlops-2627q1-mds-upc/MLOPS_Recommenditos/blob/main/docs/docs/requirements.md}{requirements} FR-14, FR-15, NFR-01, NFR-11; \href{https://github.com/mlops-2627q1-mds-upc/MLOPS_Recommenditos/blob/main/docs/docs/specification.md}{specification} FR-14 and FR-15; \href{https://github.com/mlops-2627q1-mds-upc/MLOPS_Recommenditos/blob/main/docs/docs/project-brief.md}{project brief} section 3.2 (\texttt{ES} new-market drift scenario) and section 5 (target architecture); \hyperref[EDN-03]{EDN-03}, \hyperref[EDN-08]{EDN-08}; \href{https://github.com/mlops-2627q1-mds-upc/MLOPS_Recommenditos/pull/19}{PR \#19}.},
}

% Generated by tools/edn_latex.py from reports/edn.md.
% Do not edit: change reports/edn.md and run `make edn`.
\ednentry{
  id = {EDN-13},
  title = {Keep the \texorpdfstring{\texttt{/\allowbreak{}feedback}}{/feedback} endpoint, but reachable only from inside the Compose network},
  date = {2026-09-23},
  milestone = {M6: Monitoring},
  topic = {API Design, ML System Design, Monitoring},
  participants = {@lukas2510},
  decision = {\texttt{POST /\allowbreak{}feedback} (FR-10) stays in the component and is no longer \texttt{[open]}. It takes a request ID and an observed price, stored as a delayed label, and FR-14 joins those labels to the prediction log to report the real error (MdAPE) and the empirical interval coverage per window, not only input drift. The endpoint is not publicly reachable: a reverse proxy in front of the API refuses \texttt{/\allowbreak{}feedback} from outside, so only the replay job inside the Compose network reaches it; the \texttt{X-\allowbreak{}API-\allowbreak{}Key} check stays behind that as a second layer in case the proxy rule is ever wrong. Because the key therefore never crosses the public internet, NFR-09 no longer requires TLS, and the previously \texttt{[proposed]} TLS item is dropped.},
  rationale = {\emph{Alternatives considered:}
\begin{itemize}[leftmargin=*, topsep=2pt]
\item \textbf{Option A: keep \texttt{/\allowbreak{}feedback} as a publicly reachable endpoint, protected by the API key alone.} The variant the requirements page described before this decision.
\newline Pros: one fewer component in the deployment; the endpoint behaves like a real public feedback channel.
\newline Cons: the key would travel in cleartext unless the VM gets TLS, and the FIB Virtech VM sits behind NAT with a forwarded port, so an HTTP-01 Let's Encrypt challenge on port 80 of a domain we control is most likely not possible; that left the whole point of adding the key resting on an unverified assumption about the VM.
\item \textbf{Option B: drop \texttt{/\allowbreak{}feedback}; the drift job reads the \texttt{ES} prices directly from the versioned artefact and joins them to the prediction log by row ID.}
\newline Pros: no write path, no key, no TLS question, no shared store needed for feedback; least code.
\newline Cons: the labels never enter the system through an interface, so the delayed-label mechanism exists only inside one batch job; the architecture no longer shows how ground truth would arrive in operation, which is the part M6 actually grades (``model performance degradation, and user impact''); EDN-03 chose the \texttt{ES} holdout precisely because every held-out listing has a price, and that argument only pays off if those prices flow back in somewhere.
\item \textbf{Option C (chosen): keep the endpoint, but make it unreachable from outside (reverse-proxy rule), with the API key as a second layer.}
\newline Pros: keeps the feedback loop as a real interface, so FR-14 can report error and coverage and FR-15's retraining is triggered by a measured degradation rather than by an input-distribution change alone; removes the TLS dependency entirely instead of leaving it pending on a VM capability nobody has confirmed; the attack surface of the only writing endpoint drops to zero from the internet; the reverse proxy is a component the course demo already suggests (nginx in front of \texttt{uvicorn}).
\newline Cons: one more component in \texttt{docker compose} and one proxy rule that has to be right, which is why the key is kept behind it and the deployment smoke test asserts from outside the VM that \texttt{/\allowbreak{}feedback} is refused.
\end{itemize}
\par
\emph{Why:} The endpoint's existence could not stay \texttt{[open]}, because FR-14, FR-15, NFR-03 and NFR-09 all depended on it. On whether to keep it, the decisive argument is what M6 is graded on: without labels the monitoring can only show that the inputs changed, while with them it can show that the error rises and the nominal 90\,\% intervals lose their coverage, which is also the stronger trigger for FR-15's retraining. The marginal cost is small, since the storage that dominates the work is needed for FR-13's prediction log anyway. On exposure, option C removes an unverified assumption (TLS on the Virtech VM) rather than carrying it as a pending decision, and it keeps the honest framing that traffic and labels are replayed from the \texttt{ES} holdout, not produced by real users.},
  ai = {Information seeking, Alternative generation, Alternative assessment, Recommendation},
  response = {Accepted},
  assessment = {During the review of PR \#19, AI flagged that FR-10 carried \texttt{[open]} and \texttt{[decided]} markers in the same row and that four other requirements silently depended on it, then laid out options A to C with the M6 rubric wording and the likely NAT setup of the Virtech VM as the deciding factors, and recommended keeping the endpoint while making it internal. Asked to explain the recommendation in more detail, AI also named the counter-argument itself (the labels are simulated, so a critic could call the endpoint a facade) and the mitigation (state plainly in the report that traffic and labels are replayed). Lukas reviewed the reasoning and accepted the recommendation.},
  aievidence = {Claude Code session on 2026-09-23 while reviewing PR \#19: AI recommended keeping \texttt{/\allowbreak{}feedback} and combining it with internal-only exposure; Lukas asked ``erkläre mir FR-10 genauer und warum du das eine vorschlägst'', AI explained the delayed-label mechanism, the dependencies and the three options; Lukas replied ``ok ich mag deine empfehlung ändere es so''.},
  evidence = {\href{https://github.com/mlops-2627q1-mds-upc/MLOPS_Recommenditos/blob/main/docs/docs/specification.md}{Specification} FR-10, FR-14, NFR-09; \href{https://github.com/mlops-2627q1-mds-upc/MLOPS_Recommenditos/blob/main/docs/docs/project-brief.md}{project brief} sections 3.2 and 5; \hyperref[EDN-03]{EDN-03}, \hyperref[EDN-12]{EDN-12}; \href{https://github.com/mlops-2627q1-mds-upc/MLOPS_Recommenditos/pull/19}{PR \#19}.},
}

% Generated by tools/edn_latex.py from reports/edn.md.
% Do not edit: change reports/edn.md and run `make edn`.
\ednentry{
  id = {EDN-14},
  title = {NFR-11's drift control is an i.i.d. sample, not a seller-grouped one},
  note = {\textbf{Amended by EDN-26} (2026-09-29): the window is 1,000 requests, not 100, and \texttt{model} is excluded from the comparison.
The entry below stated a window and a control clause that were measured in two different runs and never together.
\textbf{Evidence regenerated on 2026-09-29.} The numbers below are the 2026-09-23 run and are left as they were measured.
\texttt{nfr11\_\allowbreak{}results.\allowbreak{}json} was re-run after EDN-22 removed the listings registered after the reference date, which moves the
scoped, deduplicated count from 105,431 to 105,405 and the \texttt{ES} rows the API accepts from 5,981 to 5,979.},
  date = {2026-09-23},
  milestone = {M6: Monitoring},
  topic = {Monitoring, Performance Criteria},
  participants = {@lukas2510},
  decision = {NFR-11 is rewritten after measuring it on the real dataset. The drift job compares each window against a reference of 10,000 listings sampled from the training data. The \texttt{ES} replay must be flagged within its first 100 requests (was 1,000) and must name at least three drifted features other than \texttt{country\_\allowbreak{}code}. The control is an i.i.d. sample of held-out listings, of which at most 1 in 20 windows may be flagged; a seller-grouped sample is explicitly not a valid control.},
  rationale = {\emph{Alternatives considered:}
\begin{itemize}[leftmargin=*, topsep=2pt]
\item \textbf{Option A: keep the seller-grouped test split as the control, as NFR-11 originally said.}
\newline Pros: the strongest possible claim, namely that the detector stays quiet even on sellers it has never seen; reuses the split the evaluation protocol already produces.
\newline Cons: measured impossible, not merely hard. 200 of 200 seller-grouped control windows of 1,000 rows were flagged, at every reference size from 2,000 to 76,142 rows. Held-out dealers are genuinely a different distribution: their mean effect size is larger than Spain's on \texttt{make} (0.117 vs 0.070), \texttt{gears} (0.130 vs 0.034), \texttt{mileage\_\allowbreak{}km\_\allowbreak{}raw} (0.108 vs 0.091) and \texttt{age\_\allowbreak{}years} (0.113 vs 0.099), and equal on \texttt{model} (0.245 vs 0.244). Keeping this control would make NFR-11 permanently red for a reason that has nothing to do with monitoring quality.
\item \textbf{Option B: keep the seller-grouped control but add an effect-size floor on top of the p-value.}
\newline Pros: would have kept the stronger claim and is a standard remedy for over-powered statistical tests.
\newline Cons: measured not to separate the two cases. At a floor of 0.10, \texttt{ES} and the control both have 5 features above it; at 0.15 it is 2 against 1. Rejected on evidence, not on principle.
\item \textbf{Option C: drop the control clause and only require that \texttt{ES} is flagged.}
\newline Pros: simplest; the \texttt{ES} half passes comfortably.
\newline Cons: a detector that flags everything would pass. The control is what makes the requirement say anything about the monitoring at all.
\item \textbf{Option D (chosen): keep the control, but define it as an i.i.d. sample of held-out listings, and cap the reference at 10,000 rows.}
\newline Pros: measured to work. The same detector gives a 1 to 10\,\% false-alarm rate on i.i.d. windows (2\,\% at a 10,000-row reference, against a nominal 5\,\%) while still flagging \texttt{ES} in 100\,\% of trials, so it separates the two cases the requirement is about. The reference cap keeps the false-alarm rate from growing with the training set, since power against irrelevant differences rises with reference size (1\,\% at 2,000 rows, 10\,\% at 78,330).
\newline Cons: a weaker claim than option A. It states that the detector is quiet on in-distribution traffic, not that it is quiet on unseen dealers, and the requirement now has to say which sample counts as the control, which is one more thing to get right in the M6 harness.
\end{itemize}
\par
\emph{Why:} NFR-11 asserted a property of the data (``at least one feature other than \texttt{country\_\allowbreak{}code} drifts'', ``the test-set control does not flag'') that nobody had checked. Measured on the scoped, deduplicated dataset (105,431 listings; 5,981 \texttt{ES} rows that the API would accept under FR-04; reference 76,142), the first half holds with a large margin: 200 of 200 \texttt{ES} windows are flagged at every window size from 100 to 1,000, and still 200 of 200 when \texttt{country\_\allowbreak{}code} is excluded. Spain is separable for substantive reasons, above all listing completeness: \texttt{nr\_\allowbreak{}prev\_\allowbreak{}owners} is missing in 95.5\,\% of \texttt{ES} rows against 40.0\,\% in the reference, \texttt{nr\_\allowbreak{}seats} in 10.9\,\% against 3.0\,\%, and the mix differs in \texttt{body\_\allowbreak{}type} (station wagon 6.6\,\% against 14.8\,\%), \texttt{transmission} (semi-automatic 0.1\,\% against 4.0\,\%) and \texttt{fuel\_\allowbreak{}category} (diesel 42.1\,\% against 31.9\,\%). The second half was simply false, and the i.i.d. cross-check shows the detector is not at fault: the same code gives a near-nominal false-alarm rate as soon as the control window is drawn i.i.d. The requirement is therefore aligned with what can be demonstrated, and the finding behind it, that dealer-level shift is as large as country-level shift on this data, is recorded in the project brief because it also justifies the seller-grouped split and is a risk for SC-04.},
  ai = {Information seeking, Alternative generation, Alternative assessment, Recommendation},
  response = {Accepted},
  assessment = {While reviewing PR \#19, AI flagged NFR-11 as the one requirement that makes an unverified claim about the data rather than about the system. Asked for an in-depth feasibility check, AI downloaded the raw dataset from the Zenodo DOI, reproduced the scope, deduplication and split protocol, emulated alibi-detect's \texttt{TabularDrift} with scipy after confirming its tests, defaults and Bonferroni rule against the v0.13 source (including that it does no NaN handling, so missing indicators had to be added), and ran 200 trials per window size. It reported that the first half of the requirement passes with margin and the second half fails, ran two controlled cross-checks to separate a detector fault from a data property (i.i.d. split, reference size), and tested and then rejected its own effect-size remedy on the measurements. Lukas reviewed the evidence and chose option D.},
  aievidence = {Claude Code session on 2026-09-23 while reviewing PR \#19: Lukas asked ``NFR-11 please make an in depth analysis if this is feasible [...] we need to be really sure about this''; AI ran the measurement described above, reported the pass, the failure and the limits of the study, proposed options A to D and recommended D; Lukas replied ``yes please''.},
  evidence = {\href{https://github.com/mlops-2627q1-mds-upc/MLOPS_Recommenditos/blob/main/docs/docs/requirements.md}{requirements} NFR-11; \href{https://github.com/mlops-2627q1-mds-upc/MLOPS_Recommenditos/blob/main/docs/docs/specification.md}{specification} NFR-11; \href{https://github.com/mlops-2627q1-mds-upc/MLOPS_Recommenditos/blob/main/reports/analysis}{reports/analysis/} (\texttt{nfr11\_\allowbreak{}check.\allowbreak{}py}, \texttt{nfr11\_\allowbreak{}diag.\allowbreak{}py}, \texttt{nfr11\_\allowbreak{}results.\allowbreak{}json}, run of 2026-09-23); \href{https://github.com/mlops-2627q1-mds-upc/MLOPS_Recommenditos/blob/main/docs/docs/project-brief.md}{project brief} section 3.2; \hyperref[EDN-03]{EDN-03}; \href{https://github.com/mlops-2627q1-mds-upc/MLOPS_Recommenditos/pull/19}{PR \#19}.},
}

% Generated by tools/edn_latex.py from reports/edn.md.
% Do not edit: change reports/edn.md and run `make edn`.
\ednentry{
  id = {EDN-15},
  title = {UC1 required fields after measuring the fill rates, plus SC-06 for absent optional fields},
  note = {\textbf{Evidence corrected on 2026-09-29.} The first run of \texttt{fillrates.\allowbreak{}py} printed fill rates with one decimal, so \texttt{seller\_\allowbreak{}type}'s 99.986\,\% appeared as 100.0.
The script now prints three decimals and \texttt{fillrates\_\allowbreak{}results.\allowbreak{}txt} was regenerated. The decision is unchanged: 14 missing rows in 97,889 leave it effectively never absent.
The training-scope count moved from 97,913 to 97,889 when the script was brought in line with EDN-22, which drops the listings registered after the reference date.},
  date = {2026-09-29},
  milestone = {M1: Project Inception (requirements and success criteria; shapes M3: Quality Assurance)},
  topic = {Requirements Engineering, API Design, Performance Criteria},
  participants = {@lukas2510},
  decision = {\texttt{/\allowbreak{}predict} requires \texttt{body\_\allowbreak{}type} and \texttt{seller\_\allowbreak{}type} in addition to the core FR-01 asked for before (\texttt{make}, \texttt{model}, \texttt{registration\_\allowbreak{}date}, \texttt{mileage\_\allowbreak{}km\_\allowbreak{}raw}, \texttt{power\_\allowbreak{}kw}, \texttt{fuel\_\allowbreak{}category}, \texttt{transmission}, \texttt{country\_\allowbreak{}code}). The remaining basic-set fields stay optional. FR-01 no longer prescribes how the model handles an absent optional field; it states the observable outcome instead, and the new SC-06 of the problem specification bounds it: with each optional field masked on its own, and with all of them masked at once, the point model's MdAPE stays at or below 1.5 times its full-input MdAPE. This amends EDN-06, which set SC-01 to SC-05; because NFR-01 gates deployment on all of them, the gate now also covers partial UC1 inputs. Whether the rarely-missing optional fields need the random masking planned for the UC2 interval models is left to M2/M3, when the pipeline exists and the effect can be measured.},
  rationale = {\emph{Alternatives considered:}
\begin{itemize}[leftmargin=*, topsep=2pt]
\item \textbf{Option A: keep all eight basic-set fields optional and rely on LightGBM's native missing-value handling.}
\newline Pros: the smallest required set, so UC1 stays easy to call and clearly different from UC2; no training-time work at all.
\newline Cons: measured to be unsupported for two fields. \texttt{body\_\allowbreak{}type} is filled in 100.000\,\% of the training listings and \texttt{seller\_\allowbreak{}type} in 99.986\,\% (14 of 97,889 rows), so the model effectively never sees them absent and a request that omits them is sent to the default side of every split on that feature with no training example backing that direction. The contract would promise behaviour the training data cannot show.
\item \textbf{Option B: keep all eight optional and train the point model with the same random masking as the UC2 interval models.}
\newline Pros: friendliest contract; one masking scheme for both model families, which also keeps the point estimate and the interval bounds from drifting apart on partial inputs (FR-07 widens the bounds when they do).
\newline Cons: makes an API requirement depend on a pipeline that does not exist yet (no \texttt{dvc.\allowbreak{}yaml}, issue \#10 still open), so PR \#19 would stay blocked on a modelling decision that cannot be measured before the first report. Masking also costs some accuracy on full inputs, which are the common case for UC1.
\item \textbf{Option C: require all basic features.}
\newline Pros: best accuracy, no missing-value path at serving time at all.
\newline Cons: hostile to users. \texttt{nr\_\allowbreak{}prev\_\allowbreak{}owners} is filled in 61.0\,\% of the listings, \texttt{gears} in 62.5\,\% and \texttt{drive\_\allowbreak{}train} in 76.4\,\%, so even the sellers behind the listings often do not state them.
\item \textbf{Option D (chosen): require the two always-filled fields, keep the rest optional, and bound the outcome with SC-06.}
\newline Pros: closes the measured gap at zero cost to the user, because whoever owns the car knows its body type and, for UC1, is the seller. Separates the API contract (decidable now) from the training mechanism (decidable in M2/M3), so the requirement can be merged without guessing. SC-06 turns FR-01's previously unquantified ``per-field masked-input accuracy check'' into a threshold and puts it in the deployment gate.
\newline Cons: the required set grows from eight fields to ten. SC-06 amends EDN-06, which was already agreed, and it is relative rather than absolute, so it cannot rule out a model that is uniformly mediocre; SC-01 to SC-04 do that.
\item \textbf{Option E: adopt D, but report the masked-input error only, without gating on it.}
\newline Pros: no amendment to EDN-06; the number is still visible in MLflow and in the report.
\newline Cons: leaves NFR-01's gate blind to exactly the case FR-01 allows the client to trigger, which is how the gap arose in the first place.
\end{itemize}
\par
\emph{Why:} FR-01 carried an \texttt{[open]} marker because it mixed two questions: which fields the client must send, which is an API contract, and how the model is trained to honour the optional ones, which is a modelling mechanism. The fill rates that the second question depended on had never been measured. Measured on the scoped, deduplicated training data (97,889 listings, \texttt{ES} held out, 11 supported makes), the risk turns out to sit in exactly two fields: \texttt{body\_\allowbreak{}type} and \texttt{seller\_\allowbreak{}type} at 100\,\%, followed by \texttt{nr\_\allowbreak{}doors} (98.8\,\%), \texttt{nr\_\allowbreak{}seats} (97.0\,\%) and \texttt{cylinders\_\allowbreak{}volume\_\allowbreak{}cc} (91.1\,\%), while \texttt{nr\_\allowbreak{}prev\_\allowbreak{}owners} (61.0\,\%), \texttt{gears} (62.5\,\%) and \texttt{drive\_\allowbreak{}train} (76.4\,\%) have plenty of natural missingness. Only 30.7\,\% of training rows have every one of the six remaining optional fields present and 469 rows (0.5\,\%) have none of them, so the model sees partial rows including the extreme case, just never a row without \texttt{body\_\allowbreak{}type} or \texttt{seller\_\allowbreak{}type}. Requiring the two is therefore the cheapest correct fix, and it costs the user nothing. The rest is left to the modelling plan on purpose: the mechanism cannot be chosen on evidence before the pipeline runs, and blocking the requirements page on it would have held up the M1 deliverable for a decision that belongs to M2/M3. SC-06 is what makes that safe, and it is relative to the model's own full-input MdAPE, like SC-03 is relative to the baseline, because no reference value exists yet and an absolute threshold would have been invented rather than derived.},
  ai = {Information seeking, Alternative generation, Alternative assessment, Recommendation},
  response = {Accepted with modifications},
  assessment = {Asked which decisions had to be settled before PR \#19 could be merged, AI named FR-01's open item first and called it a merge blocker. Asked to re-examine that, it downloaded the raw dataset from the Zenodo DOI, reproduced the scope, deduplication and make filter (which matched the brief's 11 makes as a cross-check), measured the per-field fill rates, and then corrected two of its own claims: the item does not block code that exists, because there is no pipeline yet, and FR-01 already promised a masked-input check, so the real gap was that the check had no threshold and no link to NFR-01's gate. It also reframed the requirement as mixing contract and mechanism, which is what made a clean split possible. Lukas accepted the recommendation and the field list; the modification is that the mechanism question stays open on purpose instead of being decided together with the contract.},
  aievidence = {Claude Code session on 2026-09-29: Lukas asked (translated from German) ``which things urgently still need deciding in requirements.md before we can merge, and why?'', then ``analyse that again, then explain it to me and give me a reasoned recommendation on what the best way is''; AI measured the fill rates, revised its own analysis, proposed options A to E and recommended D with SC-06 in the gate; Lukas replied ``I like your recommendation, please change it that way in the PR''.},
  evidence = {\href{https://github.com/mlops-2627q1-mds-upc/MLOPS_Recommenditos/blob/main/docs/docs/specification.md}{Specification} FR-01; \href{https://github.com/mlops-2627q1-mds-upc/MLOPS_Recommenditos/blob/main/docs/docs/problem-spec.md}{problem specification} SC-06; \href{https://github.com/mlops-2627q1-mds-upc/MLOPS_Recommenditos/blob/main/reports/analysis}{reports/analysis/} (\texttt{fillrates.\allowbreak{}py}, run of 2026-09-29); \href{https://github.com/mlops-2627q1-mds-upc/MLOPS_Recommenditos/blob/main/docs/docs/project-brief.md}{project brief} section 4; \hyperref[EDN-06]{EDN-06}; \href{https://github.com/mlops-2627q1-mds-upc/MLOPS_Recommenditos/pull/19}{PR \#19}.},
}

% Generated by tools/edn_latex.py from reports/edn.md.
% Do not edit: change reports/edn.md and run `make edn`.
\ednentry{
  id = {EDN-16},
  title = {Separate \texorpdfstring{\texttt{/\allowbreak{}predict}}{/predict} and \texorpdfstring{\texttt{/\allowbreak{}price-\allowbreak{}range}}{/price-range} instead of one valuation endpoint},
  date = {2026-09-29},
  milestone = {M1: Project Inception (API contract; implemented in M4: Model Deployment)},
  topic = {API Design},
  participants = {@lukas2510},
  decision = {UC1 and UC2 keep their own endpoints, \texttt{POST /\allowbreak{}predict} and \texttt{POST /\allowbreak{}price-\allowbreak{}range} (FR-06, FR-07). They are no longer marked \texttt{[proposed]}.},
  rationale = {\emph{Alternatives considered:}
\begin{itemize}[leftmargin=*, topsep=2pt]
\item \textbf{Option A (chosen): one endpoint per use case.}
\newline Pros: each endpoint has one input contract, which matters because \texttt{/\allowbreak{}predict} requires ten fields (FR-01, EDN-15) and \texttt{/\allowbreak{}price-\allowbreak{}range} requires only \texttt{make} (FR-02); FR-16's per-endpoint OpenAPI examples and FR-01's per-field 422 stay straightforward; the prediction log (FR-13) keeps the two traffic shapes apart, so FR-14's drift job does not read a change in the UC1/UC2 mix as input drift on exactly the completeness features the \texttt{ES} scenario relies on (EDN-14); consistent with \texttt{/\allowbreak{}comparables} (FR-09), which is already its own endpoint.
\newline Cons: two request schemas to keep in step, and a client that wants the estimate and the interval calls twice.
\item \textbf{Option B: one \texttt{/\allowbreak{}valuation} endpoint returning the estimate and the interval.}
\newline Pros: one path, one call for a client that wants both.
\newline Cons: the required fields would depend on the use case inside one schema, so the two contracts would be rebuilt as a conditional rule instead of disappearing; mixes both traffic shapes in one log; leaving \texttt{/\allowbreak{}comparables} separate would be inconsistent, and merging it too would put SHAP, intervals and the comparables lookup in every request against NFR-02's latency budget.
\end{itemize}
\par
\emph{Why:} The two endpoints differ in their input contract rather than in their output, and EDN-15 widened that gap from eight required fields to ten against one. A single endpoint would therefore not remove the two contracts, only hide them in a conditional schema. The monitoring side settles it: FR-14 and NFR-11 compare the logged input distribution of a window against the training reference, and UC2 allows every optional field to be absent, so a shift in the client mix would look like drift on the same completeness features that make the \texttt{ES} replay detectable. The decision is also the reversible one: a combined endpoint can be added later as a convenience, while splitting one apart under \texttt{/\allowbreak{}v1} would be a breaking change.},
  ai = {Alternative assessment, Recommendation},
  response = {Accepted},
  assessment = {Asked to re-examine the split before merging PR \#19, AI first called it a merge blocker and then withdrew that: no test or report cites FR-06 or FR-07 yet, so the practical cost of deciding later starts with the M4 API tests, not with the merge. It found that the two endpoints do not differ in their output, since FR-07 already returns the point estimate, and contributed the monitoring argument above, which had not been written down anywhere. Lukas kept the scope deliberately small and dropped AI's further suggestions (a shared base schema stated in the requirement, a combined convenience endpoint) as implementation detail or as unnecessary for a course project.},
  aievidence = {Claude Code session on 2026-09-29: Lukas asked (translated from German) ``analyse that again, then explain it to me and give me a reasoned recommendation on what the best way is'' about the endpoint split, and after the recommendation replied ``I do not want overkill, it should stay simple, it is only a university course''.},
  evidence = {\href{https://github.com/mlops-2627q1-mds-upc/MLOPS_Recommenditos/blob/main/docs/docs/specification.md}{Specification} FR-06, FR-07, FR-09; \href{https://github.com/mlops-2627q1-mds-upc/MLOPS_Recommenditos/pull/19}{PR \#19}.},
}

% Generated by tools/edn_latex.py from reports/edn.md.
% Do not edit: change reports/edn.md and run `make edn`.
\ednentry{
  id = {EDN-17},
  title = {Deployment target is the FIB Virtech VM},
  date = {2026-09-29},
  milestone = {M4: Model Deployment},
  topic = {ML System Design, Deployment},
  participants = {@lukas2510},
  decision = {The component is deployed on the FIB Virtech VM the course provides (4 GB RAM, 20 GB disk, CPU-only, one semester). The project brief no longer lists the target as an open point, and the requirements say so once instead of assuming it in every target.},
  rationale = {\emph{Alternatives considered:}
\begin{itemize}[leftmargin=*, topsep=2pt]
\item \textbf{Option A (chosen): the FIB Virtech VM.}
\newline Pros: free and provided for the course, so no payment details and no personal cloud account for a semester project; 4 GB is enough with the budgets NFR-04 already sets (drift job scheduled in its own container, Prometheus capped by retention, logs capped by rotation, 4 GB disk kept free); every target in the requirements is written for this machine.
\newline Cons: tight on RAM, a single unattended host with no redundancy, and nobody on the team has access yet.
\item \textbf{Option B: another cloud (AWS, DigitalOcean, Oracle Always Free or similar).}
\newline Pros: more headroom, and a free ARM tier exists at one provider.
\newline Cons: either paid or tied to a private account with payment details; the grading rewards nothing for the hosting choice; NFR-02 to NFR-04 would have to be re-derived for different hardware, and an ARM tier would also mean multi-architecture images.
\end{itemize}
\par
\emph{Why:} The requirements already named this VM in eight places (NFR-02, NFR-04, NFR-05, NFR-09, NFR-12, NFR-13 and the network warning), so the decision had been made by writing rather than by deciding, while the brief still listed the target as open. The reason the brief gave for keeping it open, that the stack is tight on 4 GB, has since been planned for in NFR-04, so the open point was stale. What stays genuinely open is not the target but what the VM's network allows, which the requirements already track with a fallback per requirement. Access is a task, not a decision, and it has a date: if nobody has access at the M4a lab, we raise it with the teachers.},
  ai = {Information seeking, Alternative assessment, Recommendation},
  response = {Accepted},
  assessment = {While reviewing PR \#19, AI reported the brief and the requirements as contradicting each other on the deployment target. Asked to look again, it corrected its own framing: the \texttt{[proposed]} marker in the requirements refers to the performance numbers, not to the machine, and the requirements treat the VM as settled throughout, so the real issue was an implicit decision rather than a disagreement between two pages. It also pointed out that the brief's stated reason for leaving the point open is already covered by NFR-04. Lukas kept the change small and deliberately left the performance targets \texttt{[proposed]} until the first M4 load test.},
  aievidence = {Claude Code session on 2026-09-29: after Lukas asked (translated from German) ``analyse that again, then explain it to me and give me a reasoned recommendation on what the best way is'' about the deployment target, AI recommended fixing the VM as the target with a dated trigger for the missing access, and Lukas replied ``yes do that''.},
  evidence = {\href{https://github.com/mlops-2627q1-mds-upc/MLOPS_Recommenditos/blob/main/docs/docs/specification.md}{Specification} section 2 and NFR-04; \href{https://github.com/mlops-2627q1-mds-upc/MLOPS_Recommenditos/blob/main/docs/docs/project-brief.md}{project brief} section 5; \href{https://github.com/mlops-2627q1-mds-upc/MLOPS_Recommenditos/pull/19}{PR \#19}.},
}

% Generated by tools/edn_latex.py from reports/edn.md.
% Do not edit: change reports/edn.md and run `make edn`.
\ednentry{
  id = {EDN-18},
  title = {Unseen countries and models are accepted with a warning, not rejected},
  date = {2026-09-29},
  milestone = {M4: Model Deployment (serving contract; enables M6: Monitoring)},
  topic = {API Design},
  participants = {@lukas2510},
  decision = {FR-05 loses its \texttt{[proposed]} marker. A \texttt{country\_\allowbreak{}code} among the 8 countries of the data or a \texttt{model} that is missing from the training data is accepted, passed to the model as unknown, and reported in the response's \texttt{warnings} field.},
  rationale = {\emph{Alternatives considered:}
\begin{itemize}[leftmargin=*, topsep=2pt]
\item \textbf{Option A (chosen): accept and warn.}
\newline Pros: the \texttt{ES} replay reaches the model at all, which is what EDN-03's new-market drift scenario and all of M6 depend on; LightGBM handles an unseen category natively, so nothing has to be built; the caller still learns that the answer is extrapolated, and the warning is in the prediction log (FR-13).
\newline Cons: the API answers for inputs the model has never seen, and the estimate can be worse than the metrics suggest without the caller noticing the warning.
\item \textbf{Option B: reject with HTTP 422.}
\newline Pros: stricter, and never answers outside what was trained.
\newline Cons: would reject every \texttt{ES} listing, so the drift scenario could not be replayed and the retraining loop (FR-15) could not be demonstrated. The scope check we do want is about the make (FR-04), which stays a 422.
\end{itemize}
\par
\emph{Why:} The decision was effectively made by EDN-03, which chose Spain as the held-out new market precisely so that unseen traffic can be replayed against the running API in M6. Rejecting it would have removed the scenario the monitoring milestone is built on. Out-of-scope makes are a different case and keep their 422 (FR-04), so ``accept'' is not a blanket rule.},
  ai = {Alternative assessment, Recommendation},
  response = {Accepted},
  assessment = {AI pointed out that this item was already prejudiced by EDN-03 and cost nothing to confirm, which is why it stayed unmarked longer than necessary. No further analysis was needed.},
  aievidence = {Claude Code session on 2026-09-29: AI noted that FR-05 was (translated from German) ``in fact already prejudged by EDN-03'' and recommended lifting it to \texttt{[decided]}; Lukas replied ``do that''.},
  evidence = {\href{https://github.com/mlops-2627q1-mds-upc/MLOPS_Recommenditos/blob/main/docs/docs/requirements.md}{requirements} FR-05 and FR-04; \href{https://github.com/mlops-2627q1-mds-upc/MLOPS_Recommenditos/blob/main/docs/docs/specification.md}{specification} FR-05; \hyperref[EDN-03]{EDN-03}; \href{https://github.com/mlops-2627q1-mds-upc/MLOPS_Recommenditos/pull/19}{PR \#19}.},
}

% Generated by tools/edn_latex.py from reports/edn.md.
% Do not edit: change reports/edn.md and run `make edn`.
\ednentry{
  id = {EDN-19},
  title = {Which DagsHub repository the team uses as DVC remote and MLflow server},
  date = {2026-09-29},
  milestone = {M2: Reproducibility},
  topic = {Data Versioning, Experiment Tracking},
  participants = {@lukas2510 (Scrum Master), on behalf of the team. Pending confirmation by the full team at the next sprint review.},
  decision = {A DagsHub \textbf{organisation} owned by the team hosts the connected repository, and every team member joins it with push rights (option A).
The repository is \href{https://dagshub.com/recommenditos/MLOPS_Recommenditos}{\texttt{https:/\allowbreak{}/\allowbreak{}dagshub.\allowbreak{}com/\allowbreak{}recommenditos/\allowbreak{}MLOPS\_\allowbreak{}Recommenditos}}, owned by the \texttt{recommenditos} DagsHub organisation and connected through DagsHub's GitHub integration.
A first decision for \texttt{pauadal03/\allowbreak{}MLOPS\_\allowbreak{}Recommenditos} was taken and withdrawn on the same day, once it turned out that pauadal03 is not a member of this team: they hold read-only access to our GitHub repository, have never contributed to it, and are not on the roster.
Their DagsHub repository is a mirror of our public repository made from outside the team, so it is not a candidate.
Later the same day \texttt{mark.\allowbreak{}welf.\allowbreak{}atzberger/\allowbreak{}MLOPS\_\allowbreak{}Recommenditos} was deleted, so that candidate no longer exists either and the raw object it held is gone with it.
The practical choice is therefore between a DagsHub organisation owned by the team and a fresh repository under one team member's account.
Whether the repository stays public was decided separately, see EDN-20.},
  rationale = {\emph{Alternatives considered:}
\begin{itemize}[leftmargin=*, topsep=2pt]
\item \textbf{Option A (chosen): a DagsHub organisation owns the connected repository, every team member is a member with push rights.}
\newline Pros: nothing the project depends on hangs off one person's private account; push rights follow membership instead of manual collaborator entries, and the team is five people, so inviting everyone is trivial; visibility is set in one place; data and experiment history survive if a member leaves or cleans up their account.
\newline Cons: one-time setup (create the organisation, connect the GitHub repository, invite four people); under a \texttt{dvc add} mechanism the 548.6 MB raw object has to be pushed once more, though under \texttt{import-\allowbreak{}url} there is nothing to move.
\item \textbf{Option B: use \texttt{mark.\allowbreak{}welf.\allowbreak{}atzberger/\allowbreak{}MLOPS\_\allowbreak{}Recommenditos}, reconnected properly.} No longer available: the repository was deleted on 2026-09-29.
Pros (while it existed): it belonged to a teammate and already held the raw object; no organisation to create.
\newline Cons: its mirror was demonstrably not in sync (see the evidence below), and the likely fix was to delete and re-import it through the GitHub integration, which would have lost the pushed object anyway; the project would still have hung off one teammate's private account, with a manual collaborator entry needed per member.
\item \textbf{Option C: a fresh repository under one team member's account, connected through the GitHub integration.}
\newline Pros: one form less than an organisation; whoever creates it is admin and can add the other four without waiting on anyone.
\newline Cons: identical to option B's structural drawbacks -- the project hangs off one private account and rights are granted per person by hand. The saving over option A is a single setup form.
\item \textbf{Option D: use \texttt{pauadal03/\allowbreak{}MLOPS\_\allowbreak{}Recommenditos}.} Chosen on 2026-09-29 and withdrawn the same day.
\newline Pros: it is the only mirror verifiably connected through the GitHub integration, and adopting it would have cost nothing.
\newline Cons: ruled out on ownership, not on measurements -- pauadal03 is not on the team, has \texttt{read} permission on our GitHub repository and has never contributed to it, so the project's data and experiment history would sit on the account of someone who cannot even push to the repository they mirrored.
For the record, no project data was ever on that mirror: the raw object returns 404 there, so it holds our public git history and nothing else.
\item \textbf{Option E: keep several mirrors and let everyone push to their own.}
\newline Pros: no coordination needed.
\newline Cons: not viable -- the remote URL is committed in \texttt{.\allowbreak{}dvc/\allowbreak{}config}, so the team would overwrite each other's setting, and data and experiment history would be split across accounts.
\end{itemize}
\par
\emph{Why:} By the time the decision was taken the alternatives had collapsed: option D was disqualified on ownership and option B ceased to exist when that repository was deleted, so every remaining candidate required a fresh connect through the GitHub integration anyway.
That removed the only real argument against option A, which had been that a correctly connected repository already existed elsewhere at zero cost.
What was left is a one-form difference -- an organisation versus a repository under one member's account, for the same connect and the same invitations -- against the fact that fifteen cohort members outside the team hold read access to the GitHub repository and can mirror it at any time, as one of them already had.
The team's standing rule was to take the better-engineered option unless it is overkill for a one-semester course project; with four teammates to invite rather than the whole twenty-one-person collaborator list an earlier miscount had assumed, option A is not overkill.
Evidence gathered on 2026-09-29 against the DagsHub and GitHub APIs.
Ownership: write access to the GitHub repository is held by @lukas2510, @kadameit, @ulasawczuk, @W11W11W11 and @michudud04, plus @martinezmatias and @santidrj, the two lecturers who own the course organisation; @pauadal03 has \texttt{read}, zero pull requests and zero commits, and does not appear on the roster in \texttt{docs/\allowbreak{}docs/\allowbreak{}scrum/\allowbreak{}index.\allowbreak{}md}.
Sync: \texttt{pauadal03/\allowbreak{}MLOPS\_\allowbreak{}Recommenditos} is at \texttt{2e902e2143}, identical to GitHub \texttt{main}, carries all 7 branches and mirrors all 7 open issues, while \texttt{mark.\allowbreak{}welf.\allowbreak{}atzberger/\allowbreak{}MLOPS\_\allowbreak{}Recommenditos} is at \texttt{8b1d64fc}, is missing the \texttt{docs/\allowbreak{}data-\allowbreak{}facts-\allowbreak{}corrections} and \texttt{model-\allowbreak{}card} branches, still carries \texttt{feature/\allowbreak{}ruff-\allowbreak{}pl-\allowbreak{}and-\allowbreak{}coverage} which GitHub deleted after the merge, and mirrors 0 issues.
DagsHub only mirrors issues and pull requests for repositories connected through the GitHub integration, so the issue count is the clearest signal that the two were connected in different ways and that mark's is a plain git mirror.
Visibility: both report \texttt{private: false}; EDN-07 as first written on the issue \#3 branch assumed the remote was private, a premise found false and recorded in EDN-20.
The reason that original check misled us: anonymous access is not evidence either way, because DagsHub refuses every anonymous request, including for a known-public control repository (\texttt{DAGsHub-\allowbreak{}Official/\allowbreak{}dagshub-\allowbreak{}docs}, also \texttt{private: false}) and for a repository that does not exist.
Access: \texttt{lukas2510} was granted \texttt{push} on mark's repository on 2026-09-29, shortly before it was deleted, and has \texttt{pull} only on pauadal03's; the rest of the team had push on neither.
State at the end of 2026-09-29: \texttt{mark.\allowbreak{}welf.\allowbreak{}atzberger/\allowbreak{}MLOPS\_\allowbreak{}Recommenditos} returns 404 and so does the raw object it held; \texttt{pauadal03/\allowbreak{}MLOPS\_\allowbreak{}Recommenditos} still exists and still holds no data; a repository search returns that mirror as the only one left.
Structural point behind all of this: fifteen cohort members outside the team hold \texttt{read} on our GitHub repository, so anyone of them can create a mirror at any time without doing anything wrong. The rule the team needs is therefore not ``check who owns a mirror'' but ``the project's infrastructure lives somewhere the team owns'', which is an argument for option A independent of the measurements.
Follow-up: \href{https://github.com/mlops-2627q1-mds-upc/MLOPS_Recommenditos/pull/26}{PR \#26} now commits the organisation's URL in \texttt{.\allowbreak{}dvc/\allowbreak{}config}, and the owner still has to grant push rights to everyone.
Nothing has to be migrated: the 548.6 MB raw object that had been pushed to mark's repository is gone with that repository, so under the \texttt{dvc add} mechanism settled in EDN-25 (\href{https://github.com/mlops-2627q1-mds-upc/MLOPS_Recommenditos/pull/27}{PR \#27}) it is re-downloaded from Zenodo and pushed once to the organisation's remote.
\par
Amended on 2026-10-01, confirmed by the repository owner: all five team members now hold \textbf{write} access to the DagsHub repository, and it is held as a per-repository collaborator entry for each of them rather than as membership of the \texttt{recommenditos} organisation.
The follow-up above is therefore done, by a route this entry did not choose.
The decision is not revisited, because for a single repository the two forms are functionally equivalent: everyone can push, and the repository still belongs to the organisation rather than to a private account, which is what option A was for.
What the deviation does cost is the one property the entry gave as the reason to prefer membership -- rights would follow the person instead of being granted per repository -- so a second repository under the organisation needs its own five grants, and a member who leaves has to be removed once per repository rather than once.
Recorded rather than left to be discovered, so that whoever adds the next repository does not assume the invitations are already in place.},
  ai = {Information seeking, Alternative generation, Alternative assessment, Recommendation},
  response = {Accepted with modifications},
  assessment = {AI found the two mirrors, measured them against GitHub and recommended \texttt{pauadal03}'s on those measurements. The measurements were correct and the recommendation was wrong, because AI never checked whether that repository belonged to anyone on the team -- it ranked the candidates it had found without asking where they came from. Lukas caught it by asking who pauadal03 is, which invalidated the recommendation and the decision taken on it. Recorded here rather than silently corrected, because the failure is the instructive part: a candidate that scores best on every technical measurement can still be disqualified by a question nobody asked. AI's useful contributions were the measurements, the finding that mark's mirror is not properly connected, and the control-repository method behind EDN-20; its first recommendation was not one of them. Its second recommendation, option A, was the one the team accepted, and only after Lukas had rejected the first and asked who the account behind it belonged to.
This is the second decision in the same week reached on a premise nobody had checked, after ``the DagsHub repository is private'' in the first version of EDN-07 on the issue \#3 branch, and in both cases the check that settled it took one API call and was available before the decision rather than after it. The pattern, not either individual fix, is the thing worth carrying into the working agreements: state the premise a decision rests on, and verify it, before recording the decision.},
  aievidence = {Claude Code session on 2026-09-29, prompts translated from German: the comparison and the recommendation for pauadal03 followed ``isn't this our DagsHub repo, the one that matches our GitHub repo?'' and ``why did that repo score worst?''; the retraction followed ``I think pauadal03 is one of our lecturers -- is it bad that everything now has to go through them?''.},
  evidence = {\href{https://github.com/mlops-2627q1-mds-upc/MLOPS_Recommenditos/pull/26}{PR \#26}, \href{https://github.com/mlops-2627q1-mds-upc/MLOPS_Recommenditos/issues/10}{issue \#10}, EDN-20, EDN-25, roster in \texttt{docs/\allowbreak{}docs/\allowbreak{}scrum/\allowbreak{}index.\allowbreak{}md}.},
}

% Generated by tools/edn_latex.py from reports/edn.md.
% Do not edit: change reports/edn.md and run `make edn`.
\ednentry{
  id = {EDN-20},
  title = {The DagsHub remote stays public},
  date = {2026-09-29},
  milestone = {M2: Reproducibility},
  topic = {Data Versioning, Privacy},
  participants = {@lukas2510 (Scrum Master). Pending confirmation by the full team at the next sprint review.},
  decision = {The DagsHub repository that serves as our DVC remote and MLflow tracking server stays public.
We accept that the raw file's PII columns are readable by anyone with a DagsHub account, which follows from EDN-25 tracking the raw file with \texttt{dvc add} and pushing it to that remote.
The additional exposure this creates is small, because the identical file is already publicly downloadable from the pinned Zenodo DOI under the MIT license.
It is not nothing: hosting our own copy makes the team a publisher of that personal data in its own right and under its own name, which the earlier publication mitigates but does not undo.},
  rationale = {\emph{Alternatives considered:}
\begin{itemize}[leftmargin=*, topsep=2pt]
\item \textbf{Option A (chosen): keep the repository public and say so.}
\newline Pros: graders and supervisors can inspect code, data and experiments without being added as collaborators, which is the reason the course puts the project on DagsHub in the first place; it matches how the source dataset is already published.
\newline Cons: the raw file, which still carries \texttt{vin}, \texttt{street}, \texttt{seller\_\allowbreak{}company\_\allowbreak{}name} and coordinates, can be pulled by every logged-in DagsHub user; re-hosting it is a publishing act of our own, so ``the author already published it'' reduces the marginal risk but does not transfer the responsibility; the decision has to be revisited the moment we host data that is not already public elsewhere.
\item \textbf{Option B: make the repository private.}
\newline Pros: the rationale first given for pushing the raw file would hold as written; the PII sits behind an access wall.
\newline Cons: every grader, supervisor and teammate then needs a manual collaborator entry; it buys little real protection, because the identical file stays publicly downloadable from Zenodo either way.
\end{itemize}
\par
\emph{Why:} The question only came up because the premise of the raw-data decision turned out to be false.
The first version of EDN-07, on the issue \#3 branch, justified pushing the raw file with ``the DagsHub repo is private (anonymous access is redirected to the login page and an anonymous data download returns 401)''.
Checked on 2026-09-29: DagsHub refuses anonymous access to everything, public repositories included.
A known-public control repository (\texttt{DAGsHub-\allowbreak{}Official/\allowbreak{}dagshub-\allowbreak{}docs}, \texttt{private: false}) answers anonymous requests exactly like ours, and so does a repository that does not exist, so the observation has no discriminating power.
The authenticated API reports \texttt{private: false} for both mirrors that existed at the time.
Given a real choice between hiding a file that is already public and saying openly that it is public, Lukas chose the second: the protection gained would be nominal, while the access cost for graders and supervisors would be real.
Under EDN-19 the organisation is the team's own, so the visibility is ours to change and this decision can be revisited at any time without asking anyone outside the team.
Recorded deliberately as an accepted risk rather than a solved problem, because the two are not the same thing: the marginal exposure is small, but we are still the ones publishing personal data, and an entry that claimed the concern was spent would not survive a reviewer who cares about privacy.
State after the decision: \texttt{mark.\allowbreak{}welf.\allowbreak{}atzberger/\allowbreak{}MLOPS\_\allowbreak{}Recommenditos}, the repository that held the raw copy, was deleted on 2026-09-29, and the object returns 404. The team therefore publishes no personal data on any DagsHub remote at present, and this decision governs whatever remote EDN-19 settles on rather than an exposure that exists today.},
  ai = {Information seeking, Alternative generation, Alternative assessment},
  response = {Used as input for further analysis},
  assessment = {AI found that the privacy premise of the raw-data decision was false, using a control repository to show that anonymous refusal does not distinguish public from private, and put both ways out to Lukas without recommending either, since the trade-off is about how open the team wants to be rather than a technical question. Lukas decided to keep the repository public. The finding is what changed the outcome here; the decision itself was not AI's to make.},
  aievidence = {Claude Code session on 2026-09-29, prompts translated from German: after the finding was presented with the two options (``either mark switches it to private ... or you keep it public and write that honestly into the raw-data entry''), Lukas answered ``we keep it public''.},
  evidence = {EDN-19, EDN-25, \href{https://github.com/mlops-2627q1-mds-upc/MLOPS_Recommenditos/pull/26}{PR \#26}, \href{https://github.com/mlops-2627q1-mds-upc/MLOPS_Recommenditos/pull/27}{PR \#27}.},
}

% Generated by tools/edn_latex.py from reports/edn.md.
% Do not edit: change reports/edn.md and run `make edn`.
\ednentry{
  id = {EDN-21},
  title = {Report the upstream licence contradiction instead of resolving it},
  date = {2026-09-29},
  milestone = {M1: Project Inception},
  topic = {Data Licensing and Provenance},
  participants = {@lukas2510},
  decision = {The dataset card reports both licence statements that the Zenodo record makes about the AutoScout24 dataset, and names the contradiction, rather than picking one or asking the author to resolve it. The YAML frontmatter carries \texttt{license: mit} because that is the record's structured field, and the body explains the discrepancy.},
  rationale = {\emph{Alternatives considered:}
\begin{itemize}[leftmargin=*, topsep=2pt]
\item \textbf{Option A (chosen): report both statements and flag the contradiction.}
\newline Pros: it is the only accurate description of what the source actually says; it costs nothing, because both readings permit non-commercial academic use; it warns anyone who later reuses our work for a commercial purpose, which is exactly where the difference would start to matter; it demonstrates the provenance care the course rewards.
\newline Cons: a reader has to absorb an ambiguity rather than a single answer; the frontmatter still has to commit to one identifier, so the card is slightly inconsistent with itself unless the body is read.
\item \textbf{Option B: report MIT only.}
\newline Pros: simplest; matches the machine-readable metadata field that tooling and aggregators would read; MIT is the more permissive of the two, so it cannot under-claim our own rights.
\newline Cons: silently discards a restriction the author wrote on the same record; would mislead a future reader who wanted to use the data commercially; asserts a certainty we do not have.
\item \textbf{Option C: report the restrictive prose only.}
\newline Pros: the conservative reading, so it cannot over-claim rights.
\newline Cons: contradicts the structured licence field; ``research, educational or analytical'' is not a recognised licence, so it is not expressible in the card's frontmatter; over-restricts what may genuinely be MIT.
\item \textbf{Option D: ask the author to resolve it before finishing the card.}
\newline Pros: would produce a definitive answer and a citable clarification.
\newline Cons: blocks a sprint deliverable on a stranger's response time with no deadline; disproportionate for a university project whose use is permitted under either reading.
\end{itemize}
\par
\emph{Why:} The contradiction is in the source, not in our reading of it: the Zenodo record's structured licence field says MIT while the author's description on the same record says ``You are welcome to use this dataset for research, educational, or analytical purposes''. Since our use is non-commercial academic work, both readings permit it, so resolving the ambiguity would change nothing about what we are allowed to do. Documenting it honestly is therefore strictly better than choosing a side, and materially better than blocking the deliverable on an email. Option D was explicitly declined by the team as disproportionate for the size of the project.},
  ai = {Information seeking, Alternative generation, Alternative assessment, Recommendation},
  response = {Accepted with modifications},
  assessment = {AI found the contradiction while re-validating PR \#22, and in doing so corrected its own earlier review, which had asserted that the card's restrictive wording was simply wrong and contradicted the MIT licence recorded in the project brief. Fetching the Zenodo API record showed that both statements come from the same source, so the card's author had not made a mistake and the project brief's flat ``License: MIT'' was itself incomplete. AI laid out the four options above and recommended option A. The team accepted the recommendation but rejected the accompanying suggestion to contact the author, on the grounds that it is disproportionate for a university project.},
  aievidence = {Claude Code session on 2026-09-29 while reviewing PR \#22. AI fetched \texttt{https:/\allowbreak{}/\allowbreak{}zenodo.\allowbreak{}org/\allowbreak{}api/\allowbreak{}records/\allowbreak{}17643343} and reported that \texttt{metadata.\allowbreak{}license.\allowbreak{}id} is \texttt{mit-\allowbreak{}license} while the description text welcomes use for research, education and analysis, which reads as narrower than MIT; it then flagged the choice as EDN-worthy rather than deciding it. Lukas replied, in German, that the team should not be pedantic and should simply take the recommendation, since this is a small university-course project, which settled option A and ruled out option D.},
  evidence = {\href{https://github.com/mlops-2627q1-mds-upc/MLOPS_Recommenditos/pull/22}{PR \#22}, the Licensing section of \href{https://github.com/mlops-2627q1-mds-upc/MLOPS_Recommenditos/blob/main/docs/docs/dataset-card.md}{the dataset card}, Zenodo record \href{https://doi.org/10.5281/zenodo.17643343}{10.5281/zenodo.17643343}.},
}

% Generated by tools/edn_latex.py from reports/edn.md.
% Do not edit: change reports/edn.md and run `make edn`.
\ednentry{
  id = {EDN-22},
  title = {Drop listings registered after the age reference date},
  note = {\textbf{Evidence corrected on 2026-09-29.} The decision said preprocessing drops ``the 164 listings''. 164 is the raw-file count;
preprocessing runs after scoping, so it removes 27 of them. The alternatives below already said 26, measured after deduplication.
The decision itself is unchanged.
\par
\textbf{Description corrected on 2026-09-30.} The decision described the 27 as what survives ``the EDN-04 scope and the training price range''.
The rule order implemented in issue \#34 puts this rule above the price range, so the 27 survive the scope alone.
The number is the same either way, because all 27 are inside the range; only the description of how it was obtained was wrong.
The decision itself is unchanged.},
  date = {2026-09-29},
  milestone = {M3: Quality Assurance},
  topic = {Data Validation and Preprocessing},
  participants = {@lukas2510},
  decision = {Preprocessing drops every listing whose \texttt{registration\_\allowbreak{}date} is after the age reference date (the 2025-11-08 snapshot date in training, the request date in serving). There are 164 such listings in the raw file, of which \textbf{27} survive the EDN-04 scope, the one rule the implemented order puts above this one, and so are the ones preprocessing actually removes; all 27 are inside the training price range, which runs after this rule, so the count is the same whichever of the two goes first; 26 remain if deduplication runs first, the difference being one duplicate row. The Great Expectations suite asserts \texttt{registration\_\allowbreak{}date <=\allowbreak{} reference date} on the processed data as a hard expectation, and the same rule on the raw data with \texttt{mostly=\allowbreak{}0.\allowbreak{}99}, so a future scrape that suddenly carries many such rows is flagged instead of silently cleaned away.},
  rationale = {\emph{Alternatives considered:}
\begin{itemize}[leftmargin=*, topsep=2pt]
\item \textbf{Option A (chosen): drop the rows.}
\newline Pros: 164 of 118,382 raw rows is 0.14\,\%, and only 26 of them survive the EDN-04 scope, so nothing measurable is lost; makes \texttt{age >=\allowbreak{} 0} a genuine invariant that the pipeline, the tests and the API contract can all rely on; no special case anywhere in the feature code.
\newline Cons: throws away the other, possibly correct attributes of those listings; the rule has to live in preprocessing, so the raw expectation has to tolerate what preprocessing removes.
\item \textbf{Option B: clamp the age at zero.}
\newline Pros: keeps the rows and their remaining attributes; a single \texttt{max(0, .\allowbreak{}.\allowbreak{}.\allowbreak{})} in the feature code.
\newline Cons: invents a registration date that the source does not support, and does so for rows whose date is demonstrably wrong (the latest is 2026-11-01, almost a year after the snapshot); a clamped zero is indistinguishable from a genuinely new car, which is exactly the group EDN-04 scopes out.
\item \textbf{Option C: keep the rows with a negative age.}
\newline Pros: no data loss, no cleaning rule; gradient boosting tolerates the values.
\newline Cons: the API would have to accept a negative age too, or training and serving would disagree; removes the cheapest available sanity check on the most important feature.
\item \textbf{Option D: repair the date from \texttt{production\_\allowbreak{}year}.}
\newline Pros: would keep the rows with a defensible value.
\newline Cons: \texttt{production\_\allowbreak{}year} is only 19\,\% filled, so it repairs almost none of them; a repair rule needs its own validation for a group this small.
\end{itemize}
\par
\emph{Why:} For the rows this decision actually governs, a registration date after the snapshot is a data-entry error rather than a rare but real phenomenon: 137 of the 164 sit in a single month, January 2026, which looks like a year typo. The reading is weaker for the 114 new cars among them, where a dealer entering a planned first-registration date is plausible, but those fall outside the used-car scope before preprocessing runs; only 26 of the 164 reach it. At that share the cost of dropping them is nil, while every alternative either fabricates a value or gives up the \texttt{age >=\allowbreak{} 0} invariant. Bounding the raw data with \texttt{mostly=\allowbreak{}0.\allowbreak{}99} rather than a hard rule keeps the raw suite honest about what the source actually contains while still catching a systematic worsening, which matters because the \texttt{ES} holdout and the M6 drift scenario re-run these checks on data we have not profiled.},
  ai = {Information seeking, Alternative generation, Alternative assessment, Recommendation},
  response = {Accepted},
  assessment = {AI profiled the 164 rows, laid out the four options above and recommended dropping them, together with the split between a hard expectation on the processed data and a \texttt{mostly} bound on the raw data. Lukas accepted the recommendation because the dropped share is negligible and the invariant is worth more than the rows.},
  aievidence = {Claude Code session on 2026-09-29 while analysing \href{https://github.com/mlops-2627q1-mds-upc/MLOPS_Recommenditos/issues/25}{issue \#25}: AI presented the options as a decision question, Lukas chose ``drop them''.},
  evidence = {\href{https://github.com/mlops-2627q1-mds-upc/MLOPS_Recommenditos/issues/25}{Issue \#25}, \href{https://github.com/mlops-2627q1-mds-upc/MLOPS_Recommenditos/pull/24}{PR \#24}, \href{https://github.com/mlops-2627q1-mds-upc/MLOPS_Recommenditos/blob/main/docs/docs/problem-spec.md}{problem specification} section 2.},
}

% Generated by tools/edn_latex.py from reports/edn.md.
% Do not edit: change reports/edn.md and run `make edn`.
\ednentry{
  id = {EDN-23},
  title = {Read the condition flags as one-sided assertions instead of encoding them three-valued},
  date = {2026-09-29},
  milestone = {M3: Quality Assurance},
  topic = {Feature Encoding and Data Validation},
  participants = {@lukas2510},
  decision = {\texttt{has\_\allowbreak{}full\_\allowbreak{}service\_\allowbreak{}history}, \texttt{non\_\allowbreak{}smoking} and \texttt{is\_\allowbreak{}rental} stay plain booleans and are read as ``the seller asserted this'' versus ``the seller did not assert this'', never as ``yes'' versus ``no''. No three-valued (true / false / unknown) encoding is built and no \texttt{False} is mapped to missing. \texttt{is\_\allowbreak{}used} is excluded entirely because \texttt{offer\_\allowbreak{}type} already carries the same information reliably. The Great Expectations suite only checks that the columns are boolean and non-null.},
  rationale = {\emph{Alternatives considered:}
\begin{itemize}[leftmargin=*, topsep=2pt]
\item \textbf{Option A (chosen): keep them binary, fix the interpretation.}
\newline Pros: no pipeline, API or serialisation complexity; loses nothing, because the ambiguity is in the source and no encoding can resolve it; the risk it addresses (somebody reading a \texttt{False} as a denial) is a documentation and naming problem, so it is solved where it actually lives, in the problem spec, the dataset card and the model card.
\newline Cons: the column name still reads like a two-sided statement, so the documentation has to carry the caveat; a reader who skips it can still misread the feature.
\item \textbf{Option B: encode three-valued, or map \texttt{False} to missing.}
\newline Pros: makes the ambiguity visible in the data itself rather than only in prose; the shape a careful reviewer expects.
\newline Cons: it is not actually three-valued, because we cannot tell a real \texttt{False} from an unknown, so the third state would be empty and the encoding would misrepresent the problem as solved; for a gradient-boosted tree a column holding only \texttt{\{True, NaN\}} splits exactly like \texttt{\{True, False\}}, so the model gains no information at all; adds an extra state to the API schema, the training pipeline and the expectations for that zero gain.
\item \textbf{Option C: drop all affected flags.}
\newline Pros: the easiest position to defend; no ambiguous column ever reaches the model.
\newline Cons: discards the half of the signal that is reliable, since a \texttt{True} is a real assertion; full service history and non-smoking are exactly the kind of attributes that move a used-car price.
\item \textbf{Option D: recover the true value from \texttt{description} or the equipment lists.}
\newline Pros: would turn the flags into genuine two-sided features.
\newline Cons: needs multilingual free-text parsing, which the problem spec puts out of scope; produces labels with no ground truth to validate against.
\end{itemize}
\par
\emph{Why:} The flags are one-sided by construction: a listing form offers a checkbox, so a tick is evidence and an empty box is the absence of evidence. 14,744 rows flagged as neither used, new nor pre-registered, and 18,108 scoped rows with \texttt{is\_\allowbreak{}used =\allowbreak{} False} while \texttt{offer\_\allowbreak{}type =\allowbreak{} U}, are the direct measurement of that. Since no information exists to separate ``no'' from ``not stated'', every encoding choice carries identical information, and the one that costs nothing wins. This deliberately overrides the original proposal in issue \#25 to encode the flags three-valued.},
  ai = {Information seeking, Alternative generation, Alternative assessment, Recommendation},
  response = {Accepted},
  assessment = {The three-valued encoding was AI's own earlier proposal, written into issue \#25. On re-reading the issue, AI argued against it: it showed that \texttt{\{True, NaN\}} and \texttt{\{True, False\}} are informationally identical for the tree models we use, so the encoding would add complexity without any effect on the model, and recommended fixing the interpretation instead. Lukas accepted the reversal because the argument was concrete rather than stylistic.},
  aievidence = {Claude Code session on 2026-09-29 analysing \href{https://github.com/mlops-2627q1-mds-upc/MLOPS_Recommenditos/issues/25}{issue \#25}: AI presented the options as a decision question, Lukas chose ``keep them binary''.},
  evidence = {\href{https://github.com/mlops-2627q1-mds-upc/MLOPS_Recommenditos/issues/25}{Issue \#25}, \href{https://github.com/mlops-2627q1-mds-upc/MLOPS_Recommenditos/pull/24}{PR \#24}, \href{https://github.com/mlops-2627q1-mds-upc/MLOPS_Recommenditos/blob/main/docs/docs/problem-spec.md}{problem specification} section 4.},
}

% Generated by tools/edn_latex.py from reports/edn.md.
% Do not edit: change reports/edn.md and run `make edn`.
\ednentry{
  id = {EDN-24},
  title = {Keep the pre-registered exclusion although the flag behind it is unreliable},
  date = {2026-09-29},
  milestone = {M1: Project Inception},
  topic = {Problem Specification (model scope)},
  participants = {@lukas2510},
  decision = {The scope of EDN-04 stands: the pipeline keeps excluding listings with \texttt{is\_\allowbreak{}preregistered =\allowbreak{} True} and the API keeps rejecting pre-registered cars. Because a \texttt{False} in that flag can also mean ``not stated'' (EDN-23), an unknown residue of unmarked pre-registered listings stays in the training data; this is recorded as a known limitation in the problem spec instead of being worked around, and is to be carried into the dataset card and the model card. \texttt{is\_\allowbreak{}preregistered} is therefore not a feature: a column that defines the scope filter cannot also be a model input.},
  rationale = {\emph{Alternatives considered:}
\begin{itemize}[leftmargin=*, topsep=2pt]
\item \textbf{Option A (chosen): keep the exclusion, document the residue.}
\newline Pros: keeps the product scope of EDN-04 (``used-car price'') intact and the API rule unchanged; removes the 3,695 listings that are demonstrably pre-registered; states the limitation honestly instead of implying a clean filter.
\newline Cons: the residue cannot be quantified, only bounded by the 14,744 unflagged rows; a small amount of out-of-scope data stays in training.
\item \textbf{Option B: revise EDN-04 and keep pre-registered cars in scope, with the flag as a feature.}
\newline Pros: more data; a pre-registered car really is priced differently, so the flag carries genuine signal; no contradiction between filter and feature.
\newline Cons: reopens a settled scope decision whose deciding argument was product clarity, not accuracy (EDN-04 measured 6.8\,\% vs 6.9\,\% MdAPE, so almost nothing is at stake); the API would have to accept and explain a vehicle condition the component is not meant to price; the feature would be exactly as unreliable as the filter, because the same \texttt{False} ambiguity applies.
\item \textbf{Option C: detect unmarked pre-registered listings with a heuristic and drop them too.}
\newline Pros: would shrink the residue.
\newline Cons: a heuristic on recent registration and low mileage has no ground truth to validate against, so it would trade a known, documented residue for an unknown number of wrongly dropped ordinary used cars.
\end{itemize}
\par
\emph{Why:} An earlier profiling run suggested promoting \texttt{is\_\allowbreak{}preregistered} into the extended feature set, because it is \texttt{True} in 3,695 of the 113,708 rows that survive a scoping by \texttt{offer\_\allowbreak{}type} and \texttt{vehicle\_\allowbreak{}type}. That run simply did not apply the pre-registered filter, so the finding contradicted EDN-04 rather than extending it, and the suggestion was withdrawn. Of the remaining options, reopening the scope buys almost no accuracy (EDN-04's own measurement) while costing the product clarity that decided EDN-04 in the first place, and a heuristic would invent labels. Documenting a bounded, honest limitation is the proportionate answer.},
  ai = {Information seeking, Alternative generation, Alternative assessment, Recommendation},
  response = {Accepted},
  assessment = {AI found the contradiction between its own change in PR \#24 and EDN-04 while checking whether issue \#25 was still current, and reported that the scope filter itself, not only the feature set, is affected by the unreliable flag. It laid out the three options and recommended keeping EDN-04. Lukas accepted it, because reopening a settled scope decision for an effect EDN-04 had already measured as negligible is not worth it.},
  aievidence = {Claude Code session on 2026-09-29 analysing \href{https://github.com/mlops-2627q1-mds-upc/MLOPS_Recommenditos/issues/25}{issue \#25}: AI reported the contradiction unprompted and presented the options as a decision question; Lukas chose to keep EDN-04.},
  evidence = {\href{https://github.com/mlops-2627q1-mds-upc/MLOPS_Recommenditos/issues/25}{Issue \#25}, \href{https://github.com/mlops-2627q1-mds-upc/MLOPS_Recommenditos/pull/24}{PR \#24}, EDN-04, \href{https://github.com/mlops-2627q1-mds-upc/MLOPS_Recommenditos/blob/main/docs/docs/problem-spec.md}{problem specification} section 2.},
}

% Generated by tools/edn_latex.py from reports/edn.md.
% Do not edit: change reports/edn.md and run `make edn`.
\ednentry{
  id = {EDN-25},
  title = {Raw data acquisition: track with \texorpdfstring{\texttt{dvc add}}{dvc add} and push to our DagsHub remote},
  note = {\textbf{Amended by EDN-36} (2026-09-30, PR \#54): the \texttt{dvc add} pointer is retired, the \texttt{download} stage acquires the file from Zenodo itself, and no copy of the CSV is hosted by us any more.
This entry is kept as the record of what was decided on 2026-09-26 and of the reasoning EDN-36 revised; the rejection of \texttt{dvc import-\allowbreak{}url} is the part of it that still stands.},
  date = {2026-09-26},
  milestone = {M2: Reproducibility},
  topic = {Data Versioning},
  participants = {@W11W11W11, confirmed by @lukas2510 on 2026-09-29 when the premise below was corrected. Pending confirmation by the full team at the next sprint review.},
  decision = {Amends EDN-07, which had proposed importing the raw file from Zenodo and never pushing it to our own remote.
Download the raw file \texttt{autoscout24\_\allowbreak{}dataset\_\allowbreak{}20251108.\allowbreak{}csv} once from Zenodo, track it with \texttt{dvc add data/\allowbreak{}raw/\allowbreak{}autoscout24\_\allowbreak{}dataset\_\allowbreak{}20251108.\allowbreak{}csv} and push it to our DagsHub remote.
Teammates get it with \texttt{dvc pull}; the PII columns are removed later in preprocessing.
The remote is public (EDN-20), so this re-publishes the raw file's PII columns under the team's own name; that is accepted, with the reasoning recorded in EDN-20.},
  rationale = {\emph{Alternatives considered:}
\begin{itemize}[leftmargin=*, topsep=2pt]
\item \textbf{Option A: \texttt{dvc import-\allowbreak{}url} from Zenodo.} The \texttt{.\allowbreak{}dvc} file records the Zenodo URL and the file hash; the raw file is fetched from Zenodo and never stored on our remote.
\newline Pros: no second copy of the raw PII (\texttt{vin}, \texttt{street}, \texttt{seller\_\allowbreak{}company\_\allowbreak{}name}, coordinates); provenance is recorded in the pointer file itself.
\newline Cons: every fresh setup downloads 548.6 MB from Zenodo and depends on it being online; Zenodo sends no ETag, so \texttt{dvc update} change detection is less reliable; differs from the course demo, which uses \texttt{dvc add}.
\item \textbf{Option B (chosen): \texttt{dvc add} plus \texttt{dvc push} to DagsHub.}
\newline Pros: the workflow of the course demo and of issue \#3; one place to pull all data from, independent of Zenodo; fast pulls.
\newline Cons: we store a copy of the raw PII ourselves, on a remote that is public, so the team becomes a publisher of that personal data in its own right.
\item \textbf{Option C: \texttt{dvc add} plus push to a remote made private for this purpose.}
\newline Pros: the raw PII would sit behind an access wall.
\newline Cons: same access overhead as B, plus a manual collaborator entry for every grader and supervisor; it buys little real protection, because the identical file stays publicly downloadable from Zenodo either way (EDN-20).
\end{itemize}
\par
\emph{Why:} B follows the DVC workflow the course demo prescribes; \texttt{docs/\allowbreak{}dvc-\allowbreak{}demo.\allowbreak{}md} there teaches \texttt{dvc add} on the raw file, and its \texttt{dvc import}/\texttt{dvc get} commands are for data coming from another DVC or Git repository rather than from an arbitrary URL.
The remote is public, so pushing the raw file re-publishes its PII columns under the team's own name. We accept that because the identical file is already publicly downloadable from the pinned Zenodo DOI under the same licence, and because the marginal exposure is small while the access cost of a private remote for graders and supervisors would be real (EDN-20).
Option A was also weaker than it looked on the mechanism itself: the Zenodo URL answers a HEAD request with neither an \texttt{ETag} nor a \texttt{Content-\allowbreak{}MD5} header, and \texttt{dvc import-\allowbreak{}url} derives its change detection from one of those, so the mechanism was never verified against this source.
The MD5 of our copy matches the checksum Zenodo publishes (\texttt{b23a122cc51baf7de39f449193ff0d28}), so provenance stays verifiable either way.},
  ai = {Information seeking, Alternative generation, Alternative assessment, Recommendation},
  response = {Rejected},
  assessment = {AI (Claude Code) checked the Zenodo endpoint (size, MD5, no ETag), laid out options A to C and recommended A to keep the raw PII off our infrastructure. Mark chose B because it follows the course demo.
After the choice, AI checked the DagsHub repository's visibility, reported it as private, and recorded that as the mitigation for the main risk it had raised against B. That check was wrong: DagsHub refuses anonymous access to every repository, public ones included, so the observation had no discriminating power, and the authenticated API reports \texttt{private: false}.
The error is recorded rather than quietly fixed, because the decision stands while the reason given for it did not: B is still the right option, but on the grounds in the rationale above rather than on a privacy mitigation that never existed. Verified on 2026-09-29 against a known-public control repository and a repository that does not exist, both of which answer anonymously exactly like ours.},
  aievidence = {Claude Code session on 2026-09-26 while working on issue \#3: prompt ``How should the raw AutoScout24 file get into DVC (PII handling)?'' with options A to C and the recommendation for A; Mark chose B, reason ``because it follows the demo''.},
  evidence = {\href{https://github.com/mlops-2627q1-mds-upc/MLOPS_Recommenditos/issues/3}{issue \#3}, pointer file \texttt{data/\allowbreak{}raw/\allowbreak{}autoscout24\_\allowbreak{}dataset\_\allowbreak{}20251108.\allowbreak{}csv.\allowbreak{}dvc}, \href{https://github.com/mlops-2627q1-mds-upc/MLOPS_Recommenditos/blob/main/docs/docs/data-versioning.md}{data versioning conventions}, project brief \S{}3.3, EDN-07, EDN-19, EDN-20, EDN-36.},
}

% Generated by tools/edn_latex.py from reports/edn.md.
% Do not edit: change reports/edn.md and run `make edn`.
\ednentry{
  id = {EDN-26},
  title = {NFR-11's window is 1,000 requests, and \texorpdfstring{\texttt{model}}{model} is excluded from the drift comparison},
  note = {\textbf{Amended by EDN-29} (2026-09-29): the drift job runs at a significance level of 0.005, not 0.05.
The configuration below was measured on a scope that EDN-22 has since narrowed, and on one split; re-measured, its
control clause does not hold at 0.05.},
  date = {2026-09-29},
  milestone = {M6: Monitoring},
  topic = {Monitoring, Requirements},
  participants = {@lukas2510},
  decision = {Amends EDN-14. The drift job compares windows of \textbf{1,000} requests, not 100, against the 10,000-listing reference, and leaves \textbf{\texttt{model}} out of the comparison. The other clauses of NFR-11 are unchanged: the \texttt{ES} replay must be flagged, at least three features other than \texttt{country\_\allowbreak{}code} must be named, and at most 1 of 20 i.i.d. control windows may be flagged.},
  rationale = {\emph{Alternatives considered:}
\begin{itemize}[leftmargin=*, topsep=2pt]
\item \textbf{Option A (chosen): window 1,000 and \texttt{model} excluded.}
\newline Pros: the only configuration measured to satisfy all three clauses as written. Both changes are independently justified: three drifted features do not reliably co-occur in a 100-row window, and a 349-level chi-square against such a window is under-powered rather than merely noisy.
\newline Cons: the M6 drift demonstration needs 1,000 replayed requests instead of 100, so the monitoring section of the report shows a slower signal.
\item \textbf{Option B: window 1,000, keep \texttt{model}, relax the control clause to at most 2 in 20.}
\newline Pros: no feature is removed, so the monitoring watches everything it sees.
\newline Cons: weakens a promise instead of fixing its cause, and keeps a test whose cell counts do not support the chi-square.
\item \textbf{Option C: keep window 100 and relax both clauses (at least one changed feature, at most 3 false alarms in 20).}
\newline Pros: keeps the fast 100-request demonstration.
\newline Cons: two promises weakened at once; 3 false alarms in 20 windows is hard to defend for a monitoring requirement.
\item \textbf{Option D: keep NFR-11 as EDN-14 wrote it.}
\newline Cons: measured to fail. At window 100 with all features the control flags 5.4 of 20 windows, and the three-feature clause holds in only 81.5\,\% of trials.
\end{itemize}
\par
\emph{Why:} NFR-11 as written was never measured in one configuration: EDN-14 took its \texttt{ES} half from \texttt{nfr11\_\allowbreak{}check.\allowbreak{}py} (window 100, full reference) and its control half from \texttt{nfr11\_\allowbreak{}diag.\allowbreak{}py} (window 1,000, reference capped at 10,000), and the latter's output was never committed. Measuring both halves together (\texttt{nfr11\_\allowbreak{}model\_\allowbreak{}excluded.\allowbreak{}py}, 200 trials per cell, 10,000-listing reference, i.i.d. control) gives:
\par
{\small\renewcommand{\arraystretch}{1.15}%
\begin{tabularx}{\linewidth}{@{}>{\raggedright\arraybackslash}X>{\raggedright\arraybackslash}X>{\raggedright\arraybackslash}X>{\raggedright\arraybackslash}X@{}}
\toprule
Configuration & \texttt{ES} flagged & three features besides \texttt{country\_\allowbreak{}code} & control windows flagged of 20 \\
\midrule
all features, window 100 & 1.000 & 81.5\,\% of trials, minimum 1 & 5.4 \\
\texttt{model} excluded, window 100 & 1.000 & 79.5\,\% of trials, minimum 1 & 2.2 \\
all features, window 1,000 & 1.000 & 100\,\%, minimum 13 & 1.4 \\
\textbf{\texttt{model} excluded, window 1,000} & \textbf{1.000} & \textbf{100\,\%, minimum 12} & \textbf{0.7} \\
\bottomrule
\end{tabularx}}
\par
Neither change suffices alone. Excluding \texttt{model} roughly halves the false alarms at both window sizes, because at window 100 it alone accounts for 0.190 of the 0.270 flag rate, but it does not reach 1 in 20 and does not help the three-feature clause at all. Enlarging the window fixes the three-feature clause and most of the false alarms, but with \texttt{model} kept the control still flags 1.4 of 20. Only the combination satisfies every clause, so NFR-11 now states a configuration that has been demonstrated rather than assembled from two runs.},
  ai = {Information seeking, Alternative generation, Alternative assessment, Recommendation, Solution generation},
  response = {Accepted with modifications},
  assessment = {A reviewing agent reported that the control clause fails at window 100 and proposed stating window 1,000. A second agent argued that the window was the wrong lever, because the false alarms come from \texttt{model}, and recommended excluding it while keeping window 100; Lukas chose that option. The measurement then refuted it: the control improves from 5.4 to 2.2 flagged windows of 20 but stays above the promised 1, and the three-feature clause fails in one run in five at window 100 either way. Two claims in that recommendation were also wrong and are recorded rather than quietly dropped: \texttt{model} was described as having roughly 10,000 levels, where it has 349 in the reference (84,171 is \texttt{model\_\allowbreak{}version}, a different column), and the reviewing agent's control figure at window 1,000 was 0.6 of 20 against the 1.4 measured here, a difference that decides whether the clause holds with \texttt{model} kept. The argument that a 349-level chi-square against a 100-row window is under-powered survived the measurement; its effect was simply smaller than claimed. Lukas then chose the combination, which is the only configuration that was demonstrated to work.},
  aievidence = {Claude Code session on 2026-09-29: a peer agent's review of PR \#19 reported the control failure and proposed window 1,000; this session verified the finding by inspection, argued for excluding \texttt{model} instead, and Lukas chose that. After the measurement refuted it, this session reported the four measured configurations and recommended the combination, which Lukas chose.},
  evidence = {\href{https://github.com/mlops-2627q1-mds-upc/MLOPS_Recommenditos/blob/main/docs/docs/requirements.md}{requirements} NFR-11; \href{https://github.com/mlops-2627q1-mds-upc/MLOPS_Recommenditos/blob/main/docs/docs/specification.md}{specification} NFR-11; \href{https://github.com/mlops-2627q1-mds-upc/MLOPS_Recommenditos/blob/main/reports/analysis}{reports/analysis/} (\texttt{nfr11\_\allowbreak{}model\_\allowbreak{}excluded.\allowbreak{}py}, \texttt{nfr11\_\allowbreak{}model\_\allowbreak{}excluded\_\allowbreak{}results.\allowbreak{}json}, run of 2026-09-29); \hyperref[EDN-14]{EDN-14}; \href{https://github.com/mlops-2627q1-mds-upc/MLOPS_Recommenditos/pull/19}{PR \#19}.},
}

% Generated by tools/edn_latex.py from reports/edn.md.
% Do not edit: change reports/edn.md and run `make edn`.
\ednentry{
  id = {EDN-27},
  title = {Separate the requirements from their specification, keeping one set of FR/NFR IDs},
  date = {2026-09-29},
  milestone = {M1: Project Inception},
  topic = {Requirements Engineering},
  participants = {@lukas2510},
  decision = {\texttt{docs/\allowbreak{}docs/\allowbreak{}requirements.\allowbreak{}md} keeps the \texttt{FR-\allowbreak{}xx} and \texttt{NFR-\allowbreak{}xx} vocabulary but is rewritten at a design-free level: what must be true, for whom, with a priority and an acceptance criterion, and without endpoints, formats, status codes, libraries or deployment mechanics. Everything design-level moves to a new \texttt{docs/\allowbreak{}docs/\allowbreak{}specification.\allowbreak{}md}, which carries the same IDs, so \texttt{FR-\allowbreak{}04} there specifies \texttt{FR-\allowbreak{}04} here. Tests, the requirement-to-test matrix and all earlier EDN entries keep pointing at one unchanged ID space.},
  rationale = {\emph{Alternatives considered:}
\begin{itemize}[leftmargin=*, topsep=2pt]
\item \textbf{Option A (chosen): two levels, one ID space.}
\newline Pros: the requirements answer what the component must do and can be read by someone who does not know the architecture, which is what the feedback asked for; every design decision that was already analysed and settled (EDN-08, EDN-11 to EDN-18) survives in full, at the level where it belongs; a later design change, for example a different endpoint split or a different deployment path, no longer looks like a requirements change; the \texttt{FR-\allowbreak{}xx}/\texttt{NFR-\allowbreak{}xx} IDs that tests and twelve earlier EDN entries refer to keep working, and the requirement-to-test traceability is unaffected.
\newline Cons: two pages to keep consistent, and a reader has to know which of the two answers their question; every cross-reference in the brief, the problem specification, \texttt{data-\allowbreak{}versioning.\allowbreak{}md}, \texttt{CONTRIBUTING.\allowbreak{}md} and the EDN evidence links had to be repointed to the level it actually meant.
\item \textbf{Option B: rewrite the requirements as user stories (\texttt{US-\allowbreak{}xx}) with a specification behind them.}
\newline Pros: the most explicit way to show the user's intent, and the standard form in agile requirements engineering, which fits the Scrum setup of the course.
\newline Cons: introduces a second ID space on top of \texttt{FR-\allowbreak{}xx}/\texttt{NFR-\allowbreak{}xx}, so every test and every earlier EDN entry would reach its requirement through one more hop; the course material and the report template speak of functional and non-functional requirements, so the team would be renaming the artefact the report has to deliver. The team tried this form first and rejected it for these reasons; the abstraction it was meant to bring is achieved by the rewrite in option A without the renaming.
\item \textbf{Option C: keep one page and only soften its wording.}
\newline Pros: no new file, no link changes.
\newline Cons: the design decisions would still sit inside the requirements, only less visibly, so the criticism would be answered in form and not in substance; and there would be no place left for the endpoint-level detail that M4 needs.
\end{itemize}
\par
\emph{Why:} The lecturer's feedback on the M1 requirements was that they read as a specification: too technical, and fixing design where they should fix intent. The diagnosis holds. The page stated endpoints, HTTP status codes, header names, container mechanics and image tags, which meant that changing any of those looked like changing a requirement, and that a reader could not tell what the component owes its users from what we happened to choose. The detail itself is not the problem, since it is what the tests and the M4 deployment work are built on, so it was moved rather than dropped. Keeping a single ID space across both levels was the deciding factor against option B: the traceability the report has to show runs from a requirement ID to a test, and adding a second vocabulary on top of that would have cost a hop and gained nothing the rewrite does not already give.},
  ai = {Alternative generation, Alternative assessment, Recommendation, Solution generation},
  response = {Accepted with modifications},
  assessment = {Lukas relayed the feedback and asked for a more abstract version, with the explicit constraint that the existing analysis must not be lost. AI confirmed the criticism against the actual page and proposed the split into two levels, which the team accepted. Its first execution went further than asked and replaced the requirements with user stories; Lukas corrected this, because the course material, the report template and the existing IDs all speak of functional and non-functional requirements, and the stories had been meant as orientation only. AI then reworked the page to the agreed form. The abstraction of each requirement and the acceptance criteria are its proposal and were reviewed by the team.},
  aievidence = {Claude Code session on 2026-09-29 (prompts were in German, translated here): Lukas relayed the feedback as ``we got feedback on the requirements. The lecturer said they are too technical, they are more of a specification. He would rather have them like user stories, more abstract. Can you do that, but please also write a spec so the work was not for nothing'', and after AI delivered user stories corrected it with ``user stories were only meant as orientation for you. We still want to call them functional and non-functional requirements. But as said they should not be so technical, rather more abstract, and should not fix architecture decisions yet''.},
  evidence = {\href{https://github.com/mlops-2627q1-mds-upc/MLOPS_Recommenditos/blob/main/docs/docs/requirements.md}{Requirements}; \href{https://github.com/mlops-2627q1-mds-upc/MLOPS_Recommenditos/blob/main/docs/docs/specification.md}{Specification}; \href{https://github.com/mlops-2627q1-mds-upc/MLOPS_Recommenditos/pull/19}{PR \#19}.},
}

% Generated by tools/edn_latex.py from reports/edn.md.
% Do not edit: change reports/edn.md and run `make edn`.
\ednentry{
  id = {EDN-28},
  title = {One module per DVC stage, deviating from the flat Cookiecutter layout},
  date = {2026-09-29},
  milestone = {M2: Reproducibility},
  topic = {Project Structure},
  participants = {@lukas2510},
  decision = {The pipeline code is one module per DVC stage, \texttt{recommenditos/\allowbreak{}data/\allowbreak{}\{download\_\allowbreak{}raw\_\allowbreak{}dataset,preprocess,gx\_\allowbreak{}context\_\allowbreak{}configuration,validate\_\allowbreak{}data,split\_\allowbreak{}data,build\_\allowbreak{}features\}.\allowbreak{}py} and \texttt{recommenditos/\allowbreak{}modeling/\allowbreak{}\{train,evaluate\}.\allowbreak{}py}, instead of the flat Cookiecutter Data Science layout of \texttt{dataset.\allowbreak{}py}, \texttt{features.\allowbreak{}py} and \texttt{modeling/\allowbreak{}train.\allowbreak{}py} that the template generates. This is the structure the course demo repository uses.},
  rationale = {\emph{Alternatives considered:}
\begin{itemize}[leftmargin=*, topsep=2pt]
\item \textbf{Option A (chosen): one module per stage, as in the teachers' demo.}
\newline Pros: it is the layout the course demonstrates, so it needs no defending against the graders' own example; each \texttt{dvc.\allowbreak{}yaml} stage declares exactly its own module in \texttt{deps}, so a code change reruns that stage and no other; and each ticket owns one file, which is what lets five people build a sequential pipeline in parallel without colliding.
\newline Cons: a deviation from the prescribed template, which has to be justified in the report.
\item \textbf{Option B: keep the flat Cookiecutter layout as generated.}
\newline Pros: zero deviation from the prescribed template.
\newline Cons: download, preprocessing and the split would share \texttt{dataset.\allowbreak{}py}, so three people would edit one file in the same week; \texttt{deps} would be coarse, so any edit to that file reruns every stage that lists it, which weakens exactly the reproducibility property the milestone is about.
\item \textbf{Option C: keep the Cookiecutter file names and split only where a collision forces it.}
\newline Pros: a smaller deviation.
\newline Cons: neither the template nor the demo, so the report would have to justify a third structure that follows nothing; and the split boundary would be decided by whoever hits the collision first rather than by the pipeline.
\end{itemize}
\par
\emph{Why:} Cookiecutter Data Science is prescribed ``justified deviations allowed'', and the teachers' own demo repository deviates in exactly this way, which makes it the best-supported choice rather than a liberty we take. The technical argument is the DVC dependency graph: a stage's \texttt{deps} should name the code that stage actually runs, so that \texttt{dvc repro} reruns the minimum. The flat layout cannot express that, because one file backs several stages. The organisational argument is the sprint 2 cut: the pipeline is sequential, so parallel work is only possible if each ticket owns its own files, and one module per stage gives that for free.},
  ai = {Information seeking, Alternative generation, Alternative assessment, Recommendation},
  response = {Accepted},
  assessment = {Asked whether the sprint plan matched the course demo repository, AI read the demo's file tree and its \texttt{dvc.\allowbreak{}yaml} rather than relying on the existing notes, and reported that the demo splits one module per stage while our repository still carried the flat template. It connected that to the parallelisation problem the sprint was being cut around, laid out the three options with the trade-offs above and recommended A. Lukas reviewed the comparison and chose A. The contribution was useful mainly because it checked the demo directly instead of arguing from the template's documentation.},
  aievidence = {Claude Code session on 2026-09-29 during sprint 2 planning: Lukas gave the demo repository's URL and asked whether the planned milestones and tickets fit it; AI fetched its tree and \texttt{dvc.\allowbreak{}yaml}, reported what to adopt and what we deliberately add, and presented the three layout options; Lukas chose the demo layout.},
  evidence = {\href{https://github.com/mlops-2627q1-mds-upc/MLOPS_Recommenditos/blob/main/docs/docs/scrum/sprints.md}{sprint 2 planning notes}; \href{https://github.com/mlops-2627q1-mds-upc/MLOPS_Recommenditos/issues/32}{issue \#32}; \href{https://github.com/mlops-2627q1-mds-upc/MLOPS_Recommenditos/blob/main/references/course-demos.md}{references/course-demos.md}; \href{https://github.com/mlops-2627q1-mds-upc/MLOPS_Recommenditos/pull/45}{PR \#45}.},
}

% Generated by tools/edn_latex.py from reports/edn.md.
% Do not edit: change reports/edn.md and run `make edn`.
\ednentry{
  id = {EDN-29},
  title = {The drift job runs at a significance level of 0.005, not 0.05},
  date = {2026-09-29},
  milestone = {M6: Monitoring},
  topic = {Monitoring, Requirements},
  participants = {@lukas2510},
  decision = {Amends EDN-26. The drift job tests each feature at a Bonferroni-corrected threshold of \texttt{0.\allowbreak{}005 /\allowbreak{} number of features} instead of the conventional \texttt{0.\allowbreak{}05 /\allowbreak{} number of features}. Everything else stays: windows of 1,000 requests, a 10,000-listing reference, \texttt{model} excluded, an i.i.d. control of which at most 1 of 20 windows may be flagged.},
  rationale = {\emph{Alternatives considered:}
\begin{itemize}[leftmargin=*, topsep=2pt]
\item \textbf{Option A (chosen): lower the significance level to 0.005.}
\newline Pros: measured to hold with margin. Across three independent i.i.d. splits the control flags between 0 and 0.6 windows of 20 against a promise of at most 1, while the \texttt{ES} replay is still flagged in every trial of every split with at least twelve features besides \texttt{country\_\allowbreak{}code}. It fixes the cause rather than the symptom: the requirement and the detector's threshold were coupled, and nobody had noticed.
\newline Cons: a lower threshold means less sensitivity to genuinely subtle drift, which is the monitoring's real job. We have no measurement of that sensitivity, so the cost is real but unquantified.
\item \textbf{Option B: significance level 0.01.}
\newline Pros: the more conventional step down, and it does hold: 0 to 1.0 flagged windows of 20 across the three splits.
\newline Cons: the maximum sits exactly on the promise, so the requirement would again depend on which split is drawn, which is the failure mode being fixed.
\item \textbf{Option C: keep 0.05 and relax the control clause to at most 3 of 20.}
\newline Pros: no change to the detector; the measured maximum is 2.0, so 3 holds.
\newline Cons: weakens a monitoring promise to fit a threshold nobody chose deliberately. Three false alarms in twenty windows is hard to defend in the report, and the underlying coupling would stay hidden.
\item \textbf{Option D: keep the requirement as EDN-26 wrote it.}
\newline Cons: measured not to hold. On the scope EDN-22 defines, the control flags 2.3 windows of 20 at window 1,000 with \texttt{model} excluded.
\end{itemize}
\par
\emph{Why:} EDN-26's configuration was measured once, on one i.i.d. split, before EDN-22 narrowed the scope. Re-running it after that change gave 2.3 flagged control windows of 20 instead of 0.7, which 27 removed listings cannot explain: the i.i.d. split is drawn from a permutation whose length depends on the row count, so the second run used a different control sample. The estimate was never stable, and the reason is structural. The drift job flags a window when any feature falls below \texttt{P\_\allowbreak{}VAL /\allowbreak{} n\_\allowbreak{}features}, so by Bonferroni's construction the family-wise false-alarm rate \textbf{is} \texttt{P\_\allowbreak{}VAL}. At \texttt{P\_\allowbreak{}VAL =\allowbreak{} 0.\allowbreak{}05} that is 5\,\%, which is exactly the ``at most 1 of 20'' the requirement promises, so NFR-11 was asking the test to perform at its own theoretical bound with zero tolerance for estimation error. Lowering the level decouples the two. 0.005 was chosen over 0.01 because at 0.01 the measured maximum is exactly 1.0 of 20, which would leave the requirement depending on the draw again. The \texttt{ES} side is unaffected: its p-values are in the order of 1e-300, so no threshold in this range changes whether it is detected.
\par
Control windows flagged of 20, three splits, window 1,000, \texttt{model} excluded:
\par
{\small\renewcommand{\arraystretch}{1.15}%
\begin{tabularx}{\linewidth}{@{}>{\raggedright\arraybackslash}X>{\raggedright\arraybackslash}X>{\raggedright\arraybackslash}X>{\raggedright\arraybackslash}X>{\raggedright\arraybackslash}X@{}}
\toprule
\texttt{P\_\allowbreak{}VAL} & min & max & mean & \texttt{ES} clause holds in every split \\
\midrule
0.05 & 0.0 & 2.0 & 1.0 & yes \\
0.01 & 0.0 & 1.0 & 0.4 & yes \\
\textbf{0.005} & \textbf{0.0} & \textbf{0.6} & \textbf{0.2} & \textbf{yes} \\
0.001 & 0.0 & 0.6 & 0.2 & yes \\
\bottomrule
\end{tabularx}}
\par
Window 100 fails the ``at least three changed properties'' clause at every level tested, so EDN-26's window decision stands on its own.},
  ai = {Information seeking, Alternative generation, Alternative assessment, Recommendation, Solution generation},
  response = {Accepted},
  assessment = {AI re-ran the committed evidence after EDN-22 landed, found that EDN-26's configuration no longer held, and did not simply restate the number. It identified that the Bonferroni construction makes the family-wise false-alarm rate equal to the significance level, which is what the control clause promises, and wrote a measurement that separates split-to-split variability from the threshold by computing each trial's p-values once and evaluating them at several levels. It recommended 0.005 over 0.01 on the grounds that 0.01's maximum sits exactly on the promise. Lukas delegated the choice and accepted the recommendation. This is the third time a measurement has refuted a written form of NFR-11, which is itself the argument for measuring a requirement before committing to it rather than after.},
  aievidence = {Claude Code session on 2026-09-29: after the scope correction AI reported that EDN-26's configuration measured 2.3 flagged control windows of 20 instead of 0.7, explained the Bonferroni coupling, ran \texttt{nfr11\_\allowbreak{}alpha\_\allowbreak{}sweep.\allowbreak{}py} over three splits and four significance levels, and recommended 0.005; Lukas replied ``ohne team mache das was du recommendest''.},
  evidence = {\href{https://github.com/mlops-2627q1-mds-upc/MLOPS_Recommenditos/blob/main/docs/docs/specification.md}{specification} NFR-11; \href{https://github.com/mlops-2627q1-mds-upc/MLOPS_Recommenditos/blob/main/reports/analysis}{reports/analysis/} (\texttt{nfr11\_\allowbreak{}alpha\_\allowbreak{}sweep.\allowbreak{}py}, \texttt{nfr11\_\allowbreak{}alpha\_\allowbreak{}sweep\_\allowbreak{}results.\allowbreak{}json}, \texttt{nfr11\_\allowbreak{}alpha\_\allowbreak{}sweep\_\allowbreak{}results.\allowbreak{}txt}, run of 2026-09-29); \hyperref[EDN-26]{EDN-26}; \hyperref[EDN-22]{EDN-22}; \href{https://github.com/mlops-2627q1-mds-upc/MLOPS_Recommenditos/pull/47}{PR \#47}.},
}

% Generated by tools/edn_latex.py from reports/edn.md.
% Do not edit: change reports/edn.md and run `make edn`.
\ednentry{
  id = {EDN-30},
  title = {Pipeline configuration lives in \texorpdfstring{\texttt{params.\allowbreak{}yaml}}{params.yaml}, not in \texorpdfstring{\texttt{config.\allowbreak{}py}}{config.py}},
  date = {2026-09-29},
  milestone = {M2: Reproducibility},
  topic = {Project Structure, Reproducibility},
  participants = {@lukas2510},
  decision = {Everything configurable about the pipeline -- the seed, the age reference date, the scope bounds, the split ratios, the feature sets, the model variants and the SC-01 to SC-06 thresholds -- lives in \texttt{params.\allowbreak{}yaml} and is declared per stage under \texttt{params:} in \texttt{dvc.\allowbreak{}yaml}. Paths stay in \texttt{recommenditos/\allowbreak{}config.\allowbreak{}py}.},
  rationale = {\emph{Alternatives considered:}
\begin{itemize}[leftmargin=*, topsep=2pt]
\item \textbf{Option A (chosen): a \texttt{params.\allowbreak{}yaml} declared per stage.}
\newline Pros: \texttt{dvc repro} reruns exactly the stages a change affects, because DVC hashes the declared keys; \texttt{dvc params diff} shows what changed between two commits, which is evidence the report can cite; \texttt{dvc exp} becomes usable without a code change; and a value used by two stages cannot drift, because there is one copy of it.
\newline Cons: one more file to keep in step with the code, and a parameter typo is caught at run time rather than by the interpreter.
\item \textbf{Option B: Python constants in \texttt{config.\allowbreak{}py}, as the course demo does.}
\newline Pros: it is exactly what the demo shows, so it needs no defending; typos are import errors; no YAML parsing.
\newline Cons: DVC cannot see a constant, so a changed hyperparameter reruns nothing and \texttt{dvc repro} reports the pipeline as up to date while the code says otherwise -- which is the opposite of what the milestone is graded on. \texttt{dvc exp} cannot sweep it, and the demo's own \texttt{config.\allowbreak{}py} also reads an absolute \texttt{ROOT} from \texttt{.\allowbreak{}env} and raises at import without one.
\item \textbf{Option C: both, with \texttt{config.\allowbreak{}py} reading \texttt{params.\allowbreak{}yaml}.}
\newline Pros: one import surface for the stages.
\newline Cons: the indirection hides which stage reads which key, which is precisely the information \texttt{dvc.\allowbreak{}yaml} needs in order to rerun the minimum.
\end{itemize}
\par
\emph{Why:} The whole point of M2 is that a change reruns what it should and nothing else. Option B cannot deliver that for anything except code, and parameters are where most changes will happen once the stage tickets land. Paths were deliberately left out of \texttt{params.\allowbreak{}yaml}: they are not experiment variables, and deriving them from \texttt{\_\allowbreak{}\_\allowbreak{}file\_\allowbreak{}\_\allowbreak{}} keeps the project free of machine-specific absolute paths.},
  ai = {Alternative generation, Alternative assessment, Recommendation, Solution generation},
  response = {Accepted},
  assessment = {AI read the demo repository's \texttt{dvc.\allowbreak{}yaml} and \texttt{config.\allowbreak{}py} directly rather than from our notes, reported that the demo declares neither \texttt{params} nor \texttt{metrics} on any stage, and measured the consequence in a throwaway DVC repository: with a coarse \texttt{params:} declaration a single changed hyperparameter retrained every variant, and with a per-item declaration it retrained one. Lukas reviewed the measurement and accepted the recommendation.},
  aievidence = {Claude Code session on 2026-09-29 while building the pipeline skeleton: AI built a scratch DVC repository, ran \texttt{dvc repro} and \texttt{dvc status} under both declaration styles, and pasted the outputs showing one stage rerunning instead of three.},
  evidence = {\href{https://github.com/mlops-2627q1-mds-upc/MLOPS_Recommenditos/blob/main/params.yaml}{\texttt{params.\allowbreak{}yaml}}; \href{https://github.com/mlops-2627q1-mds-upc/MLOPS_Recommenditos/blob/main/dvc.yaml}{\texttt{dvc.\allowbreak{}yaml}}; \href{https://github.com/mlops-2627q1-mds-upc/MLOPS_Recommenditos/issues/32}{issue \#32}.},
}

% Generated by tools/edn_latex.py from reports/edn.md.
% Do not edit: change reports/edn.md and run `make edn`.
\ednentry{
  id = {EDN-31},
  title = {The processed-data contract is a hand-written \texorpdfstring{\texttt{schema.\allowbreak{}py}}{schema.py}, not Pandera},
  date = {2026-09-29},
  milestone = {M2: Reproducibility (shapes M3: Quality Assurance)},
  topic = {Data Validation, Testing Strategy},
  participants = {@lukas2510},
  decision = {\texttt{recommenditos/\allowbreak{}schema.\allowbreak{}py} holds the structural contract -- column names, dtypes and nullability for the raw, interim and processed frames -- as frozen dataclasses with a \texttt{validate} and a \texttt{conform} method, written by hand rather than with a validation library. Value rules stay in the Great Expectations suites. Every stage validates after reading and conforms before writing.},
  rationale = {\emph{Alternatives considered:}
\begin{itemize}[leftmargin=*, topsep=2pt]
\item \textbf{Option A (chosen): a hand-written \texttt{schema.\allowbreak{}py}.}
\newline Pros: no dependency, so the same module can later be imported by the API, which NFR-04 caps at 1 GB and which already has Pydantic for its own request schemas; we control the failure message, and it reports every problem at once with the offending row positions; and \texttt{conform} casts to the declared dtypes before each write, which is what stops Parquet's dtype drift between stages -- a library that only validates cannot do that. Great Expectations already gives us a declarative vocabulary for value rules, so a second declarative framework would overlap it.
\newline Cons: about 250 lines we own and test ourselves.
\item \textbf{Option B: Pandera.}
\newline Pros: battle-tested, declarative, less code to own, good lazy error reports as a tidy DataFrame.
\newline Cons: a third validation system next to Great Expectations and Pydantic; it would travel into the API image; and it does not check the datetime unit at all, so a \texttt{datetime64[ns]} declaration would be documentation rather than a guarantee -- exactly the failure this contract exists to prevent.
\item \textbf{Option C: Great Expectations alone, with no structural contract in code.}
\newline Pros: one tool, and Data Docs for free.
\newline Cons: the contract would exist only in a generated \texttt{gx/\allowbreak{}} store whose config carries a fresh UUID per object on every run, so it is awkward to diff and keep honest in git. It also arrives too late: the stage tickets need something to code against in week one, and Great Expectations is a later ticket.
\end{itemize}
\par
\emph{Why:} Structure and value are different questions. A structural break is a bug in our code and should fail the stage that caused it; a value break is a change in the data and belongs in an expectation suite with a \texttt{mostly=\allowbreak{}} tolerance. Splitting them that way let the contract land before Great Expectations does, which is what unblocks the parallel work. The decisive technical point is \texttt{conform}: Parquet preserves whatever dtype it is handed rather than normalising it, so without an explicit cast at each boundary the frames drift apart and the stage that notices is never the stage that caused it. Because \texttt{conform} selects only the contract's columns, the PII columns NFR-08 forbids are kept out structurally rather than by every stage remembering to drop them.},
  ai = {Information seeking, Alternative generation, Alternative assessment, Recommendation, Solution generation},
  response = {Accepted},
  assessment = {AI ran the compatibility experiments rather than arguing from documentation: it installed Pandera, Great Expectations and pyarrow against our pandas 3.0.6 and measured what survives a Parquet round trip, finding that an \texttt{object} string column silently becomes \texttt{str} across the boundary (so a schema declaring \texttt{object} passes upstream and fails downstream), that Pandera ignores the datetime unit entirely, and that the Great Expectations store regenerates with different UUIDs on every run. Those three measurements are what decided the option, and none of them was predictable from the documentation. Lukas was given the three options with these findings and chose A.},
  aievidence = {Claude Code session on 2026-09-29: AI ran the probes in a scratch environment, reported the before/after dtype table for the Parquet round trip and the sixteen declared-versus-actual datetime combinations Pandera accepts, and recommended the hand-written contract; Lukas chose it.},
  evidence = {\href{https://github.com/mlops-2627q1-mds-upc/MLOPS_Recommenditos/blob/main/recommenditos/schema.py}{\texttt{recommenditos/\allowbreak{}schema.\allowbreak{}py}}; \href{https://github.com/mlops-2627q1-mds-upc/MLOPS_Recommenditos/blob/main/tests/test_schema.py}{\texttt{tests/\allowbreak{}test\_\allowbreak{}schema.\allowbreak{}py}}; \href{https://github.com/mlops-2627q1-mds-upc/MLOPS_Recommenditos/issues/32}{issue \#32}.},
}

% Generated by tools/edn_latex.py from reports/edn.md.
% Do not edit: change reports/edn.md and run `make edn`.
\ednentry{
  id = {EDN-32},
  title = {Split proportions 60/10/10/20 and the project seed},
  date = {2026-09-29},
  milestone = {M2: Reproducibility},
  topic = {Evaluation Protocol},
  participants = {@lukas2510},
  decision = {The non-\texttt{ES} listings are split 60\,\% train, 10\,\% validation, 10\,\% calibration and 20\,\% test, grouped by \texttt{seller\_\allowbreak{}group\_\allowbreak{}id}, with the project seed \texttt{20251108} (the snapshot date). Both are pinned in \texttt{params.\allowbreak{}yaml}.},
  rationale = {\emph{Alternatives considered:}
\begin{itemize}[leftmargin=*, topsep=2pt]
\item \textbf{Option A (chosen): 60/10/10/20.}
\newline Pros: keeps the 20\,\% test share that the model card's SC-04 analysis and \texttt{reports/\allowbreak{}analysis/\allowbreak{}make\_\allowbreak{}support\_\allowbreak{}results.\allowbreak{}txt} already assume, so the published evidence stays true without a re-run. Calibration and validation get about 9,800 rows each, which is ample for conformal calibration and for early stopping.
\newline Cons: the smallest training set of the three options, about 58,700 rows.
\item \textbf{Option B: 70/10/10/10.}
\newline Pros: the most training data.
\newline Cons: halves the test set. \texttt{make\_\allowbreak{}support\_\allowbreak{}results.\allowbreak{}txt} gives the test share below which each supported make falls under SC-04's 500-row bar: Volvo at 14.9\,\% and Suzuki at 13.0\,\%. At a 10\,\% test share both drop out, so SC-04 would silently stop checking two of the eleven supported makes and the model card would have to be corrected.
\item \textbf{Option C: 64/8/8/20.}
\newline Pros: the same 20\,\% test share with slightly more training data.
\newline Cons: a less round number to explain in the report, for about 3,900 extra training rows.
\end{itemize}
\par
\emph{Why:} The test share is the binding constraint, not the training share. SC-04 only checks a segment once it holds 500 test rows, so the proportions decide how many makes the quality gate actually covers, and the model card already states which four of the eleven supported makes fall short at 20\,\%. Choosing anything below about 15\,\% would quietly enlarge that set, which is the kind of change that is invisible until someone re-reads the analysis. The seed is the snapshot date rather than a random number so that it is recognisable in a log and obviously not tuned.},
  ai = {Information seeking, Alternative assessment, Recommendation},
  response = {Accepted},
  assessment = {AI found the coupling the ticket did not mention: it read \texttt{make\_\allowbreak{}support\_\allowbreak{}results.\allowbreak{}txt}, noticed that its sensitivity table was computed at an assumed 20\,\% test share while the four-way split was still unpinned, and pointed out that pinning a smaller test share would invalidate a number already published in the model card. It laid out the three options against that constraint and recommended A. Lukas accepted.},
  aievidence = {Claude Code session on 2026-09-29: asked to pin the split, AI reported the per-make sensitivity thresholds from the committed analysis output and framed the three options around the 15\,\% floor they imply.},
  evidence = {\href{https://github.com/mlops-2627q1-mds-upc/MLOPS_Recommenditos/blob/main/params.yaml}{\texttt{params.\allowbreak{}yaml}}; \href{https://github.com/mlops-2627q1-mds-upc/MLOPS_Recommenditos/blob/main/reports/analysis/make_support_results.txt}{reports/analysis/make\_support\_results.txt}; \href{https://github.com/mlops-2627q1-mds-upc/MLOPS_Recommenditos/blob/main/docs/docs/dataset-card.md}{dataset card}; \href{https://github.com/mlops-2627q1-mds-upc/MLOPS_Recommenditos/issues/32}{issue \#32}.},
}

% Generated by tools/edn_latex.py from reports/edn.md.
% Do not edit: change reports/edn.md and run `make edn`.
\ednentry{
  id = {EDN-33},
  title = {A generated synthetic fixture, and a \texorpdfstring{\texttt{download.\allowbreak{}source}}{download.source} parameter so the skeleton runs without the raw file},
  date = {2026-09-29},
  milestone = {M2: Reproducibility (shapes M3: Quality Assurance)},
  topic = {Testing Strategy, Data Acquisition},
  participants = {@lukas2510},
  decision = {\texttt{recommenditos/\allowbreak{}data/\allowbreak{}synthetic.\allowbreak{}py} generates a schema-valid stand-in for the raw snapshot from public catalogues and a seed, never by sampling the real file. The test suite builds every fixture from it, and the \texttt{download} stage produces it while \texttt{download.\allowbreak{}source} in \texttt{params.\allowbreak{}yaml} is \texttt{synthetic}, so \texttt{dvc repro} is green on a clean clone with no credentials and no 548 MB download.},
  rationale = {\emph{Alternatives considered:}
\begin{itemize}[leftmargin=*, topsep=2pt]
\item \textbf{Option A (chosen): generate, and gate the source on a parameter.}
\newline Pros: \texttt{pytest} and \texttt{dvc repro} both run with no data access, which is what lets four people build pipeline stages at the same time; no personal data is committed to Git, so NFR-08 is satisfied by construction rather than by a rule someone has to remember; every edge case the pipeline rules exist for is guaranteed present, so a stage test asserts rather than hopes; and because the source is a parameter it is recorded in \texttt{dvc.\allowbreak{}lock} and shows up in \texttt{dvc params diff}, so a pipeline still running on synthetic data cannot pass unnoticed.
\newline Cons: the generator is code we maintain, and a pipeline on \texttt{main} can produce fabricated data until the download ticket lands.
\item \textbf{Option B: sample the fixture from the real file and commit it.}
\newline Pros: real distributions for free, and no generator to write.
\newline Cons: the raw file carries PII (street, postcode, coordinates, seller company name). EDN-25 accepts that file living behind a DVC pointer on our remote, because the identical file is public on Zenodo under the same licence; a \emph{committed sample} is a different thing, because it puts those rows in the Git history of a repository the whole cohort can read and cannot be removed later. \texttt{.\allowbreak{}gitignore} also excludes everything under \texttt{data/\allowbreak{}}, so it could not be committed without weakening that rule.
\item \textbf{Option C: no fixture; \texttt{dvc repro} and the tests run on the real file.}
\newline Pros: nothing is fabricated at any point.
\newline Cons: every test and every \texttt{dvc repro} would need \texttt{dvc pull} of 548 MB and DagsHub access, so the working agreements' requirement that a ticket have a local oracle would hold for nobody until the download ticket lands -- which is the bottleneck the whole sprint was cut to avoid.
\end{itemize}
\par
\emph{Why:} The two properties that matter are that nobody is blocked and that no personal data enters the Git history, and only a generated fixture gives both. Making the source a parameter rather than a hard-coded branch is what keeps the option honest: the choice is versioned in \texttt{dvc.\allowbreak{}lock}, it is visible in a params diff, and the stage refuses any other value with a \texttt{NotImplementedError} naming the ticket that implements it, so nobody can flip it early and silently get synthetic data. The fixture deliberately keeps the real make skew and gives each seller several listings, because a seller-grouped split on a flat distribution would be indistinguishable from a random one and the split ticket's central invariant would be untestable.
\par
Corrected on 2026-09-30 while reviewing \href{https://github.com/mlops-2627q1-mds-upc/MLOPS_Recommenditos/pull/54}{PR \#54}, and recorded rather than quietly edited: ``it is visible in a params diff'' above, and the same claim in option A, is wrong. \texttt{dvc params diff} compares the params files of two Git revisions, not the lock against the workspace, so it prints nothing about a workspace running on the synthetic source. What actually makes such a run visible is the committed \texttt{dvc.\allowbreak{}lock}, which records the source every artefact was built from, and \texttt{dvc status}, which reports a workspace that disagrees with it. The decision is unaffected: making the source a parameter is still what puts the choice in the lock.},
  ai = {Alternative generation, Alternative assessment, Recommendation, Solution generation},
  response = {Accepted},
  assessment = {AI framed the trade-off as the one between the team being unblocked and fabricated data reaching \texttt{main}, and proposed the parameter as the way to keep both, rather than presenting the synthetic stub as free. It also derived the fixture's edge-case list from the individual stage tickets so that each ticket's rule has something to act on, and pointed out the skew requirement that a naive uniform generator would have missed. Lukas reviewed the three options and accepted A.
\par
The first implementation was wrong in a way worth recording, because it is the failure mode this project's working agreements warn about. Every edge-case row was built as a copy of one body row, which made most of them duplicates of it on the deduplication key and gave all of them that row's country. Preprocessing would therefore have deleted several of the cases the fixture exists to provide, and at the project's own seed the whole block inherited \texttt{ES} and would have been diverted into the holdout, so the train, validation, calibration and test frames that four tickets build against would have contained none of them. Nothing in the first round of tests could see it: they asserted that each case was present in the raw frame, which was true. A second agent, asked to attack the change rather than describe it, reproduced the fault, and the fix now carries a test that asserts each case \emph{survives} the preprocessing rules rather than merely existing. The lesson is the one already in the working agreements: an agent's claim decides nothing until it is checked, and the check has to be at the point where the value is consumed.},
  aievidence = {Claude Code session on 2026-09-29: AI presented the three options with the parallel-work and PII consequences of each and recommended A; Lukas chose it. The generated distribution was then checked against the dataset card's published make and country shares.},
  evidence = {\href{https://github.com/mlops-2627q1-mds-upc/MLOPS_Recommenditos/blob/main/recommenditos/data/synthetic.py}{\texttt{recommenditos/\allowbreak{}data/\allowbreak{}synthetic.\allowbreak{}py}}; \href{https://github.com/mlops-2627q1-mds-upc/MLOPS_Recommenditos/blob/main/tests/test_data.py}{\texttt{tests/\allowbreak{}test\_\allowbreak{}data.\allowbreak{}py}}; \hyperref[EDN-25]{EDN-25}; \href{https://github.com/mlops-2627q1-mds-upc/MLOPS_Recommenditos/issues/32}{issue \#32}.},
}

% Generated by tools/edn_latex.py from reports/edn.md.
% Do not edit: change reports/edn.md and run `make edn`.
\ednentry{
  id = {EDN-34},
  title = {The derived raw Parquet keeps its PII columns and is pushed like any other stage output},
  date = {2026-09-30},
  milestone = {M2: Reproducibility},
  topic = {Data Versioning, Personal Data},
  participants = {@lukas2510. Open to revisit at the next sprint review, like EDN-25.},
  decision = {\texttt{data/\allowbreak{}raw/\allowbreak{}listings.\allowbreak{}parquet}, the output of the \texttt{download} stage, keeps the seven PII columns of the published file and is cached and pushed to the DagsHub remote like any other stage output. No \texttt{push: false}, and no PII removal before preprocessing. Our public remote therefore holds the same personal data twice, once as the CSV of EDN-25 and once as this Parquet.},
  rationale = {\emph{Alternatives considered:}
\begin{itemize}[leftmargin=*, topsep=2pt]
\item \textbf{Option A (chosen): push it, which is DVC's default for a stage output.}
\newline Pros: nothing to build. The marginal exposure is close to nothing, because the identical data already sits on the same public remote as a CSV (EDN-25) and on Zenodo under the same licence, so anyone who wants those columns already has them. \texttt{dvc pull} stays sufficient for a complete working state, and the pipeline's first artefact is byte-identical for everyone because it is pulled rather than regenerated.
\newline Cons: we hold a second artefact containing the same personal data under our own name. A deletion request, a licence question or a scope change then touches two artefacts instead of one. EDN-25's argument was that \emph{the identical file} is already public, and a format we produced ourselves is not literally that file.
\item \textbf{Option B: \texttt{push: false} on the output.} DVC keeps it in the local cache only and each person regenerates it from the CSV.
\newline Pros: no second copy leaves the machine.
\newline Cons: it moves the pipeline's first artefact from pulled to locally produced. Parquet writing is not guaranteed byte-stable across pyarrow versions, so a teammate on a different version would see every downstream stage rerun, which is exactly the reproducibility property NFR-06 asks for and that the skeleton was just measured to have.
\item \textbf{Option C: drop the PII in \texttt{download} and build \texttt{seller\_\allowbreak{}group\_\allowbreak{}id} there.}
\newline Pros: the Parquet every stage reads would carry no PII at all, so NFR-08 would hold from the first derived artefact rather than from the interim frame, and the file would still be pushed, so reproducibility would be unaffected.
\newline Cons: the raw layer stops being a faithful copy of the published file, which is what \texttt{RAW\_\allowbreak{}SCHEMA} and the raw expectation suite of \#25 exist to describe; and the group-key hashing moves out of \#34's first step into \#33, so a ticket boundary shifts for a gain that does not change what is publicly reachable.
\item \textbf{Option D: no Parquet at all, every stage reads the CSV.}
\newline Cons: each run parses 548.6 MB of CSV and re-infers its dtypes, which is the cost the stage exists to remove.
\end{itemize}
\par
\emph{Why:} The question is marginal exposure, not storage: measured on the real file the CSV is 548.6 MB, the Parquet of all 75 columns is 215.3 MB, and the seven PII columns are 3.1 MB of it. Against a copy of the same data that is already public on our own remote and on Zenodo, none of the alternatives changes what a reader can reach; they only change what we have to build and what we risk. B risks the reproducibility we have just demonstrated, and C buys the same nothing in exposure terms while making the raw layer no longer raw. So the honest reading is that A costs the least and hides the least.
\par
What A does cost is recorded rather than argued away: we become the publisher of that personal data in two artefacts instead of one, and EDN-25's ``the identical file is already public'' does not literally cover a format we produced ourselves. The entry exists so that a later decision to restrict the remote, or a request to remove the data, finds both artefacts named in one place.
\par
NFR-08 is unaffected either way. It forbids PII in the processed data, the prediction log, the comparables and the model artefacts; the raw layer is none of those, and preprocessing removes the columns structurally, because \texttt{Schema.\allowbreak{}conform} selects only the columns the interim contract names.},
  ai = {Information seeking, Alternative generation, Alternative assessment, Recommendation},
  response = {Accepted with modifications},
  assessment = {AI found the question in the first place, while reviewing the pipeline skeleton: the \texttt{download} stage's output had inherited EDN-25's acceptance without anyone noticing that EDN-25 had only ever weighed the CSV. Its first write-up estimated the Parquet at ``roughly another 550 MB''; asked to explain the trade-off, it measured instead and reported 215.3 MB, of which the PII columns are 3.1 MB, which moved the argument off storage entirely. It also found, while explaining, a reproducibility flaw in its own earlier suggestion of \texttt{push: false} and withdrew it. It recommended C as the clean option and A as defensible; Lukas chose A as the cheapest of the two it stood behind. The modification is that the cost of A is written down here rather than treated as settled by EDN-25.},
  aievidence = {Claude Code session on 2026-09-29 and 2026-09-30: the question was raised in the review of PR \#48, the sizes were measured against the Zenodo file, and the four options were laid out with the reproducibility catch in B; Lukas replied ``mache die einfachste variante die du recommendest''.},
  evidence = {\hyperref[EDN-25]{EDN-25}; \hyperref[EDN-20]{EDN-20}; \href{https://github.com/mlops-2627q1-mds-upc/MLOPS_Recommenditos/issues/33}{issue \#33}; \href{https://github.com/mlops-2627q1-mds-upc/MLOPS_Recommenditos/pull/48}{PR \#48}.},
}

% Generated by tools/edn_latex.py from reports/edn.md.
% Do not edit: change reports/edn.md and run `make edn`.
\ednentry{
  id = {EDN-35},
  title = {The \texorpdfstring{\texttt{download}}{download} stage owns \texorpdfstring{\texttt{data/\allowbreak{}raw/\allowbreak{}}}{data/raw/}, and the published CSV is a local cache outside the DAG},
  date = {2026-09-30},
  milestone = {M2: Reproducibility},
  topic = {Data Versioning, Pipeline Design},
  participants = {@lukas2510. Open to revisit at the next sprint review, like EDN-25.},
  decision = {The \texttt{download} stage owns \texttt{data/\allowbreak{}raw/\allowbreak{}}.
Under \texttt{download.\allowbreak{}source: zenodo} it fetches the pinned file from the Zenodo record itself and writes \texttt{data/\allowbreak{}raw/\allowbreak{}listings.\allowbreak{}parquet} as its only output, so acquisition is inside the pipeline instead of a manual step beside it.
The published CSV is kept as a local cache under \texttt{data/\allowbreak{}external/\allowbreak{}} and is neither a stage \texttt{out} nor a \texttt{dep}; \texttt{download.\allowbreak{}md5} in \texttt{params.\allowbreak{}yaml} is the dependency that pins the bytes, and the stage refuses to read a file that hashes to anything else.
The standalone pointer \texttt{data/\allowbreak{}raw/\allowbreak{}autoscout24\_\allowbreak{}dataset\_\allowbreak{}20251108.\allowbreak{}csv.\allowbreak{}dvc} is deleted, which EDN-36 records as the amendment to EDN-25.},
  rationale = {\emph{Alternatives considered:}
\begin{itemize}[leftmargin=*, topsep=2pt]
\item \textbf{Option A: the pointer stays and the stage reads it (issue \#33's option A).} \texttt{download} takes the \texttt{dvc add}-tracked CSV as a \texttt{dep} and writes only the Parquet.
\newline Pros: a fresh setup gets the CSV from DagsHub rather than from Zenodo, so it is fast and independent of Zenodo's uptime; one place to pull all data from; the snapshot's hash sits in a file at the tip of \texttt{main}, so recovering it never involves the Git history.
\newline Cons: acquisition stays a manual \texttt{dvc add} outside the pipeline, so ``reproduce from a clean clone'' is not true of the pipeline alone, which is what the data-versioning page and the course demo describe; we keep re-publishing the CSV's PII columns on a public remote; and DVC hashes every \texttt{dep} to answer \texttt{dvc status}, so 548.6 MB is read before the pipeline can say whether anything changed.
\item \textbf{Option B (chosen): the stage fetches the file, and the CSV is a local cache outside the graph.}
\newline Pros: the pipeline acquires its own input, so a clean clone reproduces it with \texttt{dvc repro} and no manual step; nothing about the CSV is re-published by us; and \texttt{dvc status} stays instant, because the stage's deps are code and params only. A clean clone does not even pay the download, because \texttt{data/\allowbreak{}raw/\allowbreak{}listings.\allowbreak{}parquet} is a pushed stage output and the stage is up to date after \texttt{dvc pull}.
\newline Cons: the file the stage really reads is not in the DAG, so DVC cannot report anything about it and a reader has to know that \texttt{download.\allowbreak{}md5} is standing in for it; and the repository stops holding a tracked copy of the raw snapshot (see the consequence below).
\item \textbf{Option C: declare the CSV as a plain \texttt{out} of the stage.}
\newline Pros: the acquired file would be in the graph, cached and recoverable with \texttt{dvc checkout}.
\newline Cons: DVC removes a stage's outputs before running it, so every \texttt{dvc repro download} would re-download 548.6 MB; and a cached \texttt{out} puts the raw PII on our remote again, which is the copy EDN-36 exists to stop.
\item \textbf{Option D: declare the CSV as a \texttt{dep}.}
\newline Cons: the same 548.6 MB hash per \texttt{dvc status} as option A, plus a worse failure mode: a missing \texttt{dep} is an error DVC raises before the stage runs, so the stage could never fetch the file it depends on, and a clean clone would be stuck.
\item \textbf{Option E: declare it as \texttt{persist: true, cache: false} (\texttt{dvc stage add -\allowbreak{}-\allowbreak{}outs-\allowbreak{}persist-\allowbreak{}no-\allowbreak{}cache}, verified present in DVC 3.67.1).}
\newline Pros: this is the option the first write-up of this decision missed, and it removes the objection to option C outright. \texttt{persist: true} means DVC does not delete the file before the run, so it is fetched once and reused, and \texttt{cache: false} keeps it out of the cache and off the remote, so there is no second copy of the PII. It would also make \texttt{dvc status} honest about the CSV: the file the stage reads would be named in the graph, and a damaged or swapped local copy would be reported by DVC and not only by our own check.
\newline Cons: a non-cached out is still hashed to build the lock and to answer \texttt{dvc status}, so it costs exactly the 548.6 MB read per status call that ruled out options A and D, on every run including the many that never touch the raw layer. It is also not recoverable despite being in the graph: \texttt{dvc checkout} cannot restore a file DVC never cached, so the lock would carry a hash with nothing behind it. And it would put a second, DVC-owned copy of a fact \texttt{download.\allowbreak{}md5} already states, so re-pinning the file would mean editing the parameter and refreshing the lock instead of only the parameter.
\end{itemize}
\par
\emph{Why:} The honest reason the CSV stays out of the graph is not that DVC cannot express ``a file the stage fetches and keeps''.
It can, with \texttt{persist: true} and \texttt{cache: false} (option E), and the first version of this entry claimed a limitation that does not exist.
The reason is what that expression costs against what it buys.
\par
What it buys is a \texttt{dvc status} that mentions the CSV and reports a local copy that has changed.
The stage already has that check and gives it more context: it re-hashes the file on every run, not only after a download, and it fails with a message naming the two things that can have happened, a damaged local copy or a genuinely re-pinned upstream file, because those need different answers.
And because \texttt{download.\allowbreak{}md5} is a parameter, re-pinning is a change DVC reruns on, which is the dependency edge a graph entry would have drawn.
\par
What it costs is a 548.6 MB read on every \texttt{dvc status}, which is the cost that ruled out the \texttt{dep} shape in the first place, plus a lock entry with no recoverable object behind it.
So the property is already covered by the stage while the cost would be paid by everyone on every status call, and the file stays a cache.
\par
It lives under \texttt{data/\allowbreak{}external/\allowbreak{}} because that is the third-party slot of the project layout: a copy of a published artefact, reproducible from the DOI and the pinned MD5, with nothing for DVC to version. Keeping it out of \texttt{data/\allowbreak{}raw/\allowbreak{}} is what lets that directory belong to the pipeline alone.
\par
\textbf{Consequence, recorded rather than argued away.} After this change no \texttt{.\allowbreak{}dvc} file at the tip of the branch references the raw blob.
The blob is still on DagsHub and in local caches, and the pointer's hash is still in the Git history of the deleted file, so an earlier commit can still \texttt{dvc pull} and the snapshot is recoverable.
But recovering it without Zenodo now means somebody reading a hash out of the Git history, and a workspace-scoped \texttt{dvc gc} would delete the blob without asking.
That is why the rule never to run \texttt{dvc gc} without \texttt{-\allowbreak{}-\allowbreak{}all-\allowbreak{}commits} is written down in the data-versioning page as part of this change.},
  ai = {Information seeking, Alternative generation, Alternative assessment, Recommendation, Solution generation},
  response = {Accepted with modifications},
  assessment = {AI implemented issue \#33's option B and proposed the refinement that the CSV is a local cache rather than a stage output, with the measurements behind it: the CSV is 548.6 MB, the Parquet 215.3 MB, the stage 88 s cold and 66 s warm.
Its option table was wrong in one way that mattered: it argued the CSV could only be an \texttt{out} that is re-downloaded every run or a \texttt{dep} that is hashed on every \texttt{dvc status}, and presented the cache as the only way out, which made a DVC limitation out of a design trade-off.
An adversarial review of PR \#54, also run with AI, found \texttt{-\allowbreak{}-\allowbreak{}outs-\allowbreak{}persist-\allowbreak{}no-\allowbreak{}cache} in DVC 3.67.1 and established that a \texttt{persist: true, cache: false} out is neither deleted before the run nor pushed, and the same review found the recoverability consequence above.
Lukas kept the mechanism and required the rationale to be rewritten so that it rests on the cost of the 548.6 MB hash and on the check the stage already performs, rather than on a limitation that does not exist, and required the consequence to be recorded.},
  aievidence = {Claude Code session on 2026-09-30 implementing \href{https://github.com/mlops-2627q1-mds-upc/MLOPS_Recommenditos/issues/33}{issue \#33}, which produced the first option table in the body of \href{https://github.com/mlops-2627q1-mds-upc/MLOPS_Recommenditos/pull/54}{PR \#54}; a second Claude Code session the same day reviewing that PR adversarially, which produced the \texttt{persist} finding and the recoverability consequence; Lukas directed that the entry be corrected rather than the mechanism changed.},
  evidence = {\href{https://github.com/mlops-2627q1-mds-upc/MLOPS_Recommenditos/issues/33}{Issue \#33}; \href{https://github.com/mlops-2627q1-mds-upc/MLOPS_Recommenditos/pull/54}{PR \#54}; the \texttt{download} stage in \texttt{dvc.\allowbreak{}yaml} and the \texttt{download} block of \texttt{params.\allowbreak{}yaml}; \href{https://github.com/mlops-2627q1-mds-upc/MLOPS_Recommenditos/blob/main/docs/docs/data-versioning.md}{Data versioning}, ``Raw data''; \texttt{tests/\allowbreak{}test\_\allowbreak{}download.\allowbreak{}py}; EDN-25, EDN-34, EDN-36.},
}

% Generated by tools/edn_latex.py from reports/edn.md.
% Do not edit: change reports/edn.md and run `make edn`.
\ednentry{
  id = {EDN-36},
  title = {Retire EDN-25's \texorpdfstring{\texttt{dvc add}}{dvc add} pointer, so we no longer host a copy of the raw CSV},
  date = {2026-09-30},
  milestone = {M2: Reproducibility},
  topic = {Data Versioning, Personal Data},
  participants = {@lukas2510. Open to revisit at the next sprint review, like EDN-25 itself.},
  decision = {Amends EDN-25, which had tracked the raw CSV with \texttt{dvc add} and pushed it to our DagsHub remote.
The pointer \texttt{data/\allowbreak{}raw/\allowbreak{}autoscout24\_\allowbreak{}dataset\_\allowbreak{}20251108.\allowbreak{}csv.\allowbreak{}dvc} is deleted, the CSV is neither tracked nor pushed by us any more, and the \texttt{download} stage acquires it instead (EDN-35).
The blob the pointer referenced is left on the remote untouched, so checking out an earlier commit and running \texttt{dvc pull} still works.
What EDN-25 decided about \texttt{dvc import-\allowbreak{}url} is unchanged and was independently confirmed while implementing the stage.},
  rationale = {\emph{Alternatives considered:}
\begin{itemize}[leftmargin=*, topsep=2pt]
\item \textbf{Option A (chosen): delete the pointer with \texttt{dvc remove}, leave the blob on the remote.}
\newline Pros: only one mechanism owns \texttt{data/\allowbreak{}raw/\allowbreak{}} at a time, which is what the tracking-granularity rule requires and what DVC enforces anyway, since it refuses an output that overlaps a tracked path; every earlier commit stays reproducible, because the blob stays; and we stop re-publishing the CSV's PII columns going forward.
\newline Cons: nothing at the tip references the blob, so recovering that snapshot without Zenodo means reading the hash out of the Git history of the deleted pointer, and a workspace-scoped \texttt{dvc gc} would delete it silently.
\item \textbf{Option B: keep the pointer beside the stage as a record of the hash.}
\newline Pros: the snapshot's hash stays visible at the tip, and \texttt{dvc pull} keeps fetching the CSV for anyone who wants it.
\newline Cons: two owners of one path, which is what the granularity rule forbids and what DVC rejects outright once a stage declares anything under \texttt{data/\allowbreak{}raw/\allowbreak{}}; and we would go on re-publishing the PII columns for no gain over the pinned Zenodo DOI.
\item \textbf{Option C: delete the pointer and the blob (\texttt{dvc gc -\allowbreak{}-\allowbreak{}cloud}), so no copy of the CSV is left anywhere of ours.}
\newline Pros: the second copy of the personal data would actually be gone, which is what EDN-07 originally wanted.
\newline Cons: every earlier commit becomes unreproducible without Zenodo, and it removes an exposure that the pinned DOI provides publicly anyway. It trades real recoverability for no real privacy.
\end{itemize}
\par
\emph{Why:} EDN-25 chose \texttt{dvc add} because the course demo teaches it for a raw input that somebody acquires by hand.
Once the pipeline acquires the file, that premise is gone: the demo's own pattern is that a stage's outputs are tracked by the pipeline rather than by \texttt{dvc add}, and the data-versioning page had already said the manual pointer would be replaced by the \texttt{download} stage once the pipeline existed.
Keeping the pointer was rejected because two owners of one path is exactly what the granularity rule exists to prevent, and deleting the blob was rejected because recoverability is worth more than removing a copy of data that is public at its source.
\par
What EDN-25 got right about the mechanism survives untouched. The Zenodo URL answers a HEAD request with neither an \texttt{ETag} nor a \texttt{Content-\allowbreak{}MD5} header, and \texttt{dvc import-\allowbreak{}url} derives its change detection from one of those, so that mechanism was never verifiable against this source. Implementing the stage measured a second instance of the same gap: Zenodo serves this file chunked and sends no \texttt{Content-\allowbreak{}Length} either.
\par
The personal-data consequence changes shape rather than going away. We stop hosting the CSV, and we host the Parquet derived from it with the same columns, which is what EDN-34 weighed; that takes effect when \texttt{download.\allowbreak{}source} flips to \texttt{zenodo} and the lock is refreshed (issue \#57).},
  ai = {Information seeking, Alternative assessment, Solution generation},
  response = {Accepted},
  assessment = {AI retired the pointer as the implementation of issue \#33's option B and, rather than assuming the consequence, checked the remote: the blob is still there (\texttt{files/\allowbreak{}md5/\allowbreak{}b2/\allowbreak{}3a12.\allowbreak{}.\allowbreak{}.\allowbreak{}}, HTTP 200, \texttt{content-\allowbreak{}length: 548610318}), so an earlier commit can still \texttt{dvc pull}, and nothing was deleted from the remote.
An adversarial review of the same PR found what that leaves behind: no pointer at the tip references the blob, \texttt{dvc gc} and \texttt{dvc gc -\allowbreak{}-\allowbreak{}cloud} both default to workspace scope and would delete it, and \texttt{dvc gc} appeared nowhere in the repository, so nobody was warned.
Lukas accepted the mechanism and required the rule about \texttt{dvc gc -\allowbreak{}-\allowbreak{}all-\allowbreak{}commits} to be documented as part of the same change.},
  aievidence = {Claude Code session on 2026-09-30 implementing \href{https://github.com/mlops-2627q1-mds-upc/MLOPS_Recommenditos/issues/33}{issue \#33} and verifying the remote by authenticated request; a second Claude Code session the same day reviewing \href{https://github.com/mlops-2627q1-mds-upc/MLOPS_Recommenditos/pull/54}{PR \#54}, which reported the \texttt{dvc gc} exposure.},
  evidence = {\href{https://github.com/mlops-2627q1-mds-upc/MLOPS_Recommenditos/issues/33}{Issue \#33}; \href{https://github.com/mlops-2627q1-mds-upc/MLOPS_Recommenditos/pull/54}{PR \#54}; the deleted pointer \texttt{data/\allowbreak{}raw/\allowbreak{}autoscout24\_\allowbreak{}dataset\_\allowbreak{}20251108.\allowbreak{}csv.\allowbreak{}dvc}; \href{https://github.com/mlops-2627q1-mds-upc/MLOPS_Recommenditos/blob/main/docs/docs/data-versioning.md}{Data versioning}, ``Raw data'' and ``Never run \texttt{dvc gc} without \texttt{-\allowbreak{}-\allowbreak{}all-\allowbreak{}commits}''; EDN-07, EDN-20, EDN-25, EDN-34, EDN-35.},
}

% Generated by tools/edn_latex.py from reports/edn.md.
% Do not edit: change reports/edn.md and run `make edn`.
\ednentry{
  id = {EDN-37},
  title = {The seller group key stays an unsalted hash, and the residual risk is disclosed},
  date = {2026-09-30},
  milestone = {M3: Quality Assurance},
  topic = {Data Validation and Preprocessing, Privacy},
  participants = {@lukas2510, deciding while reviewing PR \#56.},
  decision = {\texttt{seller\_\allowbreak{}group\_\allowbreak{}id} stays what it is: an unsalted SHA-256 of the seller's company name, or of \texttt{country\_\allowbreak{}code|zip|city} for a private seller, truncated to 16 hex characters.
No pepper and no surrogate id.
What changes is what we claim about it: it is the split's grouping key, it keeps the seller's name, street, zip and city out of every column of every artefact, and it is \textbf{not} an anonymisation measure.
The code, the interim contract and the dataset card say so, and no document describes the hash as one-way any more.},
  rationale = {\emph{Alternatives considered:}
\begin{itemize}[leftmargin=*, topsep=2pt]
\item \textbf{Option A (chosen): keep the hash unsalted and disclose what it does not protect.}
\newline Pros: the split stays reproducible from a clean clone with no secret, which is what \texttt{pytest} and \texttt{dvc repro} depend on today; nothing to build; the claim we make becomes one that can be verified, namely that no PII column reaches an artefact.
\newline Cons: an attacker who holds the published file can map every id back to the name or the location it came from, so for a private seller \texttt{zip} and \texttt{city} are recoverable from an artefact that is supposed to be free of them; we carry that statement in the dataset card rather than being able to say the ids are opaque.
\item \textbf{Option B: key the hash with a secret pepper from the gitignored \texttt{.\allowbreak{}env}.}
\newline Pros: the dictionary attack stops working; the ids become opaque to anyone without the secret, which is what the old wording already claimed.
\newline Cons: the split's group assignment then depends on a secret that is not in the repository, so \texttt{dvc repro} on a clean clone produces a different split, \texttt{dvc.\allowbreak{}lock} stops matching for anyone who lacks the pepper, and every test that asserts a grouped split has to be given one. It also protects only against an attacker who does not already have the source file, which is a public download.
\item \textbf{Option C: drop \texttt{seller\_\allowbreak{}group\_\allowbreak{}id} from the published artefacts and keep the grouping internal to \texttt{split}.}
\newline Pros: the id would not leave the pipeline, so nothing published could be inverted at all.
\newline Cons: \texttt{split} needs the key from \texttt{preprocess} because only \texttt{preprocess} sees \texttt{seller\_\allowbreak{}company\_\allowbreak{}name}, so the column has to cross an artefact boundary; and the key is the evidence for the property the split exists to have, so a test, a review or a drift analysis could no longer check that no seller appears in two sets.
\end{itemize}
\par
\emph{Why:} The measurement decided it.
A dictionary built from the file's own \texttt{seller\_\allowbreak{}company\_\allowbreak{}name}, \texttt{country\_\allowbreak{}code}, \texttt{zip} and \texttt{city} columns inverts 17,141 of 17,141 dealer ids and 13,576 of 13,576 private-seller ids in about 50 ms, so the truncated digest offers no protection worth the name: a dictionary attack does not care how short the output is, and the private-seller key space is small enough to enumerate even without the file.
That is what makes option B's cost the deciding factor rather than its benefit.
A pepper would buy protection only against someone who does not hold a file that anyone can download from the pinned Zenodo DOI, and it would pay for it with the credential-free reproducibility the fixture, the test suite and \texttt{dvc repro} on a clean clone are built on (NFR-06).
The dealer behind a listing is in any case recoverable from the public file by joining on \texttt{make}, \texttt{model}, \texttt{price}, \texttt{mileage\_\allowbreak{}km\_\allowbreak{}raw} and \texttt{registration\_\allowbreak{}date}, all of which our processed data publishes, so the group key is not the weakest link.
\par
The \texttt{zip} and \texttt{city} recovery is accepted on the same ground, and named rather than argued away.
They are two of the seven columns \texttt{params.\allowbreak{}preprocess.\allowbreak{}pii\_\allowbreak{}columns} removes, and for a private seller the group key is exactly their hash, so an artefact that has no location column still carries a value that a public file turns back into one.
We accept it because the same public file already holds those columns next to the listing itself, and because our own remote holds the raw file too (EDN-20, EDN-25), so the group key adds no exposure that is not already there.
It is written down so that the moment any of those premises changes -- a non-public source, a private remote, a dataset we scrape ourselves -- this entry is the place that says the hash was never the control.
\par
NFR-08 keeps its teeth either way, because it governs columns: no PII column appears in the processed data, the prediction log, the comparables or the model artefacts, and \texttt{Schema.\allowbreak{}conform} enforces that structurally. What this entry corrects is the stronger claim the code and the contract had grown around it.},
  ai = {Information seeking, Alternative generation, Alternative assessment, Recommendation},
  response = {Accepted},
  assessment = {The claim had been escalating rather than weakening: \texttt{preprocess.\allowbreak{}py} said ``the hash is one-way'', \texttt{schema.\allowbreak{}py} said ``never reversible to the name'', and the description of PR \#56 called the rule airtight, none of which anybody had tested. An adversarial review of that PR built the dictionary from the published file and reported the two counts above, together with the timing and the observation that the truncation is irrelevant to a dictionary attack. It laid out the three options with the reproducibility cost of a pepper and the artefact-boundary cost of dropping the column, and recommended keeping the hash and correcting the claims. Lukas accepted that, and the reasoning that settled it is his: the split must stay reproducible without secrets, and the dealer identity is already recoverable by joining on the columns we publish, so the hash is a grouping key and nothing more.},
  aievidence = {Claude Code session on 2026-09-30, reviewing PR \#56 against the real snapshot: AI was asked to attack the stage's own claims, measured the inversion of all 30,717 group ids, reported that the one-wayness claim was false, and presented keeping the hash, peppering it and dropping the column as the three options; Lukas decided to keep the unsalted hash and to state the residual risk instead.},
  evidence = {\href{https://github.com/mlops-2627q1-mds-upc/MLOPS_Recommenditos/blob/main/recommenditos/data/preprocess.py}{\texttt{recommenditos/\allowbreak{}data/\allowbreak{}preprocess.\allowbreak{}py}} (\texttt{hash\_\allowbreak{}seller\_\allowbreak{}group}); \href{https://github.com/mlops-2627q1-mds-upc/MLOPS_Recommenditos/blob/main/recommenditos/schema.py}{\texttt{recommenditos/\allowbreak{}schema.\allowbreak{}py}} (\texttt{seller\_\allowbreak{}group\_\allowbreak{}id}); \href{https://github.com/mlops-2627q1-mds-upc/MLOPS_Recommenditos/blob/main/docs/docs/dataset-card.md}{dataset card}; \hyperref[EDN-20]{EDN-20}; \hyperref[EDN-25]{EDN-25}; \href{https://github.com/mlops-2627q1-mds-upc/MLOPS_Recommenditos/issues/34}{issue \#34}; \href{https://github.com/mlops-2627q1-mds-upc/MLOPS_Recommenditos/pull/56}{PR \#56}.},
}

% Generated by tools/edn_latex.py from reports/edn.md.
% Do not edit: change reports/edn.md and run `make edn`.
\ednentry{
  id = {EDN-39},
  title = {The split's size gate is four rules with an unconditional floor, not one bound derived from the realised data},
  date = {2026-09-30},
  milestone = {M2: Reproducibility (shapes M3: Quality Assurance)},
  topic = {Evaluation Protocol, Testing Strategy},
  participants = {@lukas2510},
  decision = {The \texttt{split} stage judges its realised sizes against four named rules: every set non-empty; every set at least 75\,\% of its configured share; no seller group above 5\,\% of the pool; and the share's scatter inside a bound derived from the group sizes at 5.0 sigma, capped at 15 points. The failure names every rule that fired. \texttt{DISPERSION\_\allowbreak{}SIGMAS} is chosen from a stated false-alarm budget -- at most 1 seed in 1,000 may turn the stage red on a pool that satisfies the concentration limit -- measured over 2,000 seeds on three pools in \href{https://github.com/mlops-2627q1-mds-upc/MLOPS_Recommenditos/blob/main/reports/analysis/split_gate.py}{\texttt{reports/\allowbreak{}analysis/\allowbreak{}split\_\allowbreak{}gate.\allowbreak{}py}}.},
  rationale = {\emph{Alternatives considered:}
\begin{itemize}[leftmargin=*, topsep=2pt]
\item \textbf{Option A: one bound derived from the realised group sizes, at four sigma.} What the first implementation of the stage did.
\newline Pros: the bound is a property of the data rather than a number someone liked, it tightens as the pool gains sellers, and the derivation is algebraically correct for independent group assignment.
\newline Cons: measured unable to detect the failure it exists for. The bound takes \texttt{sum(n\_\allowbreak{}i\textasciicircum{}2)} from the \emph{realised} sizes, so it widens exactly when one dealer dominates and the split is least trustworthy. One group holding 12\,\% of a 10,000-row pool, dumped into validation, realises 20.8\,\% against a 10\,\% ratio and passes, because the bound has become 14.4 points. A call with three empty sets passes. A \texttt{group\_\allowbreak{}key} naming \texttt{make} instead of the seller realises train at 38.7\,\% and passes against a 89.9-point bound, and \texttt{country\_\allowbreak{}code} at 70.8\,\% against 101 points. Four sigma was also wrong in the other direction: the comment claimed 200 seeds reach at worst 3.5 sigma, but seeds 1 to 400 reach 3.92 and 2,000 seeds reach 4.25, so a routine \texttt{dvc exp} seed sweep would have turned the stage red on a sound split.
\item \textbf{Option B: a fixed percentage bound.} Rejected in the first implementation and still rejected.
\newline Pros: nothing about it can be widened by the data.
\newline Cons: to hold on the 1,865-row synthetic fixture pool it would have to allow about 6.8 points, and that width accepts a calibration set at a third of its intended size on the real snapshot. One number cannot serve a 99,326-row pool with 28,435 groups and a 1,865-row one with 614.
\item \textbf{Option C (chosen): keep the derived bound as a scatter check, cap it, and add an unconditional floor plus a declared concentration limit.}
\newline Pros: each rule catches what it is for, and the rule that catches an undersized set cannot be widened by any property of the realised data, because it is a fraction of the \emph{configured} ratio. Concentration now fires the gate instead of widening it, which is what makes a mis-set \texttt{group\_\allowbreak{}key} detectable at all. The cap stops the derived quantity from becoming vacuous. The failure message names the rule, so the reader is sent to the right cause. All three exploits above now fail, and the real snapshot's own split still passes with the sizes the dataset card states.
\newline Cons: four numbers to justify instead of one, and the floor is a statement about a pool of the real snapshot's granularity -- on the much smaller fixture pool the shares scatter far more, so the seed-sweep test there sweeps a looser floor and says why. Three of the four rules are module constants rather than \texttt{params.\allowbreak{}yaml} entries, so they are not swept by \texttt{dvc exp}.
\item \textbf{Option D: derive the bound from a declared maximum group share instead of the realised sizes.}
\newline Pros: would answer the objection directly, with one rule instead of three.
\newline Cons: measured not to work. The available closed form, \texttt{sum(n\_\allowbreak{}i\textasciicircum{}2) <=\allowbreak{} max(n\_\allowbreak{}i) * N}, assumes every group is the maximum size and gives a 26.8-point bound at a 5\,\% declared share, which is looser than what it replaces. Modelling the tail as equal-sized groups instead understates the real dispersion (0.0346 against a realised 0.0379) and would fail sound splits at seed 206. The declared expectation is therefore used as its own rule, where it is exact, rather than as an input to a bound it cannot model.
\end{itemize}
\par
\emph{Why:} The defect was not the formula, which is correct for what it describes, but what it was allowed to do: a single bound that both licensed and was derived from the concentration it should have flagged. Separating the two is what makes each rule checkable. The floor answers ``is this set the size we asked for'', the concentration limit answers ``is the assumption behind the scatter bound still true'', and the scatter bound answers ``is this more movement than independent assignment explains''. Only the third is a statement about noise, and only the third should be derived from the data.
\par
The constants come from measurement rather than from taste, because the first version showed what taste produces: 4.0 was justified by a sweep that stopped at 200 seeds and was wrong by 2,000. The budget is stated so the number can be re-checked: at most 1 in 1,000 seeds may false-alarm. Over 6,000 seed-pool trials the worst case is 4.25 sigma, so 5.0 gives a measured false-alarm rate of 0 with 0.75 sigma of margin. The floor at 0.75 sits below the smallest share any of 2,000 seeds produced on the real snapshot, 0.820, and rejects the calibration set at 0.58 of its intended size that option A accepted. The concentration limit at 5\,\% sits above the real snapshot's largest dealer, 3.41\,\%, and far below the 12\,\%, 34.6\,\% and 100\,\% of the three exploits.
\par
One thing is deliberately written into the comment rather than left implied: ``five sigma'' is not a normal-theory quantile here and must not be read as one. The largest dealer alone accounts for 80.7\,\% of the variance of any set's share, so the distribution is dominated by a single Bernoulli and a Gaussian tail probability would be meaningless. The multiple is an empirical quantile of a measured sweep, and the comment says so.},
  ai = {Alternative generation, Alternative assessment, Solution generation, Recommendation},
  response = {Accepted with modifications},
  assessment = {An agent asked to attack the stage rather than describe it found that the gate the pull request presented as its main safety mechanism could not detect an undersized set, and confirmed each case by execution rather than by argument: the dominant-group dump, the three empty sets, and both wrong \texttt{group\_\allowbreak{}key} values, plus a 13-mutation scorecard showing that deleting the \texttt{check\_\allowbreak{}ratios} call in \texttt{main} survived the entire test suite. It also found the constant's justification wrong in both directions, which is the part that mattered most, because the stage would have failed a sound split at the next seed sweep.
\par
The modification is option D. The review asked for the bound to be derived from a declared expectation about concentration; that was tried and measured, and it is either looser than what it replaces or tight enough to fail sound splits, so the declared expectation became its own rule instead. Recording that is the point: the reviewer's prescription was not simply adopted, it was tested and the part that did not survive measurement was replaced by something that did.},
  aievidence = {Claude Code adversarial review of PR \#53 on 2026-09-30, with a 13-mutation scorecard and a reproduction of each exploit; the fixing session re-ran every exploit against the new gate and measured the 2,000-seed sweep that chose the constants.},
  evidence = {\href{https://github.com/mlops-2627q1-mds-upc/MLOPS_Recommenditos/blob/main/reports/analysis/split_gate.py}{\texttt{reports/\allowbreak{}analysis/\allowbreak{}split\_\allowbreak{}gate.\allowbreak{}py}} and its committed output; \href{https://github.com/mlops-2627q1-mds-upc/MLOPS_Recommenditos/blob/main/recommenditos/data/split_data.py}{\texttt{recommenditos/\allowbreak{}data/\allowbreak{}split\_\allowbreak{}data.\allowbreak{}py}}; \href{https://github.com/mlops-2627q1-mds-upc/MLOPS_Recommenditos/blob/main/tests/test_split.py}{\texttt{tests/\allowbreak{}test\_\allowbreak{}split.\allowbreak{}py}}; \hyperref[EDN-32]{EDN-32}; \href{https://github.com/mlops-2627q1-mds-upc/MLOPS_Recommenditos/issues/35}{issue \#35}; \href{https://github.com/mlops-2627q1-mds-upc/MLOPS_Recommenditos/pull/53}{PR \#53}.},
}

% Generated by tools/edn_latex.py from reports/edn.md.
% Do not edit: change reports/edn.md and run `make edn`.
\ednentry{
  id = {EDN-40},
  title = {The \texorpdfstring{\texttt{ES}}{ES} holdout stays row-selected, so a cross-border dealer is reported rather than moved},
  date = {2026-09-30},
  milestone = {M2: Reproducibility (shapes M6: Monitoring)},
  topic = {Evaluation Protocol, Monitoring},
  participants = {@lukas2510},
  decision = {The \texttt{split} stage selects the holdout per listing, not per seller, so ``every \texttt{ES} listing is held out'' (EDN-03) takes precedence over ``no seller appears in two sets'' (EDN-14) for a dealer that lists in \texttt{ES} and elsewhere. The stage counts any such seller and reports it as a warning with row counts, the disjointness tests include the holdout instead of iterating the split names only, and the count is pinned by a test.},
  rationale = {\emph{Alternatives considered:}
\begin{itemize}[leftmargin=*, topsep=2pt]
\item \textbf{Option A: hold out whole sellers that touch \texttt{ES}.} Both invariants would then hold without exception.
\newline Pros: the holdout would be seller-disjoint from the training sets, so an \texttt{ES} replay could not contain a dealer the model trained on, which is the property EDN-14's finding says matters.
\newline Cons: it changes what the holdout \emph{is}. EDN-03 defines it as the \texttt{ES} listings and prices its cost as exactly those rows, so the artefact would gain non-\texttt{ES} rows that M6 replays as Spanish traffic, and its size would no longer be the 6,079 rows EDN-03 and the dataset card state. On the real snapshot it would also buy nothing: 0 sellers list both inside and outside \texttt{ES}, because \texttt{hash\_\allowbreak{}seller\_\allowbreak{}group} keys a dealer by its company name and no name occurs on both sides.
\item \textbf{Option B (chosen): keep the row selection, detect the conflict and report it with counts.}
\newline Pros: EDN-03's definition of the holdout is untouched, and on the real snapshot the resolution costs nothing measurable. The violation becomes visible where it does occur rather than being silently excluded from the tests: the fixtures the tests run on have 91 such sellers at 2,000 raw rows and 898 at 20,000, and no test would have noticed, because both disjointness tests iterated the split names and skipped the holdout.
\newline Cons: a weaker invariant than option A. The stage does not guarantee seller-disjointness across all five artefacts, only across the four split sets, and a future snapshot with cross-border dealer names would need this decision revisited rather than being handled automatically.
\item \textbf{Option C: leave it as it was, selected per row and unmentioned.}
\newline Cons: this is what the pull request did, and it is the reason the conflict went unnoticed. Two documented invariants cannot both hold, and nothing said which one wins or what it costs.
\end{itemize}
\par
\emph{Why:} The two rules genuinely conflict and one of them has to lose, so the only bad answer is not saying which. EDN-03 is the authority on what the holdout contains, and its whole purpose is that the replayed rows are a market: rows from other countries in it would change what the drift measurement measures, and the 6,079-row figure is already cited in the dataset card and in EDN-14's count of the \texttt{ES} rows the API would accept. Against that, option A's benefit on the real snapshot is exactly zero, because no dealer name crosses the border there.
\par
What makes option B acceptable rather than a shrug is that the cost is now measured and visible. EDN-14's finding is that dealer-level shift is as large as country-level shift, so a replay containing dealers the model trained on understates drift and confounds both NFR-11's flagging and FR-15's retrain comparison. That is a real limitation of the synthetic fixture as a drift rehearsal, and it is now a reported number with a test pinning it, so if a future snapshot does grow cross-border dealer names the stage says so on every run instead of quietly violating the invariant.},
  ai = {Information seeking, Alternative generation, Alternative assessment, Recommendation},
  response = {Accepted},
  assessment = {The adversarial review of PR \#53 found the conflict, which the pull request never mentioned, and measured that it occurs 91 and 898 times in the two fixtures the tests run on while occurring 0 times on the real snapshot -- the combination that explains why it was invisible. It also identified the mechanism, that both disjointness tests iterate the split names and therefore exclude the holdout by construction. It offered both resolutions and accepted either, provided the invariant became visible. Lukas chose the row selection, because EDN-03 owns the definition of the holdout and option A would have changed a figure already published for no gain on the real data.},
  aievidence = {Claude Code adversarial review of PR \#53 on 2026-09-30, which measured the cross-border counts in both fixtures and on the real snapshot and named the conflict between EDN-03 and EDN-14.},
  evidence = {\href{https://github.com/mlops-2627q1-mds-upc/MLOPS_Recommenditos/blob/main/recommenditos/data/split_data.py}{\texttt{recommenditos/\allowbreak{}data/\allowbreak{}split\_\allowbreak{}data.\allowbreak{}py}} (\texttt{cross\_\allowbreak{}holdout\_\allowbreak{}sellers}); \href{https://github.com/mlops-2627q1-mds-upc/MLOPS_Recommenditos/blob/main/tests/test_split.py}{\texttt{tests/\allowbreak{}test\_\allowbreak{}split.\allowbreak{}py}}; \href{https://github.com/mlops-2627q1-mds-upc/MLOPS_Recommenditos/blob/main/reports/analysis/split_gate.py}{\texttt{reports/\allowbreak{}analysis/\allowbreak{}split\_\allowbreak{}gate.\allowbreak{}py}}; \hyperref[EDN-03]{EDN-03}; \hyperref[EDN-14]{EDN-14}; \href{https://github.com/mlops-2627q1-mds-upc/MLOPS_Recommenditos/issues/35}{issue \#35}; \href{https://github.com/mlops-2627q1-mds-upc/MLOPS_Recommenditos/pull/53}{PR \#53}.},
}

% Generated by tools/edn_latex.py from reports/edn.md.
% Do not edit: change reports/edn.md and run `make edn`.
\ednentry{
  id = {EDN-41},
  title = {\texorpdfstring{\texttt{model\_\allowbreak{}version}}{model\_version} becomes the leading token of the normalised trim plus a frequency floor},
  date = {2026-09-30},
  milestone = {M2: Reproducibility (feature engineering; feeds M3: Quality Assurance)},
  topic = {Feature Encoding},
  participants = {@lukas2510},
  decision = {\texttt{model\_\allowbreak{}version} enters the extended feature set as the first \texttt{features.\allowbreak{}model\_\allowbreak{}version\_\allowbreak{}tokens} tokens of its normalised form -- case and accents folded, separators collapsed -- keeping only the levels above \texttt{features.\allowbreak{}model\_\allowbreak{}version\_\allowbreak{}min\_\allowbreak{}frequency} of the training rows and treating everything else as missing. Both numbers stay parameters so the experiment ladder can sweep them. The two finer variants measured below are rejected for now, and the justification in the code says what the token actually is rather than what it was assumed to be.},
  rationale = {\emph{Alternatives considered:}
\begin{itemize}[leftmargin=*, topsep=2pt]
\item \textbf{Option A (chosen): normalise, keep the leading token, then a frequency floor.}
\newline Pros: 279 levels covering 85.9\,\% of the training rows, from a column with 80,332 distinct raw values, at the cost of one regular expression. Both knobs are parameters, so the ladder measures whether the coarseness costs accuracy rather than anyone arguing about it. The level set is an artefact the API loads, so a request is normalised exactly as a training row was.
\newline Cons: coarse, and unevenly so. The leading token names the body style for Audi and Porsche (\texttt{avant}, \texttt{coupe}) and the engine letter for the German premium models that are most of the data: \texttt{d}, \texttt{911} and \texttt{e} are the three most frequent levels and 12.4\,\% of the training rows lead with an engine letter, which largely repeats \texttt{fuel\_\allowbreak{}category} while discarding the trim a buyer cares about (\texttt{M Sport}, \texttt{Touring}). Inside one make and model the merging is heavy -- Porsche 992 level \texttt{911} covers 743 distinct raw trims, Audi A6 \texttt{avant} 473, BMW 320 \texttt{d} 331 -- and 329 pairs of levels survive where one is a prefix of the other, so \texttt{(\textquotesingle{}20d\textquotesingle{}, \textquotesingle{}320d\textquotesingle{})} are separate levels for the same car. 14.1\,\% of the rows have no level at all.
\item \textbf{Option B: prefix the token with make and model, so a level is \texttt{BMW|320|d}.}
\newline Pros: reads as the obvious fix for the prefix-pair problem, and it does reduce those from 329 pairs to 90.
\newline Cons: measured, it does not fix the merging it was meant to fix and it costs a lot. The worst merge inside one make and model stays at exactly 743, because Porsche 992 level \texttt{911} is already inside one make and model, and the number of groups holding over 50 raw trims is unchanged at 203. Meanwhile each qualified level is rarer, so the same frequency floor drops far more of them: coverage falls from 85.9\,\% to 62.8\,\%, and the median merge rises from 3 to 47.
\item \textbf{Option C: keep two tokens when the first is an engine letter, one otherwise.}
\newline Pros: the only variant that touches the specific weakness. BMW 320's worst merge falls from 331 to 200, BMW 520's from 327 to 130, Mercedes C 220's from 74 to 35, and the groups over 50 raw trims fall from 203 to 179.
\newline Cons: not clearly better overall. It does nothing for the two largest merges (Porsche 992 \texttt{911} at 743, Audi A6 \texttt{avant} at 473, neither an engine letter), it costs 3 points of coverage (85.9\,\% to 82.9\,\%) and takes the level count from 279 to 318, and 1,797 training rows (2.9\,\%) lose their level entirely. It also hard-codes a per-make letter list into a rule that is otherwise a swept parameter, so it cannot be turned off from \texttt{params.\allowbreak{}yaml}.
\item \textbf{Option D: two tokens throughout, which is option A at \texttt{model\_\allowbreak{}version\_\allowbreak{}tokens: 2}.}
\newline Pros: the best merging of the four -- worst merge 358 rather than 743, only 64 groups over 50 trims -- and it needs no code at all, because it is a parameter change.
\newline Cons: coverage collapses to 44.0\,\%, so more than half the column is missing. Kept as the parameter value the ladder will try, not as the default.
\item \textbf{Option E: normalise only, with no token cut.}
\newline Pros: loses nothing.
\newline Cons: unusable as a level set. 78,738 distinct values survive normalisation over 105,405 listings, and a frequency floor on the full string keeps 7 levels covering 0.8\,\% of the rows.
\item \textbf{Option F: hash the trim into a fixed number of buckets.}
\newline Pros: bounded level count with no vocabulary to version.
\newline Cons: destroys the SHAP explainability FR-08 promises, because a bucket has no meaning to show a user.
\end{itemize}
\par
\emph{Why:} A is the cheapest thing that produces a usable level set, and the alternatives were measured rather than argued. B and C both looked better on the face of it and neither survived the measurement: B leaves the merging it targets untouched while costing 23 points of coverage, and C buys a real improvement on 12.4\,\% of the rows by making the other 87.6\,\% slightly worse and by hard-coding a rule the ladder cannot sweep. The decision that actually matters is that both knobs are parameters, so the question ``is the coarse trim good enough'' is answered by the experiment ladder with an MdAPE rather than by anyone's judgement here, and D is already the configuration it will compare against. What had to change regardless of the outcome is the justification: the code claimed the leading token ``is the one that names the trim'', which is true for Audi and Porsche body styles and false for the German premium engine codes that dominate the data, and a comment that claims a property the data does not have is worse than no comment.},
  ai = {Information seeking, Alternative generation, Alternative assessment, Recommendation},
  response = {Accepted with modifications},
  assessment = {The normalisation and the frequency floor were AI's own proposal in the \texttt{features} ticket, and the overstated justification was AI's wording. An adversarial review, also by AI, measured the collisions on the real column and showed the claim was false for most of the data; it proposed B and C as cheap improvements and asked for them to be implemented only if they measured better. They were then measured, and both were rejected on their own numbers -- which is the part worth recording, because the suggestion was plausible enough that it would have been accepted without the measurement. Lukas accepted keeping A with the corrected justification.},
  aievidence = {Claude Code session on 2026-09-30: an adversarial review of \href{https://github.com/mlops-2627q1-mds-upc/MLOPS_Recommenditos/pull/55}{PR \#55} reported that ``the docstring's claim that the leading token is the one that names the trim is true for Audi and Porsche body styles and false for the German premium engine codes, which are the bulk of the data'', and asked for the two variants to be evaluated and implemented ``only if it is clearly better on the real data and you can show the numbers''. The numbers are in \texttt{reports/\allowbreak{}analysis/\allowbreak{}extended\_\allowbreak{}features\_\allowbreak{}results.\allowbreak{}txt}.},
  evidence = {\href{https://github.com/mlops-2627q1-mds-upc/MLOPS_Recommenditos/blob/main/reports/analysis/extended_features.py}{\texttt{reports/\allowbreak{}analysis/\allowbreak{}extended\_\allowbreak{}features.\allowbreak{}py}} and its results file (2026-09-30); \href{https://github.com/mlops-2627q1-mds-upc/MLOPS_Recommenditos/blob/main/docs/docs/model-card.md}{model card} ``Feature space''; \href{https://github.com/mlops-2627q1-mds-upc/MLOPS_Recommenditos/blob/main/docs/docs/problem-spec.md}{problem specification} section 4; \hyperref[EDN-15]{EDN-15}; \href{https://github.com/mlops-2627q1-mds-upc/MLOPS_Recommenditos/issues/36}{issue \#36}; \href{https://github.com/mlops-2627q1-mds-upc/MLOPS_Recommenditos/pull/55}{PR \#55}.},
}

% Generated by tools/edn_latex.py from reports/edn.md.
% Do not edit: change reports/edn.md and run `make edn`.
\ednentry{
  id = {EDN-42},
  title = {A categorical value the training rows never saw becomes missing, and the holdout's country is allowed to vanish},
  date = {2026-09-30},
  milestone = {M2: Reproducibility (feature encoding; the holdout feeds M6: Monitoring)},
  topic = {Feature Encoding, Monitoring},
  participants = {@lukas2510},
  decision = {Every categorical feature leaves the \texttt{features} stage as a \texttt{category} over the levels the training split holds, for every split and for a request alike. A value outside those levels becomes missing rather than a level of its own. The accepted consequence is that the \texttt{ES} holdout's \texttt{country\_\allowbreak{}code} is missing in \textbf{every row} of its feature matrix, because \texttt{ES} is held out by construction and is therefore not a training level. This is documented rather than worked around: the stage warns, naming any column a written split never fills, a test pins the holdout case as intended, and the problem specification states it.},
  rationale = {\emph{Alternatives considered:}
\begin{itemize}[leftmargin=*, topsep=2pt]
\item \textbf{Option A (chosen): training levels only, an unseen value is missing, and the holdout's empty column is documented.}
\newline Pros: a code means the same car in training, in test, in the holdout and at request time, which is the only way a fitted model can be applied to anything but its training frame. It is exactly what \hyperref[EDN-18]{EDN-18} already decided for the serving path -- an unseen country is accepted and passed to the model as unknown -- so the pipeline and the API agree without a second rule. The \texttt{ES} matrix being country-blind is the new-market scenario stated honestly: the model has never seen that market, which is the whole point of \hyperref[EDN-03]{EDN-03}. Costs nothing to implement.
\newline Cons: a column that is empty in every row looks exactly like a defect, and the holdout is not replayed until M6, months after anyone remembers why. Mitigated by the warning, the test and the specification sentence rather than by changing the encoding.
\item \textbf{Option B: give the holdout country a level of its own.}
\newline Pros: the column is no longer empty, so nothing looks broken.
\newline Cons: hands the model a code it was never fitted on, which is worse than missing: a tree splits on it with no training example behind the branch. It also renumbers every other level relative to the fitted model unless the level is appended, and appending a level nothing was fitted on is a contract the artefact cannot honour. And it would make the drift scenario measure something other than a new market.
\item \textbf{Option C: drop \texttt{country\_\allowbreak{}code} from the feature sets, so the holdout matrix has no such column.}
\newline Pros: no empty column anywhere.
\newline Cons: throws away a feature the problem specification names for every other split, where it is filled in 99.98\,\% of rows, in order to tidy up one holdout frame. Country is a real price driver across seven markets.
\item \textbf{Option D: let each frame decide its own levels.}
\newline Pros: no frame ever has an empty categorical.
\newline Cons: the failure this whole mechanism exists to prevent. Every split would get different codes, so a model fitted on \texttt{train} could not be applied to \texttt{test} at all, and a request could not be scored. It is also a leak: validation, test and the holdout would be deciding the feature space.
\end{itemize}
\par
\emph{Why:} The encoding rule was effectively settled by EDN-02 (native categoricals) and EDN-18 (accept an unseen value, do not reject it); what was open is whether the holdout's fully empty \texttt{country\_\allowbreak{}code} is a defect to fix or a consequence to document. It is a consequence, and the two ways of ``fixing'' it are both worse: option B feeds the model an unfitted code and option C deletes a real feature. The actual defect was that nothing said so -- the reasoning existed only in a pull-request description, which does not survive a squash merge -- so the fix is the warning, the test and one sentence in the specification. Whether the model should instead be fitted and gated only on the makes and markets the API serves is a separate question that is deliberately not settled here.},
  ai = {Alternative assessment, Recommendation},
  response = {Accepted},
  assessment = {An adversarial review found the empty column by running the stage on the real snapshot and checking every column of every split, and established that \texttt{country\_\allowbreak{}code} is the only column fully null in the holdout and not also fully null in training. It read the situation as faithful to the drift scenario rather than as a bug, and recommended documenting it in three places rather than changing the encoding, explicitly declining to drop the column or to add a level because those are modelling changes nobody had asked for. That framing is what made this an entry rather than a patch.},
  aievidence = {Claude Code session on 2026-09-30, adversarial review of \href{https://github.com/mlops-2627q1-mds-upc/MLOPS_Recommenditos/pull/55}{PR \#55}: ``this is faithful to the drift scenario rather than a defect, because \texttt{ES} is held out by construction, so the model genuinely has never seen that market, and EDN-18 already requires an unseen value to be accepted rather than rejected. What is unacceptable is that nothing says so.''},
  evidence = {\href{https://github.com/mlops-2627q1-mds-upc/MLOPS_Recommenditos/blob/main/recommenditos/data/build_features.py}{\texttt{recommenditos/\allowbreak{}data/\allowbreak{}build\_\allowbreak{}features.\allowbreak{}py}} (\texttt{\_\allowbreak{}warn\_\allowbreak{}about\_\allowbreak{}unobserved\_\allowbreak{}columns}); \href{https://github.com/mlops-2627q1-mds-upc/MLOPS_Recommenditos/blob/main/tests/test_features.py}{\texttt{tests/\allowbreak{}test\_\allowbreak{}features.\allowbreak{}py}} (\texttt{test\_\allowbreak{}the\_\allowbreak{}holdout\_\allowbreak{}country\_\allowbreak{}is\_\allowbreak{}missing\_\allowbreak{}throughout\_\allowbreak{}its\_\allowbreak{}matrix\_\allowbreak{}on\_\allowbreak{}purpose}); \href{https://github.com/mlops-2627q1-mds-upc/MLOPS_Recommenditos/blob/main/docs/docs/problem-spec.md}{problem specification} section 2; \hyperref[EDN-03]{EDN-03}; \hyperref[EDN-18]{EDN-18}; \href{https://github.com/mlops-2627q1-mds-upc/MLOPS_Recommenditos/pull/55}{PR \#55}.},
}

% Generated by tools/edn_latex.py from reports/edn.md.
% Do not edit: change reports/edn.md and run `make edn`.
\ednentry{
  id = {EDN-43},
  title = {The persisted contract carries a categorical's level list, and refuses to infer one},
  date = {2026-09-30},
  milestone = {M2: Reproducibility},
  topic = {Data Validation, Feature Encoding},
  participants = {@lukas2510},
  decision = {\texttt{Column} gains a \texttt{levels} field that a \texttt{category} cannot be declared without and that no other dtype may carry, so the artefact the \texttt{features} stage writes is lossless for the only data-dependent dtype it produces. \texttt{Schema.\allowbreak{}validate} compares the level list, not only \texttt{str(dtype)}. \texttt{Schema.\allowbreak{}conform} applies the declared levels instead of inferring them from the frame in front of it, and \texttt{build\_\allowbreak{}features.\allowbreak{}FeatureSpace} is the single public way for a consumer to load the contract and the vocabulary and cast a frame against them.},
  rationale = {\emph{Alternatives considered:}
\begin{itemize}[leftmargin=*, topsep=2pt]
\item \textbf{Option A (chosen): put the levels on \texttt{Column}, in the artefact, and make the cast apply them.}
\newline Pros: the contract becomes self-sufficient -- a consumer holding only the loaded \texttt{Schema} and a frame has a correct cast, which is what \texttt{train} (\#37) and \texttt{evaluate} (\#39) both need. The invariant is structural rather than remembered: \texttt{Column} refuses a \texttt{category} without levels, so no code path can produce a contract that says \texttt{\textquotesingle{}category\textquotesingle{}} and leaves the levels to be guessed. \texttt{validate} now catches a reordered level set, which changes every code in the column while the dtype string stays the same.
\newline Cons: the contract file grows by the level lists, including 279 trims, so the extended set's artefact is larger and noisier to read. Nullable columns keep no \texttt{levels} key, which keeps the noise to the 15 columns that need it.
\item \textbf{Option B: leave the levels in the vocabulary and make the private applier public.}
\newline Pros: no change to the contract at all, and the vocabulary already holds them.
\newline Cons: two objects have to be loaded and kept in step for a cast that is conceptually one contract's business, and \texttt{validate} still cannot check the levels because the schema still does not know them. It also leaves the trap in place: \texttt{conform} would keep silently inferring levels for anyone who used it.
\item \textbf{Option C: document that \texttt{.\allowbreak{}cat.\allowbreak{}codes} is unsafe and leave the mechanism alone.}
\newline Pros: no code change; the pipeline has no active bug today, because pyarrow preserves the dictionary order through Parquet and LightGBM remaps categories by value.
\newline Cons: relies on every future consumer reading the warning. \#37's \texttt{ridge} variant encodes the categoricals itself and \#39's masking sweep rebuilds frames per masked field, and \texttt{.\allowbreak{}cat.\allowbreak{}codes} is the obvious idiom for both, so the trap would be walked into twice within the milestone, each time producing a model scored on a different numbering with no error anywhere.
\end{itemize}
\par
\emph{Why:} The absence was reproduced three ways before it was fixed: \texttt{validate} accepted a frame with reversed levels, \texttt{conform} turned a 50-row subset of 25 levels into 5, and a one-row serving frame for \texttt{make=\allowbreak{}\textquotedbl{}BMW\textquotedbl{}} came back with code 0 where the training code is 3. None of it broke anything today, which is exactly why it was worth fixing now rather than after \#37 and \#39 had each written their own cast. Option B was the smaller change and was rejected because it leaves \texttt{conform} -- the function every stage already calls -- silently wrong; a contract that cannot check the only data-dependent part of itself is not the boundary this module claims to be.},
  ai = {Information seeking, Alternative assessment, Solution generation},
  response = {Accepted with modifications},
  assessment = {The lossy artefact was AI's own design in the \texttt{features} ticket, and an adversarial review by AI found it, reproduced all three consequences by execution and named the two tickets that would walk into it. The modification is the shape of the fix: the review suggested either putting the levels in the contract or having the schema take them from the vocabulary, and the first was chosen because it is the only one that also lets \texttt{validate} check them. The \texttt{Column} invariant and the single \texttt{FeatureSpace} entry point were added on top, so that the property holds by construction rather than by review.},
  aievidence = {Claude Code session on 2026-09-30, adversarial review of \href{https://github.com/mlops-2627q1-mds-upc/MLOPS_Recommenditos/pull/55}{PR \#55}: the review reported that ``\texttt{Schema.\allowbreak{}conform} INFERS a categorical's levels from the data'' and ``\texttt{Vocabulary} does hold the levels, but the only function that applies them is private, so a consumer holding the loaded \texttt{Schema} and a frame has no correct way to cast'', with a reproduction of each case.},
  evidence = {\href{https://github.com/mlops-2627q1-mds-upc/MLOPS_Recommenditos/blob/main/recommenditos/schema.py}{\texttt{recommenditos/\allowbreak{}schema.\allowbreak{}py}}; \href{https://github.com/mlops-2627q1-mds-upc/MLOPS_Recommenditos/blob/main/tests/test_schema.py}{\texttt{tests/\allowbreak{}test\_\allowbreak{}schema.\allowbreak{}py}}; \hyperref[EDN-02]{EDN-02}; \hyperref[EDN-31]{EDN-31}; \href{https://github.com/mlops-2627q1-mds-upc/MLOPS_Recommenditos/issues/37}{issue \#37}; \href{https://github.com/mlops-2627q1-mds-upc/MLOPS_Recommenditos/issues/39}{issue \#39}; \href{https://github.com/mlops-2627q1-mds-upc/MLOPS_Recommenditos/pull/55}{PR \#55}.},
}

% Generated by tools/edn_latex.py from reports/edn.md.
% Do not edit: change reports/edn.md and run `make edn`.
\ednentry{
  id = {EDN-44},
  title = {The traceability gate is milestone-scoped, and its expected-coverage set is a validated YAML file},
  date = {2026-09-30},
  milestone = {M3: Quality Assurance},
  topic = {Testing Strategy, CI/CD},
  participants = {Lukas},
  decision = {NFR-07's requirement-to-test matrix is gated per milestone rather than over the whole table, and which requirement owes evidence when lives in \texttt{tools/\allowbreak{}expected\_\allowbreak{}coverage.\allowbreak{}yaml}, a schema-validated file the generator checks, rather than in a \texttt{Due} column of the specification's tables. Milestone names are validated against M1 to M6 and the enforced set is read off that known order, never off the order the file lists them in.},
  rationale = {\emph{Alternatives considered:}
\begin{itemize}[leftmargin=*, topsep=2pt]
\item \textbf{Option A (chosen): a small validated YAML file, gate enforced per milestone up to \texttt{current}.}
\newline Pros: the due date of a requirement is a parameter of a CI gate, not a statement about the system, so it does not belong in a requirement document; the generator can enforce a strict schema (every documented requirement booked exactly once, no unknown ID, no unknown or out-of-order milestone) where a table cell cannot be checked at all; bumping \texttt{current} is a one-line, dated, reviewable diff; the file carries the reasoning per milestone in comments.
\newline Cons: a second place to look; the gate is only as honest as the bookings, so a requirement can be deferred by moving one line, which is why the M3 membership and the verification route of every already-due requirement are pinned by a test.
\item \textbf{Option B: a \texttt{Due} column in the specification's tables.}
\newline Pros: everything in one document, visible next to the requirement it belongs to.
\newline Cons: mixes a CI parameter into a document that is supposed to describe the system; \texttt{current} would still have needed a home in prose; the specification's tables are already the widest thing in the docs; nothing can validate a table cell, so a typo or a missing cell would be invisible.
\item \textbf{Option C: no completeness gate at all, matrix published for a human to read.}
\newline Pros: never a red build nobody can fix; zero maintenance.
\newline Cons: leaves NFR-07's promise to a review that nothing reminds anyone to do, which is the state that let the matrix go unimplemented until this PR.
\item \textbf{Option D: enforce completeness over the whole table from the start.}
\newline Pros: no bookkeeping, no bookings to argue about.
\newline Cons: before M4 there is no API, so fourteen of the sixteen functional requirements cannot have a test; the build would be red from the first commit until M6, and a gate that is always red stops being read.
\end{itemize}
\par
\emph{Why:} A gate is only useful if someone acts on it, which rules out D, and only trustworthy if something enforces it, which rules out C. Between A and B, the deciding argument is that the generator can refuse a malformed expected-coverage set and cannot refuse a malformed table cell: the file makes the gate tighten by itself, because a requirement added to the documents fails the build until someone decides when its evidence is due. The decision rests on two specification amendments, both recorded in this PR: NFR-07 said CI ``does not enforce completeness'' and now describes this gate, and the ``Verified by'' legend claimed \textbf{[automated]} means a test already carries the marker, which was false for fourteen of sixteen functional entries and is now described as the route to the evidence with the due date pointing at \texttt{tools/\allowbreak{}expected\_\allowbreak{}coverage.\allowbreak{}yaml}. Neither amendment changes what the system must do; both make the documents describe what CI actually does. An adversarial review of the first implementation found three ways to drop a requirement out of the gate with a one-line edit, one of which -- appending an out-of-order milestone block, because the enforced set was positional in the file -- was not a reviewable statement at all; validating the names against M1 to M6 and reading the order off that closes it, and the other two are now pinned by a test so the argument for loosening the gate has to be written down.},
  ai = {Alternative generation, Alternative assessment, Recommendation, Solution generation},
  response = {Accepted with modifications},
  assessment = {AI laid out A to D with the trade-offs and recommended A, and implemented it. An adversarial AI review of that implementation then found, by execution rather than by reading, that the gate could be loosened three ways without a reviewable statement, and that the enforced set was positional in the YAML. That finding is the modification: the validated milestone order and the tests that pin the gate's parameters were not in the first design and are the reason the decision is defensible as written.},
  aievidence = {Claude Code sessions on 2026-09-30: the options were laid out and A implemented for issue \#40; a second, adversarial review session reproduced each loosening route by running the generator against a modified file and reported the exit codes, and the fixes were then verified the same way.},
  evidence = {\href{https://github.com/mlops-2627q1-mds-upc/MLOPS_Recommenditos/issues/40}{issue \#40}; \href{https://github.com/mlops-2627q1-mds-upc/MLOPS_Recommenditos/pull/51}{PR \#51}; \texttt{tools/\allowbreak{}expected\_\allowbreak{}coverage.\allowbreak{}yaml}; \texttt{DUE\_\allowbreak{}AT\_\allowbreak{}M3} in \texttt{tests/\allowbreak{}test\_\allowbreak{}requirement\_\allowbreak{}matrix.\allowbreak{}py}.},
}

% Generated by tools/edn_latex.py from reports/edn.md.
% Do not edit: change reports/edn.md and run `make edn`.
\ednentry{
  id = {EDN-45},
  title = {A \texorpdfstring{\texttt{req}}{req} marker is verification only where the specification says a test is the evidence},
  date = {2026-09-30},
  milestone = {M3: Quality Assurance},
  topic = {Testing Strategy},
  participants = {Lukas},
  decision = {The matrix reports four statuses, not two, and the specification's ``Verified by'' tag decides which one a \texttt{req} marker earns: a marker on an \textbf{[automated]} entry is \emph{verified by a test}, a marker on a \textbf{[manual]} entry is \emph{named by a test} and does not stand in for the drill its cell names. A marker on a \textbf{[manual]} entry whose cell names nothing does not satisfy the gate.},
  rationale = {\emph{Alternatives considered:}
\begin{itemize}[leftmargin=*, topsep=2pt]
\item \textbf{Option A (chosen): the specification's tag decides, and ``named by a test'' is a status of its own.}
\newline Pros: a fully checkable rule, no judgement in the tool; catches the real failure the review found, where NFR-06's determinism and MLflow provenance were reported as covered by a test whose whole body asserts that two files exist; closes the cheapest way to fake coverage, which was to tag an entry \textbf{[manual]}, leave its cell empty and put a marker on any test; makes the matrix state plainly that no status means a person agreed the evidence is enough.
\newline Cons: one more concept for a reader; a marker on a \textbf{[manual]} entry with named evidence still satisfies the gate, so the distinction changes the report and not the build in that case.
\item \textbf{Option B: keep one \texttt{covered} status, fix only the two wrong markers.}
\newline Pros: smallest diff; the two specific lies are gone.
\newline Cons: the mechanism that produced them stays, so the next marker on a \textbf{[manual]} entry reports as verified again; and the matrix keeps asserting something a tool cannot know.
\item \textbf{Option C: require the specification cell to name the test, and count a marker only when the named test exists.}
\newline Pros: the strongest claim of the three, and would also catch a marker on the wrong test.
\newline Cons: most cells name a category and not a test (``API test''), so it would need every cell rewritten before the API exists; ``does this cell name a test'' is a judgement a parser cannot make without a heuristic, and a heuristic in the middle of the traceability evidence is worse than a coarser rule that is exact.
\item \textbf{Option D: drop the automated/manual distinction from the gate and treat any marker as a claim needing human sign-off.}
\newline Pros: honest about the tool's limits.
\newline Cons: gives up the one thing the tool can check, and turns the gate back into a reminder.
\end{itemize}
\par
\emph{Why:} The matrix is what the report cites for NFR-07, so its worst failure mode is reporting a requirement as verified when nothing verified it, and that is exactly what happened: NFR-06 and NFR-01 were both green while their evidence was partly or wholly absent. A tool cannot read a test body and judge whether it verifies a criterion, so the honest move is to let the document that already says which evidence counts decide, and to say out loud in the generated table that a status is about the route being walked and not about a person agreeing it is enough. C would claim more than the documents can back today; B leaves the mechanism in place. The same reasoning is why the human review stays in NFR-07 rather than being replaced by this gate: NFR-01 still reads \emph{verified by a test} while two of its six success criteria have no thresholds, and only a person reading the table can notice that.},
  ai = {Alternative generation, Alternative assessment, Recommendation, Solution generation},
  response = {Accepted with modifications},
  assessment = {An adversarial AI review found the flaw and proved it by execution, naming the test whose body asserts that two files exist and the covering test whose own comment says four of the six criteria are still null. A second AI session generated A to D, recommended A over C on the grounds that C needs a heuristic, and implemented it. The modification is in the markers rather than the tool: AI's first instinct was to strip NFR-01's markers so the requirement reads as uncovered, which would have thrown away real evidence and turned the M3 gate red with nothing to fix; instead the marker moved onto the test that keeps SC-04 to SC-06 from ever counting as met, and the residual over-claim is stated in the pull request rather than hidden.},
  aievidence = {Claude Code adversarial review of PR \#51 on 2026-09-30, which reproduced the finding by running the generator and reading the covering tests; the follow-up session laid out A to D and verified each fix by re-running the reproduction.},
  evidence = {\href{https://github.com/mlops-2627q1-mds-upc/MLOPS_Recommenditos/pull/51}{PR \#51}; \texttt{Entry.\allowbreak{}status} and \texttt{Entry.\allowbreak{}has\_\allowbreak{}the\_\allowbreak{}promised\_\allowbreak{}evidence} in \texttt{tools/\allowbreak{}requirement\_\allowbreak{}matrix.\allowbreak{}py}; the status table in \texttt{CONTRIBUTING.\allowbreak{}md}; \href{https://github.com/mlops-2627q1-mds-upc/MLOPS_Recommenditos/issues/39}{issue \#39} for the SC-04 to SC-06 thresholds.},
}

% Generated by tools/edn_latex.py from reports/edn.md.
% Do not edit: change reports/edn.md and run `make edn`.
\ednentry{
  id = {EDN-46},
  title = {DagsHub credentials live in two gitignored stores, and the token is entered twice},
  date = {2026-09-30},
  milestone = {M2: Reproducibility},
  topic = {Credential handling, developer onboarding},
  participants = {Lukas},
  decision = {A contributor's DagsHub token is stored twice, in two gitignored files: in \texttt{.\allowbreak{}env} as \texttt{MLFLOW\_\allowbreak{}TRACKING\_\allowbreak{}PASSWORD}, which MLflow reads from the environment, and in \texttt{.\allowbreak{}dvc/\allowbreak{}config.\allowbreak{}local} via \texttt{dvc remote modify origin -\allowbreak{}-\allowbreak{}local password}, which is the only place DVC's HTTP remote reads a password from.
Nothing derives one store from the other, and \texttt{.\allowbreak{}env.\allowbreak{}template} deliberately does not list \texttt{DAGSHUB\_\allowbreak{}USERNAME} or \texttt{DAGSHUB\_\allowbreak{}USER\_\allowbreak{}TOKEN}.
\texttt{recommenditos/\allowbreak{}config.\allowbreak{}py} loads \texttt{<repo>/\allowbreak{}.\allowbreak{}env} by name rather than searching upwards, so the credentials of a checkout are its own.},
  rationale = {\emph{Alternatives considered:}
\begin{itemize}[leftmargin=*, topsep=2pt]
\item \textbf{Option A (chosen): two stores, the token pasted into each, and a template that lists only what the project reads.}
\newline Pros: no machinery to maintain; both stores are already gitignored, so no token can reach a commit; each tool reads its credentials the way it documents; the template cannot mislead, because every variable in it is read and a test asserts that.
\newline Cons: the token is entered twice, and the reader has to be told why or it looks like an oversight; the \texttt{dvc remote modify} form puts the token in shell history and in \texttt{ps}.
\item \textbf{Option B: one source of truth in \texttt{.\allowbreak{}env}, plus a script that writes \texttt{.\allowbreak{}dvc/\allowbreak{}config.\allowbreak{}local} from it.}
\newline Pros: one paste; one file to rotate.
\newline Cons: a piece of project machinery whose whole job is to write a secret to disk on its own initiative, which is a worse thing to own than a second paste; another step that can fail or drift; nothing in the course requires it.
\item \textbf{Option C: \texttt{.\allowbreak{}env} only, relying on an environment variable for DVC.}
\newline Pros: would be the simplest of all, if it existed.
\newline Cons: it does not. Measured rather than assumed: with \texttt{MLFLOW\_\allowbreak{}*}, \texttt{DAGSHUB\_\allowbreak{}USERNAME} and \texttt{DAGSHUB\_\allowbreak{}USER\_\allowbreak{}TOKEN} all exported, \texttt{dvc status -\allowbreak{}c} still fails with \texttt{configuration error -\allowbreak{} HTTP \textquotesingle{}basic\textquotesingle{} authentication require both \textquotesingle{}user\textquotesingle{} and \textquotesingle{}password\textquotesingle{}}, and DVC 3.67.1's HTTP remote schema offers exactly \texttt{user}, \texttt{password} and \texttt{ask\_\allowbreak{}password} and no environment route at all.
\item \textbf{Option D: \texttt{ask\_\allowbreak{}password true}, prompting instead of storing.}
\newline Pros: nothing on disk; nothing in shell history.
\newline Cons: \texttt{dvc\_\allowbreak{}http} calls \texttt{getpass}, which needs a terminal, so it prompts on every \texttt{dvc pull} and breaks any unattended \texttt{dvc repro} or CI job. Documented as the better choice on a shared machine, not as our default.
\end{itemize}
\par
\emph{Why:} Option C would have been the right answer and is not available, which is what makes the double entry a property of DVC rather than a choice.
Between A and B, the property that matters is that no token can reach a commit, and both stores already have it; B buys one fewer paste at the price of owning a secret-writing script.
Listing \texttt{DAGSHUB\_\allowbreak{}*} in the template was actively harmful: nothing reads those names, so filling them in configures nothing while reading as though it configured DVC, which is the confusion the getting-started page exists to prevent.
Scoping \texttt{.\allowbreak{}env} to the repository belongs to the same decision: a bare \texttt{load\_\allowbreak{}dotenv()} searches upwards, so a clone nested under another checkout inherited that one's token and looked configured when it was not.},
  ai = {Information seeking, Alternative assessment, Solution generation},
  response = {Accepted with modifications},
  assessment = {AI established by execution that DVC has no environment-variable route, which is the fact the whole decision rests on, and it wrote the two-store walkthrough. An adversarial review of that work, also by AI, then found three things the first pass had asserted rather than checked: that the \texttt{.\allowbreak{}env} scoping was not repository-local, so the review's own ``fresh clone'' verification had in fact been running on the parent checkout's credentials; that \texttt{mlflow.\allowbreak{}db} was neither gitignored nor prevented; and that the documented newcomer command printed a traceback where the page promised a message. All three were reproduced before being fixed. The modification is that the template lost the \texttt{DAGSHUB\_\allowbreak{}*} block, which the first pass had defended as documentation, once the review showed nothing reads it.},
  aievidence = {Claude Code sessions on 2026-09-30: the DVC environment-variable question was settled by running \texttt{dvc status -\allowbreak{}c} with and without the variables exported; the review findings were each reproduced before any fix, including a probe showing \texttt{find\_\allowbreak{}dotenv} resolving to a \texttt{.\allowbreak{}env} three directories above a credential-free worktree.},
  evidence = {\href{https://github.com/mlops-2627q1-mds-upc/MLOPS_Recommenditos/issues/42}{issue \#42}; \href{https://github.com/mlops-2627q1-mds-upc/MLOPS_Recommenditos/pull/50}{PR \#50}; \hyperref[EDN-19]{EDN-19}; NFR-15 in \href{https://github.com/mlops-2627q1-mds-upc/MLOPS_Recommenditos/blob/main/docs/docs/specification.md}{the specification}, which held this as part of NFR-09 when this entry was written (\hyperref[EDN-66]{EDN-66}).},
}

% Generated by tools/edn_latex.py from reports/edn.md.
% Do not edit: change reports/edn.md and run `make edn`.
\ednentry{
  id = {EDN-47},
  title = {The serving runtime is \texorpdfstring{\texttt{[project] dependencies}}{[project] dependencies}, and everything else is a PEP 735 dependency group},
  date = {2026-10-05},
  milestone = {M3: Quality Assurance (shapes M5: Model Packaging)},
  topic = {Containerization, Dependency management},
  participants = {@lukas2510},
  decision = {\texttt{pyproject.\allowbreak{}toml} no longer has one flat dependency list.
\texttt{[project] dependencies} holds the serving runtime only, meaning what it takes to load a bundle from \texttt{models/\allowbreak{}<variant>/\allowbreak{}} and price a listing: joblib, lightgbm, loguru, numpy, pandas, pyarrow, python-dotenv, pyyaml, scikit-learn and typer.
Everything else is in \texttt{[dependency-\allowbreak{}groups]}: \texttt{pipeline} (dvc, matplotlib, mlflow, requests, tqdm), \texttt{notebook} (ipython, jupyterlab, notebook), \texttt{docs} (mkdocs), \texttt{test} (coverage, pytest, pytest-cov) and \texttt{dev} (nbstripout, pip, pre-commit, ruff).
\texttt{[tool.\allowbreak{}uv] default-\allowbreak{}groups =\allowbreak{} \textquotedbl{}all\textquotedbl{}} keeps a plain \texttt{uv sync} and \texttt{uv run} installing every group, so no documented command changed.
The API image will install the runtime alone with \texttt{uv sync -\allowbreak{}-\allowbreak{}locked -\allowbreak{}-\allowbreak{}no-\allowbreak{}default-\allowbreak{}groups}.
FastAPI, uvicorn and pydantic join the runtime with the API in M4; Great Expectations, CodeCarbon and \texttt{shap} (EDN-11) go into groups.
\texttt{make test-\allowbreak{}serving}, run by CI's \texttt{Serving runtime} job, which is a required check, builds the runtime set plus the \texttt{test} group from the lock and runs \texttt{tests/\allowbreak{}test\_\allowbreak{}serving.\allowbreak{}py} from it, without \texttt{tests/\allowbreak{}conftest.\allowbreak{}py}.
A serving module that imports a group's package therefore fails the pull request instead of M5's image build.},
  rationale = {\emph{Alternatives considered:}
\begin{itemize}[leftmargin=*, topsep=2pt]
\item \textbf{Option A (chosen): PEP 735 dependency groups, with the runtime in \texttt{[project] dependencies} and \texttt{default-\allowbreak{}groups =\allowbreak{} \textquotedbl{}all\textquotedbl{}}.}
\newline Pros: one lock file and one resolution for everything, so the full set's versions do not move and the image installs a strict subset of exactly what CI tested; uv supports it natively (\texttt{-\allowbreak{}-\allowbreak{}group}, \texttt{-\allowbreak{}-\allowbreak{}only-\allowbreak{}group}, \texttt{-\allowbreak{}-\allowbreak{}no-\allowbreak{}default-\allowbreak{}groups}); with \texttt{default-\allowbreak{}groups =\allowbreak{} \textquotedbl{}all\textquotedbl{}}, a contributor's \texttt{uv sync}, CI's \texttt{uv sync -\allowbreak{}-\allowbreak{}locked} and every \texttt{uv run} behave as before; groups are not published as package metadata, which suits a course project that nobody installs from an index.
\newline Cons: the default is the opposite of what the image needs, so the image and the serving check have to remember \texttt{-\allowbreak{}-\allowbreak{}no-\allowbreak{}default-\allowbreak{}groups}, and a \texttt{uv run} without \texttt{-\allowbreak{}-\allowbreak{}no-\allowbreak{}sync} in that environment would quietly install every group into it.
Two tests in the check close that gap: one fails when the environment holds anything beyond the runtime and \texttt{test} closures read from \texttt{uv.\allowbreak{}lock}, so the check cannot pass vacuously; the other runs the serving path in a fresh interpreter and fails when it imported a module only the \texttt{test} group installs, because the check's environment has pytest and five packages it pulls in, which the image will not.
\item \textbf{Option B: optional extras (\texttt{[project.\allowbreak{}optional-\allowbreak{}dependencies]}).}
\newline Pros: older and more widely known than dependency groups, and every installer has supported them for years.
\newline Cons: \texttt{uv sync} leaves extras out unless it is given \texttt{-\allowbreak{}-\allowbreak{}all-\allowbreak{}extras} or \texttt{-\allowbreak{}-\allowbreak{}extra}, so every documented \texttt{uv sync} (getting-started, CONTRIBUTING, the README, the pipeline docs and every CI job) would either have to change or would silently install less than it does today.
Extras are also published package metadata, an install-time choice for the users of a library, and nobody installs this project as a library.
\item \textbf{Option C: keep one flat list and swap \texttt{mlflow} for \texttt{mlflow-\allowbreak{}skinny}.}
\newline Pros: a one-line change; MLflow is the dependency people usually suspect of bloat.
\newline Cons: measured while reviewing \#42, it does not do what it promises.
\texttt{mlflow-\allowbreak{}skinny} depends on fastapi, uvicorn, starlette and databricks-sdk as well, so none of them leaves.
What it does drop (alembic, docker, flask, flask-cors, graphene, gunicorn, waitress, sqlalchemy, huey, skops, aiohttp) is 39 MB of a 915 MB \texttt{.\allowbreak{}venv}, which is noise next to jupyterlab and the rest of the training stack.
Dropping sqlalchemy and alembic also breaks the two tests that exercise a real MLflow store offline, which are the only place the suite does so without credentials.
And the image would still contain the notebooks, the docs and the tooling.
\item \textbf{Option D: keep one flat list, and give the image its own hand-written requirements file.}
\newline Pros: no change to \texttt{pyproject.\allowbreak{}toml}.
\newline Cons: a second dependency list that nothing ties to \texttt{uv.\allowbreak{}lock}, so the image would either resolve its own versions, and then it is not running what CI tested, or need its versions copied by hand and kept in step forever.
A serving import of a training package would still first fail in the image build.
\item \textbf{Option E: a separate project or uv workspace member for the API.}
\newline Pros: the hardest boundary, with its own \texttt{pyproject.\allowbreak{}toml}.
\newline Cons: the API imports \texttt{recommenditos} itself (\texttt{model.\allowbreak{}py}, \texttt{build\_\allowbreak{}features}), so it would depend on the package anyway and the split would end up inside it again.
A workspace means two members to version and document, which is more structure than a five-person, one-semester project needs.
\end{itemize}
\par
\emph{Why:} The issue framed the lever correctly: what matters is which dependencies reach the image at all, not which MLflow package is installed.
Option C was ruled out by measurement before this work started.
Between A and B the deciding difference is the default.
Extras would have changed every documented \texttt{uv sync}, or quietly installed less with it, while groups with \texttt{default-\allowbreak{}groups =\allowbreak{} \textquotedbl{}all\textquotedbl{}} leave every existing command exactly as it was and only ask the image to opt out.
\par
The runtime set comes from evidence, not from a guess.
An import hook traced a real \texttt{load\_\allowbreak{}model} and \texttt{predict\_\allowbreak{}eur} on every committed bundle, in the full environment and in the runtime one (\texttt{reports/\allowbreak{}analysis/\allowbreak{}serving\_\allowbreak{}imports.\allowbreak{}py}).
\texttt{model.\allowbreak{}py} reaches \texttt{build\_\allowbreak{}features}, \texttt{split\_\allowbreak{}data}, \texttt{pipeline}, \texttt{schema} and \texttt{config}, and those import nine of the ten packages above themselves.
The tenth, pyarrow, is imported by pandas, whose \texttt{read\_\allowbreak{}parquet} needs it but which declares it only as an extra.
The other packages that got imported belong to those ten: their own requirements, such as narwhals, which both lightgbm and scikit-learn require, scipy and threadpoolctl.
The exceptions are optional imports that are skipped when the package is absent: tqdm from \texttt{config.\allowbreak{}py}, psutil from joblib and charset\_normalizer from numpy.
All three appear in the full environment's trace and are missing from the runtime one's, which still serves.
Nothing on the serving path imports MLflow, DVC or anything else from the training stack, so no refactor was needed.
Three small packages stay in the runtime on purpose.
typer is there because every stage module declares its command line when it is imported, and the serving seam imports two of them.
It costs 6.7 MB with rich and Pygments, and splitting \texttt{build\_\allowbreak{}features} to avoid it would break EDN-28's one module per stage and the \texttt{deps} that \texttt{dvc.\allowbreak{}yaml} names.
pyyaml and python-dotenv are there because \texttt{pipeline.\allowbreak{}py} and \texttt{config.\allowbreak{}py} import them, and together they are 3.0 MB.
tqdm went the other way, into \texttt{pipeline}: \texttt{config.\allowbreak{}py} sends loguru through \texttt{tqdm.\allowbreak{}write} with colour forced on whenever tqdm can be imported, which would put escape codes into every container log line.
Without it, the API logs to loguru's default stderr sink, which colours only a terminal.
\texttt{test} is a group of its own rather than part of \texttt{dev}, because the serving check installs it on top of the runtime, and every extra package there would be one a serving import could lean on unnoticed.
matplotlib is in \texttt{pipeline}, not \texttt{notebook}, because a DVC stage draws figures too: issue \#38's \texttt{compare-\allowbreak{}energy} stage writes \texttt{reports/\allowbreak{}figures/\allowbreak{}energy-\allowbreak{}vs-\allowbreak{}error.\allowbreak{}png}, and an install of the pipeline alone has to be able to run it.
\par
Measured on 2026-10-05 by building both environments from \texttt{uv.\allowbreak{}lock} (\texttt{reports/\allowbreak{}analysis/\allowbreak{}runtime\_\allowbreak{}footprint.\allowbreak{}py}): the runtime is \textbf{23 distributions, the project included, and 421 MB of site-packages as installed (519 MB with bytecode)}.
The full environment is \textbf{248 distributions and 815 MB (1,054 MB)}.
So the runtime is 52\,\% of the full environment as installed and 49\,\% with bytecode, and the full one, once compiled, would exceed NFR-04's 1 GB cap with its site-packages alone, before a base image or a model.
The runtime leaves 579 MB as installed, or 481 MB compiled, for the base image, FastAPI, the model and the comparables index.
The largest runtime packages are pyarrow (157 MB), scipy (114 MB), numpy (57 MB), pandas (39 MB), scikit-learn (31 MB) and lightgbm (10 MB).
pyarrow is therefore the first place M5 should look if the image runs tight, but it is not a free saving: \texttt{b0}'s bundle carries a Parquet lookup table, and pandas 3 backs its \texttt{str} dtype with pyarrow whenever it is installed, so dropping it changes how every string column of a request is stored.
The four committed bundles price their first five test rows identically in both environments, so the runtime set serves the same numbers rather than merely importing.
\par
The full set's resolution did not move: \texttt{uv export -\allowbreak{}-\allowbreak{}all-\allowbreak{}groups} lists the same 252 packages at the same versions before and after, the project itself aside, and the only lock change is the project's own entry.
The guarantee is made permanent rather than measured once.
Adding \texttt{import mlflow} to \texttt{model.\allowbreak{}py} made \texttt{make test-\allowbreak{}serving} fail with \texttt{ModuleNotFoundError: No module named \textquotesingle{}mlflow\textquotesingle{}}, while the full suite would still have passed.
Adding \texttt{import packaging.\allowbreak{}version} there failed it as well, naming \texttt{packaging} as a module only the \texttt{test} group installs.
A module-level \texttt{import codecarbon} in \texttt{tests/\allowbreak{}conftest.\allowbreak{}py}, which is what issue \#38's PR does, leaves it green, because the check never loads that file.},
  ai = {Information seeking, Alternative generation, Alternative assessment, Solution generation},
  response = {Accepted},
  assessment = {Lukas chose dependency groups over extras before the implementation started, and the \texttt{mlflow-\allowbreak{}skinny} measurement had already been taken while reviewing \#42.
What AI contributed is the evidence that the split is correct and stays correct.
It traced the serving import graph by running it rather than by reading the imports, which is what separated the three optional imports from real requirements.
It built both environments from the lock, priced the committed bundles in each, checked that the lock's resolution was unchanged, and placed every dependency with a stated reason.
That includes the three that did not obviously belong anywhere: typer kept in the runtime rather than splitting the stage modules, tqdm moved out of it because of the forced-colour log sink, and \texttt{test} separated from \texttt{dev}.
It then proved the CI check catches a training import by injecting one, and added the guard that fails the check when anything beyond the runtime and the test group is installed, so a lost \texttt{-\allowbreak{}-\allowbreak{}no-\allowbreak{}default-\allowbreak{}groups} cannot make the job pass vacuously.
An independent AI review of the pull request then found three ways the check could still be wrong, and reproduced each: the check loaded \texttt{tests/\allowbreak{}conftest.\allowbreak{}py}, so \#38's conftest import of CodeCarbon would have turned it red with no serving module at fault; the test group's own packages could have stood in for a missing runtime dependency; and a change to the \texttt{Makefile} alone did not trigger the job.
Each is fixed and covered by a break-it check; the review also caught that narwhals had been called optional when lightgbm and scikit-learn both require it.},
  aievidence = {Claude Code sessions on 2026-10-05 implementing issue \#61 and reviewing PR \#69: the import trace and the footprint run committed under \texttt{reports/\allowbreak{}analysis/\allowbreak{}}, the before/after \texttt{uv export -\allowbreak{}-\allowbreak{}all-\allowbreak{}groups} comparison, the injected imports that turned \texttt{make test-\allowbreak{}serving} red, and the review's reproduction of the conftest failure by merging \#68's head into this branch.},
  evidence = {\href{https://github.com/mlops-2627q1-mds-upc/MLOPS_Recommenditos/issues/61}{issue \#61}; \href{https://github.com/mlops-2627q1-mds-upc/MLOPS_Recommenditos/pull/69}{PR \#69}; \texttt{pyproject.\allowbreak{}toml}; \href{https://github.com/mlops-2627q1-mds-upc/MLOPS_Recommenditos/blob/main/reports/analysis/runtime_footprint.py}{\texttt{reports/\allowbreak{}analysis/\allowbreak{}runtime\_\allowbreak{}footprint.\allowbreak{}py}} and \href{https://github.com/mlops-2627q1-mds-upc/MLOPS_Recommenditos/blob/main/reports/analysis/runtime_footprint_results.txt}{its output}; \href{https://github.com/mlops-2627q1-mds-upc/MLOPS_Recommenditos/blob/main/reports/analysis/serving_imports.py}{\texttt{reports/\allowbreak{}analysis/\allowbreak{}serving\_\allowbreak{}imports.\allowbreak{}py}} and \href{https://github.com/mlops-2627q1-mds-upc/MLOPS_Recommenditos/blob/main/reports/analysis/serving_imports_results.txt}{its output}; \texttt{tests/\allowbreak{}test\_\allowbreak{}serving.\allowbreak{}py}; \texttt{make test-\allowbreak{}serving} and the \texttt{Serving runtime} job in \texttt{.\allowbreak{}github/\allowbreak{}workflows/\allowbreak{}ci.\allowbreak{}yml}; the Dependencies section of \href{https://github.com/mlops-2627q1-mds-upc/MLOPS_Recommenditos/blob/main/CONTRIBUTING.md}{CONTRIBUTING.md}; NFR-04 in \href{https://github.com/mlops-2627q1-mds-upc/MLOPS_Recommenditos/blob/main/docs/docs/specification.md}{the specification}; \hyperref[EDN-08]{EDN-08}, \hyperref[EDN-11]{EDN-11}, \hyperref[EDN-28]{EDN-28}.},
}

% Generated by tools/edn_latex.py from reports/edn.md.
% Do not edit: change reports/edn.md and run `make edn`.
\ednentry{
  id = {EDN-48},
  title = {\texorpdfstring{\texttt{split}}{split} records the supported-make list and stays a lossless partition; the downstream stages apply it},
  note = {\textbf{Amended by EDN-67} (2026-10-05): \texttt{features} is now the one stage that applies the list, to all five frames and before the vocabulary is built, and \texttt{train} and \texttt{evaluate} check it rather than filtering again.
All four guarantees listed below are discharged; \texttt{docs/\allowbreak{}docs/\allowbreak{}pipeline.\allowbreak{}md} records where under ``The supported makes''.},
  date = {2026-09-30},
  milestone = {M2: Reproducibility (shapes M3: Quality Assurance)},
  topic = {Evaluation Protocol, Data Preparation},
  participants = {@lukas2510, decided by the repository owner},
  decision = {The \texttt{split} stage computes the supported-make list (EDN-05) and writes it as \texttt{data/\allowbreak{}processed/\allowbreak{}supported\_\allowbreak{}makes.\allowbreak{}json}, but it does \textbf{not} remove the rows of unsupported makes: every interim row lands in exactly one of the five artefacts. The \texttt{features}, \texttt{train} and \texttt{evaluate} stages read that artefact and apply it, so the model is fitted and the SC-01 to SC-06 gate measured only on the makes the API serves. \texttt{dvc.\allowbreak{}yaml} gives \texttt{features} the artefact as a dependency, so changing \texttt{split.\allowbreak{}min\_\allowbreak{}listings\_\allowbreak{}per\_\allowbreak{}make} reruns the matrices, the models and the metrics.},
  rationale = {\emph{Alternatives considered:}
\begin{itemize}[leftmargin=*, topsep=2pt]
\item \textbf{Option A: filter in \texttt{split}, and write the removed rows to a sixth artefact.}
\newline Pros: the frames leaving \texttt{split} are already in scope, so no downstream stage can forget the filter, which is the failure this decision is most exposed to. The skeleton pointed this way: \texttt{schema.\allowbreak{}py} said ``the supported-make filter and the ES holdout belong to \texttt{split}'', \texttt{preprocess.\allowbreak{}py} said the filter ``is deliberately NOT here [...] which happens in \texttt{split}'', and every ``training scope'' figure in the docs is post-filter.
\newline Cons: it makes \texttt{split} lossy, so ``the split loses no row'' stops being a testable invariant and becomes a claim about two artefacts plus a sixth one. EDN-05's threshold could no longer be revisited without re-running the split, which is the irreversible direction: a downstream stage can always apply a list, but it cannot recover rows the split threw away. The sixth artefact also needs its own \texttt{dvc.\allowbreak{}yaml} output and its own contract, which is outside issue \#35.
\item \textbf{Option B: record the list, and measure the metrics over every make anyway.}
\newline Pros: nothing to wire; the split stays lossless and no stage has to remember anything.
\newline Cons: it reports a number the product cannot deliver. The gate would be computed over 1,437 listings whose make the API refuses with a 422 (FR-04), and \texttt{evaluate}'s per-make segments would include makes that are out of scope, so SC-04 would check a segment the API never answers.
\item \textbf{Option C (chosen): record the list in \texttt{split}, and restrict the model input and the metrics downstream.}
\newline Pros: keeps both properties that matter. \texttt{split} stays a lossless partition, so the threshold stays revisitable and no row disappears without an artefact saying where it went; and the reported population is the served population, because the stages that build model input and measure the gate apply the list. It is also the reversible direction, and \texttt{dvc.\allowbreak{}lock} ties the artefact's hash to the stages that consume it.
\newline Cons: the guarantee now lives in three stages rather than in one, so it has to be written down and tested rather than being true by construction. That is the cost this entry accepts, and \texttt{docs/\allowbreak{}docs/\allowbreak{}pipeline.\allowbreak{}md} carries the four obligations so they are not only in a squash-merge commit message.
\end{itemize}
\par
\emph{Why:} The authority on what the split does is problem-spec section 5, the evaluation protocol, and it removes nothing but \texttt{ES}: ``The remaining listings are split into train, validation, calibration and test sets grouped by seller''. Section 2 scopes the \emph{model}, not the split, and EDN-05 already places the rejection at the API: ``the pipeline computes the list and the API rejects other makes''. The repository owner read the two documents the same way and confirmed option C.
\par
The measurements say the choice is cheap in both directions. Applying the list moves every realised share by at most \textbf{0.08 pp} (99,326 rows to 97,889), so restricting downstream does not disturb the split proportions the gate of EDN-39 checks. Of the \textbf{6,079} holdout rows, \textbf{5,979} are a supported make and \textbf{100} are not, which matches EDN-14's count of the \texttt{ES} rows the API would accept under FR-04, so the holdout needs the same filter and the drift and retrain work is unaffected by where the filter sits.
\par
One framing from the original write-up is corrected here, because it would have gone into the report wrong. The 1,437 out-of-scope rows would \textbf{not} inflate the reported metrics: rare makes are harder to price, so a pooled figure computed over them is if anything pessimistic. The real problem is a different one, and it is about comparability rather than optimism: the reported population would not be the served population, and it would not be comparable to the reference values in problem-spec section 8, which are all post-filter. That is why the metrics are restricted even though leaving them unrestricted would not flatter the model.
\par
What the chosen option owes, and what \texttt{docs/\allowbreak{}docs/\allowbreak{}pipeline.\allowbreak{}md} records under ``The supported makes'', is four guarantees from \#36, \#37 and \#39, none of them optional: the filter applied to train, validation, calibration, test \textbf{and} the \texttt{ES} holdout; applied \textbf{before} any training-derived vocabulary or statistic is computed, since fitting on rows the API refuses puts unservable makes into the model's own inputs; \texttt{evaluate}'s per-make segments restricted to the listed makes; and a test asserting that the filtered frames contain only supported makes, so the guarantee is checked rather than intended.
How each was discharged: the third by \texttt{evaluate} in \#39 (\texttt{check\_\allowbreak{}make\_\allowbreak{}levels\_\allowbreak{}are\_\allowbreak{}served}); the first and second together by \texttt{features} in \#63 (\hyperref[EDN-67]{EDN-67}), which replaced the train-and-validation filter \#37 had put in \texttt{train} and the test filter \#39 had put in \texttt{evaluate} with checks that fail instead of filtering; and the fourth by a test for each of the three, every one of which first asserts that its input holds an unsupported make.
\par
Two stale claims were a consequence of this decision, and are now applied: the comments in \texttt{recommenditos/\allowbreak{}schema.\allowbreak{}py} and \texttt{recommenditos/\allowbreak{}data/\allowbreak{}preprocess.\allowbreak{}py} said the supported-make filter belongs to, or happens in, \texttt{split}. Both now say that \texttt{split} computes the list and the stages building model input apply it.},
  ai = {Information seeking, Alternative generation, Alternative assessment, Recommendation},
  response = {Accepted},
  assessment = {AI found that the skeleton contained evidence for both readings and that the two documents behind them disagree, rather than implementing one and moving on: \texttt{schema.\allowbreak{}py} and \texttt{preprocess.\allowbreak{}py} place the filter in \texttt{split}, while problem-spec section 5 and EDN-05 place it at the API. It laid out the options with the irreversibility argument that decided it, and flagged the choice as an EDN candidate before acting on it. It also measured the consequences instead of asserting them, which is what produced the 0.08 pp share shift and the 5,979 of 6,079 holdout figure, and it corrected its own earlier framing that the out-of-scope rows would inflate the metrics once it checked the direction of the effect. Lukas confirmed option C.},
  aievidence = {Claude Code sessions on 2026-09-30 while implementing and then adversarially reviewing issue \#35: the first flagged the two readings and asked for a decision rather than choosing silently; the adversarial review verified the numbers and found that the obligation existed only in the pull request body, with issues \#36, \#37 and \#39 mentioning the make list nowhere at all.},
  evidence = {\href{https://github.com/mlops-2627q1-mds-upc/MLOPS_Recommenditos/blob/main/recommenditos/data/split_data.py}{\texttt{recommenditos/\allowbreak{}data/\allowbreak{}split\_\allowbreak{}data.\allowbreak{}py}}; the \texttt{supported\_\allowbreak{}makes.\allowbreak{}json} dependency on the \texttt{features} stage in \href{https://github.com/mlops-2627q1-mds-upc/MLOPS_Recommenditos/blob/main/dvc.yaml}{\texttt{dvc.\allowbreak{}yaml}}; \href{https://github.com/mlops-2627q1-mds-upc/MLOPS_Recommenditos/blob/main/docs/docs/pipeline.md}{pipeline docs}, ``The supported makes'', for the four downstream guarantees; \href{https://github.com/mlops-2627q1-mds-upc/MLOPS_Recommenditos/blob/main/reports/analysis/split_gate.py}{\texttt{reports/\allowbreak{}analysis/\allowbreak{}split\_\allowbreak{}gate.\allowbreak{}py}} for the 0.08 pp and 5,979 measurements; \hyperref[EDN-05]{EDN-05}; \hyperref[EDN-14]{EDN-14}; \hyperref[EDN-39]{EDN-39}; \href{https://github.com/mlops-2627q1-mds-upc/MLOPS_Recommenditos/issues/35}{issue \#35}; \href{https://github.com/mlops-2627q1-mds-upc/MLOPS_Recommenditos/pull/53}{PR \#53}.},
}

% Generated by tools/edn_latex.py from reports/edn.md.
% Do not edit: change reports/edn.md and run `make edn`.
\ednentry{
  id = {EDN-49},
  title = {No bias correction on the log-to-euro inverse transform},
  date = {2026-09-30},
  milestone = {M3: Quality Assurance},
  topic = {Modelling Approach, Evaluation Protocol},
  participants = {@lukas2510},
  decision = {\texttt{predict\_\allowbreak{}eur} is \texttt{exp} of the model's log-space prediction, with no Duan smearing factor and no lognormal variance correction. The target functional is the conditional \textbf{median} price, and that is what every metric this project reports measures.},
  rationale = {\emph{Alternatives considered:}
\begin{itemize}[leftmargin=*, topsep=2pt]
\item \textbf{Option A (chosen): naive \texttt{exp}.}
\newline Pros: it is the conditional median of price, which is exactly what MdAPE, the +/-20\,\% share and the conformal intervals are about; it needs no statistic fitted on the training residuals, so nothing extra has to be persisted, versioned or kept in step between \texttt{train}, \texttt{evaluate} and the API; and it is what the model card already told readers the model does.
\newline Cons: it is \emph{not} an unbiased estimate of the mean price, so anyone reading the estimate as an expected sale value is reading it wrong. The model card and the response schema say ``typical price'' for that reason.
\item \textbf{Option B: Duan's smearing estimator.} Multiply by the mean of \texttt{exp(residual)} over the training rows.
\newline Pros: the textbook correction for a log-linear model, non-parametric, and it is what a reviewer who knows the retransformation problem will ask about.
\newline Cons: it targets the conditional mean, so it shifts every prediction upward, away from the typical asking price. Measured on the real snapshot's 19,665 test rows it makes every variant worse: factor 1.0258 and MdAPE 9.52\,\% to 9.83\,\% for \texttt{b1}; 1.0052 and 6.83\,\% to 6.86\,\% for \texttt{lgbm-\allowbreak{}basic}; 1.0028 and 6.32\,\% to 6.37\,\% for \texttt{lgbm-\allowbreak{}extended}. It also adds a fitted constant to the artefact that has to travel with the model and be applied identically in three places.
\item \textbf{Option C: lognormal variance correction,} multiply by \texttt{exp(sigma\textasciicircum{}2 /\allowbreak{} 2)}.
\newline Pros: closed form, no residual pass needed.
\newline Cons: everything option B has, plus a distributional assumption the residuals do not have to satisfy; it is strictly weaker than smearing, which estimates the same quantity without the assumption.
\end{itemize}
\par
\emph{Why:} The correction answers a question nobody here asks. Fitting with squared error on \texttt{log(price)} estimates \texttt{E[log P | x]}, so \texttt{exp} of it is the geometric mean, which for the conditional distribution of a price is the median. Every number in problem-spec section 6 and every criterion in section 8 is median- or quantile-flavoured: MdAPE is the primary metric, SC-02 is a share within a band, and SC-05's intervals are conformal and therefore calibrated after the fact on held-out data, so they need no correction either. A mean-targeting correction would make each of those numbers worse, and the measurements say it does, on all three variants that have an inverse transform at all.
\par
\texttt{b0} is the reason the decision is visible in the code rather than implicit. The median baseline is fitted on \texttt{price} directly, not on \texttt{log\_\allowbreak{}price}, because a median commutes with a monotone transform: \texttt{exp(median(log price))} and \texttt{median(price)} are the same number, and the readable one is the one worth storing in a lookup table a person can audit. So B0 needs no inverse transform, B1 and the two LightGBM variants share one, and the seam puts it inside the model -- \texttt{predict\_\allowbreak{}eur} returns euros and \texttt{predict\_\allowbreak{}log\_\allowbreak{}price} is defined as its log -- so \texttt{evaluate} and the API cannot come to disagree about it.
\par
The honest limit of this entry: the effect is small. The largest degradation measured is 0.32 pp, on the interpretable baseline rather than on the deployed candidate. It is recorded because the retransformation problem is a standard question about a log-target model, the answer has numbers behind it, and the alternative would have quietly cost accuracy.},
  ai = {Information seeking, Alternative generation, Alternative assessment, Recommendation},
  response = {Accepted with modifications},
  assessment = {AI raised the retransformation question unprompted, named the two standard corrections and argued from the target functional rather than from convention, which is the argument that decides it. Its first write-up carried the effect sizes from a different run (0.5 to 2.2 pp) as if they were ours; those were re-measured on the real snapshot and are an order of magnitude smaller, so the entry now claims what this project measured and says that the effect is small. The direction of the effect -- worse on every variant -- reproduced.},
  aievidence = {Claude Code session on 2026-09-30 implementing issue \#37: the smearing factor and the paired MdAPE per variant were computed on the real snapshot in a scratchpad run of the whole chain, with the numbers reproduced in the model card.},
  evidence = {\href{https://github.com/mlops-2627q1-mds-upc/MLOPS_Recommenditos/blob/main/recommenditos/modeling/model.py}{\texttt{recommenditos/\allowbreak{}modeling/\allowbreak{}model.\allowbreak{}py}}, \texttt{Model.\allowbreak{}predict\_\allowbreak{}eur} and \texttt{MedianBaselineModel}; \href{https://github.com/mlops-2627q1-mds-upc/MLOPS_Recommenditos/blob/main/docs/docs/model-card.md}{model card}, Training Procedure, Training; \texttt{tests/\allowbreak{}test\_\allowbreak{}model.\allowbreak{}py::test\_\allowbreak{}predict\_\allowbreak{}log\_\allowbreak{}price\_\allowbreak{}is\_\allowbreak{}exactly\_\allowbreak{}the\_\allowbreak{}log\_\allowbreak{}of\_\allowbreak{}predict\_\allowbreak{}eur}; \href{https://github.com/mlops-2627q1-mds-upc/MLOPS_Recommenditos/issues/37}{issue \#37}.},
}

% Generated by tools/edn_latex.py from reports/edn.md.
% Do not edit: change reports/edn.md and run `make edn`.
\ednentry{
  id = {EDN-50},
  title = {One-hot encoding for the Ridge baseline only, against EDN-02's ``no one-hot''},
  date = {2026-09-30},
  milestone = {M3: Quality Assurance},
  topic = {Feature Engineering, Modelling Approach},
  participants = {@lukas2510},
  decision = {The \texttt{b1} Ridge variant one-hot encodes its categoricals, with \texttt{handle\_\allowbreak{}unknown=\allowbreak{}\textquotedbl{}infrequent\_\allowbreak{}if\_\allowbreak{}exist\textquotedbl{}} and \texttt{min\_\allowbreak{}category\_\allowbreak{}rows: 5} as the threshold below which levels share one column. This is confined to \texttt{b1}: the LightGBM variants and the CatBoost challenger take the categoricals natively, as \texttt{category} codes over the levels the \texttt{features} stage fixed, and no one-hot column exists anywhere near them.},
  rationale = {\emph{Alternatives considered:}
\begin{itemize}[leftmargin=*, topsep=2pt]
\item \textbf{Option A (chosen): one-hot, confined to B1.}
\newline Pros: it is the standard encoding for a linear model, it needs no target information, and it keeps B1 able to do the job it exists for -- telling a Porsche from a Dacia as part of an interpretable depreciation baseline. \texttt{handle\_\allowbreak{}unknown=\allowbreak{}\textquotedbl{}infrequent\_\allowbreak{}if\_\allowbreak{}exist\textquotedbl{}} is also what makes an unseen level an answer rather than an error, which EDN-18 requires of the serving path.
\newline Cons: it is the encoding EDN-02 rules out, so a reader who takes that rule as project-wide sees a violation. The column count grows with the level count: 503 columns at real scale before any folding, against 16 input features.
\item \textbf{Option B: target encoding with cross-fitting.}
\newline Pros: one column per categorical however many levels it has, and it usually beats one-hot for a linear model on high-cardinality data.
\newline Cons: it encodes the target into a feature, so it leaks unless it is cross-fitted, and cross-fitting means a second resampling scheme inside the stage, a second thing to persist and a second thing that can differ between training and serving. For a baseline whose purpose is to be simple and auditable that is the wrong trade, and the leak is the kind of defect this project would only find by not finding it.
\item \textbf{Option C: drop the categoricals from B1 and fit on the numerics alone.}
\newline Pros: no encoding decision at all, and the ladder still has a linear baseline.
\newline Cons: it destroys the baseline. A depreciation model that cannot see the make is not a weaker version of B1, it is a different and much worse thing, and SC-03's comparison against B0 would be against a baseline nobody would have built.
\item \textbf{Option D: read EDN-02 as project-wide and drop B1 from the ladder.}
\newline Pros: no apparent contradiction to explain.
\newline Cons: problem-spec section 7 names B1 as a baseline, and the interpretable linear reference is what makes the tree models' gain legible. Removing a baseline to preserve the letter of a rule about a different model family is the wrong direction.
\end{itemize}
\par
\emph{Why:} EDN-02 is a decision about the \emph{model family}, and ``categoricals stay categorical, no one-hot'' is stated there as one of the reasons gradient boosting was chosen over the alternatives -- LightGBM and CatBoost handle them natively, so the pipeline does not have to blow up the matrix. It is not a constraint on how a linear baseline encodes its inputs, because a linear model has no native handling to use: the choice is one-hot, target encoding or nothing. Recording this keeps the next reader from either ``fixing'' B1 to match the rule or reading the rule as broken.
\par
Two details of the encoding are decisions in their own right and are measured rather than assumed.
\par
Absence is a level. Categoricals are mapped to a literal \texttt{\textquotedbl{}\_\allowbreak{}\_\allowbreak{}missing\_\allowbreak{}\_\allowbreak{}\textquotedbl{}} value before encoding, deterministically at fit and at predict time, so an absent value gets a coefficient like any other level instead of failing the fit or being dropped. That is the same treatment EDN-15 asks for, expressed in the only way a linear model can express it.
\par
\texttt{min\_\allowbreak{}category\_\allowbreak{}rows} is an absolute row count rather than a share, so that what the threshold means does not change as the training set grows. Measured on the real snapshot (60,378 training rows, MdAPE on the 19,665 test rows): no folding gives 503 encoded columns at 9.48\,\%, 5 gives 404 at 9.52\,\%, 30 gives 294 at 9.71\,\%, 100 gives 187 at 10.61\,\% and 1,000 gives 68 at 13.45\,\%. So folding costs accuracy monotonically, and the reason to fold is not accuracy.
\par
What the reason is, read off the fitted encoder rather than assumed: at 5, folding touches \textbf{2 of the 8 categoricals}. \texttt{model} folds 100 of its 438 levels into one infrequent column and \texttt{fuel\_\allowbreak{}category} folds 1. The other six have no infrequent group at all, and \texttt{make} and \texttt{body\_\allowbreak{}type} have no \texttt{\_\allowbreak{}\_\allowbreak{}missing\_\allowbreak{}\_\allowbreak{}} level either, because the real training rows never leave those two empty. An unseen or an absent value in one of those six therefore encodes as all zeros rather than into a group.
\par
EDN-18 is satisfied either way, and checked rather than assumed: masking each of the eight in turn, and giving each an unseen value in turn, the real \texttt{b1} answers with a finite positive price all sixteen times. It is simply satisfied by the all-zero encoding for six of the eight -- a coherent encoding of ``no level applies'' -- and not by the infrequent group the threshold produces. The threshold is what stops \texttt{handle\_\allowbreak{}unknown=\allowbreak{}\textquotedbl{}infrequent\_\allowbreak{}if\_\allowbreak{}exist\textquotedbl{}} from being pointless for \texttt{model}, which is the one column where rare levels are the norm, and nothing more.
\par
So the justification for 5 is what is left: folding is nearly free there, 0.04 pp, and it stops a level seen in a handful of rows from getting a coefficient fitted on those rows. At real scale that is 47 \texttt{model} levels seen exactly once, among the 100 the threshold folds. An infrequent group does not need 5 to exist -- 2 already gives one, folding 48 levels for 0.02 pp -- so 5 is a round number inside a flat region rather than a boundary, and the value is a parameter precisely so the ladder can revisit it.
\par
One claim that did \textbf{not} survive measurement and is recorded so nobody repeats it: the design this work followed justified a much larger threshold with a small-data pathology, a mid-range car with an unseen model and country predicted at 94,323 EUR. That did not reproduce here at either scale. On the fixture the same probe gives 13,993 EUR unfolded and 15,413 EUR at a threshold of 30; at real scale, 32,332 EUR and 33,104 EUR. The threshold is justified by the infrequent group and by the one-row-coefficient argument, not by a pathology this project observed.},
  ai = {Information seeking, Alternative generation, Alternative assessment, Recommendation},
  response = {Accepted with modifications},
  assessment = {AI spotted the apparent conflict with EDN-02 before writing the estimator and flagged it as an EDN candidate rather than resolving it silently, and its reading -- that EDN-02 is about the tree family -- is the one the entry records. It was overruled on the threshold: it recommended 30 on the strength of a measurement from another run and asserted the value was free at real scale. Sweeping it here showed folding is not free and that its pathology did not reproduce, so the value is 5 and the entry says what the evidence actually supports.},
  aievidence = {Claude Code session on 2026-09-30 implementing issue \#37: the threshold sweep was run on the fixture and then on the real snapshot, and the ``unseen model and country'' probe was re-measured at both scales before the recommended value was changed.},
  evidence = {\href{https://github.com/mlops-2627q1-mds-upc/MLOPS_Recommenditos/blob/main/recommenditos/modeling/model.py}{\texttt{recommenditos/\allowbreak{}modeling/\allowbreak{}model.\allowbreak{}py}}, \texttt{RidgeModel}; the \texttt{min\_\allowbreak{}category\_\allowbreak{}rows} comment in \href{https://github.com/mlops-2627q1-mds-upc/MLOPS_Recommenditos/blob/main/params.yaml}{\texttt{params.\allowbreak{}yaml}}; \href{https://github.com/mlops-2627q1-mds-upc/MLOPS_Recommenditos/blob/main/docs/docs/model-card.md}{model card}, Training Procedure, Training; \texttt{tests/\allowbreak{}test\_\allowbreak{}model.\allowbreak{}py::test\_\allowbreak{}an\_\allowbreak{}unseen\_\allowbreak{}category\_\allowbreak{}is\_\allowbreak{}treated\_\allowbreak{}as\_\allowbreak{}missing\_\allowbreak{}not\_\allowbreak{}an\_\allowbreak{}error}; \hyperref[EDN-02]{EDN-02}; \hyperref[EDN-18]{EDN-18}; \href{https://github.com/mlops-2627q1-mds-upc/MLOPS_Recommenditos/issues/37}{issue \#37}.},
}

% Generated by tools/edn_latex.py from reports/edn.md.
% Do not edit: change reports/edn.md and run `make edn`.
\ednentry{
  id = {EDN-51},
  title = {Mean fill plus a per-feature missingness indicator for the Ridge numerics, which is not imputation},
  date = {2026-09-30},
  milestone = {M3: Quality Assurance},
  topic = {Feature Engineering, Modelling Approach},
  participants = {@lukas2510},
  decision = {The \texttt{b1} Ridge variant's numeric branch is the union of two paths: a mean-filled, standardised copy of every numeric feature, and \texttt{MissingIndicator(features=\allowbreak{}\textquotedbl{}all\textquotedbl{})}, which emits one indicator column per feature whether or not that feature was ever missing at fit time. No other estimator fills anything, and no artefact of the pipeline is imputed: the fill exists only inside B1's fitted pipeline.},
  rationale = {\emph{Alternatives considered:}
\begin{itemize}[leftmargin=*, topsep=2pt]
\item \textbf{Option A (chosen): mean fill plus an indicator for every feature.}
\newline Pros: information-preserving for a linear model. For a row where feature \texttt{j} is absent the contribution is \texttt{beta\_\allowbreak{}j * mean\_\allowbreak{}j + gamma\_\allowbreak{}j}, and \texttt{gamma\_\allowbreak{}j} is free to absorb whatever the absence is worth, so the fill value is a numerically neutral placeholder rather than a guess at the value. Nothing about the missingness is destroyed, which is what EDN-15 is protecting. \texttt{features=\allowbreak{}\textquotedbl{}all\textquotedbl{}} also makes SC-06 honest: the criterion masks a field and measures what it costs, and it can only do that if masking a field the training rows happened to fill still produces an indicator.
\newline Cons: a reader who sees \texttt{SimpleImputer} in the pipeline will conclude the project imputes, which is exactly the misreading this entry exists to prevent. It also doubles the numeric column count.
\item \textbf{Option B: \texttt{SimpleImputer} alone, or \texttt{SimpleImputer(add\_\allowbreak{}indicator=\allowbreak{}True)} at its default.}
\newline Pros: less to explain, fewer columns.
\newline Cons: this is the option that silently breaks SC-06, and it is the default, which is why it is worth an entry. scikit-learn emits an indicator only for features that were missing \textbf{at fit time}, so a column complete in the training split is mean-filled with no indicator the moment the criterion masks it -- the model then reports the training mean for that field, the prediction barely moves, and SC-06 measures the model's own imputation rather than the cost of the missing input. Without the indicator at all, an absent value is indistinguishable from an average one, which is imputation in the sense EDN-15 rejects.
\item \textbf{Option C: drop the rows or the columns with missing values.}
\newline Pros: no fill anywhere, so no explaining to do.
\newline Cons: dropping rows would fit B1 on a different and much smaller population than the other three variants, so the ladder would stop being a comparison. Dropping columns would throw away \texttt{nr\_\allowbreak{}prev\_\allowbreak{}owners} and \texttt{gears}, which are missing in a third and a quarter of the training rows respectively, and the missingness itself carries signal (EDN-15).
\item \textbf{Option D: give up on B1 taking numerics with gaps and use a model that handles them natively.}
\newline Pros: no encoding at all.
\newline Cons: that model is LightGBM, which is ladder steps 3 and 4. B1's value is that it is a \emph{linear} reference.
\end{itemize}
\par
\emph{Why:} EDN-15's rule is that a missing value is never imputed, because the missingness itself carries signal. Ridge cannot consume NaN, so B1 either encodes the gaps or does not exist. The construction above is the encoding under which nothing is lost: the pair \texttt{(filled value, indicator)} is a bijection with \texttt{(value, present)} for a present value and carries the absence explicitly for an absent one, so the linear model has exactly the information the tree models get from a NaN split. The fill is a placeholder chosen to be numerically harmless after standardisation, not an estimate of the missing value, and the entry exists because those two look identical in the code.
\par
\texttt{features=\allowbreak{}\textquotedbl{}all\textquotedbl{}} is the part a test has to hold, because it is invisible from the outside. It is also invisible on the synthetic fixture: every numeric column there has a missing value in the training split, so the default emits the same set of indicators and a test written against the fixture cannot tell the two apart. That was found by breaking the implementation on purpose and watching the test suite pass. The test now fills one column completely before fitting and fails with the default.
\par
What that test constructs is \textbf{not} the situation on the real snapshot, and the entry says so rather than borrowing the stronger claim. Of the 8 numeric features of the \texttt{basic} set -- the only set \texttt{b1} uses -- \textbf{none} is complete in the real training split: the emptiest is \texttt{age\_\allowbreak{}years} with 1 gap in 60,378 rows and the fullest \texttt{nr\_\allowbreak{}prev\_\allowbreak{}owners} with 22,697. So scikit-learn's default would emit the same eight indicators here, and \texttt{features=\allowbreak{}\textquotedbl{}all\textquotedbl{}} changes nothing in the shipped configuration. It is insurance, not a measured fix, and what it insures against is real: 136 of the \texttt{extended} set's 147 numeric features are complete at fit time, so a Ridge on \texttt{extended} -- which the ladder may well want, since B1 is the interpretable reference for the winning feature set -- would lose 136 of its 147 indicators to the default, and SC-06 would measure the model's own imputation for every one of them. The guard is kept because it costs nothing and because the configuration it protects is one change away.
\par
The other half of the same discipline is what the estimator does \emph{not} do. \texttt{SimpleImputer(keep\_\allowbreak{}empty\_\allowbreak{}features=\allowbreak{}True)} keeps an all-missing column in place rather than dropping it, so the matrix shape does not depend on which columns the training rows happened to fill; and nothing outside B1's pipeline fills anything, so no artefact on disk and no other variant carries a filled value.},
  ai = {Alternative generation, Alternative assessment, Recommendation, Solution generation},
  response = {Accepted},
  assessment = {AI identified that the obvious construction would make SC-06 measure its own imputation, which is a subtle failure that would have produced a passing criterion and a wrong conclusion, and it named \texttt{features=\allowbreak{}\textquotedbl{}all\textquotedbl{}} as the fix and flagged the whole construction as an EDN candidate because of how it reads. It also proposed the test for it. The test it proposed was then shown, by mutating the implementation, to pass either way on this fixture, and it was replaced with one that fills a column first; that correction came from the mutation exercise rather than from the design.},
  aievidence = {Claude Code session on 2026-09-30 implementing issue \#37: the mutation battery run before the pull request recorded \texttt{MissingIndicator(features=\allowbreak{}\textquotedbl{}all\textquotedbl{})} to \texttt{MissingIndicator()} as a surviving mutation, which is what produced the replacement test.},
  evidence = {\href{https://github.com/mlops-2627q1-mds-upc/MLOPS_Recommenditos/blob/main/recommenditos/modeling/model.py}{\texttt{recommenditos/\allowbreak{}modeling/\allowbreak{}model.\allowbreak{}py}}, \texttt{RidgeModel} and \texttt{\_\allowbreak{}ridge\_\allowbreak{}pipeline}; \texttt{tests/\allowbreak{}test\_\allowbreak{}model.\allowbreak{}py::test\_\allowbreak{}the\_\allowbreak{}ridge\_\allowbreak{}emits\_\allowbreak{}a\_\allowbreak{}missing\_\allowbreak{}indicator\_\allowbreak{}for\_\allowbreak{}every\_\allowbreak{}numeric\_\allowbreak{}feature} and \texttt{::test\_\allowbreak{}a\_\allowbreak{}numeric\_\allowbreak{}feature\_\allowbreak{}complete\_\allowbreak{}in\_\allowbreak{}training\_\allowbreak{}still\_\allowbreak{}gets\_\allowbreak{}an\_\allowbreak{}indicator}; \href{https://github.com/mlops-2627q1-mds-upc/MLOPS_Recommenditos/blob/main/docs/docs/model-card.md}{model card}, Training Procedure, Training; \hyperref[EDN-15]{EDN-15}; \href{https://github.com/mlops-2627q1-mds-upc/MLOPS_Recommenditos/issues/37}{issue \#37}.},
}

% Generated by tools/edn_latex.py from reports/edn.md.
% Do not edit: change reports/edn.md and run `make edn`.
\ednentry{
  id = {EDN-52},
  title = {\texorpdfstring{\texttt{predict\_\allowbreak{}eur}}{predict\_eur} bounds every prediction to the training price range},
  date = {2026-09-30},
  milestone = {M3: Quality Assurance},
  topic = {Modelling Approach, API Design},
  participants = {@lukas2510, decided by the repository owner},
  decision = {\texttt{Model.\allowbreak{}predict\_\allowbreak{}eur} clips the estimator's log-space output to \texttt{log} of the observed minimum and maximum \texttt{price} of that variant's own training rows, before the exponential, and records both bounds in \texttt{models/\allowbreak{}<variant>/\allowbreak{}model.\allowbreak{}json} under \texttt{training.\allowbreak{}price\_\allowbreak{}min\_\allowbreak{}eur} and \texttt{training.\allowbreak{}price\_\allowbreak{}max\_\allowbreak{}eur}. The clip is in the base class and \texttt{predict\_\allowbreak{}eur} is final, so no estimator can bypass it, and a bundle whose recorded range cannot bound a price is refused on load.},
  rationale = {\emph{Alternatives considered:}
\begin{itemize}[leftmargin=*, topsep=2pt]
\item \textbf{Option A: return the estimator's output unbounded.}
\newline Pros: nothing hidden; whatever the model believes is what the caller sees, which is the right default for a component that reports its own errors.
\newline Cons: the interface cannot then promise what \texttt{evaluate} and the API rely on. \texttt{exp} overflows to \texttt{inf} above about 710 in log space and underflows to exactly \texttt{0.\allowbreak{}0} below -746, and both are reachable from a linear model with no bounded link. An \texttt{inf} prediction makes MdAPE \texttt{nan}, and \texttt{nan <=\allowbreak{} 0.\allowbreak{}09} is \texttt{False}, so the gate would record a failure with no indication that the number was not a number; a \texttt{0.\allowbreak{}0} prediction breaks the ``strictly positive'' promise the API's response schema rests on.
\item \textbf{Option B (chosen): clip to the training price range, in log space, before the exponential.}
\newline Pros: makes ``finite and strictly positive'' a property of the arithmetic rather than something hoped for, because the bounds are finite and positive by construction. Bounding before the exponential rather than after means the overflow is unreachable instead of merely clipped, so no \texttt{RuntimeWarning} and no \texttt{inf} ever exists. The bound is the range the training data supports, which is the only range about which the model has evidence. And because both edges are recorded in the bundle, the share of predictions sitting exactly on an edge can be reported, so the clip cannot hide a pathology it is covering up.
\newline Cons: a clipped prediction is wrong in a way that no longer looks wrong. It is a plausible number, so a reader cannot tell it from an estimate unless someone reports the share on the bound -- which is why that reporting is part of the decision rather than a nice-to-have.
\item \textbf{Option C: clip to a quantile range of the training prices, say the 0.1st to the 99.9th percentile.}
\newline Pros: a tighter and more defensible band than the extremes, and less sensitive to one collector car in the training rows.
\newline Cons: it would move real predictions. The bound stops being a guard rail and becomes a silent recalibration of the top and bottom of the distribution, which is a modelling decision dressed up as a safety check, and it would bias the metrics for exactly the segments SC-04 is watching.
\item \textbf{Option D: raise on a prediction outside the training range.}
\newline Pros: the loudest possible signal, and no wrong number is ever served.
\newline Cons: it turns a model's extrapolation into a 500 from the API for a request that is in scope, which contradicts FR-05's promise that an in-scope car is always answered. It would also make \texttt{evaluate} unable to finish a run over the test set, so the one place the behaviour can be measured would be the place that cannot run.
\end{itemize}
\par
\emph{Why:} The measurement is what settled it, and it also corrected the reason the bound was originally proposed. The design this work followed justified the clip with a Ridge returning \textbf{487,417 EUR} for a make-only request. That number would \textbf{not} be bounded: on the real snapshot the training range is \textbf{500 to 1,814,750 EUR}, so 487,417 is comfortably inside it and the clip would never fire. Repeating the make-only probe here gives 6,529 EUR for B1, nowhere near the bound.
\par
What the clip does catch, measured over the 19,665 test rows of the real snapshot, is a genuine extrapolation: B1's unbounded prediction leaves the range \textbf{once}, at \textbf{10,635,538 EUR}. None of the other three variants leaves it at all, with unbounded maxima of 1,549,910 EUR for B0, 1,055,373 for \texttt{lgbm-\allowbreak{}basic} and 731,089 for \texttt{lgbm-\allowbreak{}extended}. So the clip is load-bearing for exactly one variant on roughly one row in twenty thousand -- rare, and worth having anyway, because the alternative is one impossible number per twenty thousand answers in a product whose whole claim is a credible price.
\par
The honest limit is stated in the code and in the model card rather than left for a reader to discover: bounding that 10.6 million turns it into 1,814,750 EUR, which is the price of the most expensive car in the training rows and not a sensible estimate for the car that was asked about. The bound makes the number representable, not right. The reserved \texttt{predict\_\allowbreak{}interval\_\allowbreak{}eur} of SC-05 is where a caller will eventually see that such a prediction is not to be trusted.
\par
One guard came out of writing this down. \texttt{predict\_\allowbreak{}eur} promises a strictly positive value and keeps that promise \emph{through} the bound, so a recorded range that cannot bound a price -- a lower edge of 0, which would clip to \texttt{log(0)} and exponentiate to exactly \texttt{0.\allowbreak{}0} -- has to be refused rather than served. \texttt{preprocess.\allowbreak{}price\_\allowbreak{}min\_\allowbreak{}eur: 500} makes it unreachable from the pipeline; the guard exists for a hand-edited bundle, and it is what keeps the promise unconditional instead of conditional on a parameter somewhere else in the system.},
  ai = {Alternative generation, Alternative assessment, Recommendation, Solution generation},
  response = {Accepted with modifications},
  assessment = {AI proposed the clip and the float-limit argument for applying it in log space, which is the part that makes the interface's promise structural. Its justifying measurement did not hold up: the 487,417 EUR case it cited is inside the training range and would never be clipped, which was found by measuring the range instead of assuming it. Re-measuring over the whole test split produced the case that does justify the clip, and the entry and the code now claim that one and explicitly disclaim the other. AI also proposed the test; the test was strengthened to replace the estimator with one returning an absurd log price, so that it fails when the clip is removed rather than depending on how far a particular fit happens to extrapolate.},
  aievidence = {Claude Code session on 2026-09-30 implementing issue \#37: the unbounded predictions of all four variants were computed over the real test split, and the make-only probe was re-run at both data scales before the justification in the entry was rewritten.},
  evidence = {\href{https://github.com/mlops-2627q1-mds-upc/MLOPS_Recommenditos/blob/main/recommenditos/modeling/model.py}{\texttt{recommenditos/\allowbreak{}modeling/\allowbreak{}model.\allowbreak{}py}}, \texttt{Model.\allowbreak{}predict\_\allowbreak{}eur} and the range guard in \texttt{Model.\allowbreak{}\_\allowbreak{}\_\allowbreak{}init\_\allowbreak{}\_\allowbreak{}}; \texttt{tests/\allowbreak{}test\_\allowbreak{}model.\allowbreak{}py::test\_\allowbreak{}a\_\allowbreak{}prediction\_\allowbreak{}that\_\allowbreak{}would\_\allowbreak{}overflow\_\allowbreak{}is\_\allowbreak{}bounded\_\allowbreak{}to\_\allowbreak{}the\_\allowbreak{}training\_\allowbreak{}range}, \texttt{::test\_\allowbreak{}a\_\allowbreak{}training\_\allowbreak{}price\_\allowbreak{}range\_\allowbreak{}that\_\allowbreak{}cannot\_\allowbreak{}bound\_\allowbreak{}a\_\allowbreak{}price\_\allowbreak{}is\_\allowbreak{}refused} and the range assertions in \texttt{::test\_\allowbreak{}predict\_\allowbreak{}eur\_\allowbreak{}contract}; \href{https://github.com/mlops-2627q1-mds-upc/MLOPS_Recommenditos/blob/main/docs/docs/model-card.md}{model card}, Training Procedure, Training; \href{https://github.com/mlops-2627q1-mds-upc/MLOPS_Recommenditos/pull/62}{PR \#62}; \href{https://github.com/mlops-2627q1-mds-upc/MLOPS_Recommenditos/issues/37}{issue \#37}.},
}

% Generated by tools/edn_latex.py from reports/edn.md.
% Do not edit: change reports/edn.md and run `make edn`.
\ednentry{
  id = {EDN-53},
  title = {\texorpdfstring{\texttt{train.\allowbreak{}num\_\allowbreak{}threads}}{train.num\_threads} is pinned to 1, trading fit speed for a machine-independent artefact},
  date = {2026-09-30},
  milestone = {M3: Quality Assurance},
  topic = {Reproducibility, Modelling Approach},
  participants = {@lukas2510, decided by the repository owner},
  decision = {\texttt{params.\allowbreak{}yaml} carries \texttt{train.\allowbreak{}num\_\allowbreak{}threads: 1} and \texttt{train} passes it to LightGBM as \texttt{n\_\allowbreak{}jobs}. It is never left at LightGBM's default of \texttt{-\allowbreak{}1}, which means every core. The reproducibility it buys is a \texttt{booster.\allowbreak{}txt} that is byte-identical on any machine. What it costs is narrower than first recorded: measured against a pin of 2 or 4 threads it costs about 1.5x the fit, and measured against the default it actually \emph{saves} time and, more to the point, removes an enormous variance.},
  rationale = {\emph{Alternatives considered:}
\begin{itemize}[leftmargin=*, topsep=2pt]
\item \textbf{Option A: leave \texttt{n\_\allowbreak{}jobs} at its default of \texttt{-\allowbreak{}1}.}
\newline Pros: nothing to configure. Not, as it turns out, the fastest fit: see the rationale.
\newline Cons: LightGBM writes the thread count into \texttt{booster.\allowbreak{}txt} as a \texttt{[num\_\allowbreak{}threads: N]} line, so the artefact -- and therefore its DVC hash and \texttt{dvc.\allowbreak{}lock} -- depends on the core count of whoever ran \texttt{dvc repro}. Two contributors with different laptops would produce different artefacts from identical data and parameters, and \texttt{dvc status} could never say the model was unchanged. On top of that, on an 8-core machine the default is both slower and wildly variable at real scale, and it makes the test suite's cost unpredictable, which matters because the suite fits all four variants on 744 rows.
\item \textbf{Option B (chosen): pin to 1.}
\newline Pros: makes \texttt{booster.\allowbreak{}txt} byte-identical across thread counts, which is a stronger claim than NFR-06's ``metrics within 0.1 percentage points'' and is the one the reproducibility test asserts. It also removes oversubscription if the four independent \texttt{train@\allowbreak{}*} stages are ever run at once, and it is the value with by far the smallest spread between repeated fits at either data scale -- which is what makes a reported fit time worth reporting.
\newline Cons: it is measurably slower than a \textbf{pinned} 2 or 4. On the 60,378-row training split \texttt{lgbm-\allowbreak{}extended} fits in a median 18.3 s at 1 thread against 12.2 s at 4, so option C would buy back about a third of the fit per variant. It is not slower than the default it actually replaces.
\item \textbf{Option C: pin to a fixed larger number, 2 or 4.}
\newline Pros: the same machine independence as option B, because the value is pinned rather than discovered, with most of the speed back. This, and not option A, is the alternative that the pin genuinely costs something against.
\newline Cons: it is only machine-independent on machines that \emph{have} that many cores. LightGBM silently uses what is available, so a 2-core CI runner asked for 4 would write a different \texttt{booster.\allowbreak{}txt} than an 8-core laptop, and the property would hold right up until the first machine that breaks it -- the worst kind of guarantee. 1 is the only value every machine can honour.
\end{itemize}
\par
\emph{Why:} The measurement that decides the direction is that the \textbf{trees are already thread-independent and only the recorded setting is not}. Fitting \texttt{lgbm-\allowbreak{}basic} on the fixture in separate processes at 1, 2, 4 and 8 threads gives the same \texttt{best\_\allowbreak{}iteration} (98) and a byte-identical model file once the \texttt{[num\_\allowbreak{}threads: N]} line is normalised away, at every thread count; the md5 of the whole file differs at every one of them. So pinning the value does not make the model reproducible -- it already is -- it makes the \emph{file} reproducible, which is what DVC hashes.
\par
That also corrects a claim worth not repeating. The design this work followed said \texttt{deterministic=\allowbreak{}True} and \texttt{force\_\allowbreak{}row\_\allowbreak{}wise=\allowbreak{}True} are what keep the trees stable across thread counts. They are not, at this data size: with both flags \textbf{off}, the normalised file is still byte-identical across 1, 2, 4 and 8 threads. What the flags do change is \emph{which} trees are built -- the normalised md5 differs between the flags on and off -- so they are not free of consequence, they are simply not what makes the result thread-independent here. They are kept as insurance documented by LightGBM, and the code and the model card now say that rather than claiming a measurement that does not exist.
\par
The cost side was measured rather than assumed, and the first version of this entry stated it against a configuration the default never produces. The full real-scale curve, medians of three fits per thread count on the 60,378-row training split, on an 8-core machine with nothing else running:
\par
{\small\renewcommand{\arraystretch}{1.15}%
\begin{tabularx}{\linewidth}{@{}>{\raggedright\arraybackslash}X>{\raggedright\arraybackslash}X>{\raggedright\arraybackslash}X>{\raggedright\arraybackslash}X>{\raggedright\arraybackslash}X@{}}
\toprule
Variant & 1 thread & 2 & 4 & 8 (the default here) \\
\midrule
\texttt{lgbm-\allowbreak{}basic} & 7.1 s (7.0 to 7.2) & 5.6 s (5.3 to 6.8) & 6.8 s (6.4 to 7.4) & 17.7 s (16.5 to 36.5) \\
\texttt{lgbm-\allowbreak{}extended} & 18.3 s (17.9 to 18.9) & 12.8 s (12.7 to 13.1) & 12.2 s (11.3 to 14.3) & 29.2 s (12.6 to 38.7) \\
\bottomrule
\end{tabularx}}
\par
Two things follow, and the second is the correction. More threads do help, so the design's ``there is no performance reason to use more'' is wrong: 2 to 4 threads are consistently the fastest, and the pin costs about 1.5x the fit for \texttt{lgbm-\allowbreak{}extended} and 1.27x for \texttt{lgbm-\allowbreak{}basic} against them. But the alternative this decision actually rejects is not a pin of 2 or 4, it is LightGBM's \textbf{default of every core}, and against that the pin is free or a gain: 18.3 s against 29.2 s for \texttt{lgbm-\allowbreak{}extended}, 7.1 s against 17.7 s for \texttt{lgbm-\allowbreak{}basic}. So the trade is with option C, not with option A.
\par
The spread matters more than the median. At 1 thread the three fits of a variant land within 0.2 s and 1.0 s of each other; at 8 they span 20 s and 26 s, and the fastest \texttt{lgbm-\allowbreak{}extended} fit of the whole sweep (12.6 s) and the slowest (38.7 s) are both at 8 threads. A number that varies threefold between identical runs is not a fit time anyone can report, which is a reproducibility argument for the pin quite apart from the artefact's hash.
\par
The earlier entry attributed the reversal at 8 threads to the 744-row fixture, where the direction is even sharper: 0.33 s at 1 thread, 0.27 s at 2, 0.48 s at 4 and a median of \textbf{28.58 s at 8}, with one fit of 143.71 s, on a machine already at load 13 to 19. That reading was too narrow. The reversal is not about the row count but about the work per thread: it happens at 60,378 rows too, and on \texttt{lgbm-\allowbreak{}basic} -- 16 features -- more decisively than on \texttt{lgbm-\allowbreak{}extended}'s 162, which is where the extra threads have something to divide. Since the test suite fits every variant on the fixture, on whatever CI runner picks up the job, the pin is what keeps the suite's cost predictable as well.
\par
Absolute seconds here are machine- and load-dependent and are not comparable with the ladder table's one-off fit times in the model card; the ratios within one sweep are the measurement.},
  ai = {Information seeking, Alternative assessment, Recommendation, Solution generation},
  response = {Accepted with modifications},
  assessment = {AI identified that LightGBM records the thread count in the artefact, which is the fact the whole decision rests on, and proposed pinning it. Three of its supporting claims were wrong and were corrected by measuring. That \texttt{deterministic=\allowbreak{}True} and \texttt{force\_\allowbreak{}row\_\allowbreak{}wise=\allowbreak{}True} are what make the trees thread-independent: they are not at this size, the trees are identical with the flags off, though the flags do change which trees are built. That pinning to 1 costs nothing at real scale: it costs about 1.5x the fit against a pin of 2 to 4. And then the correction to that correction -- the entry's own second version stated the cost against 2 to 4 as if that were what the default does, and put eight threads at real scale under ``not measured''. Measuring it reversed the comparison that matters: against the default the pin is free or a gain. The entry now names which alternative the cost is against, and the code comments were rewritten to match.},
  aievidence = {Claude Code session on 2026-09-30 implementing issue \#37: \texttt{booster.\allowbreak{}txt} was hashed with the \texttt{[num\_\allowbreak{}threads: N]} line normalised away across four thread counts and both settings of the determinism flags. The thread sweep was then redone during the review of pull request \#62, as medians of three fits at 1, 2, 4 and 8 threads on both LightGBM variants at real scale on an idle machine, after a first attempt was discarded for having shared the machine with a test run -- which is also how the size of the variance at 8 threads came to light.},
  evidence = {\href{https://github.com/mlops-2627q1-mds-upc/MLOPS_Recommenditos/blob/main/recommenditos/modeling/model.py}{\texttt{recommenditos/\allowbreak{}modeling/\allowbreak{}model.\allowbreak{}py}}, \texttt{LightGBMModel.\allowbreak{}fit}; the \texttt{num\_\allowbreak{}threads} comment in \href{https://github.com/mlops-2627q1-mds-upc/MLOPS_Recommenditos/blob/main/params.yaml}{\texttt{params.\allowbreak{}yaml}}; \texttt{tests/\allowbreak{}test\_\allowbreak{}model.\allowbreak{}py::test\_\allowbreak{}num\_\allowbreak{}threads\_\allowbreak{}is\_\allowbreak{}pinned\_\allowbreak{}from\_\allowbreak{}params}, which trains at a value the project does not use so a hard-coded 1 would fail it, and \texttt{::test\_\allowbreak{}the\_\allowbreak{}payload\_\allowbreak{}of\_\allowbreak{}a\_\allowbreak{}refit\_\allowbreak{}is\_\allowbreak{}byte\_\allowbreak{}identical}; \href{https://github.com/mlops-2627q1-mds-upc/MLOPS_Recommenditos/blob/main/docs/docs/model-card.md}{model card}, Training Procedure, Determinism; \href{https://github.com/mlops-2627q1-mds-upc/MLOPS_Recommenditos/blob/main/docs/docs/requirements.md}{requirements} NFR-06 and NFR-10; \href{https://github.com/mlops-2627q1-mds-upc/MLOPS_Recommenditos/issues/37}{issue \#37}.},
}

% Generated by tools/edn_latex.py from reports/edn.md.
% Do not edit: change reports/edn.md and run `make edn`.
\ednentry{
  id = {EDN-54},
  title = {Tracking degrades silently when it is not configured, with an environment variable to make it fatal},
  date = {2026-09-30},
  milestone = {M3: Quality Assurance},
  topic = {CI/CD, Experiment Tracking},
  participants = {@lukas2510, decided by the repository owner},
  decision = {A pipeline stage opens its MLflow run through \texttt{optional\_\allowbreak{}run}, which yields a real run handle when the server is configured and reachable and a \texttt{DisabledRun} that logs nothing when it is not -- whether because nothing is configured, because the credentials are refused or because the server cannot be reached. The stage runs to completion either way and records \texttt{mlflow.\allowbreak{}tracking\_\allowbreak{}mode} in \texttt{model.\allowbreak{}json}. Setting \texttt{RECOMMENDITOS\_\allowbreak{}REQUIRE\_\allowbreak{}TRACKING=\allowbreak{}1} turns any of those failures back into a failed stage. There is never a fallback to a local tracking store.},
  rationale = {\emph{Alternatives considered:}
\begin{itemize}[leftmargin=*, topsep=2pt]
\item \textbf{Option A: fail the stage whenever tracking is not available,} which is what \texttt{tracked\_\allowbreak{}run} does and what EDN-46's helper was built for.
\newline Pros: a run that was not recorded never happened, which is the strongest possible reading of NFR-06; and the failure is the cheapest one, before the fit.
\newline Cons: CI has no DagsHub credentials and a fresh clone has none either, so \texttt{pytest} and \texttt{dvc repro} would both fail on the one path every contributor and every pull request takes. The suite's whole premise, stated in \texttt{tests/\allowbreak{}conftest.\allowbreak{}py}, is that it runs on a clean clone with no credentials and no network; this option contradicts it.
\item \textbf{Option B: fall back to a local tracking store when the server is not configured.}
\newline Pros: the run is recorded somewhere, so nothing is lost and the code path is the same.
\newline Cons: since MLflow 3.16 an unset tracking URI resolves to \texttt{sqlite:/\allowbreak{}/\allowbreak{}/\allowbreak{}\$PWD/\allowbreak{}mlflow.\allowbreak{}db}, and for a DVC stage \texttt{\$PWD} is the repository root. The fallback therefore drops a database into the repository on every run of every contributor, gitignored so nothing would ever point it out. \texttt{tests/\allowbreak{}conftest.\allowbreak{}py} already guards \texttt{mlflow.\allowbreak{}db} and \texttt{mlruns/\allowbreak{}} as forbidden paths for exactly this reason, and writing the guard for \texttt{models/\allowbreak{}} caught a real instance of it: a local store puts its \emph{artifacts} in \texttt{.\allowbreak{}/\allowbreak{}mlruns} too.
\item \textbf{Option C (chosen): degrade silently, with \texttt{RECOMMENDITOS\_\allowbreak{}REQUIRE\_\allowbreak{}TRACKING=\allowbreak{}1} as the escape hatch.}
\newline Pros: every contributor and CI can train, which is a property of the whole pipeline rather than of this stage. And the runs whose numbers are cited can be made to fail loudly, so ``the runs are on the server'' is something a person can verify instead of something a reader has to take on trust.
\newline Cons: two behaviours instead of one, and the quiet one is the default. A contributor who \emph{meant} to record a run and mistyped a variable gets a warning line in a long stage log and a model that looks finished. The mitigation is that \texttt{model.\allowbreak{}json} carries \texttt{tracking\_\allowbreak{}mode: \textquotedbl{}disabled\textquotedbl{}} and a null \texttt{run\_\allowbreak{}id}, so the artefact itself says it was not recorded, and \texttt{evaluate} says so again when it finds no run to resume.
\item \textbf{Option D: degrade silently with no escape hatch.}
\newline Pros: one behaviour, nothing to remember.
\newline Cons: it makes the report's claim unfalsifiable. Nobody could demonstrate that a given set of runs reached the server, because the successful and the skipped path are indistinguishable from the outside.
\end{itemize}
\par
\emph{Why:} The two halves of the decision answer two different audiences and neither can be dropped.
\par
Silence is right for the default audience, which is CI and a contributor on a fresh clone. Training with no credentials is not a degraded mode for this project; it is the normal mode for everything except the handful of runs that produce numbers for the report, and a gate that turns red for everyone stops being read.
\par
The escape hatch is right for the other audience, which is whoever produces those numbers. It was not hypothetical: the live verification of this stage against the team's DagsHub server was run with \texttt{RECOMMENDITOS\_\allowbreak{}REQUIRE\_\allowbreak{}TRACKING=\allowbreak{}1} precisely so that a misconfiguration could not let the check pass while logging nothing. Without it the verification would have proved that the stage finishes, which it does in either case, rather than that the runs landed.
\par
The rule that there is never a local fallback is what makes ``no tracking URI means tracking is off'' a statement with content, and it is the same reasoning \texttt{recommenditos/\allowbreak{}config.\allowbreak{}py} already applies to loading \texttt{.\allowbreak{}env} from a named path rather than searching upwards.
\par
One consequence was measured rather than assumed, and it is recorded in EDN-57: detecting ``the server is not reachable'' is not free, because MLflow retries. That is why the retry budget is pinned, and the two decisions have to be read together.},
  ai = {Alternative generation, Alternative assessment, Recommendation, Solution generation},
  response = {Accepted},
  assessment = {AI identified that the existing helper's fail-fast behaviour, which EDN-46 chose deliberately for a credential check, is the wrong behaviour for a pipeline stage, and that the obvious repair -- a local fallback -- would put a database in the repository root under MLflow 3.16's changed default. It proposed the null-object handle so the branch exists once rather than in front of every log call, and the environment variable so the quiet default stays falsifiable. It also used the variable in its own live verification rather than only testing it, which is what showed the hatch earns its place.},
  aievidence = {Claude Code session on 2026-09-30 implementing issue \#37: the DagsHub verification was run with \texttt{RECOMMENDITOS\_\allowbreak{}REQUIRE\_\allowbreak{}TRACKING=\allowbreak{}1} and the throwaway experiment deleted afterwards; the \texttt{mlruns} leak was found by the new \texttt{models/\allowbreak{}} guard in \texttt{tests/\allowbreak{}conftest.\allowbreak{}py} failing.},
  evidence = {\href{https://github.com/mlops-2627q1-mds-upc/MLOPS_Recommenditos/blob/main/recommenditos/tracking.py}{\texttt{recommenditos/\allowbreak{}tracking.\allowbreak{}py}}, \texttt{optional\_\allowbreak{}run}, \texttt{Run} and \texttt{DisabledRun}; \texttt{tests/\allowbreak{}test\_\allowbreak{}model.\allowbreak{}py::test\_\allowbreak{}a\_\allowbreak{}tracking\_\allowbreak{}failure\_\allowbreak{}does\_\allowbreak{}not\_\allowbreak{}fail\_\allowbreak{}training}, \texttt{::test\_\allowbreak{}require\_\allowbreak{}tracking\_\allowbreak{}makes\_\allowbreak{}a\_\allowbreak{}failure\_\allowbreak{}fatal}, \texttt{::test\_\allowbreak{}tracking\_\allowbreak{}is\_\allowbreak{}disabled\_\allowbreak{}without\_\allowbreak{}a\_\allowbreak{}tracking\_\allowbreak{}uri} and \texttt{::test\_\allowbreak{}require\_\allowbreak{}tracking\_\allowbreak{}refuses\_\allowbreak{}a\_\allowbreak{}model\_\allowbreak{}with\_\allowbreak{}no\_\allowbreak{}run\_\allowbreak{}to\_\allowbreak{}resume}; the \texttt{FORBIDDEN\_\allowbreak{}PATHS} guard in \href{https://github.com/mlops-2627q1-mds-upc/MLOPS_Recommenditos/blob/main/tests/conftest.py}{\texttt{tests/\allowbreak{}conftest.\allowbreak{}py}}; \hyperref[EDN-46]{EDN-46}; \hyperref[EDN-57]{EDN-57}; \href{https://github.com/mlops-2627q1-mds-upc/MLOPS_Recommenditos/issues/37}{issue \#37}.},
}

% Generated by tools/edn_latex.py from reports/edn.md.
% Do not edit: change reports/edn.md and run `make edn`.
\ednentry{
  id = {EDN-55},
  title = {\texorpdfstring{\texttt{evaluate}}{evaluate} resumes the run \texorpdfstring{\texttt{train}}{train} created instead of opening its own},
  date = {2026-09-30},
  milestone = {M3: Quality Assurance},
  topic = {Experiment Tracking, Evaluation Protocol},
  participants = {@lukas2510, decided by the repository owner},
  decision = {\texttt{train@\allowbreak{}<variant>} creates the MLflow run and writes its id into \texttt{models/\allowbreak{}<variant>/\allowbreak{}model.\allowbreak{}json}. \texttt{evaluate} reads that id and reopens the same run with \texttt{resume\_\allowbreak{}run}, appending the test metrics and the gate verdict. One experiment holds exactly one run per variant, and no stage opens a run of its own except the stage that first produces one.},
  rationale = {\emph{Alternatives considered:}
\begin{itemize}[leftmargin=*, topsep=2pt]
\item \textbf{Option A: one run per DVC stage, joined by a shared tag.} This is what the course demo does and what its notes warn about.
\newline Pros: the simplest possible wiring; each stage owns what it logs and nothing is shared; no stage depends on another's identifier.
\newline Cons: four variants become eight runs, and the thing anyone actually wants -- the hyperparameters next to the MdAPE they produced -- is in two different rows that a reader has to join by hand. MLflow's compare view compares runs, not pairs of runs, so the one view the rubric asks for (``help you and others compare, understand and extend your results'') cannot be built from it. It also doubles again as soon as \#38 adds energy figures.
\item \textbf{Option B: nested runs, a parent per pipeline state with four children.}
\newline Pros: MLflow supports it natively and it expresses the ladder's structure.
\newline Cons: the parent has no natural owner. The four \texttt{train@\allowbreak{}<variant>} stages are independent DVC stages that can run in any order or alone, so whichever one created the parent would be doing it on behalf of the others, and \texttt{dvc repro train@\allowbreak{}lgbm-\allowbreak{}basic} on its own would either create a parent with one child or reuse a stale one. MLflow's compare view already handles four siblings, so the structure buys nothing it does not cost.
\item \textbf{Option C (chosen): \texttt{train} creates the run, \texttt{evaluate} resumes it.}
\newline Pros: one run per variant carries its hyperparameters, its train and validation loss, its model artefact, \#38's energy figures and \#39's test metrics and gate verdict, so the comparable view is the experiment's own table with no joining. It is also the natural place for the artefact, because the run that holds the metrics holds the model they were measured on.
\newline Cons: \texttt{evaluate} now depends on an identifier produced by another stage, so a model trained with tracking off has no run to append to and its metrics are dropped. That is visible rather than silent -- \texttt{model.\allowbreak{}json} says \texttt{tracking\_\allowbreak{}mode: \textquotedbl{}disabled\textquotedbl{}} and \texttt{run\_\allowbreak{}id: null}, and \texttt{resume\_\allowbreak{}run} logs which variant it could not resume -- and \texttt{RECOMMENDITOS\_\allowbreak{}REQUIRE\_\allowbreak{}TRACKING=\allowbreak{}1} makes it fatal for the runs that matter (EDN-54). A resumed run's recorded duration is also the training run's, not the evaluation's, so run duration in the UI means ``how long the fit took''.
\end{itemize}
\par
\emph{Why:} The demo's own notes flag one-run-per-stage as a wart, and the rubric asks for results that can be compared and understood. Resumption is the only one of the three options where the comparison needs no post-processing, and the cost it adds -- a cross-stage identifier -- is small and is carried in an artefact both stages already depend on.
\par
Two details were raised by the \#39 side of the interface and are part of the decision rather than of the implementation. \texttt{resume\_\allowbreak{}run} must be a context manager that ends the run it opened, because \texttt{evaluate} is one stage covering all four variants in one process and a run left open would collect the next variant's metrics. And a metric a stage could not measure is dropped rather than logged, because SC-04 to SC-06 are \texttt{null} until their machinery exists and MLflow has no representation for ``not measured'' -- so the absence of the key is the record, which is consistent with the metrics artefact recording \texttt{null} and never a fabricated pass.
\par
Verified against the team's real DagsHub server rather than only against a local store: four runs for four variants, each with 9 to 12 parameters, the fit metrics, the four tags this stage sets and the three provenance tags the seam sets, and the whole model bundle uploaded under \texttt{model/\allowbreak{}}. Resuming one of them appended the test metrics to it, left its status \texttt{FINISHED}, correctly skipped a \texttt{null} criterion, and \textbf{did not create a fifth run}. Filtering the experiment by \texttt{tags.\allowbreak{}git\_\allowbreak{}commit} returned exactly those four, which is the comparable view, and the filter string is written into the model card so a grader can reproduce it.
\par
One property of MLflow that the decision does not fix, and that the model card states: there is no upsert by run name, so a second \texttt{dvc repro} of one variant creates a second run rather than replacing the first. The earlier run stays as history, which is what tracking is for, and the view is filtered by the commit of the state being looked at.},
  ai = {Alternative generation, Alternative assessment, Recommendation, Solution generation},
  response = {Accepted},
  assessment = {The resumption design came out of a negotiation between the agents working on issues \#37 and \#39 rather than from one of them, which is also where the context-manager requirement came from -- the \#39 side pointed out that one stage covering four variants would otherwise leak a run across variants. AI then verified the result against the real server instead of asserting it, including the run count before and after resuming, which is the assertion that distinguishes this option from option A.},
  aievidence = {Claude Code sessions on 2026-09-30 designing and implementing issues \#37 and \#39: the seam was agreed between the two over three rounds, and the live DagsHub verification of the run count, the parameters, the tags and the uploaded artefacts was run in a throwaway experiment that was deleted afterwards.},
  evidence = {\href{https://github.com/mlops-2627q1-mds-upc/MLOPS_Recommenditos/blob/main/recommenditos/tracking.py}{\texttt{recommenditos/\allowbreak{}tracking.\allowbreak{}py}}, \texttt{resume\_\allowbreak{}run}; \href{https://github.com/mlops-2627q1-mds-upc/MLOPS_Recommenditos/blob/main/recommenditos/modeling/train.py}{\texttt{recommenditos/\allowbreak{}modeling/\allowbreak{}train.\allowbreak{}py}} and \href{https://github.com/mlops-2627q1-mds-upc/MLOPS_Recommenditos/blob/main/recommenditos/modeling/evaluate.py}{\texttt{recommenditos/\allowbreak{}modeling/\allowbreak{}evaluate.\allowbreak{}py}}; \texttt{tests/\allowbreak{}test\_\allowbreak{}model.\allowbreak{}py::test\_\allowbreak{}one\_\allowbreak{}variant\_\allowbreak{}is\_\allowbreak{}one\_\allowbreak{}run\_\allowbreak{}that\_\allowbreak{}evaluate\_\allowbreak{}appends\_\allowbreak{}to}; \href{https://github.com/mlops-2627q1-mds-upc/MLOPS_Recommenditos/blob/main/docs/docs/model-card.md}{model card}, Training Procedure, Experiment tracking; \hyperref[EDN-54]{EDN-54}; \href{https://github.com/mlops-2627q1-mds-upc/MLOPS_Recommenditos/issues/37}{issue \#37}; \href{https://github.com/mlops-2627q1-mds-upc/MLOPS_Recommenditos/issues/39}{issue \#39}.},
}

% Generated by tools/edn_latex.py from reports/edn.md.
% Do not edit: change reports/edn.md and run `make edn`.
\ednentry{
  id = {EDN-56},
  title = {The model bundle carries its own provenance, so \texorpdfstring{\texttt{models/\allowbreak{}}}{models/} is deliberately not byte-reproducible},
  date = {2026-09-30},
  milestone = {M3: Quality Assurance},
  topic = {Reproducibility, Data Versioning},
  participants = {@lukas2510, decided by the repository owner},
  decision = {\texttt{models/\allowbreak{}<variant>/\allowbreak{}model.\allowbreak{}json} carries \texttt{trained\_\allowbreak{}at} and \texttt{mlflow.\allowbreak{}run\_\allowbreak{}id} alongside the hyperparameters, the feature list and the training statistics. Both change between two otherwise identical fits, so the model directory is \textbf{not} byte-reproducible by design. The reproducibility claim NFR-06 rests on is made about the predictions and about the estimator's own payload file instead, and both are asserted.},
  rationale = {\emph{Alternatives considered:}
\begin{itemize}[leftmargin=*, topsep=2pt]
\item \textbf{Option A (chosen): provenance inside \texttt{model.\allowbreak{}json}.}
\newline Pros: the bundle is self-describing wherever it ends up. EDN-08 bakes it into the API image at build time, so at serving time there is no MLflow client, no \texttt{dvc.\allowbreak{}lock} and no repository -- and \texttt{GET /\allowbreak{}health} still has to be able to say which model version it serves and where its numbers came from. It is also what lets \texttt{evaluate} resume the right run (EDN-55) without a second artefact to keep in step.
\newline Cons: \texttt{dvc repro train@\allowbreak{}<variant>} always produces a new directory hash even when the model is identical, so \texttt{evaluate} always reruns after a retrain, and \texttt{dvc status} can never say ``the model is unchanged''. A reviewer comparing two pipeline states cannot tell a real change from a re-run by the hash alone.
\item \textbf{Option B: keep \texttt{models/\allowbreak{}} byte-stable and move the provenance to a second, non-cached output.}
\newline Pros: the model directory's hash then means exactly ``the model changed'', which is the property a reviewer and \texttt{dvc status} both want, and a re-run that changes nothing reruns nothing downstream.
\newline Cons: it needs a second \texttt{outs} entry per variant, and it puts the run id in a file \texttt{evaluate} and the API would each have to find beside the model rather than in it -- a cross-stage dependency for \#39 and a second file the API image has to be told to copy. It also splits one artefact that is restored together into two that might not be.
\item \textbf{Option C: drop the provenance and rely on MLflow for it.}
\newline Pros: \texttt{models/\allowbreak{}} becomes byte-stable at no structural cost, and MLflow already records the commit, the data version and the parameters of every run.
\newline Cons: it works only where MLflow is reachable. A bundle obtained with \texttt{dvc pull} on a machine with no credentials -- which is the normal case for this project, see EDN-54 -- could then not say which run produced it, and neither could the running API. It would also break the resumption of EDN-55 entirely, because the id would live only on the server whose runs it is meant to identify.
\end{itemize}
\par
\emph{Why:} This is a trade against NFR-06 and it is worth being exact about which half of NFR-06 it trades. NFR-06 asks for two things: that the documented steps from a clean clone reproduce the same splits and metrics within 0.1 percentage points, and that \textbf{every training run records which code, which data version and which parameters produced it}. The provenance fields serve the second half directly. What they cost is a property NFR-06 never asks for, which is that the output directory hash is stable.
\par
The reproducibility that is claimed is therefore claimed about the right things and is tested. Two fits of the same variant on the same rows give bit-identical predictions for all four variants, and a byte-identical payload file -- \texttt{booster.\allowbreak{}txt}, \texttt{pipeline.\allowbreak{}joblib} or \texttt{lookup.\allowbreak{}parquet}, which is the whole of the fitted model. \texttt{model.\allowbreak{}json} is excluded from that comparison explicitly, in the test and in its docstring, so nobody reads the exclusion as an oversight.
\par
The cost is real and is accepted rather than argued away: every \texttt{dvc repro} of a \texttt{train} stage reruns \texttt{evaluate}, even when the model is identical. At the measured scale that is about three seconds for the evaluation, against the 43 seconds of the fit that preceded it, so it changes nothing about how the pipeline is used. If the evaluation ever becomes expensive -- the SC-06 masking sweep of \#39 is the candidate, at roughly 25 predict calls per variant -- option B becomes worth revisiting, and this entry is the record of what it would buy.},
  ai = {Alternative generation, Alternative assessment, Recommendation},
  response = {Accepted},
  assessment = {AI raised the trade as an explicit decision rather than letting a timestamp drift into an artefact unremarked, which is the failure mode this repository has already had once with a test run's validation summary. It named the alternative that preserves byte stability and the cost of that alternative, and it separated the reproducibility that is claimed from the reproducibility that is not, then wrote the test so that the exclusion of \texttt{model.\allowbreak{}json} is visible in the assertion instead of being silently absent.},
  aievidence = {Claude Code session on 2026-09-30 implementing issue \#37: two fits of each variant were compared byte for byte in a scratchpad before the test was written, and the payload files were confirmed identical while \texttt{model.\allowbreak{}json} differed.},
  evidence = {\href{https://github.com/mlops-2627q1-mds-upc/MLOPS_Recommenditos/blob/main/recommenditos/modeling/model.py}{\texttt{recommenditos/\allowbreak{}modeling/\allowbreak{}model.\allowbreak{}py}}, \texttt{fit\_\allowbreak{}variant} and \texttt{Model.\allowbreak{}save}; \texttt{tests/\allowbreak{}test\_\allowbreak{}model.\allowbreak{}py::test\_\allowbreak{}the\_\allowbreak{}payload\_\allowbreak{}of\_\allowbreak{}a\_\allowbreak{}refit\_\allowbreak{}is\_\allowbreak{}byte\_\allowbreak{}identical} and \texttt{::test\_\allowbreak{}training\_\allowbreak{}is\_\allowbreak{}reproducible}; \href{https://github.com/mlops-2627q1-mds-upc/MLOPS_Recommenditos/blob/main/docs/docs/requirements.md}{requirements} NFR-06; \hyperref[EDN-08]{EDN-08}; \hyperref[EDN-55]{EDN-55}; \href{https://github.com/mlops-2627q1-mds-upc/MLOPS_Recommenditos/issues/37}{issue \#37}.},
}

% Generated by tools/edn_latex.py from reports/edn.md.
% Do not edit: change reports/edn.md and run `make edn`.
\ednentry{
  id = {EDN-57},
  title = {The MLflow retry budget is pinned to three requests},
  date = {2026-09-30},
  milestone = {M3: Quality Assurance},
  topic = {Experiment Tracking, CI/CD},
  participants = {@lukas2510, decided by the repository owner},
  decision = {\texttt{recommenditos/\allowbreak{}tracking.\allowbreak{}py} sets \texttt{MLFLOW\_\allowbreak{}HTTP\_\allowbreak{}REQUEST\_\allowbreak{}MAX\_\allowbreak{}RETRIES} to \texttt{3} with \texttt{setdefault}, rather than leaving MLflow's default of 5. A caller who wants MLflow's own budget can still ask for it by setting the variable.},
  rationale = {\emph{Alternatives considered:}
\begin{itemize}[leftmargin=*, topsep=2pt]
\item \textbf{Option A: leave MLflow's default of 5 retries.}
\newline Pros: no deviation to explain, and the most patience with a flaky connection.
\newline Cons: measured against a closed port, giving up takes \textbf{246 s}. Because tracking degrades rather than failing (EDN-54), that is 246 seconds a stage spends before training normally, once per stage. Over the four-variant ladder it is \textbf{16.4 minutes of waiting on a server that is not there}, on top of the 67 seconds the fits actually take.
\item \textbf{Option B (chosen): 3 retries.}
\newline Pros: measured at \textbf{13.7 s} to give up, so the four-variant ladder loses under a minute instead of a quarter of an hour, and three attempts still ride out a connection that drops a packet or two rather than being down. \texttt{setdefault} keeps the decision overridable in the one case where patience is worth more than latency, which is a large artefact upload over a bad line.
\newline Cons: a deviation from the library's default, so an upload that would have succeeded on the fourth attempt now fails and that run's artefact is missing while its metrics are present. It is a warning in the log rather than a failure, which is the same quiet path EDN-54 chose and inherits the same mitigation.
\item \textbf{Option C: 0 or 1 retries.}
\newline Pros: essentially free to detect an unreachable server -- measured at 0.00 s and 0.12 s -- which is what the tests use.
\newline Cons: it makes a single dropped packet lose a run's tracking. For the runs whose numbers are cited that is the opposite of what is wanted, and those are exactly the runs that upload a 6.8 MB booster.
\item \textbf{Option D: probe the server with a short bounded request before configuring MLflow.}
\newline Pros: the fastest possible detection, with MLflow's own budget left intact for the calls that matter.
\newline Cons: a hand-rolled health check that duplicates what the client already does, has to know DagsHub's URL shape, and can disagree with the client about whether the server is usable. More code and one more thing that can be wrong.
\end{itemize}
\par
\emph{Why:} The numbers make the choice, and they are steeply non-linear because MLflow backs off exponentially. Measured against \texttt{http:/\allowbreak{}/\allowbreak{}127.\allowbreak{}0.\allowbreak{}0.\allowbreak{}1:1/\allowbreak{}}, a port that is closed so the connection is refused rather than left hanging: 5 retries take 246.48 s, 3 take 13.68 s, 2 take 4.33 s, 1 takes 0.12 s and 0 takes 0.00 s. Three is the knee of that curve -- an order of magnitude cheaper than the default while still retrying.
\par
One thing this entry deliberately does \textbf{not} claim, because it would be an overstatement. NFR-10 budgets 15 minutes for ``training the chosen configuration'', singular. At the default budget the chosen configuration alone would take about 4.8 minutes -- 246 s of waiting plus a 43 s fit -- which is inside NFR-10, although 85\,\% of it would be spent on a server that is not there. It is the \textbf{whole ladder} that would exceed the figure, at about 17.5 minutes for four variants, and the ladder is not what the requirement measures. The honest statement is that the default makes \texttt{dvc repro} unusable in practice and distorts every wall-clock number anyone reports, not that it breaks NFR-10.
\par
The variable is set beside the existing \texttt{MLFLOW\_\allowbreak{}DISABLE\_\allowbreak{}AGENT\_\allowbreak{}HINT} and for the same reason: both are MLflow environment settings that this project has a view on, and one place that sets them is better than each contributor discovering them. Both use \texttt{setdefault}, so neither takes the choice away.},
  ai = {Information seeking, Alternative assessment, Recommendation},
  response = {Accepted with modifications},
  assessment = {The 246-second figure came out of a measurement AI ran because a test against an unreachable server had not returned, rather than from the design, which had listed ``does a bogus URI fail fast?'' as an open question. AI swept the budget to find the knee instead of picking a number. Its first write-up claimed the default would eat NFR-10's training budget; checking the requirement's wording showed it budgets one configuration rather than the ladder, and the entry now says what the numbers support and explicitly disclaims the stronger version.},
  aievidence = {Claude Code session on 2026-09-30 implementing issue \#37: the retry sweep against a closed port produced 246.48 s, 13.68 s, 4.33 s, 0.12 s and 0.00 s for 5, 3, 2, 1 and 0 retries, and the same measurement is why the tests pin the budget to 0.},
  evidence = {\href{https://github.com/mlops-2627q1-mds-upc/MLOPS_Recommenditos/blob/main/recommenditos/tracking.py}{\texttt{recommenditos/\allowbreak{}tracking.\allowbreak{}py}}, the \texttt{MLFLOW\_\allowbreak{}HTTP\_\allowbreak{}REQUEST\_\allowbreak{}MAX\_\allowbreak{}RETRIES} default; \texttt{tests/\allowbreak{}test\_\allowbreak{}model.\allowbreak{}py::\_\allowbreak{}point\_\allowbreak{}at\_\allowbreak{}an\_\allowbreak{}unreachable\_\allowbreak{}server}; \href{https://github.com/mlops-2627q1-mds-upc/MLOPS_Recommenditos/blob/main/docs/docs/requirements.md}{requirements} NFR-10; \hyperref[EDN-54]{EDN-54}; \href{https://github.com/mlops-2627q1-mds-upc/MLOPS_Recommenditos/issues/37}{issue \#37}.},
}

% Generated by tools/edn_latex.py from reports/edn.md.
% Do not edit: change reports/edn.md and run `make edn`.
\ednentry{
  id = {EDN-58},
  title = {An unmeasured success criterion blocks the gate, and the artefact says how close the model is anyway},
  date = {2026-09-30},
  milestone = {M3: Quality Assurance},
  topic = {Testing Strategy, Model Evaluation (the NFR-01 gate)},
  participants = {@lukas2510},
  decision = {A criterion the \texttt{evaluate} stage could not measure is recorded as \texttt{null}, never as a pass and never as a failure, and \texttt{null} blocks NFR-01's gate. Two criteria are in that state today: SC-05, because no model exposes \texttt{predict\_\allowbreak{}interval\_\allowbreak{}eur} and the UC2 conformal intervals are a later ticket, and SC-04 whenever no segment level reaches the 500-row minimum. Alongside the blocking verdict, \texttt{metrics.\allowbreak{}json} carries two leaves that say what is outstanding and how close the candidates are without it: \texttt{criteria\_\allowbreak{}not\_\allowbreak{}measured} names the criteria, and \texttt{n\_\allowbreak{}variants\_\allowbreak{}passing\_\allowbreak{}measurable} counts the variants where nothing that \emph{could} be measured failed. Each criterion also carries a structured status rather than only prose: \texttt{sc05\_\allowbreak{}status: \textquotedbl{}not\_\allowbreak{}measured\textquotedbl{}}, \texttt{sc05\_\allowbreak{}capability\_\allowbreak{}checked: \textquotedbl{}predict\_\allowbreak{}interval\_\allowbreak{}eur\textquotedbl{}}, \texttt{sc05\_\allowbreak{}reason: \textquotedbl{}.\allowbreak{}.\allowbreak{}.\allowbreak{}\textquotedbl{}}, and \texttt{sc04\_\allowbreak{}note} naming the largest level and its row count.},
  rationale = {\emph{Alternatives considered:}
\begin{itemize}[leftmargin=*, topsep=2pt]
\item \textbf{Option A (chosen): block, with a structured \texttt{null} and a second, clearly named verdict.}
\newline Pros: NFR-01's word is ``every'', and this is the only reading of it that survives an unbuilt criterion. A reader of the artefact can tell ``we did not measure it'' from ``the model missed it'' without parsing English, which is what EDN-12 needs, since promotion is a human decision and the human has to see what is outstanding. \texttt{n\_\allowbreak{}variants\_\allowbreak{}passing\_\allowbreak{}measurable} keeps the first delivery's metrics file informative: \texttt{n\_\allowbreak{}variants\_\allowbreak{}passing: 0} says nothing about the model, and this says how close the ladder is.
\newline Cons: two leaves and a status string per criterion that a one-criterion-at-a-time reader does not need, and a second verdict that a careless reader could quote as if it were the gate.
\item \textbf{Option B: block, with \texttt{null} and nothing else.}
\newline Pros: the narrowest possible \texttt{metrics.\allowbreak{}json}, which is what the existing comment in \texttt{dvc.\allowbreak{}yaml} was written to defend.
\newline Cons: the first delivery's artefact then reads ``0 of 4 variants pass'' with no explanation anywhere in it, and the reason is a ticket that has not been written rather than anything about the models. That is the number the report would have to cite, and it would have to be explained in prose beside the artefact instead of in it.
\item \textbf{Option C: exclude an unmeasurable criterion from the gate and pass on the rest.}
\newline Pros: the gate would produce a deployable candidate now, which is what a demo wants.
\newline Cons: unacceptable. It would let a model pass NFR-01 without any interval ever being calibrated, which is precisely the promise the requirement makes. It also gets easier to do the more criteria are outstanding, so the gate would weaken exactly when it matters most.
\item \textbf{Option D (for SC-04 specifically): treat ``every level satisfies P'' as met when no level qualifies.}
\newline Pros: vacuously true, and it is what \texttt{all([])} returns, so it is what the obvious implementation does.
\newline Cons: it reports a check nobody ran as a check that passed. On the 2,000-row test fixture no level reaches 500 rows, so this option would have reported SC-04 as met for every variant in every CI run.
\end{itemize}
\par
\emph{Why:} Option C is the only one that is actually wrong, and the choice between A and B is about whether the artefact has to be readable on its own. It does: EDN-12 makes promotion a pull-request review, and the thing the reviewer reads is \texttt{dvc metrics diff}, which is long-format and lists only the leaves that changed. A row reading \texttt{criteria\_\allowbreak{}not\_\allowbreak{}measured: sc05 -\allowbreak{}> \textquotedbl{}\textquotedbl{}} is the single most useful line that diff can carry, and it does not exist under option B. Option D was rejected on a measurement rather than on taste: at the committed fixture size zero segment levels clear the 500-row bar, so the vacuous reading would have been the one CI exercised, and an implementation of SC-04 that simply returned ``passed'' would have kept the suite green. The test suite answers that by running \texttt{evaluate} a second time with \texttt{sc04\_\allowbreak{}min\_\allowbreak{}segment\_\allowbreak{}rows} lowered to 50, which produces a real verdict on the fixture; growing \texttt{FIXTURE\_\allowbreak{}ROWS} from 2,000 to about 20,000 so the project's own bar is exercised was considered and rejected, because it multiplies the cost of the session fixture across the whole suite for a property the parameterised run already covers more precisely, and the fixture size is itself a recorded decision (EDN-33).},
  ai = {Alternative generation, Alternative assessment, Recommendation, Solution generation},
  response = {Accepted with modifications},
  assessment = {AI laid out the four readings of an unmeasurable criterion and recommended A, with the argument that option C weakens the gate more the more criteria are outstanding. The part that was worth more than the recommendation was the measurement behind option D: AI ran the split and the segment cut at three data scales before writing any of the stage, found that no level clears 500 rows on the pytest fixture, and concluded that the vacuous reading would be the one CI exercised and that a test asserting \texttt{sc04\_\allowbreak{}passed is None} there would be satisfied by an implementation that returned \texttt{None} unconditionally. That is the trap the second \texttt{evaluate} run in \texttt{tests/\allowbreak{}test\_\allowbreak{}pipeline.\allowbreak{}py} exists for, and it came out of the measurement rather than out of the design. The modification was to the leaf count: an earlier proposal also put \texttt{sc04\_\allowbreak{}worst\_\allowbreak{}segment} and the per-field sweep into \texttt{metrics.\allowbreak{}json}, which was moved to \texttt{reports/\allowbreak{}metrics/\allowbreak{}<variant>.\allowbreak{}json} to keep \texttt{dvc metrics show} at 24 columns.},
  aievidence = {Claude Code sessions on 2026-09-30 designing and then implementing issue \#39. The design session measured the split and segment sizes at 416, 3,695 and 22,551 test rows and reported ``SC-04's 500-row rule is not measurable at all on the 2,000-row pytest fixture; zero levels qualify, the largest is \texttt{seller\_\allowbreak{}type=\allowbreak{}Dealer} at 345 rows''; it also validated the proposed \texttt{metrics.\allowbreak{}json} shape against \texttt{dvc metrics show} and \texttt{dvc metrics diff} in a throwaway DVC repository before the shape was fixed.},
  evidence = {\href{https://github.com/mlops-2627q1-mds-upc/MLOPS_Recommenditos/blob/main/recommenditos/modeling/evaluate.py}{\texttt{recommenditos/\allowbreak{}modeling/\allowbreak{}evaluate.\allowbreak{}py}}, \texttt{evaluate\_\allowbreak{}sc04}, \texttt{evaluate\_\allowbreak{}sc05}, \texttt{gate\_\allowbreak{}passed}, \texttt{gate\_\allowbreak{}blocked\_\allowbreak{}by} and \texttt{\_\allowbreak{}check\_\allowbreak{}every\_\allowbreak{}criterion\_\allowbreak{}is\_\allowbreak{}decided}; \texttt{tests/\allowbreak{}test\_\allowbreak{}model.\allowbreak{}py::test\_\allowbreak{}sc04\_\allowbreak{}is\_\allowbreak{}not\_\allowbreak{}measured\_\allowbreak{}rather\_\allowbreak{}than\_\allowbreak{}vacuously\_\allowbreak{}met\_\allowbreak{}when\_\allowbreak{}no\_\allowbreak{}level\_\allowbreak{}qualifies}, \texttt{::test\_\allowbreak{}an\_\allowbreak{}unmeasured\_\allowbreak{}criterion\_\allowbreak{}blocks\_\allowbreak{}the\_\allowbreak{}gate\_\allowbreak{}rather\_\allowbreak{}than\_\allowbreak{}passing\_\allowbreak{}it}, \texttt{::test\_\allowbreak{}a\_\allowbreak{}criterion\_\allowbreak{}added\_\allowbreak{}to\_\allowbreak{}the\_\allowbreak{}list\_\allowbreak{}without\_\allowbreak{}a\_\allowbreak{}measurement\_\allowbreak{}fails\_\allowbreak{}the\_\allowbreak{}stage}; \texttt{tests/\allowbreak{}test\_\allowbreak{}pipeline.\allowbreak{}py::test\_\allowbreak{}sc04\_\allowbreak{}is\_\allowbreak{}measured\_\allowbreak{}as\_\allowbreak{}soon\_\allowbreak{}as\_\allowbreak{}a\_\allowbreak{}segment\_\allowbreak{}level\_\allowbreak{}is\_\allowbreak{}large\_\allowbreak{}enough}; \href{https://github.com/mlops-2627q1-mds-upc/MLOPS_Recommenditos/blob/main/docs/docs/problem-spec.md}{problem specification} section 8; \hyperref[EDN-06]{EDN-06}, \hyperref[EDN-12]{EDN-12}, \hyperref[EDN-33]{EDN-33}; \href{https://github.com/mlops-2627q1-mds-upc/MLOPS_Recommenditos/issues/39}{issue \#39}.},
}

% Generated by tools/edn_latex.py from reports/edn.md.
% Do not edit: change reports/edn.md and run `make edn`.
\ednentry{
  id = {EDN-59},
  title = {SC-06 masks API input fields, not feature columns, and an absent field's value follows EDN-23},
  date = {2026-09-30},
  milestone = {M3: Quality Assurance},
  topic = {Model Evaluation (SC-06), Feature Encoding},
  participants = {@lukas2510},
  decision = {SC-06's masking sweep is defined over the \textbf{API input fields} of the evaluated variant's own feature set, not over the model's feature columns. Every input field is mapped to the column or columns it produced, so \texttt{registration\_\allowbreak{}date} covers the derived \texttt{age\_\allowbreak{}years} and one equipment field covers every multi-hot column built from it; the optional set is then that mapping minus \texttt{evaluate.\allowbreak{}required\_\allowbreak{}input\_\allowbreak{}fields}, which is 6 fields for the basic set and 23 for the extended one. What ``masked'' means is decided per field from the contract rather than being one missing value for everything, and it follows EDN-23: an assertion flag becomes \texttt{False}, an equipment multi-hot column becomes \texttt{False}, a categorical becomes a null code over the declared levels, and a number becomes \texttt{NaN}. Any other encoding raises, naming the column and its dtype. The mask is applied to the feature matrix rather than to the raw frame, and the artefact records which fields were masked so the approximation is visible.
\par
\textbf{Scope, after the review of pull request \#62:} what an omitted equipment list is was re-opened there and is recorded by \hyperref[EDN-61]{EDN-61}, not here. It settled as \textbf{the empty list, every derived column \texttt{False}}, which is the value this entry was already written against and the one its tests assert, so nothing measured here changed. This entry therefore cites EDN-61 for that value and keeps what is its own: SC-06 defined over input fields rather than feature columns, the mapping from an input field to the columns it produced, the per-dtype rule for the other three kinds of column, and the record below of why nullable-boolean masking was rejected for the three condition flags. \texttt{Model.\allowbreak{}mask\_\allowbreak{}absent} is the single implementation of the rule and the masking sweep calls it.},
  rationale = {\emph{Alternatives considered:}
\begin{itemize}[leftmargin=*, topsep=2pt]
\item \textbf{Option A (chosen): input fields, mapped to their columns, masked in the feature matrix, absence per EDN-23.}
\newline Pros: it measures the request FR-01 actually permits. A caller cannot leave out one equipment \emph{item}, so masking one multi-hot column measures something the API can never be asked. The mapping also makes the sweep cost stable across the encoding: 33 input fields and 23 optional ones for the extended set whether the equipment columns number 16 (fixture) or 133 (real snapshot), where a column-wise sweep would have grown from 24 predict calls to 162. And it removes the defect the naive reading has: \texttt{age\_\allowbreak{}years} is not in \texttt{required\_\allowbreak{}input\_\allowbreak{}fields}, so a plain set difference over feature columns makes the single most important feature in the model look optional and SC-06 would have reported the cost of hiding a car's age.
\newline Cons: it needs one explicit mapping in \texttt{params.\allowbreak{}yaml} (\texttt{input\_\allowbreak{}field\_\allowbreak{}to\_\allowbreak{}feature}), which is a line somebody has to maintain, and the mask is applied to the matrix rather than by re-running the feature build, so for a field whose encoding is not one-to-one it is an approximation of the serving path.
\item \textbf{Option B: mask feature columns, one at a time.}
\newline Pros: no mapping, no parameter, and the sweep is derivable from \texttt{model.\allowbreak{}features} alone.
\newline Cons: measures requests the API cannot make, masks \texttt{age\_\allowbreak{}years} as if it were optional, and its cost and its result both change when issue \#36's encoding changes, so the criterion would not be comparable across feature-set revisions.
\item \textbf{Option C: mask the raw processed frame and re-run the feature build per scenario.}
\newline Pros: literally the serving path, so exact for every field including the equipment ones.
\newline Cons: needs a reusable transform the \texttt{features} stage does not expose yet, and it is 25 feature builds per variant instead of 25 predict calls. Worth revisiting when the API needs that same seam.
\item \textbf{Option D: mask an assertion flag to \texttt{pd.\allowbreak{}NA} on a nullable \texttt{boolean} column, so the criterion measures absence rather than a denial.}
\newline Pros: the shape a careful reviewer expects, and it distinguishes ``not asserted'' from ``asserted false''.
\newline Cons: it is the encoding EDN-23 examined and rejected, for the reason that in this data there is no third state to encode: a listing form offers a checkbox, a tick is evidence and an empty box is the absence of evidence, so \texttt{False} \textbf{is} absence. It also does not work: the three flags are non-nullable in the persisted contract, so \texttt{Schema.\allowbreak{}conform} refuses a null one and \texttt{predict\_\allowbreak{}eur} raises rather than predicting. Reversing EDN-23 for one criterion would have made the pipeline and the API disagree about what an omitted flag is.
\end{itemize}
\par
\emph{Why:} The field-level definition is the one that matches the requirement: FR-01 is a statement about a request, so a criterion about ``a missing optional field'' has to be measured one request at a time. The two assertions that make the mapping safe are what justify keeping it in code rather than in a document: every required input field must map to at least one column of the trained model, and no column may be claimed by two fields, so a params edit that breaks either fails the stage instead of quietly changing which fields the criterion covers. Option D was chosen first and then withdrawn on re-reading EDN-23; the deciding facts were that the flags are declared non-nullable, so \texttt{pd.\allowbreak{}NA} raises rather than being measured, and that the regression test for the alternative encoding demonstrates what would have happened instead: assigning \texttt{NaN} into a \texttt{bool} column on pandas 3 replaces it with a float64 column of \texttt{NaN}, on which \texttt{.\allowbreak{}all()} returns \texttt{True}, so a consumer reading it as a flag sees ``asserted for every row'' and the column is destroyed rather than masked. The mask therefore derives its value from the contract's dtype and refuses an encoding it does not recognise, so an encoding choice nobody anticipated fails loudly instead of masking the wrong way. Option A's approximation is recorded rather than hidden: for every field of the basic set the encoding is one-to-one, so A and C are identical there, and the difference is confined to the equipment fields.},
  ai = {Alternative generation, Alternative assessment, Recommendation, Solution generation},
  response = {Accepted with modifications},
  assessment = {AI found the \texttt{age\_\allowbreak{}years} defect in the written criterion before any code was written, by computing the naive set difference and noticing it yields seven optional fields for the basic set where EDN-15 states six; two independent sessions found it separately, which is the cross-check that it is a defect in the specification rather than a reading of it. AI also proposed option D and wrote it into the design as settled, and the reversal to option A came from checking it against EDN-23 and against the persisted contract, where it fails outright. That is the part of the contribution that was rejected, and it is worth recording as rejected: the argument for \texttt{pd.\allowbreak{}NA} was about what a reviewer would expect to see rather than about this data, and EDN-23 had already answered it with measurements. The per-dtype rule, the two mapping assertions and the ``raise on an encoding you do not recognise'' rule were accepted as proposed.},
  aievidence = {Claude Code sessions on 2026-09-30 on issues \#37 and \#39. The \#39 design session reported ``the naive set difference yields 7 optional fields for \texttt{basic} where EDN-15 says 6, and the spurious one is \texttt{age\_\allowbreak{}years}, the single most important feature in the model'', and separately measured the mapping against a simulated post-\#36 matrix to show the optional-field count is stable at 23. The same session recommended the nullable-boolean masking of option D; the implementation session checked it against EDN-23 and against \texttt{Schema.\allowbreak{}\_\allowbreak{}null\_\allowbreak{}problems}, found that it raises \texttt{SchemaError}, and implemented \texttt{False}.},
  evidence = {\href{https://github.com/mlops-2627q1-mds-upc/MLOPS_Recommenditos/blob/main/recommenditos/modeling/evaluate.py}{\texttt{recommenditos/\allowbreak{}modeling/\allowbreak{}evaluate.\allowbreak{}py}}, \texttt{input\_\allowbreak{}field\_\allowbreak{}columns}, \texttt{optional\_\allowbreak{}input\_\allowbreak{}fields}, \texttt{sc06\_\allowbreak{}scenarios} and \texttt{masking\_\allowbreak{}sweep}; \href{https://github.com/mlops-2627q1-mds-upc/MLOPS_Recommenditos/blob/main/recommenditos/modeling/model.py}{\texttt{recommenditos/\allowbreak{}modeling/\allowbreak{}model.\allowbreak{}py}}, \texttt{Model.\allowbreak{}mask\_\allowbreak{}absent} and \texttt{Model.\allowbreak{}columns\_\allowbreak{}of}, which this entry defers to for the value and calls for the expansion; \texttt{params.\allowbreak{}yaml}, \texttt{evaluate.\allowbreak{}input\_\allowbreak{}field\_\allowbreak{}to\_\allowbreak{}feature}; \texttt{tests/\allowbreak{}test\_\allowbreak{}model.\allowbreak{}py::test\_\allowbreak{}the\_\allowbreak{}optional\_\allowbreak{}fields\_\allowbreak{}of\_\allowbreak{}the\_\allowbreak{}basic\_\allowbreak{}set\_\allowbreak{}are\_\allowbreak{}the\_\allowbreak{}six\_\allowbreak{}of\_\allowbreak{}edn15}, \texttt{::test\_\allowbreak{}an\_\allowbreak{}equipment\_\allowbreak{}field\_\allowbreak{}covers\_\allowbreak{}every\_\allowbreak{}multi\_\allowbreak{}hot\_\allowbreak{}column\_\allowbreak{}it\_\allowbreak{}produced}, \texttt{::test\_\allowbreak{}every\_\allowbreak{}feature\_\allowbreak{}column\_\allowbreak{}is\_\allowbreak{}claimed\_\allowbreak{}by\_\allowbreak{}exactly\_\allowbreak{}one\_\allowbreak{}input\_\allowbreak{}field} and \texttt{::test\_\allowbreak{}mask\_\allowbreak{}absent\_\allowbreak{}is\_\allowbreak{}the\_\allowbreak{}one\_\allowbreak{}answer\_\allowbreak{}to\_\allowbreak{}what\_\allowbreak{}an\_\allowbreak{}omitted\_\allowbreak{}field\_\allowbreak{}is}; \href{https://github.com/mlops-2627q1-mds-upc/MLOPS_Recommenditos/blob/main/docs/docs/problem-spec.md}{problem specification} section 8 and \href{https://github.com/mlops-2627q1-mds-upc/MLOPS_Recommenditos/blob/main/docs/docs/specification.md}{specification} FR-01, both reworded by this decision; \hyperref[EDN-15]{EDN-15}, \hyperref[EDN-23]{EDN-23}; \href{https://github.com/mlops-2627q1-mds-upc/MLOPS_Recommenditos/issues/39}{issue \#39}.},
}

% Generated by tools/edn_latex.py from reports/edn.md.
% Do not edit: change reports/edn.md and run `make edn`.
\ednentry{
  id = {EDN-60},
  title = {Which segments may gate is enforced in code that fails, not by a rule in a document},
  date = {2026-09-30},
  milestone = {M3: Quality Assurance},
  topic = {Model Evaluation (SC-04), Testing Strategy},
  participants = {@lukas2510},
  decision = {The three rules that decide which segment levels SC-04 is allowed to gate on are checks that fail the stage rather than conventions a reader has to remember. A module constant names the segments that condition on the target (\texttt{price\_\allowbreak{}bucket}), and adding one to \texttt{evaluate.\allowbreak{}segments} raises; the row builder reads that constant itself, so no caller can produce a price-bucket row that counts. A segment over an input field FR-01 does not require also raises, because the rule that excludes a \texttt{\textquotedbl{}(missing)\textquotedbl{}} level is only correct for a required field. And a per-make level naming a make the model was not fitted on raises, so EDN-48's third obligation is a checked property rather than an intention. Every excluded level stays in the exported table with its metrics, its row count and the name of the rule that excluded it.},
  rationale = {\emph{Alternatives considered:}
\begin{itemize}[leftmargin=*, topsep=2pt]
\item \textbf{Option A (chosen): a constant plus three refusals, and every exclusion recorded in the table.}
\newline Pros: the rule lives where it is applied, so it cannot be forgotten by someone editing \texttt{params.\allowbreak{}yaml} a month from now, and each refusal names the reason rather than the symptom. Keeping the excluded rows in the table with \texttt{excluded\_\allowbreak{}because} means the report can show that \texttt{Volkswagen} had 81 test rows and \texttt{fuel\_\allowbreak{}category=\allowbreak{}(missing)} had 5, which is the difference between a breakdown a reader can audit and one that silently omits what did not fit.
\newline Cons: more code than a comment, and three error paths that only fire on a configuration nobody has written yet.
\item \textbf{Option B: document the rules in the problem specification and rely on review.}
\newline Pros: no code, and the specification is where the criterion is defined anyway.
\newline Cons: the failure it has to prevent is a one-line \texttt{params.\allowbreak{}yaml} edit that produces a passing gate. SC-04's worst qualifying level for \texttt{lgbm-\allowbreak{}extended} is 17.8\,\%; its worst price bucket is 22.3\,\%, so adding \texttt{price\_\allowbreak{}bucket} to the segment list changes the number the gate blocks on while looking like a reporting improvement. A rule that is only written down is checked exactly as often as somebody reads it.
\item \textbf{Option C: drop the price buckets from the stage entirely.}
\newline Pros: nothing to exclude, so nothing to enforce.
\newline Cons: they are the diagnostic that shows \emph{why} they may not gate. Measured on the real snapshot, their profile is U-shaped (22.3\,\% under 5,000 EUR and 6.3\,\% over 80,000 EUR against 5.5\,\% in the 40,000 to 80,000 EUR bucket), which is the regression-to-the-mean artefact of conditioning on the target. Without the table the exclusion is an assertion; with it, it is evidence.
\item \textbf{Option D: exclude a \texttt{\textquotedbl{}(missing)\textquotedbl{}} level for every segment, without checking that the field is required.}
\newline Pros: one rule, no second check.
\newline Cons: it would be wrong the first time somebody segments by an optional field. For a required field the absence cannot occur at serving time, because FR-01 answers such a request with a 422; for an optional field it is a case the API answers every day, and excluding it would hide exactly the rows worth looking at.
\end{itemize}
\par
\emph{Why:} All three rules share one shape: a statement about what the criterion means that would otherwise be maintained by discipline. The M3 rubric rewards exactly this distinction, and the measurement makes it concrete rather than stylistic: the price buckets would change SC-04's verdict if they gated, so the difference between option A and option B is a number in the report. Option D was rejected on the same grounds as the price buckets: it is right today only because all five configured segments happen to be required fields, and an accidental correctness is what the check turns into a guaranteed one.},
  ai = {Alternative generation, Alternative assessment, Recommendation, Solution generation},
  response = {Accepted with modifications},
  assessment = {AI proposed the target-conditioned constant and the ``absence of a required field'' exclusion, and measured the price-bucket profile that shows why the exclusion matters rather than asserting that it does. Two modifications were made to what it proposed. Its construction passed an \texttt{eligible=\allowbreak{}False} flag into the row builder for the price-bucket rows, which leaves the exclusion dependent on the caller passing the flag; the flag was removed and the builder now reads the segment's own name, so the exclusion holds however the rows are built. And it did not propose the third check at all: the requirement that a segment be a required input field came from asking what makes the \texttt{\textquotedbl{}(missing)\textquotedbl{}} exclusion correct, which turned out to be an assumption nothing enforced.},
  aievidence = {Claude Code session on 2026-09-30 designing issue \#39, which measured the price-bucket breakdown at 22,551 test rows and reported ``the U shape is the regression-to-the-mean artefact of conditioning on the target, not a fairness finding, and the 490k MAE in the top bucket would dominate any aggregate''; and the implementation session on the same day, which replaced the \texttt{eligible} flag with the name-based rule while reducing the row builder's argument count.},
  evidence = {\href{https://github.com/mlops-2627q1-mds-upc/MLOPS_Recommenditos/blob/main/recommenditos/modeling/evaluate.py}{\texttt{recommenditos/\allowbreak{}modeling/\allowbreak{}evaluate.\allowbreak{}py}}, \texttt{TARGET\_\allowbreak{}CONDITIONED\_\allowbreak{}SEGMENTS}, \texttt{criterion\_\allowbreak{}segments}, \texttt{\_\allowbreak{}sc04\_\allowbreak{}eligibility} and \texttt{check\_\allowbreak{}make\_\allowbreak{}levels\_\allowbreak{}are\_\allowbreak{}served}; \texttt{tests/\allowbreak{}test\_\allowbreak{}model.\allowbreak{}py::test\_\allowbreak{}a\_\allowbreak{}target\_\allowbreak{}conditioned\_\allowbreak{}segment\_\allowbreak{}cannot\_\allowbreak{}become\_\allowbreak{}a\_\allowbreak{}criterion}, \texttt{::test\_\allowbreak{}a\_\allowbreak{}segment\_\allowbreak{}over\_\allowbreak{}an\_\allowbreak{}optional\_\allowbreak{}input\_\allowbreak{}field\_\allowbreak{}is\_\allowbreak{}refused}, \texttt{::test\_\allowbreak{}a\_\allowbreak{}price\_\allowbreak{}bucket\_\allowbreak{}row\_\allowbreak{}is\_\allowbreak{}reported\_\allowbreak{}and\_\allowbreak{}cannot\_\allowbreak{}count}, \texttt{::test\_\allowbreak{}sc04\_\allowbreak{}refuses\_\allowbreak{}rows\_\allowbreak{}from\_\allowbreak{}a\_\allowbreak{}segment\_\allowbreak{}it\_\allowbreak{}may\_\allowbreak{}not\_\allowbreak{}gate}, \texttt{::test\_\allowbreak{}a\_\allowbreak{}per\_\allowbreak{}make\_\allowbreak{}level\_\allowbreak{}the\_\allowbreak{}api\_\allowbreak{}would\_\allowbreak{}reject\_\allowbreak{}is\_\allowbreak{}refused}; \texttt{reports/\allowbreak{}metrics/\allowbreak{}segments.\allowbreak{}csv}; \href{https://github.com/mlops-2627q1-mds-upc/MLOPS_Recommenditos/blob/main/docs/docs/model-card.md}{model card}, the per-segment breakdown; \hyperref[EDN-48]{EDN-48}; \href{https://github.com/mlops-2627q1-mds-upc/MLOPS_Recommenditos/issues/39}{issue \#39}.},
}

% Generated by tools/edn_latex.py from reports/edn.md.
% Do not edit: change reports/edn.md and run `make edn`.
\ednentry{
  id = {EDN-61},
  title = {An omitted equipment list is the empty list, and the seam owns the masking},
  date = {2026-10-01},
  milestone = {M3: Quality Assurance},
  topic = {Feature Encoding, API Design, Evaluation Protocol},
  participants = {@lukas2510, decided by the repository owner},
  decision = {A request that omits one of the four equipment lists is read as an \textbf{empty} list, so every multi-hot column that list produced reaches the estimators as \texttt{False}. It is not read as absent. \texttt{Model.\allowbreak{}mask\_\allowbreak{}absent} in \texttt{recommenditos/\allowbreak{}modeling/\allowbreak{}model.\allowbreak{}py} is the single implementation of what an omitted field is, for the API, for SC-06's masking sweep and for the tests, and \texttt{Model.\allowbreak{}columns\_\allowbreak{}of} is the single implementation of which columns a field became.},
  rationale = {\emph{Alternatives considered:}
\begin{itemize}[leftmargin=*, topsep=2pt]
\item \textbf{Option A (chosen): omitted means the empty list, every item \texttt{False}.}
\newline Pros: it is a state the model was fitted on, and well represented. The raw field is always a list and \texttt{\textquotedbl{}[]\textquotedbl{}} when empty: over the 105,405 scoped listings there is \textbf{not one null} in any of the four equipment fields, while the literal \texttt{\textquotedbl{}[]\textquotedbl{}} is 5.18\,\% of them for \texttt{equipment\_\allowbreak{}comfort}, 8.03\,\% for \texttt{equipment\_\allowbreak{}entertainment}, 8.01\,\% for \texttt{equipment\_\allowbreak{}extra} and 6.44\,\% for \texttt{equipment\_\allowbreak{}safety} -- and 2.76\,\% of the training rows have all four empty at once, every one of the 133 multi-hot columns \texttt{False}. The estimators can only respond to what they were fitted on, so the encoding that lands inside the training distribution is the one that produces a prediction anybody can reason about. It is also the reading a listing form suggests: a seller who ticks nothing has an empty list.
\newline Cons: it is an assertion the request did not make. A Porsche whose comfort list was simply not sent is scored as a Porsche with no comfort equipment, and for a car whose price partly reflects its equipment that is a real, if small, downward pull. The entry accepts that rather than hiding it.
\item \textbf{Option B: omitted means absent, \texttt{pd.\allowbreak{}NA} across the list's columns.}
\newline Pros: it is what EDN-15 asks for in general -- a missing value is never imputed, because the missingness carries signal -- and reading an omission as ``none of these'' is an imputation. It is also what \texttt{build\_\allowbreak{}features.\allowbreak{}encode\_\allowbreak{}equipment} produces from a null source cell, so the serving path would need no special case to agree with it, and the feature contract names both states in so many words: ``an empty list is every item False, a null list is null throughout''.
\newline Cons: the state exists in the contract and not in the data. Every training row has a list, so a model asked to score an all-\texttt{pd.\allowbreak{}NA} equipment column is extrapolating off its own training distribution, with no fitted behaviour to appeal to at all -- LightGBM would route it by a default direction learned from a split that never saw a missing value, and a Ridge on \texttt{extended} would fire 133 missingness indicators whose coefficients were fitted on columns that were never missing. The semantic argument is about what the source \emph{could} express; this one is about what the model can answer.
\item \textbf{Option C: leave it to each caller, as it was.}
\newline Pros: no new code.
\newline Cons: this is where the entry comes from. The documentation said both things at once -- \texttt{\_\allowbreak{}align}'s docstring said an absent equipment list must arrive as NaN rather than \texttt{False}, \texttt{predict\_\allowbreak{}eur}'s said EDN-23 fixes it as empty -- and the test module had a private helper of its own implementing one of them. Worse, the one-liner a caller reaches for, \texttt{frame[column] =\allowbreak{} np.\allowbreak{}nan}, is wrong twice and loud once: for the three non-nullable condition flags it fails the contract's null check, so the caller finds out, and for a nullable equipment column it upcasts to float64, survives \texttt{conform} and comes back as \texttt{pd.\allowbreak{}NA} -- silently choosing option B.
\end{itemize}
\par
\emph{Why:} No measurement decides between A and B, and establishing \emph{that} is most of the work this entry did. Measured on the real snapshot's \texttt{lgbm-\allowbreak{}extended} over all 19,665 test rows, comparing predictions rather than metrics, the two rules are \textbf{bit-identical}: \texttt{equipment\_\allowbreak{}comfort} 6.8975\,\% MdAPE either way, \texttt{equipment\_\allowbreak{}entertainment} 6.3865\,\%, \texttt{equipment\_\allowbreak{}extra} 6.8543\,\%, \texttt{equipment\_\allowbreak{}safety} 6.4888\,\%, and all 23 optional fields masked at once 10.0082\,\%, with a maximum absolute difference of 0.0 EUR in every scenario. So the choice had to be made on what the two states mean and on which one the model has seen, not on a number.
\par
Why the equivalence holds, and when it would stop. \texttt{\_\allowbreak{}align} maps \texttt{False} to \texttt{0.\allowbreak{}0} and \texttt{pd.\allowbreak{}NA} to \texttt{NaN}, and for columns that are overwhelmingly \texttt{False} LightGBM sends both down the same branch -- which is EDN-23's argument that \texttt{\{True, NaN\}} splits like \texttt{\{True, False\}}, here measured over 133 columns instead of three. It is a property of the \textbf{tree models}, not of the seam and not of the decision. B1's numeric branch is a mean fill plus \texttt{MissingIndicator(features=\allowbreak{}\textquotedbl{}all\textquotedbl{})} (EDN-51), so an \texttt{pd.\allowbreak{}NA} equipment column would fire an indicator where a \texttt{False} one would not, and the two rules would diverge. \texttt{b1} is \texttt{basic}-only today and \texttt{basic} carries no equipment column, so nothing in the shipped ladder measures that -- which is precisely why the rule is fixed now, on semantics, rather than left resting on an equivalence that happens to hold for the two variants that currently exist.
\par
The API consequence is worth stating explicitly, because it is the mistake this entry is most likely to have to prevent, and it is the opposite of the one the first draft warned about: an omitted equipment list must reach the model as \texttt{False} throughout, so a request model that leaves the field \texttt{None} and passes it on has silently chosen option B. Defaulting it to \texttt{[]} is correct here.
\par
One thing this entry does \textbf{not} claim: that the seam can refuse the wrong masking. It cannot, and that is why a shared helper is the answer rather than a guard. An equipment column of all \texttt{pd.\allowbreak{}NA} is a legal frame the contract allows, and a numeric column of \texttt{0.\allowbreak{}0} is indistinguishable from a genuine zero, so there is nothing for \texttt{predict\_\allowbreak{}eur} to detect. What is left is to have one correct implementation, use it everywhere, and make the wrong value impossible to reach by accident.},
  ai = {Information seeking, Alternative generation, Alternative assessment, Recommendation, Solution generation},
  response = {Rejected},
  assessment = {The question was found by an adversarial review of pull request \#62, which spotted the contradiction and the masking trap and recommended option A. A second AI pass argued the other way from the contract wording, EDN-15 and \texttt{encode\_\allowbreak{}equipment}, recommended option B and implemented it. The repository owner overruled that on the ground neither pass had weighed: the training data contains no null equipment column, so option B hands the model a state it was never fitted on, while option A hands it one that 2.76\,\% of its training rows are. The semantic argument was about what the source could express, the distributional one about what the model can answer, and the second decides it. AI also produced a wrong supporting number for option B -- a 0.09\,\% metric gap between the rules -- which turned out to be its own measurement error and is corrected in this entry. What survived from AI is the shared helper and the \texttt{columns\_\allowbreak{}of} expansion, which are the durable part and are unaffected by the direction.},
  aievidence = {Claude Code session on 2026-09-30 reviewing pull request \#62 and on 2026-10-01 applying the owner's ruling. The 0.09\,\% gap AI first reported was traced to its own probe script masking a second column, \texttt{has\_\allowbreak{}full\_\allowbreak{}service\_\allowbreak{}history}, in the ``absent'' frame and not in the ``empty'' one: the real figures are 6.3297\,\% for that flag alone, 6.8975\,\% for \texttt{equipment\_\allowbreak{}comfort} under either rule, and 6.8982\,\% for the two together, which is the number that was mistaken for an equipment effect. Re-measured per field on the full test split, the two rules are bit-identical.},
  evidence = {\href{https://github.com/mlops-2627q1-mds-upc/MLOPS_Recommenditos/blob/main/recommenditos/modeling/model.py}{\texttt{recommenditos/\allowbreak{}modeling/\allowbreak{}model.\allowbreak{}py}}, \texttt{Model.\allowbreak{}mask\_\allowbreak{}absent}, \texttt{Model.\allowbreak{}columns\_\allowbreak{}of} and \texttt{Model.\allowbreak{}\_\allowbreak{}align}; \href{https://github.com/mlops-2627q1-mds-upc/MLOPS_Recommenditos/blob/main/recommenditos/data/build_features.py}{\texttt{recommenditos/\allowbreak{}data/\allowbreak{}build\_\allowbreak{}features.\allowbreak{}py}}, \texttt{\_\allowbreak{}equipment\_\allowbreak{}columns} and \texttt{encode\_\allowbreak{}equipment}; \texttt{tests/\allowbreak{}test\_\allowbreak{}model.\allowbreak{}py::test\_\allowbreak{}mask\_\allowbreak{}absent\_\allowbreak{}is\_\allowbreak{}the\_\allowbreak{}one\_\allowbreak{}answer\_\allowbreak{}to\_\allowbreak{}what\_\allowbreak{}an\_\allowbreak{}omitted\_\allowbreak{}field\_\allowbreak{}is} and \texttt{::test\_\allowbreak{}predict\_\allowbreak{}eur\_\allowbreak{}does\_\allowbreak{}not\_\allowbreak{}raise\_\allowbreak{}with\_\allowbreak{}every\_\allowbreak{}optional\_\allowbreak{}field\_\allowbreak{}absent}; \hyperref[EDN-15]{EDN-15}; \hyperref[EDN-23]{EDN-23}; \hyperref[EDN-51]{EDN-51}; \href{https://github.com/mlops-2627q1-mds-upc/MLOPS_Recommenditos/pull/62}{pull request \#62}; \href{https://github.com/mlops-2627q1-mds-upc/MLOPS_Recommenditos/issues/39}{issue \#39}.},
}

% Generated by tools/edn_latex.py from reports/edn.md.
% Do not edit: change reports/edn.md and run `make edn`.
\ednentry{
  id = {EDN-62},
  title = {\texorpdfstring{\texttt{lgbm-\allowbreak{}basic}}{lgbm-basic} is the candidate, and SC-04's miss on cars over 20 years is accepted rather than worked around},
  date = {2026-09-30},
  milestone = {M3: Quality Assurance},
  topic = {Modelling Approach, Model Evaluation (the NFR-01 gate)},
  participants = {@lukas2510, decided by the repository owner},
  decision = {Two decisions the measured criteria forced, recorded together because the same measurement produced both. \textbf{First:} \texttt{lgbm-\allowbreak{}basic} is the candidate the model card and the report put forward, not \texttt{lgbm-\allowbreak{}extended}, although \texttt{lgbm-\allowbreak{}extended} has the lower pooled MdAPE (6.32\,\% against 6.83\,\%). \texttt{lgbm-\allowbreak{}extended} fails SC-06 at 1.584x and \texttt{lgbm-\allowbreak{}basic} passes at 1.208x, and on a request carrying only the ten fields FR-01 requires -- which the contract lets any caller send -- \texttt{lgbm-\allowbreak{}basic} answers at 8.25\,\% against \texttt{lgbm-\allowbreak{}extended}'s 10.01\,\%. The candidate is a documented judgement, not a computed one: \texttt{metrics.\allowbreak{}json}'s \texttt{best\_\allowbreak{}variant} stays the lowest MdAPE overall, because that field means what it says, and \texttt{deployable\_\allowbreak{}variant} stays null while any criterion blocks. \textbf{Second:} SC-04's failure on \texttt{age\_\allowbreak{}bucket=\allowbreak{}over 20} (17.12\,\% for the candidate against a 15\,\% bound, over 784 test rows) is recorded as a missed criterion. The threshold is not moved, no upper age bound is added to the scope, and the segment is not dropped from the breakdown.},
  rationale = {\emph{Alternatives considered:}
\begin{itemize}[leftmargin=*, topsep=2pt]
\item \textbf{Option A (chosen): put \texttt{lgbm-\allowbreak{}basic} forward, and accept SC-04's miss with the reason recorded.}
\newline Pros: it is the only reading under which SC-06 does any work. A criterion exists to change a decision, and this is the decision it changes. The candidate meets every criterion that is measurable except SC-04, which no variant meets, so the choice costs nothing in criteria met and 0.51 pp of pooled MdAPE on a fully described car. It is also the cheaper model by every other measure: 16 columns against 162, 17.8 s of fit against 43 s, 6.3 MB of payload against 6.8 MB, and no 90-second \texttt{features} stage for the extended matrix. Accepting SC-04's miss keeps the bar where EDN-06 set it and turns the failure into the finding it is.
\newline Cons: the report has to explain why the model with the better headline number is not the one put forward, which is more words than ``we picked the best''. And the candidate is now a statement in a document rather than a field in an artefact, so a reader of \texttt{metrics.\allowbreak{}json} alone does not see it.
\item \textbf{Option B: keep \texttt{lgbm-\allowbreak{}extended} as the candidate and note SC-06 as a known limitation.}
\newline Pros: the lowest MdAPE is the number a reader looks for first, and the criteria are advisory until a release actually happens.
\newline Cons: it makes SC-06 decorative. The criterion was added by EDN-15 precisely to catch a model whose accuracy depends on optional inputs, it caught one, and overriding it the first time it fires teaches everyone reading the report that the gate is a formality. It would also put forward the \emph{worse} model for the request FR-01 guarantees, by 1.76 pp.
\item \textbf{Option C: serve \texttt{lgbm-\allowbreak{}extended} when the request is complete and \texttt{lgbm-\allowbreak{}basic} when it is not.}
\newline Pros: strictly the best point estimate for every request, and the 0.51 pp is not thrown away.
\newline Cons: two models in the image, two sets of SHAP baselines, two rows in every metrics table, and a routing rule that has to be explained to a user whose estimate changes when they add an optional field. It is also not this ticket's decision to take: it is an API design question, and FR-02's answer may turn out to be an interval rather than a second model. Left open for the deployment milestone rather than ruled out.
\item \textbf{Option D (for SC-04): add an upper age bound to the training scope, the way EDN-04 bounds price.}
\newline Pros: the criterion would pass, and there is precedent in the same document.
\newline Cons: it would pass by removing the rows it fails on, which is the shape of a criterion adjusted to its result. The cars over 20 years old are 784 of 19,665 test rows and they are cars people ask about; declaring them out of scope is a product decision about who the component serves, and it cannot be justified by a criterion the model missed.
\item \textbf{Option E (for SC-04): raise the threshold to 20\,\%, or report the criterion on a subset.}
\newline Pros: it would make the gate green.
\newline Cons: it is the same objection without the honesty of a stated scope change. SC-04's 15\,\% came from a reference run and a deliberate margin, and nothing has been learned about the \emph{bar} -- only about the model.
\end{itemize}
\par
\emph{Why:} The two halves of this entry rest on one measurement, and the measurement is what makes the choice defensible rather than tasteful. Both LightGBM variants were swept over every optional field on the real snapshot: each of \texttt{lgbm-\allowbreak{}extended}'s 23 optional fields costs at most 9.2\,\% of its MdAPE on its own, but all 23 at once cost 58.4\,\%, and that all-at-once scenario is not a hypothetical -- it is exactly a request carrying only the required fields, which FR-01 promises to answer. Comparing the two variants \emph{on that request} is comparing them on a request the product has to serve, and there the ladder's conclusion reverses. So the honest statement of what ladder step 4 measured is conditional: the extended features are worth 0.51 pp when the caller fills them in and cost 1.76 pp when the caller does not, and for a component that lets a caller omit 23 of its 33 inputs the second half is the one that decides.
\par
SC-04 is the same measurement read the other way. The problem specification named two possible remedies for the over-20-years segment, ``the extended features or a scope change for classic cars'', and the extended features have now been measured against it: 17.78\,\% against the basic set's 17.12\,\%, so they make that segment marginally \emph{worse} while improving the pooled figure. That rules the first remedy out on evidence, which is a finding about the features rather than an argument about the bar, and it is worth recording as such. The second remedy is a product decision nobody has taken, so the criterion is reported as missed. Nothing is unblocked by softening it in any case: SC-05 has no measurement until the UC2 conformal intervals exist, so NFR-01's gate blocks whatever SC-04 says, and a threshold edited under those circumstances would have no effect except on the report's credibility.},
  ai = {Information seeking, Alternative generation, Alternative assessment, Recommendation},
  response = {Accepted},
  assessment = {The finding is AI's, and it was not the one anybody was looking for. Two sessions measured SC-06 independently -- the implementer of issue \#39 and the reviewer of issue \#37 -- and agreed to four decimal places on every variant, which is why the numbers were trusted. The implementer then noticed what the all-at-once scenario \emph{is}, namely a required-fields-only request, and that comparing the two variants on it inverts the ladder's ranking; that step is what turned a criterion result into a decision. AI recommended option A for both halves and raised options C, D and E as things it would not decide on its own, which is the right boundary: which model the API serves and which cars the product covers are not conclusions to draw from a metrics file. Two of its numbers were corrected before they were used: an expected single-field failure on \texttt{has\_\allowbreak{}full\_\allowbreak{}service\_\allowbreak{}history} (predicted at 1.399x) measured at 1.002x on real data and was reported as not reproducing rather than as confirmed, and \texttt{b1}'s figures were re-measured after it turned out the available models had been fitted at a stale \texttt{min\_\allowbreak{}category\_\allowbreak{}rows}.},
  aievidence = {Claude Code sessions on 2026-09-30 implementing and reviewing issues \#37 and \#39. The masking sweep over 19,665 real test rows produced 1.208x for \texttt{lgbm-\allowbreak{}basic} and 1.584x for \texttt{lgbm-\allowbreak{}extended}, reproduced independently by both sessions, and \texttt{reports/\allowbreak{}metrics/\allowbreak{}masked-\allowbreak{}inputs.\allowbreak{}csv} carries the per-field table behind those two numbers. The owner was given options A to E with these measurements and chose A.},
  evidence = {\href{https://github.com/mlops-2627q1-mds-upc/MLOPS_Recommenditos/blob/main/docs/docs/model-card.md}{model card}, Results and the experiment ladder; \texttt{reports/\allowbreak{}metrics/\allowbreak{}masked-\allowbreak{}inputs.\allowbreak{}csv} and \texttt{reports/\allowbreak{}metrics/\allowbreak{}segments.\allowbreak{}csv}; \href{https://github.com/mlops-2627q1-mds-upc/MLOPS_Recommenditos/blob/main/docs/docs/problem-spec.md}{problem specification} section 8; \hyperref[EDN-06]{EDN-06}, \hyperref[EDN-15]{EDN-15}, \hyperref[EDN-59]{EDN-59}; \href{https://github.com/mlops-2627q1-mds-upc/MLOPS_Recommenditos/issues/39}{issue \#39}, \href{https://github.com/mlops-2627q1-mds-upc/MLOPS_Recommenditos/pull/65}{PR \#65}.},
}

% Generated by tools/edn_latex.py from reports/edn.md.
% Do not edit: change reports/edn.md and run `make edn`.
\ednentry{
  id = {EDN-63},
  title = {\texorpdfstring{\texttt{metrics.\allowbreak{}json}}{metrics.json} carries the data source and the fitted estimator, at the cost of five more columns},
  date = {2026-10-01},
  milestone = {M3: Quality Assurance},
  topic = {Experiment Tracking, Release Gate},
  participants = {@lukas2510},
  decision = {\texttt{metrics.\allowbreak{}json} gains a \texttt{data\_\allowbreak{}source} leaf, the resolved \texttt{download.\allowbreak{}source}, and an \texttt{estimator} leaf per variant, the model actually fitted; each \texttt{reports/\allowbreak{}metrics/\allowbreak{}<variant>.\allowbreak{}json} carries both as well. That widens the DVC metrics file from 24 leaves to 29, and \texttt{evaluate} declares \texttt{download.\allowbreak{}source} under \texttt{params:} so the recorded value cannot disagree with the lock.},
  rationale = {\emph{Alternatives considered:}
\begin{itemize}[leftmargin=*, topsep=2pt]
\item \textbf{Option A (chosen): the provenance goes into \texttt{metrics.\allowbreak{}json} itself, five leaves wide.}
\newline Pros: \texttt{dvc metrics show} and \texttt{dvc metrics diff} are where a promotion reviewer reads a result (EDN-12), so the one place a number can be quoted from is the one place that now says what produced it. \texttt{dvc metrics diff} lists a changed leaf, so a source or estimator that changes between two commits shows up beside the metric it moved.
\newline Cons: the file was already past the width a terminal table keeps, and this is five more columns of something that is not a measurement.
\item \textbf{Option B: the provenance goes only into the per-variant records under \texttt{reports/\allowbreak{}metrics/\allowbreak{}}.}
\newline Pros: costs nothing in the metrics table; the per-variant record already held \texttt{estimator}.
\newline Cons: it answers the question in the file nobody reads first. The defect this exists to prevent was a reader quoting \texttt{dvc metrics show}, and a second file they have to open is exactly the step that was skipped.
\item \textbf{Option C: an MLflow tag on each run and nothing in the artefacts.}
\newline Pros: the experiment is where the ladder is compared, and a tag is free.
\newline Cons: tracking is optional by construction -- it degrades to disabled with no credentials -- so the one artefact that is always written would be the one that says nothing. It also puts the evidence behind a login, while \texttt{metrics.\allowbreak{}json} is in the pull request diff.
\item \textbf{Option D: nothing, and rely on \texttt{dvc.\allowbreak{}lock} recording the source.}
\newline Pros: the lock already records it, so the information exists.
\newline Cons: it exists in a file whose \texttt{download} stage entry nobody reads to interpret an MdAPE, and it is a different file from the one holding the number. That is the state that let a constant-median stub's figure on generated rows sit on \texttt{main} as \texttt{best\_\allowbreak{}mdape} for a week.
\end{itemize}
\par
\emph{Why:} The honest version of the trade-off is that the metrics file is already too wide to read in a terminal and this makes it wider, so the question is only whether provenance earns five of those columns. It does, because the failure it prevents happened: \texttt{metrics.\allowbreak{}json} on \texttt{main} carried one MdAPE under all four variant names, computed by a constant-median stub on a synthetic fixture, and nothing in the artefact said either thing. The width argument also has a limit that matters here: the file is already read with \texttt{dvc metrics show -\allowbreak{}-\allowbreak{}md}, where five more columns cost nothing a reader notices, while \texttt{dvc metrics diff} is long-format and lists only what changed, so the new leaves are invisible until they are the answer.
Declaring \texttt{download.\allowbreak{}source} under the stage's \texttt{params:} is the part that makes it a guarantee rather than a label. Without it the stage would read a key DVC does not watch, and a recorded value could drift from the lock; with it, the convention \texttt{dvc.\allowbreak{}yaml} states for every other stage holds for this one too.
One consequence was found while implementing it and is worth recording: the end-to-end test fixture ran the stages against the project's own \texttt{params.\allowbreak{}yaml}, so after the flip to \texttt{zenodo} the fixture's artefacts would have claimed \texttt{data\_\allowbreak{}source: zenodo} over generated rows -- the exact mislabel the field exists to prevent. The fixture now runs against a params copy pinned to \texttt{synthetic}, and a test asserts the artefacts say so, which is what keeps the new field from being decorative.},
  ai = {Alternative generation, Alternative assessment, Solution generation},
  response = {Accepted},
  assessment = {Issue \#57 had already identified the defect and asked for the source and the estimator in \texttt{metrics.\allowbreak{}json}, so the problem was not AI's to find. What AI contributed is the shape and the two things that make it hold: declaring \texttt{download.\allowbreak{}source} under the stage's \texttt{params:} rather than reading an undeclared key, and noticing that the test fixture would start mislabelling its own artefacts as soon as the project's source flipped. The second one is the kind of defect this project has repeatedly produced -- a claim the code does not back -- and it would have shipped as a passing test suite.},
  aievidence = {Claude Code session on 2026-10-01 implementing issue \#57: the four options were laid out against the existing docstring's 24-leaf constraint, and the fixture's mislabelling was found by asking what \texttt{params[\textquotedbl{}download\textquotedbl{}][\textquotedbl{}source\textquotedbl{}]} would resolve to inside the test run rather than inside the pipeline.},
  evidence = {\href{https://github.com/mlops-2627q1-mds-upc/MLOPS_Recommenditos/blob/main/recommenditos/modeling/evaluate.py}{\texttt{recommenditos/\allowbreak{}modeling/\allowbreak{}evaluate.\allowbreak{}py}}, \texttt{gate\_\allowbreak{}summary} and \texttt{\_\allowbreak{}evaluate\_\allowbreak{}variant}; \texttt{tests/\allowbreak{}test\_\allowbreak{}pipeline.\allowbreak{}py::test\_\allowbreak{}the\_\allowbreak{}metrics\_\allowbreak{}artefacts\_\allowbreak{}say\_\allowbreak{}which\_\allowbreak{}data\_\allowbreak{}and\_\allowbreak{}which\_\allowbreak{}estimator\_\allowbreak{}produced\_\allowbreak{}them}; \texttt{metrics.\allowbreak{}json}; \hyperref[EDN-12]{EDN-12}, \hyperref[EDN-33]{EDN-33}; \href{https://github.com/mlops-2627q1-mds-upc/MLOPS_Recommenditos/issues/57}{issue \#57}.},
}

% Generated by tools/edn_latex.py from reports/edn.md.
% Do not edit: change reports/edn.md and run `make edn`.
\ednentry{
  id = {EDN-64},
  title = {Pynblint runs from its own locked environment, and a wrapper that reads its JSON report is the gate},
  date = {2026-10-01},
  milestone = {M3: Quality Assurance},
  topic = {Static Analysis, CI/CD},
  participants = {@lukas2510, decided by the repository owner},
  decision = {Pynblint 0.1.6 runs from a dedicated tool environment in \texttt{tools/\allowbreak{}pynblint-\allowbreak{}env/\allowbreak{}}, with its own \texttt{pyproject.\allowbreak{}toml} (\texttt{pynblint=\allowbreak{}=\allowbreak{}0.\allowbreak{}1.\allowbreak{}6}, \texttt{click<8.\allowbreak{}2}, Python 3.12) and its own hash-verified \texttt{uv.\allowbreak{}lock}, invoked as \texttt{uv run -\allowbreak{}-\allowbreak{}project tools/\allowbreak{}pynblint-\allowbreak{}env -\allowbreak{}-\allowbreak{}locked}.
A stdlib-only wrapper, \texttt{tools/\allowbreak{}notebook\_\allowbreak{}lint.\allowbreak{}py}, runs inside that environment, lints each notebook, reads Pynblint's JSON report and is what fails the build, because Pynblint itself exits 0 whatever it finds.
\texttt{make notebook-\allowbreak{}lint} is the one command, and the CI job \texttt{Notebook lint (Pynblint)} runs exactly that target.},
  rationale = {\emph{Alternatives considered:}
\begin{itemize}[leftmargin=*, topsep=2pt]
\item \textbf{Option A (chosen): a locked tool environment, plus a wrapper that is the gate.}
\newline Pros: all 63 packages are pinned and their hashes verified, so the tool cannot change under us; \texttt{-\allowbreak{}-\allowbreak{}locked} fails when the \texttt{pyproject.\allowbreak{}toml} and the lock drift apart; the root \texttt{uv.\allowbreak{}lock} is untouched and the directory is not a workspace member; and the lock is a file Dependabot can read.
The wrapper runs Pynblint with an allowlisted environment in an empty working directory, so no stray variable or \texttt{.\allowbreak{}pynblint} file can reconfigure it, and it fails loudly when Pynblint crashes instead of reading a crash as a clean run.
\newline Cons: two more files, one of them a generated 75 KB lock, a wrapper of its own with tests, and a deviation from the issue's literal ``via \texttt{uvx pynblint}''.
\item \textbf{Option B: \texttt{uvx -\allowbreak{}-\allowbreak{}python 3.\allowbreak{}12 -\allowbreak{}-\allowbreak{}with \textquotesingle{}click=\allowbreak{}=\allowbreak{}8.\allowbreak{}1.\allowbreak{}8\textquotesingle{} pynblint=\allowbreak{}=\allowbreak{}0.\allowbreak{}1.\allowbreak{}6}.}
\newline Pros: one line, and the closest to what the issue asked for.
\newline Cons: about 59 transitive dependencies float, and that is exactly how the tool broke: a click release nobody chose crashed it, so CI can turn red on a pull request that changed nothing near it, and a run cannot be reproduced later.
\item \textbf{Option C: the same, plus \texttt{-\allowbreak{}-\allowbreak{}exclude-\allowbreak{}newer <date>}.}
\newline Pros: the whole resolution is frozen in time, still in one line.
\newline Cons: no hash verification, a magic date in the Makefile, and no lock file for the dependency graph or Dependabot to see.
\item \textbf{Option D: a hashed requirements file passed to \texttt{uvx -\allowbreak{}c}.}
\newline Pros: looks like a lock without a second project.
\newline Cons: \texttt{uvx} silently ignores the hashes; a deliberately tampered hash still ran with exit 0, while \texttt{uv pip install -\allowbreak{}r} on the same file failed with ``Hash mismatch''.
\item \textbf{Option E: a pre-commit hook with \texttt{additional\_\allowbreak{}dependencies}.}
\newline Pros: feedback at commit time, and CI's existing lint job would run it.
\newline Cons: the dependencies float as in option B, and the raw hook reported ``Passed'' on a notebook with findings, so it needs the same wrapper anyway.
\item \textbf{Option F: a dependency group in the root project, kept apart with uv's \texttt{conflicts}.}
\newline Pros: one lock file for everything.
\newline Cons: it ties our lock to an unmaintained tool's \texttt{typer<0.\allowbreak{}13} and \texttt{ipython<9}, and every resolution of the project would carry that weight.
\item \textbf{For the gate: \texttt{jq -\allowbreak{}e} over the JSON in the Makefile, or the wrapper run in the project environment.}
\newline Pros: \texttt{jq} needs no code; the project environment would let the wrapper use typer and loguru like \texttt{tools/\allowbreak{}requirement\_\allowbreak{}matrix.\allowbreak{}py}.
\newline Cons: \texttt{jq} is not on every contributor's machine and prints no reasons; the project environment would make a seconds-long CI job install mlflow, lightgbm and DVC just to lint.
\end{itemize}
\par
\emph{Why:} The decision rests on three measured facts about Pynblint 0.1.6, each reproduced rather than read from its documentation.
First, it does not run as the issue proposed: it pins \texttt{typer<0.\allowbreak{}13} but not click, and typer 0.12 crashes on click 8.2 and newer (released 2025-05-10) with \texttt{TypeError: Secondary flag is not valid for non-\allowbreak{}boolean flag}, on Python 3.12, 3.13 and 3.14 alike.
A tool that broke once through a dependency nobody pinned should not be run in a way that lets it break again, which rules out options B and E and makes the lock worth its two files.
Second, it exits 0 with findings: a notebook with eleven findings still returned 0 in both file and directory mode, so a CI step that only runs it can never fail and something has to read the report.
Third, its configuration leaks in from outside: its settings are read from environment variables with no prefix, so an ordinary \texttt{INCLUDE=\allowbreak{}/\allowbreak{}usr/\allowbreak{}include} crashed a run with a validation error, and a \texttt{.\allowbreak{}pynblint} file in the current directory reconfigures it silently.
The wrapper is the smallest thing that answers all three: it pins nothing itself, it isolates the process, and its exit code is the finding count.
It is stdlib only so that it can run in the tool environment, which keeps the CI job free of the project's own dependencies.},
  ai = {Information seeking, Alternative generation, Alternative assessment, Recommendation, Solution generation},
  response = {Accepted},
  assessment = {AI did the experiments the decision rests on: it reproduced the click crash and pinned it to click rather than to the Python version, measured the exit code with findings, found that \texttt{uvx -\allowbreak{}c} ignores hashes by tampering with one, and timed every isolation option cold and warm.
It also found that the issue's own premise, ``run it isolated via \texttt{uvx pynblint}'', did not work as written.
AI then laid out options A to F with a recommendation, and the owner chose A.
AI also wrote the wrapper and its tests, test-first, and two further AI review passes, one on the code and one on the documents, found defects that were fixed before the pull request was opened.
A mutation run showed that no test noticed when the gate stopped passing the exclusions to Pynblint; an enforced repository rule would have been reported as checked although it never runs on a single notebook; and a report of an unexpected shape exited 1, which reads as findings, instead of 2.
Each now has a test.
An independent AI review of the pull request then found two more: the Pynblint answers the tests replay were recorded once and never checked against the real tool, so a release that stopped reporting findings would have passed every test, and a report that could not be written exited 1, which reads as findings.
A test that runs the real Pynblint from the tool environment and a fix for the exit code followed.
The owner reviews the result in the pull request.},
  aievidence = {Claude Code session on 2026-10-01 implementing issue \#41: a research run on Pynblint 0.1.6 in a scratch directory (rule catalogue, configuration, exit codes, isolation options with timings), then an architecture pass that turned the measurements into options with pros and cons; the owner was given the isolation and gate options with these measurements and chose the locked environment with the wrapper.
Claude Code review session on 2026-10-05 of PR \#67, which ran the gate on deliberately broken copies of the notebook and checked the tests' recorded answers against the real Pynblint.},
  evidence = {\href{https://github.com/mlops-2627q1-mds-upc/MLOPS_Recommenditos/blob/main/tools/pynblint-env/pyproject.toml}{\texttt{tools/\allowbreak{}pynblint-\allowbreak{}env/\allowbreak{}pyproject.\allowbreak{}toml}}, \href{https://github.com/mlops-2627q1-mds-upc/MLOPS_Recommenditos/blob/main/tools/notebook_lint.py}{\texttt{tools/\allowbreak{}notebook\_\allowbreak{}lint.\allowbreak{}py}}, \texttt{tests/\allowbreak{}test\_\allowbreak{}notebook\_\allowbreak{}lint.\allowbreak{}py}, the \texttt{Notebook lint (Pynblint)} job in \href{https://github.com/mlops-2627q1-mds-upc/MLOPS_Recommenditos/blob/main/.github/workflows/ci.yml}{\texttt{.\allowbreak{}github/\allowbreak{}workflows/\allowbreak{}ci.\allowbreak{}yml}}; \hyperref[EDN-65]{EDN-65}; \href{https://github.com/mlops-2627q1-mds-upc/MLOPS_Recommenditos/issues/41}{issue \#41}, \href{https://github.com/mlops-2627q1-mds-upc/MLOPS_Recommenditos/pull/67}{PR \#67}.},
}

% Generated by tools/edn_latex.py from reports/edn.md.
% Do not edit: change reports/edn.md and run `make edn`.
\ednentry{
  id = {EDN-65},
  title = {The Pynblint policy: enforce what can fire on a committed notebook, and show execution by running it},
  date = {2026-10-01},
  milestone = {M3: Quality Assurance},
  topic = {Static Analysis, Testing Strategy},
  participants = {@lukas2510, decided by the repository owner},
  decision = {Of Pynblint 0.1.6's 22 rules, 13 are enforced at their default thresholds and 9 are excluded, each with its reason, in the table in \href{https://github.com/mlops-2627q1-mds-upc/MLOPS_Recommenditos/blob/main/docs/docs/notebooks.md\#pynblint-rules}{Notebooks}.
That table is the configuration: the gate reads it, and refuses to run unless it classifies exactly the rules the installed Pynblint has.
The three rules that read execution counts are excluded, because nbstripout removes those counts on every commit, and that a notebook runs top to bottom is shown by running it with \texttt{make notebook-\allowbreak{}run} instead.
The five repository-level rules are not run, each with the reason and with what enforces the same property here.
Pynblint lints every notebook git would commit, one file at a time, and two rules of our own require every notebook to sit directly in \texttt{notebooks/\allowbreak{}} under the Cookiecutter name \texttt{<number>.\allowbreak{}<version>-\allowbreak{}<initials>-\allowbreak{}<description>.\allowbreak{}ipynb}, and to carry none of the flags that make nbstripout keep its outputs.
There is no warning tier and no per-notebook waiver: a finding fails the build, and a rule that is wrong for us is excluded in the table, with its reason, in a reviewed pull request.},
  rationale = {\emph{Alternatives considered:}
\begin{itemize}[leftmargin=*, topsep=2pt]
\item \textbf{Option A (chosen): exclude the execution rules, keep nbstripout as it is, and evidence execution by running the notebook.}
\newline Pros: nbstripout 0.9.1 sets every \texttt{execution\_\allowbreak{}count} to null, so on a committed notebook \texttt{non-\allowbreak{}executed-\allowbreak{}notebook} always fires, \texttt{non-\allowbreak{}linear-\allowbreak{}execution} can never fire, and \texttt{non-\allowbreak{}executed-\allowbreak{}cells} fires only on an empty code cell and then flags every other code cell; enforcing them would either keep the build red forever or claim checks that do not happen.
Stripping also keeps rows of a dataset with personal data out of git, which is more than diff hygiene.
The profiling notebook's summary cell fails when a dataset card figure does not reproduce, so the drill checks the card as well as the code.
\newline Cons: execution is not gated in CI; \texttt{make notebook-\allowbreak{}run} is a local drill, because the notebook needs the 548 MB source file.
\item \textbf{Option B: \texttt{nbstripout -\allowbreak{}-\allowbreak{}keep-\allowbreak{}count}, so the three rules have something to read.}
\newline Pros: three more live rules.
\newline Cons: execution counts churn in every diff, and they record what one kernel once did, not that a fresh kernel runs the notebook: an edited cell that was never re-run keeps its old count.
Every edit would need a full re-run on the source file before committing, and an editor that writes no counts trips the rules.
\item \textbf{Option C: scan the repository root and run the repository rules.}
\newline Pros: the literal reading of ``notebook and repository QA''.
\newline Cons: \texttt{repository-\allowbreak{}not-\allowbreak{}versioned} looks for a \texttt{.\allowbreak{}git} directory and fires in every git worktree, which this project uses for parallel work; locally the scan walks \texttt{.\allowbreak{}venv} and other worktrees and lints other branches' stale notebook copies; \texttt{test-\allowbreak{}coverage-\allowbreak{}data-\allowbreak{}not-\allowbreak{}available} depends on whether the tests ran in that checkout.
Each repository rule is about a property something else here already enforces more strongly, which the table says rule by rule.
\item \textbf{Option D: scan the \texttt{notebooks/\allowbreak{}} directory.}
\newline Pros: one command, no file list.
\newline Cons: three structural false positives to exclude, untracked scratch files get linted, and a notebook committed anywhere else escapes the scan.
\item \textbf{Option E: the policy in a commented \texttt{.\allowbreak{}pynblint} file.}
\newline Pros: Pynblint's own configuration format.
\newline Cons: it is read only from the current directory and is overridden by environment variables, its reasons are comments nothing checks, and a misspelt rule name is accepted silently, so an exclusion with a typo excludes nothing and says nothing.
\item \textbf{Option F: raise \texttt{cell-\allowbreak{}too-\allowbreak{}long} or \texttt{notebook-\allowbreak{}too-\allowbreak{}long} so the first notebook passes.}
\newline Pros: no restructuring.
\newline Cons: it is the threshold adjusted to its result that EDN-62 rejects; the reason would be the notebook, not the rule.
\item \textbf{Option G: report the excluded rules as non-blocking warnings.}
\newline Pros: nothing is hidden.
\newline Cons: Pynblint has no severities, and a warning tier is exactly the ``findings nobody acts on'' the issue warns about.
\item \textbf{Option H: execute the notebook in CI, with the source file cached under its MD5.}
\newline Pros: the execution evidence that options A and B argue about would be automated, and because the summary cell fails on a mismatch, every pull request would check the dataset card against the data.
\newline Cons: a cold cache downloads 548 MB from Zenodo, so a third party's server would decide whether our CI is green, which is a flaky test by construction; the job needs the full project environment, about a minute and 2.8 GB of RAM; and GitHub evicts a cache unused for seven days, so the download recurs.
It stays open: if the raw data ever comes from our own DVC remote instead, most of the cons go away.
\end{itemize}
\par
\emph{Why:} A rule is enforced where it can fire on what we commit, and excluded where it cannot, with the reason written next to it.
That principle is what keeps the gate honest in both directions: nothing it reports is noise, and nothing it lists as checked is unchecked.
It also bit on the first run: the dataset card's profiling notebook, as first drafted, had five code cells over 30 lines, and they were restructured into shorter cells with precomputed values rather than the threshold being raised, so the notebook stayed within the 50-cell limit as well.
On the stripped copy that is committed, the only rule that fired was \texttt{non-\allowbreak{}executed-\allowbreak{}notebook}, which is the measurement behind excluding the execution rules.
Making the documentation table the configuration follows the pattern of the requirement matrix, which reads the requirements and the specification rather than a copy of them: the reasons cannot drift from what the gate does, and a rule a future Pynblint adds or renames stops the build until someone decides about it.
The table names every rule, enforced or excluded, and the gate refuses one that does not classify exactly the installed rules, enforces nothing, or enforces a repository rule that would never run, so a misspelt slug, a rule a new release adds or an empty policy stops the gate rather than silently changing what is checked.},
  ai = {Information seeking, Alternative generation, Alternative assessment, Recommendation, Solution generation},
  response = {Accepted},
  assessment = {AI read the Pynblint 0.1.6 source and tested every rule on purpose-built notebooks, stripped and unstripped, which is how the interaction with nbstripout was found rule by rule rather than assumed; it also found the blind spots the table now states as caveats.
It proposed the policy, the docs-as-configuration design and the location rule, and laid out options A to G; the owner chose A for the execution rules and the docs table as the policy.
The notebook itself was drafted by AI, and its first Pynblint run found five code cells over 30 lines; they were restructured rather than the threshold raised, so the policy applied to AI-written code as to any other.
An AI review of the documents then found two holes in the first version of this policy, both closed before the pull request was opened: nothing stopped a \texttt{keep\_\allowbreak{}output} flag from carrying outputs past nbstripout, and \texttt{make notebook-\allowbreak{}run}, the stand-in for the excluded execution rules, passed even when a dataset card figure did not reproduce.},
  aievidence = {Claude Code session on 2026-10-01 implementing issue \#41: a rule-by-rule experiment on stripped and unstripped sample notebooks, an architecture pass that produced the policy table and options A to G, the owner's choice among them, and a document review that added option H and the two fixes above.},
  evidence = {\href{https://github.com/mlops-2627q1-mds-upc/MLOPS_Recommenditos/blob/main/docs/docs/notebooks.md}{Notebooks}, the rule table; \href{https://github.com/mlops-2627q1-mds-upc/MLOPS_Recommenditos/blob/main/notebooks/1.0-lh-dataset-card-profiling.ipynb}{\texttt{notebooks/\allowbreak{}1.\allowbreak{}0-\allowbreak{}lh-\allowbreak{}dataset-\allowbreak{}card-\allowbreak{}profiling.\allowbreak{}ipynb}}; \texttt{tests/\allowbreak{}test\_\allowbreak{}notebook\_\allowbreak{}lint.\allowbreak{}py}; \href{https://github.com/mlops-2627q1-mds-upc/MLOPS_Recommenditos/blob/main/docs/docs/specification.md}{specification} NFR-07; \hyperref[EDN-64]{EDN-64}; \href{https://github.com/mlops-2627q1-mds-upc/MLOPS_Recommenditos/issues/41}{issue \#41}, \href{https://github.com/mlops-2627q1-mds-upc/MLOPS_Recommenditos/pull/67}{PR \#67}.},
}

% Generated by tools/edn_latex.py from reports/edn.md.
% Do not edit: change reports/edn.md and run `make edn`.
\ednentry{
  id = {EDN-66},
  title = {NFR-06 and NFR-09 are each split in two, and each ID keeps the half its citations mean},
  date = {2026-10-05},
  milestone = {M3: Quality Assurance},
  topic = {Requirements Engineering, Testing Strategy},
  participants = {@lukas2510, decided by the repository owner},
  decision = {NFR-06 and NFR-09 each promised two things that come due in different milestones and are verified by different evidence, so each is split in two, and every half has its own ID, acceptance criterion, verification route and milestone.
The existing ID is narrowed to the half its citations mean, and the other half gets the next free ID.
\textbf{NFR-06} keeps reproducibility: from a clean clone, \texttt{dvc pull} and \texttt{dvc repro} reproduce the splits and metrics within \ensuremath{\pm}0.1 percentage points; \textbf{[manual]}, the clean-clone drill of \href{https://github.com/mlops-2627q1-mds-upc/MLOPS_Recommenditos/blob/main/docs/docs/pipeline.md}{The DVC pipeline}, due M3.
\textbf{NFR-14} takes run provenance: every training run records its commit, its data version and its parameters; \textbf{[automated]}, by the tracking tests and a \texttt{train} test that reads a real run back, due M3.
\textbf{NFR-09} keeps the running system's security: a read-only public surface, the body-size limit, no TLS, the API key from the environment, non-root containers and the vulnerability scan; \textbf{[automated]}, due M4.
\textbf{NFR-15} takes credential hygiene: no secret in any commit of the history, a scan before every merge, and a leaked secret revoked rather than only removed; \textbf{[manual]}, the CI secret scan of EDN-71, due M3.},
  rationale = {\emph{Alternatives considered:}
\begin{itemize}[leftmargin=*, topsep=2pt]
\item \textbf{Option A (chosen): narrow the existing ID to one half, and give the other half a new ID.}
\newline Pros: on \texttt{main} before the split, 66 lines cited NFR-06, 15 of them \texttt{req} markers, and 18 cited NFR-09.
Apart from the requirement rows themselves and the passages that tell NFR-06's history, eleven of NFR-06's citations mean provenance -- the tracking seam and its tests, two documentation pages and EDN-54 -- and of NFR-09's only EDN-46's evidence pointer means secret hygiene.
Every other citation, the report's included, means the half the ID keeps and stays correct without an edit, and the twelve tests marked NFR-06 for determinism and the download pin keep their marker.
No ID is reused for a different requirement: a narrowed ID promises a subset of what it promised before and nothing new.
\newline Cons: a reader holding an older text can still read NFR-06 as including provenance, so the citations of the moved halves had to be found and repointed.
Earlier entries written under the combined meaning, EDN-45, EDN-55 and EDN-56, are records of their date and are not rewritten; EDN-56 in fact already argued the two halves of NFR-06 separately, which is the strongest sign they were two requirements.
\item \textbf{Option B: retire NFR-06 and NFR-09 and issue four new IDs.}
\newline Pros: no ID ever changes meaning, not even by narrowing.
\newline Cons: every one of those 84 citations becomes a pointer to a requirement that no longer exists, including three sections of the LaTeX report; the fifteen NFR-06 markers fail the per-test check until each is rewritten; and the matrix needs every requirement to have a table row, so a retired ID could only survive in prose, where nothing checks it.
\item \textbf{Option C: suffixes, NFR-06a and NFR-06b.}
\newline Pros: the family stays visible in the name.
\newline Cons: \texttt{tools/\allowbreak{}requirement\_\allowbreak{}matrix.\allowbreak{}py} accepts \texttt{(?:FR|NFR)-\allowbreak{}\textbackslash{}d\{2,\}} and nothing else, so the tool would have to change first; and a bare ``NFR-06'' in an existing citation would become ambiguous between the parent and its halves, which is the ambiguity the split exists to remove.
\item \textbf{Option D: keep one ID each, and let the ``Verified by'' cell name both routes.}
\newline Pros: no ID change at all.
\newline Cons: it is the state issue \#60 describes.
The parser requires exactly one of \textbf{[automated]} and \textbf{[manual]}, and the gate books an ID to one milestone, so NFR-09's hygiene half sat in M4 behind an API that does not exist yet while the token it protects was already on every contributor's disk, and NFR-06's tests could only ever make it ``named by a test''.
\end{itemize}
\par
\emph{Why:} The constraint is that a requirement ID is never reused for a different requirement, because the report and the EDN cite them; between A and B the question is only which citations break.
In A it is a bounded, listed set, the ones that mean the half that moved; in B it is all of them.
The halves get their milestones from when their evidence can exist, not from the requirement they came from: NFR-14 is due in M3 because the tracking seam (PR \#50) and the \texttt{train} stage's parameter logging (PR \#62) are both merged, and NFR-15 is due in M3 because the risk has existed since issue \#42 and PR \#50 put a DagsHub token into every contributor's \texttt{.\allowbreak{}env} and \texttt{.\allowbreak{}dvc/\allowbreak{}config.\allowbreak{}local} (EDN-46).
NFR-09 stays booked to M4 as decided with the split, although its image scan, its Dockerfile review and the check from outside the VM are M5 work; \texttt{tools/\allowbreak{}expected\_\allowbreak{}coverage.\allowbreak{}yaml} says so, so the delivery review notices if they are still missing then.
NFR-14 is \textbf{[automated]} because what a run records can be read back from a run; NFR-15 is \textbf{[manual]} because its evidence is a CI job and not a Pytest test, the route NFR-07 already uses for the CI run it names, and making it \textbf{[automated]} would have needed a test that only stands in for the scan.
Giving NFR-14 a route of its own also showed that it was not met: \texttt{train.\allowbreak{}num\_\allowbreak{}threads} is declared as a parameter of the \texttt{train} stage in \texttt{dvc.\allowbreak{}yaml} and changes the artefact (EDN-53), yet the run carried it only inside the \texttt{model.\allowbreak{}json} artefact, not as an MLflow parameter, so a comparison of runs by parameter could not see it, and the only test spot-checked two parameters.
So a test now reads the keys \texttt{dvc.\allowbreak{}yaml} declares for each \texttt{train@\allowbreak{}<variant>} out of \texttt{dvc.\allowbreak{}yaml} itself and compares them with the parameters of a real run; it failed on the code before the change, on \texttt{num\_\allowbreak{}threads} for all four variants, and passes since the stage logs it.
\texttt{train.\allowbreak{}mlflow\_\allowbreak{}experiment} is the one declared key recorded as what it is, the experiment the run belongs to, and the test checks it there rather than as a repeated parameter.
\texttt{train.\allowbreak{}py} and \texttt{recommenditos/\allowbreak{}tracking.\allowbreak{}py}, whose three comments also move from NFR-06 to NFR-14, are dependencies of the four \texttt{train} stages and of \texttt{evaluate}, so the change reruns them on the real snapshot.
The re-run on 2026-10-05 moved no number: \texttt{metrics.\allowbreak{}json}, the segment table and the masking sweep are unchanged, and inside each model bundle only \texttt{model.\allowbreak{}json} changed, by the provenance EDN-56 puts there, while every payload file kept its hash; the four new runs carry all of their stage's declared parameters, \texttt{num\_\allowbreak{}threads} among them.},
  ai = {Information seeking, Alternative generation, Alternative assessment, Recommendation, Solution generation},
  response = {Accepted with modifications},
  assessment = {The owner decided to split both requirements, which half comes due when, and that the hygiene half is verified by a secret scan; the issue and its comment are his.
How to split the IDs he delegated, under the constraint that no ID is reused.
AI counted every citation of both IDs and sorted them by the half each one means, which is the evidence that decides between A and B, read the matrix's ID pattern to rule out C, and chose A.
It also found the unlogged \texttt{num\_\allowbreak{}threads} while writing NFR-14's route, and its first recommendation was to leave that and the \texttt{tracking.\allowbreak{}py} comments to a later change that retrains anyway.
An independent AI review of the pull request rejected that, because it left NFR-14 reported as verified by a test while its parameter clause was unmet, and the fix was made in this pull request, test first, which is the modification.},
  aievidence = {Claude Code session on 2026-10-05 implementing issue \#60: \texttt{grep -\allowbreak{}rn \textquotedbl{}NFR-\allowbreak{}0[69]\textquotedbl{}} over the documents, the report, the tests, the tools and the package, each hit classified by the half it means, and \texttt{git grep} counts on \texttt{main} for the figures above; the requirement matrix rebuilt after the split, with NFR-14 verified by five tests and NFR-15 verified by hand; \texttt{test\_\allowbreak{}a\_\allowbreak{}run\_\allowbreak{}records\_\allowbreak{}every\_\allowbreak{}parameter\_\allowbreak{}its\_\allowbreak{}stage\_\allowbreak{}declares} run against the unchanged \texttt{train.\allowbreak{}py} first, failing for all four variants on \texttt{num\_\allowbreak{}threads} alone.},
  evidence = {\href{https://github.com/mlops-2627q1-mds-upc/MLOPS_Recommenditos/issues/60}{issue \#60} and its comment on NFR-06; \href{https://github.com/mlops-2627q1-mds-upc/MLOPS_Recommenditos/pull/70}{PR \#70}; \href{https://github.com/mlops-2627q1-mds-upc/MLOPS_Recommenditos/blob/main/docs/docs/requirements.md}{requirements} and \href{https://github.com/mlops-2627q1-mds-upc/MLOPS_Recommenditos/blob/main/docs/docs/specification.md}{specification} NFR-06, NFR-09, NFR-14 and NFR-15; \texttt{tools/\allowbreak{}expected\_\allowbreak{}coverage.\allowbreak{}yaml}; \texttt{DUE\_\allowbreak{}AT\_\allowbreak{}M3} in \texttt{tests/\allowbreak{}test\_\allowbreak{}requirement\_\allowbreak{}matrix.\allowbreak{}py}; \hyperref[EDN-27]{EDN-27} for the one ID space, \hyperref[EDN-44]{EDN-44}, \hyperref[EDN-45]{EDN-45}, \hyperref[EDN-46]{EDN-46}, \hyperref[EDN-56]{EDN-56}, \hyperref[EDN-71]{EDN-71}.},
}

% Generated by tools/edn_latex.py from reports/edn.md.
% Do not edit: change reports/edn.md and run `make edn`.
\ednentry{
  id = {EDN-67},
  title = {\texorpdfstring{\texttt{features}}{features} applies the supported-make list to every frame, the \texorpdfstring{\texttt{ES}}{ES} holdout included, before the vocabulary is built},
  date = {2026-10-05},
  milestone = {M3: Quality Assurance (feature engineering; the holdout feeds M6: Monitoring)},
  topic = {Feature Encoding, Evaluation Protocol, Monitoring},
  participants = {@lukas2510, decided by the repository owner},
  decision = {The \texttt{features} stage restricts all five frames it reads -- train, validation, calibration, test and the \texttt{ES} holdout -- to the makes \texttt{data/\allowbreak{}processed/\allowbreak{}supported\_\allowbreak{}makes.\allowbreak{}json} names, and does so before any training-derived vocabulary or statistic is computed.
That discharges the first and second of the four guarantees EDN-48 left to the downstream stages, in one place.
The holdout matrix therefore holds 5,979 of the 6,079 \texttt{ES} rows: the 100 it loses are exactly the listings the API answers with a 422 under FR-04, so the matrix is the holdout as the model receives it.
It is not what M6 replays: a request cannot be rebuilt from a matrix, because the holdout's \texttt{country\_\allowbreak{}code} is empty in every row (EDN-42) and \texttt{registration\_\allowbreak{}date} has become \texttt{age\_\allowbreak{}years}, both fields FR-01 requires.
The replay reads \texttt{data/\allowbreak{}processed/\allowbreak{}holdout\_\allowbreak{}es.\allowbreak{}parquet}, which \texttt{split} still writes with all 6,079 rows, and has to filter to the supported makes itself or count the 100 refused requests as FR-04's 422s.
As a consequence \texttt{train} and \texttt{evaluate} no longer filter at all; they check the property with \texttt{check\_\allowbreak{}served\_\allowbreak{}makes\_\allowbreak{}only} and fail when it does not hold.},
  rationale = {\emph{Alternatives considered:}
\begin{itemize}[leftmargin=*, topsep=2pt]
\item \textbf{Option A (chosen): filter every frame in \texttt{features}, the holdout included, before the vocabulary.}
\newline Pros: the rule gets one owner, at the only point where it still decides the feature space, so the \texttt{make} levels (11 rather than 25), the \texttt{model} levels (437 rather than 561), the equipment threshold and the \texttt{model\_\allowbreak{}version} floor are all shares of the population the product serves.
Obligations 1 and 2 close together, and the UC2 conformal calibration, which has no ticket yet, inherits a served population instead of having to remember the filter.
Any model-side use of the holdout matrix -- scoring it offline, or comparing encoded inputs for drift -- is on the 5,979 rows the API would pass to the model, the population EDN-14 measured NFR-11 on.
\newline Cons: the holdout matrix no longer holds every \texttt{ES} listing, so a model-side analysis that wants the refused rows has to start from \texttt{data/\allowbreak{}processed/\allowbreak{}holdout\_\allowbreak{}es.\allowbreak{}parquet}, which \texttt{split} still writes losslessly.
The vocabulary moves, so every model and every number in the model card had to be re-measured.
\item \textbf{Option B: restrict only the training rows the vocabulary is fitted on, and leave the holdout matrix unfiltered for the M6 ticket to decide.}
\newline Pros: the smallest change to the holdout matrix, which would keep every \texttt{ES} row.
\newline Cons: it leaves obligation 1 open for calibration and the holdout, so the guarantee stays spread over stages and tickets that do not exist yet, which is the cost EDN-48 accepted and this ticket exists to remove.
And the raw holdout matrix would not even be raw: once the vocabulary is fitted on served rows, the 100 unserved rows are encoded with no \texttt{make} level at all, so the matrix would hold rows the model never receives, in an encoding that cannot say what they were.
Nor would it keep anything for M6: the replay cannot be built from a matrix under either option, and the split frame keeps the refused rows under both.
\end{itemize}
\par
\emph{Why:} The restriction belongs where the feature space is decided, and the holdout matrix belongs inside it rather than beside it.
The matrix is how the model sees the holdout, and the API refuses an unsupported make with a 422 before the model is called (FR-04), so what the model can ever receive from the \texttt{ES} listings is the 5,979 rows, not the 6,079 -- the population EDN-14 measured NFR-11 on.
What M6 replays is a separate question this decision does not settle.
The replay is built from the split frame, because a matrix cannot be turned back into requests, and it must either apply the same list or count FR-04's 422s for the 100 refused rows; either way the drift and error figures it reports are about the 5,979 the model answers.
That is also why option B's apparent advantage, keeping refused requests replayable, is not one: they stay in \texttt{data/\allowbreak{}processed/\allowbreak{}holdout\_\allowbreak{}es.\allowbreak{}parquet} under either option, so the choice made here stays reversible.
\par
\textbf{Interaction with EDN-42, checked on the real run.} The filter acts on rows by \texttt{make} and touches no other encoding, so the holdout's \texttt{country\_\allowbreak{}code} is still missing in every one of its 5,979 rows, because \texttt{ES} is still not a training level, and the stage's warning about a column a split never fills still names \texttt{country\_\allowbreak{}code} and nothing else.
What changes is the rest of the holdout matrix: before, its 100 unserved rows were encoded with make levels the model was never meant to serve; now every holdout row is one of the 11 training makes and none has a missing make, so \texttt{country\_\allowbreak{}code} is the only feature of the holdout matrix the model sees as unknown, which is the new-market signal EDN-42 chose to keep.
\par
\textbf{What it moved, measured on the real snapshot (2026-10-05, \texttt{reports/\allowbreak{}analysis/\allowbreak{}served\_\allowbreak{}vocabulary.\allowbreak{}py} and the committed \texttt{dvc.\allowbreak{}lock}).} The \texttt{make} levels go from 25 to 11, the \texttt{model} levels from 561 to 437, and the \texttt{model\_\allowbreak{}version} levels from 279 covering 85.9\,\% of the training rows to 278 covering 86.5\,\% of the served ones; the equipment vocabulary is the same 133 of 146 items either way, so neither matrix gains or loses a column.
On the ladder, three of the four variants are unchanged to full precision: the \texttt{b0} and \texttt{b1} payloads are byte-identical, and \texttt{lgbm-\allowbreak{}basic} builds the same trees, with only the category codes inside its split bitsets renumbered.
\texttt{lgbm-\allowbreak{}extended} is the one that moved, through the one \texttt{model\_\allowbreak{}version} level it lost: MdAPE 6.318\,\% to 6.325\,\%, within 20\,\% 90.40\,\% to 90.29\,\%, SC-04's worst segment 17.78\,\% to 17.86\,\%, SC-06 1.584x to 1.583x, and its best iteration is now the 1,000th rather than the 996th, so early stopping fires for neither LightGBM variant inside the 1,000-tree budget.
No criterion changed its verdict for any variant, and EDN-62's choice of candidate holds: \texttt{lgbm-\allowbreak{}extended} still fails SC-06, and its pooled advantage shrinks from 0.51 pp to 0.50 pp.
\par
\textbf{Consequences for the other stages.} \texttt{train} and \texttt{evaluate} used to filter their own rows; the filter is now a check, \texttt{check\_\allowbreak{}served\_\allowbreak{}makes\_\allowbreak{}only}, which refuses a matrix holding a make outside the list, a row with no make level, or a \texttt{make} level the list does not name.
A second filter would never remove a row from matrices \texttt{features} wrote, and on any other matrices it would quietly repair the rows while keeping a feature space decided over another population, so the rule keeps one owner and its consumers check a property rather than an intention.
This changes behaviour only for inconsistent inputs -- matrices from another run, or a make list rewritten without rebuilding them -- which used to be filtered silently and now fail with a message naming the state.
One state is not an inconsistency and gets a failure of its own in \texttt{features}: a supported make that no training row holds, which \texttt{split} can produce because it counts support over all four sets, so after a change to \texttt{seed} or to the threshold every seller of a make can land outside train.
The API would accept that make while the model had never seen it, and only \texttt{features} can say why, so it refuses the state by name rather than leaving \texttt{train} to report matrices from another run.
For the same reason the guard against an empty population moved into \texttt{features} and covers every frame: an empty matrix cannot even be read back, because Parquet keeps no levels for an empty categorical, so the downstream guards it replaces could no longer be reached.},
  ai = {Alternative generation, Alternative assessment, Recommendation, Solution generation},
  response = {Accepted},
  assessment = {AI generated options A and B with the holdout counts and the EDN-42 interaction, and recommended A; Lukas chose A before the implementation started.
Implementing it test first turned up two things the options had not anticipated.
Once the vocabulary is fitted on served rows, an unfiltered holdout would encode its unserved rows with no make level at all, which takes away most of what option B was meant to preserve; and an empty matrix cannot be read back through Parquet, which is why the empty-population guard moved from \texttt{train} and \texttt{evaluate} into \texttt{features}.
AI measured the accuracy effect rather than assuming it, and established why three variants did not move rather than reporting the unchanged numbers as a coincidence.
An independent review of the pull request, also by AI, then corrected two points without changing the decision.
The first version of this entry argued that the filtered matrix is what M6 replays, which it cannot be, since a request cannot be rebuilt from a matrix; the replay reads the split frame, and the reasoning above says so.
And a supported make with no training row was only warned about in \texttt{features}, so the failure surfaced in \texttt{train} with advice that looped; \texttt{features} now refuses it by name.},
  aievidence = {Claude Code session on 2026-10-05 implementing issue \#63: the two options were laid out with the 5,979-of-6,079 holdout count and the EDN-42 interaction, and Lukas chose A.
The independent review of PR \#71 on the same day found the replay argument and the untrained-make diagnosis, and reproduced the second with a training split holding only BMW against the list BMW and Audi.
The measurement behind the numbers is \texttt{reports/\allowbreak{}analysis/\allowbreak{}served\_\allowbreak{}vocabulary\_\allowbreak{}results.\allowbreak{}txt} and the ladder re-run the committed \texttt{dvc.\allowbreak{}lock} records, tracked as one MLflow run per variant.},
  evidence = {\href{https://github.com/mlops-2627q1-mds-upc/MLOPS_Recommenditos/issues/63}{issue \#63}; \href{https://github.com/mlops-2627q1-mds-upc/MLOPS_Recommenditos/pull/71}{PR \#71}; \href{https://github.com/mlops-2627q1-mds-upc/MLOPS_Recommenditos/blob/main/recommenditos/data/build_features.py}{\texttt{recommenditos/\allowbreak{}data/\allowbreak{}build\_\allowbreak{}features.\allowbreak{}py}} (\texttt{served\_\allowbreak{}rows}, \texttt{check\_\allowbreak{}served\_\allowbreak{}makes\_\allowbreak{}only}); \texttt{tests/\allowbreak{}test\_\allowbreak{}features.\allowbreak{}py::test\_\allowbreak{}every\_\allowbreak{}frame\_\allowbreak{}the\_\allowbreak{}stage\_\allowbreak{}writes\_\allowbreak{}holds\_\allowbreak{}only\_\allowbreak{}supported\_\allowbreak{}makes} and \texttt{::test\_\allowbreak{}the\_\allowbreak{}vocabulary\_\allowbreak{}is\_\allowbreak{}decided\_\allowbreak{}by\_\allowbreak{}the\_\allowbreak{}served\_\allowbreak{}training\_\allowbreak{}rows\_\allowbreak{}alone}; \href{https://github.com/mlops-2627q1-mds-upc/MLOPS_Recommenditos/blob/main/reports/analysis/served_vocabulary.py}{\texttt{reports/\allowbreak{}analysis/\allowbreak{}served\_\allowbreak{}vocabulary.\allowbreak{}py}} and its results file (2026-10-05); \href{https://github.com/mlops-2627q1-mds-upc/MLOPS_Recommenditos/blob/main/docs/docs/pipeline.md}{pipeline docs}, ``The supported makes''; \href{https://github.com/mlops-2627q1-mds-upc/MLOPS_Recommenditos/blob/main/docs/docs/model-card.md}{model card}, ``Feature space'' and the ladder; \hyperref[EDN-14]{EDN-14}, \hyperref[EDN-42]{EDN-42}, \hyperref[EDN-48]{EDN-48}.},
}

% Generated by tools/edn_latex.py from reports/edn.md.
% Do not edit: change reports/edn.md and run `make edn`.
\ednentry{
  id = {EDN-68},
  title = {The Great Expectations store is a cached build product, the suites are generated from the contract, and each rule sits on the frame where it can hold},
  date = {2026-10-05},
  milestone = {M3: Quality Assurance},
  topic = {Data Validation, Testing Strategy, Reproducibility},
  participants = {@lukas2510, decided by the repository owner},
  decision = {\texttt{configure\_\allowbreak{}gx} rebuilds the Great Expectations file context under \texttt{gx/\allowbreak{}} from scratch on every run, from \texttt{recommenditos/\allowbreak{}data/\allowbreak{}gx\_\allowbreak{}context\_\allowbreak{}configuration.\allowbreak{}py}, \texttt{schema.\allowbreak{}py} and \texttt{params.\allowbreak{}yaml}, and \texttt{gx/\allowbreak{}} is that stage's DVC-cached output: gitignored, pushed, and a dependency of \texttt{validate-\allowbreak{}data}. The context holds two dataframe assets, \texttt{raw} and \texttt{interim}, with one suite, validation definition and checkpoint each. The suites' column lists, dtypes and non-null checks are generated from \texttt{schema.\allowbreak{}py} rather than written by hand. Every rule of issue \#25 sits on the frame where it can hold: \texttt{registration\_\allowbreak{}date} at or before the reference date hard on the interim frame and with \texttt{validate.\allowbreak{}raw\_\allowbreak{}mostly} on the raw one, evaluated on the parsed date; the training price range on the interim frame only; the FR-03 mileage range as a tolerant check on both frames (EDN-72); no agreement rule between \texttt{is\_\allowbreak{}used} and \texttt{offer\_\allowbreak{}type}; and no supported-make rule at all, because \texttt{split} computes the list after this stage and the downstream stages apply it (EDN-48). No suite can pass vacuously: both require \texttt{validate.\allowbreak{}min\_\allowbreak{}rows} rows and the columns their bounds read filled up to \texttt{validate.\allowbreak{}filled\_\allowbreak{}mostly}, and \texttt{validate-\allowbreak{}data} refuses an empty suite and a run that returns fewer results than its suite holds. A failed expectation fails the stage after it has written a git-tracked summary without timestamps and the Data Docs, both under \texttt{reports/\allowbreak{}data-\allowbreak{}validation/\allowbreak{}}; a frame that breaks its contract is not handed to its suite.},
  rationale = {\emph{Alternatives considered:}
\begin{itemize}[leftmargin=*, topsep=2pt]
\item \textbf{Option A: commit \texttt{gx/\allowbreak{}} as configuration, as the course demo does.}
\newline Pros: the suites are readable as JSON on GitHub and the report can link them; it is the layout the course teaches.
\newline Cons: Great Expectations writes a fresh UUID into every object it saves, so measured on 1.23.2, two builds of the same code differ in 6 of the 9 files a context holds, and a second build in place still rewrites 3 of them. Every rerun of the stage would be a diff, two branches that both rebuilt the store would conflict, and the committed JSON says less readably what the module already says. The demo's Parquet assets also record the absolute path of its author's checkout, so its store works on one machine only.
\item \textbf{Option B: commit \texttt{gx/\allowbreak{}}, and pin every id so a rebuild is byte-identical.}
\newline Pros: committed and stable at once, and a test could assert the committed store is the one the code generates.
\newline Cons: Great Expectations overwrites an id passed to a suite, a validation definition or a checkpoint when it saves one (measured), so pinning means rewriting the files after it has written them. That is a post-processor against an internal format, and it would break silently on an upgrade.
\item \textbf{Option C: no store; build an ephemeral context inside \texttt{validate-\allowbreak{}data}.}
\newline Pros: the least code and nothing to version.
\newline Cons: drops the \texttt{configure\_\allowbreak{}gx} stage that \#25 and the course demo prescribe, leaves the suites with no inspectable form at all, and the Data Docs need stored results anyway.
\item \textbf{Option D (chosen): the store is \texttt{configure\_\allowbreak{}gx}'s cached DVC output, rebuilt from scratch.}
\newline Pros: no UUID noise in Git; \texttt{dvc pull} restores the exact store a validation ran against; declaring it as an output is what puts \texttt{configure\_\allowbreak{}gx} in the DAG, which the demo's stage is not; a rebuild cannot keep a suite or an expectation the module no longer defines, while the demo's script fails outright on its second run because the data source already exists.
\newline Cons: the suites are not readable on GitHub, so a reviewer reads the module that builds them, and the Data Docs render them after \texttt{dvc pull}. A rebuild changes the output hash even when no rule changed, so \texttt{validate-\allowbreak{}data} reruns whenever \texttt{configure\_\allowbreak{}gx} does -- which is only when the module, the contract or one of its parameters changed.
\end{itemize}
\par
\emph{Why:} The store question was decided by measurement rather than taste: an artefact that changes on every rebuild with no change in meaning is a build product, and this project already treats build products as DVC outputs (\texttt{models/\allowbreak{}}, EDN-56). The other choices follow from what the tool can and cannot do, each found by running it on our frames before building on it (\texttt{reports/\allowbreak{}analysis/\allowbreak{}gx\_\allowbreak{}probe.\allowbreak{}py}):
\par
\begin{itemize}[leftmargin=*, topsep=2pt]
\item \textbf{Dataframe assets, one checkpoint per frame.} A dataframe asset records no path, so the store works in every clone, and \texttt{validate-\allowbreak{}data} validates the frame its contract has just accepted instead of reading the 215 MB file twice. \texttt{Checkpoint.\allowbreak{}run} takes one set of batch parameters for all its validation definitions, so two frames need two checkpoints. The Data Docs are rendered once per stage run rather than by a checkpoint action, which would render the whole site once per frame.
\item \textbf{Results and Data Docs outside \texttt{gx/\allowbreak{}}.} They go to \texttt{reports/\allowbreak{}data-\allowbreak{}validation/\allowbreak{}results/\allowbreak{}} and \texttt{data-\allowbreak{}docs/\allowbreak{}}, \texttt{validate-\allowbreak{}data}'s own outputs. Written inside \texttt{gx/\allowbreak{}}, every validation would modify \texttt{configure\_\allowbreak{}gx}'s output and \texttt{dvc status} would report that stage changed after every run.
\item \textbf{Generated contract expectations.} \#25 puts the column schema and types in the suite; generating them from \texttt{schema.\allowbreak{}py} keeps it the one place the contract is written, and makes the Data Docs a complete statement of each frame. They pass whenever they are reached, because \texttt{Schema.\allowbreak{}validate} runs first, and they are weaker than it: Great Expectations compares only a column's scalar type, so it cannot tell nullable, extension or \texttt{object} dtypes from their counterparts. Measured, \texttt{datetime64[us]} passes \texttt{datetime64[ns]}, \texttt{boolean} passes \texttt{bool}, \texttt{Float64} passes \texttt{float64}, and \texttt{object} and \texttt{string[python]} pass \texttt{str}, so the declared dtype stays \texttt{schema.\allowbreak{}py}'s to check (EDN-31 found the datetime unit gap in Pandera too).
\item \textbf{The fill-rate rule checks preprocessing, not the scrape.} This was true when \#25 merged: the interim suite asserted filled the columns the raw contract declared non-null, and a gap in a scrape failed \texttt{download} first. \href{https://github.com/mlops-2627q1-mds-upc/MLOPS_Recommenditos/issues/78}{\#78} moved those fill rates into the raw suite (EDN-77).
\item \textbf{No vacuous pass.} Every column rule holds over no rows, and Great Expectations skips missing values, so an empty frame or a bounded column that has emptied out would pass every rule that reads it. Both suites therefore check the row count against \texttt{validate.\allowbreak{}min\_\allowbreak{}rows} and the fill rate of \texttt{registration\_\allowbreak{}date} and \texttt{mileage\_\allowbreak{}km\_\allowbreak{}raw} against \texttt{validate.\allowbreak{}filled\_\allowbreak{}mostly}, and \texttt{validate-\allowbreak{}data} refuses a suite with no expectations and a run with fewer results than expectations.
\item \textbf{The raw date rule on the parsed date.} Great Expectations types a bound as a number or a date: a text bound is parsed into a \texttt{date}, and comparing it with the raw file's text column raises. So the raw suite sees the published file with \texttt{registration\_\allowbreak{}date} cast by the interim contract's own cast, and every other column untouched. It therefore flags exactly the rows the interim rule then removes.
\item \textbf{What was taken out of \#25's list, and why.} An \texttt{is\_\allowbreak{}used} / \texttt{offer\_\allowbreak{}type} agreement rule would fail 18,108 of the published file's 113,708 used passenger-car rows by design (EDN-23). The supported-make list cannot be a rule here: \texttt{split} computes it after this stage, and the processed sets deliberately keep every make (EDN-48), so no frame \texttt{validate-\allowbreak{}data} sees is meant to satisfy it. Its guarantee is enforced where the list is applied, in \texttt{features} (EDN-67), and checked again by \texttt{train} and \texttt{evaluate}.
\item \textbf{How a failure gates.} Any failed expectation fails the stage, as the structural contract already did, and every frame's result is read; the demo only logs a count, and only for the first validation result of its checkpoint run. The summary and the Data Docs are written before the raise, so a failed run says which rule broke and shows a sample of the failing values.
\end{itemize}
\par
On the real snapshot every expectation of both suites passes. The raw suite's date rule finds 164 registrations after the reference date, 0.1388\,\% of the 118,161 dated rows and so inside the 1\,\% tolerance, and its mileage rule three readings above 1,000,000 km, which EDN-72 is about.
\par
One cost of the tool was found only on the real data: Great Expectations hashes every in-memory batch into a fingerprint for the result's markers, which on the raw frame doubles the stage's memory. Measured with \texttt{reports/\allowbreak{}analysis/\allowbreak{}gx\_\allowbreak{}fingerprint\_\allowbreak{}cost.\allowbreak{}py}, the stage takes 14 s and peaks at 4.4 GB with the fingerprint, and 11 s and 2.2 GB without it, the raw frame itself being most of that. \texttt{validate-\allowbreak{}data} switches the fingerprint off for its own runs, and a test fails if an upgrade makes that switch stop working.},
  ai = {Information seeking, Alternative generation, Alternative assessment, Recommendation, Solution generation},
  response = {Accepted with modifications},
  assessment = {AI implemented \#25 and proposed every choice above, with the alternatives, after testing three assumptions before building on them, each of which changed the design. Running every expectation kind on a satisfying and a breaking frame showed that Great Expectations works on pandas 3 but cannot bound a text date. Measuring what the store writes across two builds turned ``commit \texttt{gx/\allowbreak{}}'', the demo's layout and the stub's own plan, into a cached output. Checking each listed rule against the point in the DAG where it runs found that the supported-make rule could not mean anything there, which the issue, \texttt{schema.\allowbreak{}py} and \texttt{docs/\allowbreak{}docs/\allowbreak{}pipeline.\allowbreak{}md} all still assumed. Lukas confirmed the setup as proposed. An independent AI review of PR \#76 then found four things the first version got wrong, all fixed before merge: the fill-rate rule was described as catching a scrape gap that the raw contract already stops at \texttt{download}; a zero-row frame, an emptied bounded column or an empty suite passed vacuously; the type blind spots were understated as two, while \texttt{object}, \texttt{string[python]} and \texttt{Float64} pass as well; and the generated not-null checks had no test, so deleting them passed the suite.},
  aievidence = {Claude Code session on 2026-10-05 implementing issue \#25: throwaway probes on the fixture and the real snapshot before any stage code, the store measurement, and the options above presented with a recommendation, which Lukas confirmed; then an independent AI review of PR \#76, whose findings are listed above.},
  evidence = {\href{https://github.com/mlops-2627q1-mds-upc/MLOPS_Recommenditos/blob/main/recommenditos/data/gx_context_configuration.py}{\texttt{recommenditos/\allowbreak{}data/\allowbreak{}gx\_\allowbreak{}context\_\allowbreak{}configuration.\allowbreak{}py}} and \href{https://github.com/mlops-2627q1-mds-upc/MLOPS_Recommenditos/blob/main/recommenditos/data/validate_data.py}{\texttt{recommenditos/\allowbreak{}data/\allowbreak{}validate\_\allowbreak{}data.\allowbreak{}py}}, whose docstrings give each rule's reason; \href{https://github.com/mlops-2627q1-mds-upc/MLOPS_Recommenditos/blob/main/tests/test_validate_data.py}{\texttt{tests/\allowbreak{}test\_\allowbreak{}validate\_\allowbreak{}data.\allowbreak{}py}}; \href{https://github.com/mlops-2627q1-mds-upc/MLOPS_Recommenditos/blob/main/reports/analysis/gx_probe.py}{\texttt{reports/\allowbreak{}analysis/\allowbreak{}gx\_\allowbreak{}probe.\allowbreak{}py}} and \href{https://github.com/mlops-2627q1-mds-upc/MLOPS_Recommenditos/blob/main/reports/analysis/gx_fingerprint_cost.py}{\texttt{gx\_\allowbreak{}fingerprint\_\allowbreak{}cost.\allowbreak{}py}} with their results; \href{https://github.com/mlops-2627q1-mds-upc/MLOPS_Recommenditos/blob/main/reports/data-validation/summary.json}{\texttt{reports/\allowbreak{}data-\allowbreak{}validation/\allowbreak{}summary.\allowbreak{}json}}; \href{https://github.com/mlops-2627q1-mds-upc/MLOPS_Recommenditos/blob/main/docs/docs/pipeline.md}{pipeline docs}, ``Data validation''; \hyperref[EDN-22]{EDN-22}, \hyperref[EDN-23]{EDN-23}, \hyperref[EDN-31]{EDN-31}, \hyperref[EDN-48]{EDN-48}; \href{https://github.com/mlops-2627q1-mds-upc/MLOPS_Recommenditos/issues/25}{issue \#25}, \href{https://github.com/mlops-2627q1-mds-upc/MLOPS_Recommenditos/pull/76}{PR \#76}.},
}

% Generated by tools/edn_latex.py from reports/edn.md.
% Do not edit: change reports/edn.md and run `make edn`.
\ednentry{
  id = {EDN-69},
  title = {CodeCarbon measures the fit alone, offline and per process, and its figure is reported as the estimate it is},
  date = {2026-10-05},
  milestone = {M3: Quality Assurance},
  topic = {Energy Efficiency, Experiment Tracking, Evaluation Protocol},
  participants = {@lukas2510, decided by the repository owner},
  decision = {\texttt{train} measures exactly the \texttt{fit\_\allowbreak{}variant} call with CodeCarbon 3.3.1's \texttt{OfflineEmissionsTracker}, in \texttt{process} mode, against the Spanish grid (\texttt{ESP}, 174.05 g CO2eq/kWh) and with \texttt{measure\_\allowbreak{}power\_\allowbreak{}secs: 15}, all three from \texttt{params.\allowbreak{}yaml}'s \texttt{train.\allowbreak{}energy}, which the \texttt{train} stages declare. CodeCarbon's row goes to \texttt{reports/\allowbreak{}emissions/\allowbreak{}<variant>.\allowbreak{}csv} and into the MLflow run that also receives the variant's test metrics, as metrics, with the CPU power method and the machine behind them as parameters. A \texttt{compare-\allowbreak{}energy} stage joins each model, its row and its metrics into \texttt{reports/\allowbreak{}figures/\allowbreak{}energy-\allowbreak{}vs-\allowbreak{}error.\allowbreak{}\{csv,png\}}. On every machine we train on the figure is CodeCarbon's estimate, not a reading: 10 W for the RAM plus the CPU's TDP times the share of its logical CPUs the fit kept busy. We keep that estimate in the pipeline, word every comparison as one of fit times expressed in energy, within one run on one machine, and checked the estimate once against the RAPL energy counter, with read access granted for that check only.},
  rationale = {\emph{Alternatives considered:}
\begin{itemize}[leftmargin=*, topsep=2pt]
\item \textbf{How the figure is obtained, given that it is an estimate.}
\par
\begin{itemize}[leftmargin=*, topsep=2pt]
\item \textbf{Option A: keep CodeCarbon's estimate and word it as what it is.}
\newline Pros: nothing changes on anyone's machine; the same mechanism and settings everywhere, hermetic, with no network call; it is the prescribed tool used as prescribed.
\newline Cons: the energy axis adds no information beyond the fit time it is computed from, and the absolute watt-hours are model output that depends on the machine and on its load.
\item \textbf{Option B: grant read access to RAPL on Lukas's laptop permanently,} with a udev or tmpfiles rule, so that CodeCarbon reads the package counter there.
\newline Pros: a hardware counter rather than a model.
\newline Cons: RAPL is root-only on purpose since kernel 5.10, because unprivileged power readings are a side channel that has been used to extract keys (Platypus, CVE-2020-8694), and this laptop has SGX. RAPL counts the whole package and CodeCarbon does not apportion it to the process, so on a laptop several agents share -- the 1-minute load ranged from 10 to 31 on the day of the measurement -- it would charge each fit for the others. Only Lukas's runs would read it, so a teammate's \texttt{dvc repro} would write a figure of another kind into the same files. And it is a standing change to a personal machine that somebody has to remember.
\item \textbf{Option C (chosen): A, plus a one-off check of the estimate against the counter.} Read access is granted with a \texttt{chmod} that does not survive a reboot, the four fits and an idle baseline are measured with both, and access is revoked again.
\newline Pros: the report can say what the estimate is worth instead of only what it is, at no standing cost to the laptop's security, and the pipeline stays as in A.
\newline Cons: it needs Lukas's \texttt{sudo} and an otherwise idle laptop, and it validates one machine on one day, not every run.
\item \textbf{Option D: compare the ladder on fit seconds and CPU-seconds, and keep CodeCarbon's figure only as NFR-10's log.}
\newline Pros: both are direct measurements.
\newline Cons: it drops the energy comparison issue \#38 and the rubric ask for, and the comparison would show the same ratios under another name.
\end{itemize}
\item \textbf{\texttt{process} against \texttt{machine} mode.} \texttt{machine} charges the fit for the load of the whole machine: on the shared laptop, a process that only slept was charged 4.5 W of CPU in \texttt{machine} mode and 0.3 W in \texttt{process} mode. \texttt{process} reads the fit's own CPU time.
\item \textbf{The online tracker against the offline one with a pinned country.} Per CodeCarbon's source, the online tracker geolocates the machine through a call to geojs, writes its latitude and longitude into the CSV, and queries the AWS, Azure and GCP metadata endpoint, switching to the provider's regional factor when one answers, which on a GitHub runner, an Azure VM, one would. The emissions of an identical fit would then depend on where it ran, a git-tracked file would carry a contributor's location, and the test suite would make network calls. With every socket connection made to fail, the offline tracker made none.
\item \textbf{The demo's \texttt{measure\_\allowbreak{}power\_\allowbreak{}secs=\allowbreak{}1} against CodeCarbon's default of 15.} The power is sampled every second by a separate thread either way; this setting is how often the samples are integrated, and at 1 s that integration races the tracker's own shutdown. Against CodeCarbon's own model, 1 s gave 0.86x to 1.07x on a single-threaded LightGBM fit and up to 4.31x on a busy loop's RAM half, an interval counted several times over; 15 s gave 0.987x to 0.993x on the same fit.
\end{itemize}
\par
\emph{Why:} Issue \#38 was written on the assumption that CodeCarbon measures the energy of a fit, and the first thing done was to test that on the laptop the ladder is trained on, before any number was quoted. It does not. Since kernel 5.10 RAPL is root-only, so CodeCarbon falls back to its \texttt{cpu\_\allowbreak{}load} model, with the i5-10210U's TDP of 15 W from its own table and a constant 10 W for the RAM, its floor of two DIMMs for any x86 machine. The energy of a fit is then 10 W times its wall time plus 15 W / 8 logical CPUs times its CPU time, and at \texttt{train.\allowbreak{}num\_\allowbreak{}threads: 1} that is its duration times 11.75 to 11.80 W, to within 1\,\% over repeated fits. The RAM constant is 83\,\% of a LightGBM fit's figure. On a GitHub runner, an AMD EPYC 7763 with four vCPUs and no powercap interface at all, the same model charges 80 W per second, because it divides the chip's 280 W TDP by the four vCPUs the VM sees. So the figure is a fit time multiplied by a machine-specific constant, and it moves with the machine's load as well as with the work: the same \texttt{lgbm-\allowbreak{}basic} fit took 33.0 s on a laptop shared with other jobs and 8.6 s on a quiet one, and was charged accordingly.
\par
That makes the comparison the issue asked for weaker than it reads, not wrong. Between two variants fitted in one run on one machine the ratio of their energies is the ratio of their fit times, which is a fair measure of their relative cost. What it cannot be is a statement in watt-hours that holds across machines, or evidence that a variant uses energy differently from how it uses time. So the pipeline keeps the estimate (option A), every place that shows the figure says what it is -- the CSV, the \texttt{codecarbon.\allowbreak{}cpu\_\allowbreak{}power\_\allowbreak{}method} parameter of every run, the subtitle of the figure, the model card and the report -- and the one check that can say how far the estimate is from the hardware was made once (option C) rather than built into every run (option B). NFR-10's specification says the same since this decision, because ``CodeCarbon measures'' was not true of any machine we have.
\par
The RAPL check, \texttt{reports/\allowbreak{}analysis/\allowbreak{}rapl\_\allowbreak{}validation.\allowbreak{}py}, measured an idle baseline and three fits of each variant on an idle laptop with the package, core, uncore, DRAM and platform counters, CodeCarbon's forced \texttt{cpu\_\allowbreak{}load} estimate and CodeCarbon's own RAPL figure, all around the same fit. \textbf{[pending RAPL check: idle draw, estimate against the counter per variant, whether the ranking and the ratios hold.]}
\par
The rest of the design follows from treating the figure as evidence that has to be checkable. Only the fit is inside the tracker, as the demo does, which \texttt{tests/\allowbreak{}test\_\allowbreak{}model.\allowbreak{}py::test\_\allowbreak{}only\_\allowbreak{}the\_\allowbreak{}fit\_\allowbreak{}is\_\allowbreak{}measured} asserts by making the read and the write slow. Every CodeCarbon argument is passed explicitly, because it reads \texttt{\textasciitilde{}/\allowbreak{}.\allowbreak{}codecarbon.\allowbreak{}config}, \texttt{.\allowbreak{}/\allowbreak{}.\allowbreak{}codecarbon.\allowbreak{}config} and any \texttt{CODECARBON\_\allowbreak{}*} variable for every argument it is not given, so a contributor's configuration could switch on its API output or force a CPU power. CodeCarbon swallows every exception in its constructor, \texttt{start} and \texttt{stop}, so the module reads its row back by CodeCarbon's run id and by column name and fails the stage when it is missing, and refuses a country CodeCarbon would silently measure against the world average. The row is joined to its model by that run id, which \texttt{model.\allowbreak{}json} records, and \texttt{fit\_\allowbreak{}cpu\_\allowbreak{}seconds} is logged beside the energy because it is a measurement and the CPU half is proportional to it. DVC deletes a stage's outputs before running it, so the CSV CodeCarbon appends to holds exactly one row after \texttt{dvc repro} (measured: one row after three \texttt{dvc repro -\allowbreak{}-\allowbreak{}force}, three with \texttt{persist: true}).
\par
What the comparison concluded, on the ladder of \textbf{[pending final run: commit]}: \textbf{[pending final run: per-variant fit time, energy, emissions and MdAPE, and whether \texttt{lgbm-\allowbreak{}extended} earns its extra training energy]}. The whole ladder is \textbf{[pending final run]} of fitting, against NFR-10's 15-minute budget.},
  ai = {Information seeking, Alternative generation, Alternative assessment, Recommendation, Solution generation},
  response = {Accepted},
  assessment = {AI found the limitation the decision is about, by reading CodeCarbon's source and then measuring it on the laptop and on a GitHub runner, before writing the module; the issue, the course demo and the report draft all assumed a measurement. It laid out the four options with the security and the shared-machine arguments against B, recommended A with C, and Lukas chose that. It also chose the tracker's settings on measurements rather than on the demo: \texttt{process} mode, the offline tracker, and 15 s rather than the demo's 1 s once it was shown to race the tracker's shutdown. It did not grant itself RAPL access by other means, although Lukas's \texttt{docker} group would have allowed it, because that would have bypassed the permission boundary this decision is partly about.},
  aievidence = {Claude Code session on 2026-10-05 implementing issue \#38: probes of CodeCarbon 3.3.1 with every socket blocked, on four workloads in both modes and at both integration intervals, a DVC repository testing the append, and a temporary CI step on pull request \#68 that printed what CodeCarbon does on a GitHub runner; then the four options above with a recommendation, which Lukas decided as option C.},
  evidence = {\href{https://github.com/mlops-2627q1-mds-upc/MLOPS_Recommenditos/blob/main/recommenditos/modeling/energy.py}{\texttt{recommenditos/\allowbreak{}modeling/\allowbreak{}energy.\allowbreak{}py}}, whose docstring states what the figure is, and \href{https://github.com/mlops-2627q1-mds-upc/MLOPS_Recommenditos/blob/main/recommenditos/modeling/compare_energy.py}{\texttt{recommenditos/\allowbreak{}modeling/\allowbreak{}compare\_\allowbreak{}energy.\allowbreak{}py}}; \texttt{train.\allowbreak{}energy} in \href{https://github.com/mlops-2627q1-mds-upc/MLOPS_Recommenditos/blob/main/params.yaml}{\texttt{params.\allowbreak{}yaml}}; \href{https://github.com/mlops-2627q1-mds-upc/MLOPS_Recommenditos/blob/main/tests/test_energy.py}{\texttt{tests/\allowbreak{}test\_\allowbreak{}energy.\allowbreak{}py}} and \texttt{tests/\allowbreak{}test\_\allowbreak{}model.\allowbreak{}py::test\_\allowbreak{}every\_\allowbreak{}training\_\allowbreak{}run\_\allowbreak{}records\_\allowbreak{}its\_\allowbreak{}emissions\_\allowbreak{}beside\_\allowbreak{}its\_\allowbreak{}accuracy} and \texttt{::test\_\allowbreak{}only\_\allowbreak{}the\_\allowbreak{}fit\_\allowbreak{}is\_\allowbreak{}measured}; \href{https://github.com/mlops-2627q1-mds-upc/MLOPS_Recommenditos/blob/main/reports/analysis/codecarbon_validity.py}{\texttt{reports/\allowbreak{}analysis/\allowbreak{}codecarbon\_\allowbreak{}validity.\allowbreak{}py}} with its results, \href{https://github.com/mlops-2627q1-mds-upc/MLOPS_Recommenditos/blob/main/reports/analysis/codecarbon_github_runner.txt}{\texttt{codecarbon\_\allowbreak{}github\_\allowbreak{}runner.\allowbreak{}txt}} and \href{https://github.com/mlops-2627q1-mds-upc/MLOPS_Recommenditos/blob/main/reports/analysis/rapl_validation.py}{\texttt{rapl\_\allowbreak{}validation.\allowbreak{}py}} with its results; \texttt{reports/\allowbreak{}figures/\allowbreak{}energy-\allowbreak{}vs-\allowbreak{}error.\allowbreak{}png}; \href{https://github.com/mlops-2627q1-mds-upc/MLOPS_Recommenditos/blob/main/docs/docs/specification.md}{specification} NFR-10; \href{https://github.com/mlops-2627q1-mds-upc/MLOPS_Recommenditos/blob/main/docs/docs/model-card.md}{model card}, Environmental Impact; \href{https://github.com/mlops-2627q1-mds-upc/MLOPS_Recommenditos/issues/38}{issue \#38}, \href{https://github.com/mlops-2627q1-mds-upc/MLOPS_Recommenditos/pull/68}{pull request \#68}.},
}

% Generated by tools/edn_latex.py from reports/edn.md.
% Do not edit: change reports/edn.md and run `make edn`.
\ednentry{
  id = {EDN-70},
  title = {\texorpdfstring{\texttt{n\_\allowbreak{}estimators}}{n\_estimators} is a ceiling of 5,000 that early stopping stays under, not a tree count},
  date = {2026-10-05},
  milestone = {M3: Quality Assurance},
  topic = {Modelling Approach, Hyperparameters, Energy Efficiency},
  participants = {@lukas2510, decided by the repository owner},
  decision = {Both LightGBM variants get \texttt{n\_\allowbreak{}estimators: 5000}, with \texttt{early\_\allowbreak{}stopping\_\allowbreak{}rounds: 50} unchanged, so early stopping on the validation split decides how many trees each model keeps: round 1,243 for \texttt{lgbm-\allowbreak{}basic} and 1,647 for \texttt{lgbm-\allowbreak{}extended} on the real snapshot.
At 1,000 the ceiling had ended both fits before early stopping could.
\texttt{train} now records in \texttt{model.\allowbreak{}json} and in the MLflow run whether early stopping or the ceiling ended each fit (\texttt{early\_\allowbreak{}stopped}, beside \texttt{boosting\_\allowbreak{}rounds}, the rounds the fit ran), and logs a warning when it was the ceiling.
\texttt{learning\_\allowbreak{}rate} and \texttt{num\_\allowbreak{}leaves} are not changed by this entry.},
  rationale = {\emph{Alternatives considered:}
\begin{itemize}[leftmargin=*, topsep=2pt]
\item \textbf{Option A: keep 1,000.}
\newline Pros: no retrain, and the report's numbers stay where they are.
\newline Cons: the parameter, not the validation curve, decides the tree count of both variants, so every later comparison -- of the feature sets, of \texttt{learning\_\allowbreak{}rate}, of \texttt{num\_\allowbreak{}leaves} -- is judged between models cut off by a budget.
And it hides itself: \texttt{lgbm-\allowbreak{}extended} once recorded a best iteration of 996 out of 1,000, which reads like early stopping and was not.
\item \textbf{Option B: 2,000, just above where early stopping lands.}
\newline Pros: the smallest worst case of every cost that grows with the trees.
\newline Cons: \texttt{lgbm-\allowbreak{}extended} ran to round 1,697, so 2,000 leaves 303 rounds of slack.
A modest change to the data -- a new snapshot, or the Great Expectations work on the stages before \texttt{split} -- or a lower learning rate would bind it again, and only the new warning would say so.
\item \textbf{Option C (chosen): 5,000.}
\newline Pros: three times the later best round and four times the earlier, so it does not bind under any change measured so far: in a 12-point grid of \texttt{learning\_\allowbreak{}rate} and \texttt{num\_\allowbreak{}leaves} around the committed values, the slowest point (\texttt{learning\_\allowbreak{}rate} 0.025, \texttt{num\_\allowbreak{}leaves} 31) ran 4,116 rounds.
And its worst case stays inside every budget the trees drive: a fit that used all 5,000 trees would take about 80 s for \texttt{lgbm-\allowbreak{}extended} against NFR-10's 15 minutes, write a booster of about 33 MB against NFR-04's 1 GB image, and spend about 105 ms on one row's native SHAP export against NFR-02's 200 ms for the whole of \texttt{/\allowbreak{}predict}.
\newline Cons: it is a value with a stated rationale, not a measured optimum, and the latency half of that rationale is a laptop measurement while NFR-02 is judged on the course VM, so its margin is less certain than the fit-time one.
\item \textbf{Option D: effectively unbounded, 20,000 or more, as the exploratory fits used.}
\newline Pros: it could never bind in any plausible configuration.
\newline Cons: the ceiling would stop being a guard.
A configuration whose validation curve keeps creeping down -- a lower learning rate, more data -- would run to a model whose SHAP export alone exceeds NFR-02 (about 210 ms at 10,000 trees on the same laptop), with a fit time, an energy figure and an artefact size nobody chose.
\item \textbf{Option E: fix the tree count at the round early stopping chose (\texttt{n\_\allowbreak{}estimators: 1243}, early stopping off).}
\newline Pros: the same model today, and a fit that no longer reads the validation split.
\newline Cons: it bakes one snapshot's answer into a parameter.
The next retrain, which EDN-12 makes a response to drift, would get it wrong in either direction, where early stopping decides again on every retrain.
\end{itemize}
\par
\emph{Why:} The value was found by measurement rather than guessed, in two steps that never open the test split.
\texttt{reports/\allowbreak{}analysis/\allowbreak{}early\_\allowbreak{}stopping\_\allowbreak{}ceiling.\allowbreak{}py} fitted each variant once at a ceiling of 20,000 with every other setting as committed, through the stage's own \texttt{fit\_\allowbreak{}variant}, on the \texttt{train} and \texttt{validation} matrices only: early stopping ended both fits, at best rounds 1,243 and 1,647 after 1,293 and 1,697 rounds.
The ceiling was then chosen from those numbers and the per-tree costs the same script measured, and only the committed \texttt{dvc repro} read the test split, through \texttt{evaluate}.
\par
\textbf{What the extra trees are worth, which is less than issue \#64 expected.} The issue read the 1,000-tree numbers as a lower bound on what the hyperparameters could reach.
They were very nearly the bound itself.
On validation, from 1,000 trees to the round early stopping chose, L1 in log space falls by 0.2\,\% and 0.6\,\% while MdAPE stays flat (6.947\,\% to 6.956\,\%, 6.448\,\% to 6.453\,\%): L1 is a mean and MdAPE a median, so the last trees improve the tail of the errors and not the typical one.
On the test split the committed run moves MdAPE from 6.827\,\% to 6.808\,\% for \texttt{lgbm-\allowbreak{}basic} and from 6.325\,\% to 6.256\,\% for \texttt{lgbm-\allowbreak{}extended}, within 20\,\% from 89.68\,\% to 89.81\,\% and from 90.29\,\% to 90.40\,\%, SC-03 from 43.9\,\% to 44.0\,\% and from 48.0\,\% to 48.6\,\%.
SC-04's worst segment, cars over 20 years, moves from 17.12\,\% to 16.77\,\% for \texttt{lgbm-\allowbreak{}basic} and from 17.86\,\% to 18.45\,\% for \texttt{lgbm-\allowbreak{}extended}; SC-06 from 1.208x to 1.213x and from 1.583x to 1.646x.
No criterion changes its verdict for any variant, \texttt{b0} and \texttt{b1} are unchanged, and EDN-62's choice of candidate holds more firmly than before: \texttt{lgbm-\allowbreak{}extended}'s pooled advantage grows from 0.50 pp to 0.55 pp, and its disadvantage on a request carrying only the required fields from 1.77 pp to 2.04 pp (10.30\,\% against 8.26\,\%).
The training L1 falls much further than the validation L1 -- 0.0688 to 0.0655 against 0.1005 to 0.1003 for the candidate -- so the added trees mostly fit the training rows more closely, and on \texttt{lgbm-\allowbreak{}extended} that comes with a stronger dependence on its optional fields.
\par
\textbf{What they cost, which is linear in the trees.} 1.29x and 1.70x the boosting rounds, so the same factor of fit time and of the energy the fit draws; the boosters grow from 6.2 to 7.6 MB and from 6.8 to 10.9 MB; one row's native SHAP export, the FR-08 path, from 18.9 to 26.1 ms and from 22.3 to 35.4 ms on the laptop that measured it.
The committed run fits the candidate in 8.3 s and the whole ladder in 37.6 s, far inside NFR-10's 15 minutes; fits of the identical trees on 2026-10-05 and 2026-10-06 ranged from 8.0 s to 10.0 s and from 26.4 s to 32.9 s with machine load, so the rounds, not the seconds, are the measurement.
Bought at that price, 0.02 pp of MdAPE on the candidate is not much; it is accepted because the alternative is a model whose size is set by an arbitrary parameter rather than by its data, which is what makes every other comparison in the ladder trustworthy.
\par
\textbf{The issue's ``Care needed'' claims, checked in the code.} Early stopping reads the validation split and nothing else: \texttt{LightGBMModel.\allowbreak{}fit} passes \texttt{validation} as the only evaluation set, under \texttt{first\_\allowbreak{}metric\_\allowbreak{}only}, and \texttt{train.\allowbreak{}read\_\allowbreak{}matrices} reads only \texttt{train} and \texttt{validation}; the test split is read by \texttt{evaluate} alone, after training.
One claim needs a qualification rather than a correction: \texttt{early\_\allowbreak{}stopping\_\allowbreak{}rounds} bounds how many rounds the fit runs past its best one, not the optimism of choosing that best round on the same split it is scored on.
That optimism is what the untouched test split measures, and here test and validation moved together and by little.
\par
\textbf{A record that can say whether the ceiling bound.} Investigating the issue showed that \texttt{best\_\allowbreak{}iteration} alone cannot: lightgbm 4.7 reports the best round inside the budget whether or not early stopping fired, so \texttt{best\_\allowbreak{}iteration\_\allowbreak{}} is never 0 at the ceiling, as a code comment in \texttt{model.\allowbreak{}py} claimed, and the 996 of \texttt{lgbm-\allowbreak{}extended} was a fit the ceiling ended.
\texttt{model.\allowbreak{}json} therefore records the rounds run as well, and \texttt{early\_\allowbreak{}stopped} is true when the fit ran the full patience past its best round.},
  ai = {Information seeking, Alternative generation, Alternative assessment, Recommendation, Solution generation},
  response = {Accepted},
  assessment = {The method was set in the brief: raise the ceiling until early stopping fires, find where with one generous exploratory fit per variant, and commit a value with clear headroom.
AI ran the measurement, chose 5,000 by the worst-case argument above rather than by a multiple of the observed best round alone, and laid out options A to E.
Lukas confirmed 5,000 after seeing the exploratory fits, the per-tree costs and the committed run's numbers.
Two of its findings went beyond the brief.
It established that the issue's premise, that the 1,000-tree models were materially short of what the hyperparameters could reach, does not hold: the extra trees are worth 0.02 pp and 0.07 pp of test MdAPE.
And it found, by probing lightgbm 4.7's early-stopping callback, that the code comment explaining the \texttt{best\_\allowbreak{}iteration\_\allowbreak{} or num\_\allowbreak{}trees()} fallback was wrong and that the record could not tell a binding ceiling from early stopping, which is why \texttt{early\_\allowbreak{}stopped} and \texttt{boosting\_\allowbreak{}rounds} were added with tests at the \texttt{train} seam.},
  aievidence = {Claude Code session on 2026-10-05 implementing issue \#64: the exploratory fits are \texttt{reports/\allowbreak{}analysis/\allowbreak{}early\_\allowbreak{}stopping\_\allowbreak{}ceiling.\allowbreak{}py} and its results file; the committed run is the \texttt{dvc.\allowbreak{}lock} of PR \#77, tracked as one MLflow run per variant in \texttt{recommenditos-\allowbreak{}price} with \texttt{best\_\allowbreak{}iteration}, \texttt{boosting\_\allowbreak{}rounds} and \texttt{early\_\allowbreak{}stopped} logged.},
  evidence = {\href{https://github.com/mlops-2627q1-mds-upc/MLOPS_Recommenditos/issues/64}{issue \#64}; \href{https://github.com/mlops-2627q1-mds-upc/MLOPS_Recommenditos/pull/77}{PR \#77}; \href{https://github.com/mlops-2627q1-mds-upc/MLOPS_Recommenditos/blob/main/params.yaml}{\texttt{params.\allowbreak{}yaml}}, \texttt{train.\allowbreak{}variants}; \href{https://github.com/mlops-2627q1-mds-upc/MLOPS_Recommenditos/blob/main/recommenditos/modeling/model.py}{\texttt{recommenditos/\allowbreak{}modeling/\allowbreak{}model.\allowbreak{}py}}, \texttt{LightGBMModel.\allowbreak{}fit}; \texttt{tests/\allowbreak{}test\_\allowbreak{}model.\allowbreak{}py::test\_\allowbreak{}a\_\allowbreak{}lightgbm\_\allowbreak{}record\_\allowbreak{}says\_\allowbreak{}early\_\allowbreak{}stopping\_\allowbreak{}ended\_\allowbreak{}the\_\allowbreak{}fit}, \texttt{::test\_\allowbreak{}early\_\allowbreak{}stopped\_\allowbreak{}is\_\allowbreak{}decided\_\allowbreak{}by\_\allowbreak{}the\_\allowbreak{}patience\_\allowbreak{}not\_\allowbreak{}by\_\allowbreak{}the\_\allowbreak{}best\_\allowbreak{}round} and \texttt{::test\_\allowbreak{}a\_\allowbreak{}budget\_\allowbreak{}early\_\allowbreak{}stopping\_\allowbreak{}never\_\allowbreak{}reaches\_\allowbreak{}is\_\allowbreak{}recorded\_\allowbreak{}as\_\allowbreak{}a\_\allowbreak{}budget\_\allowbreak{}that\_\allowbreak{}bound}; \href{https://github.com/mlops-2627q1-mds-upc/MLOPS_Recommenditos/blob/main/reports/analysis/early_stopping_ceiling.py}{\texttt{reports/\allowbreak{}analysis/\allowbreak{}early\_\allowbreak{}stopping\_\allowbreak{}ceiling.\allowbreak{}py}} and \href{https://github.com/mlops-2627q1-mds-upc/MLOPS_Recommenditos/blob/main/reports/analysis/tuning_sweep.py}{\texttt{tuning\_\allowbreak{}sweep.\allowbreak{}py}} with their results files (2026-10-05); \href{https://github.com/mlops-2627q1-mds-upc/MLOPS_Recommenditos/blob/main/docs/docs/model-card.md}{model card}, Training and the ladder; \hyperref[EDN-53]{EDN-53}, \hyperref[EDN-62]{EDN-62}.},
}

% Generated by tools/edn_latex.py from reports/edn.md.
% Do not edit: change reports/edn.md and run `make edn`.
\ednentry{
  id = {EDN-71},
  title = {Secrets are scanned by gitleaks before every commit and over the whole history in CI, with a rule of our own for the DagsHub token},
  date = {2026-10-05},
  milestone = {M3: Quality Assurance},
  topic = {CI/CD, Security, Credential handling},
  participants = {@lukas2510, decided by the repository owner},
  decision = {NFR-15 is enforced with gitleaks 8.30.1, run twice with one configuration, \texttt{.\allowbreak{}gitleaks.\allowbreak{}toml}.
The \texttt{gitleaks} pre-commit hook scans the staged change before the commit exists.
The \texttt{Secret scan (gitleaks)} CI job scans every commit GitHub publishes for the repository -- every branch, every tag and the head of every pull request ever opened -- on every pull request and every push to \texttt{main}, with no paths filter, using the release binary pinned by version and SHA-256 rather than \texttt{gitleaks/\allowbreak{}gitleaks-\allowbreak{}action}.
\texttt{.\allowbreak{}gitleaks.\allowbreak{}toml} keeps gitleaks' default rules and adds one, \texttt{dagshub-\allowbreak{}token}, which recognises a DagsHub token by its context because its shape cannot be told from a commit SHA.
A merge commit is scanned for what it changed itself, with \texttt{git log -\allowbreak{}-\allowbreak{}remerge-\allowbreak{}diff}, so a secret written into a conflict resolution is seen as well.
Before the job scans the history it runs \texttt{tools/\allowbreak{}secret\_\allowbreak{}scan.\allowbreak{}py canary}: a throwaway history with a fake token committed in each of thirteen places at its real path, two tokens each, and two merge commits that each bring in one more, scanned with the command and the configuration the real scan uses, which fails unless all 28 are reported for the commit and the file they were planted in.},
  rationale = {\emph{Alternatives considered:}
\begin{itemize}[leftmargin=*, topsep=2pt]
\item \textbf{Option A (chosen): gitleaks, as a pre-commit hook and as a CI job over the whole history.}
\newline Pros: purpose-built, offline and MIT-licensed; the hook stops a secret before it is committed, and the job catches what the hook never saw, a commit made with \texttt{-\allowbreak{}-\allowbreak{}no-\allowbreak{}verify} or from a clone without the hook, including a secret added in one commit and deleted in the next, because it reads the history rather than the tree; the configuration is a file in the repository the course grades, reviewed like code; the default rules stay on, so the credentials the project does not have yet, such as FR-10's API key, a cloud key or a GitHub token, are covered when they arrive.
Everything GitHub publishes for the repository, 228 commits, scans in under two seconds and has no finding, with the default rules or with ours, so no secret was ever committed.
\newline Cons: gitleaks has no rule for a DagsHub token, and its generic rule misses one more often than it looks: in the seven places where the token is assigned with \texttt{=\allowbreak{}} or \texttt{:} it finds 93.0\,\% to 97.5\,\% of them, because 4.9\,\% of such tokens fall below its entropy threshold of 3.5, and in the other six places it finds none, including \texttt{uv run dvc remote modify origin -\allowbreak{}-\allowbreak{}local password <token>}, the exact command the getting-started page documents.
So the project owns a regex, and the canary exists to keep it honest.
The hook is built from source, so its first install is slow, once per machine and gitleaks version: measured on a thermally throttled laptop, 65 s with Go installed and 5 min 35 s without, when pre-commit downloads Go first, at about 145 s of CPU and 300 MB of RAM either way, leaving 47 MB or, with the downloaded Go, 329 MB in pre-commit's cache and about 180 MB in Go's build cache.
Each commit then costs about a quarter of a second of scanning.
\item \textbf{Option B: detect-secrets with a committed baseline.}
\newline Pros: a Python package that installs into our environment, with no Go toolchain; the baseline makes an existing finding auditable rather than a permanent red build.
Its default plugins also detect the DagsHub token better than gitleaks' defaults: in 11 of the 13 places, all 200 tokens in each, missing only the two bare command-line forms, \texttt{dvc remote modify .\allowbreak{}.\allowbreak{}.\allowbreak{} password <token>} and \texttt{dagshub login -\allowbreak{}-\allowbreak{}token <token>}.
\newline Cons: it scans files and, by design, not the git history, so the CI net for a commit that skipped the hook would only see the tip, and a secret added and then deleted would pass both checks; the baseline is a file the team has to keep honest, and an entry audited as ``not a secret'' is an allowlist that nobody re-reviews; and its entropy detectors flag hex and base64 by entropy alone, which on today's tree is three findings to audit -- \texttt{download.\allowbreak{}md5} in \texttt{params.\allowbreak{}yaml}, the gitleaks checksum the CI job pins, and a test string -- and grows with every hash we commit, while the MD5s of \texttt{dvc.\allowbreak{}lock} stay quiet only because it skips lock files.
\item \textbf{Option C: GitHub secret scanning with push protection.}
\newline Pros: nothing to install or maintain, and push protection blocks the push itself rather than reporting it afterwards.
\newline Cons: DagsHub is not among GitHub's supported secret scanning patterns (checked against the documentation and the \texttt{github/\allowbreak{}docs} repository on 2026-10-05), so it would not recognise the one secret the project has; a custom pattern would be another platform setting, as would the feature itself, leaving no trace in the repository the course grades and nothing a reviewer can read in a pull request.
\item \textbf{Option D: a Pytest test that greps the working tree for credential-shaped strings.}
\newline Pros: entirely ours, and the evidence would sit in the test suite, where the matrix could report NFR-15 as verified by a test.
\newline Cons: the weakest detection of the four: it sees only the tree at the commit under test, so a secret that was committed and deleted passes, and it would have to reimplement the context matching every scanner above already has; it would run only in the test job, which has a paths filter and a shallow checkout; and the matrix status it bought would describe a test standing in for the scan rather than the scan.
\end{itemize}
\par
\emph{Why:} What the requirement has to survive is a commit that skipped the hook, and the repository is public, so a secret in any commit of any pushed branch is published whether or not the tip still holds it.
That makes reading the history the property that decides, and it rules out B and D, which read files; C fails earlier, because it cannot see a DagsHub token at all.
Detection is not what decided it, and the measurement says so: detect-secrets' defaults see the token in more places than gitleaks' do.
What decides is that gitleaks' gap can be closed with one rule in a file we own, while B's history gap cannot be closed at all.
Three choices inside A follow from the same evidence.
The rule matches a 40-hex run by its context -- after a word naming a credential or one of the two services, or as the password in a URL -- with word boundaries, so the 32-hex MD5s of \texttt{dvc.\allowbreak{}lock}, \texttt{params.\allowbreak{}yaml} and the \texttt{.\allowbreak{}dvc} files and the 64-hex SHA-256s of \texttt{uv.\allowbreak{}lock} cannot fire it; with it, gitleaks finds all 2,600 planted tokens in all 13 places and still nothing in the tree or the history.
The CI job uses the release binary because \texttt{gitleaks/\allowbreak{}gitleaks-\allowbreak{}action} requires a licence key for a repository owned by an organisation, as this one is, obtained through a registration form and stored as a repository secret, which a pull request from a fork would not receive; the binary is checked against a SHA-256 pinned in the workflow, not one downloaded beside it.
And the job scans every ref, the heads of merged pull requests included, rather than only the pull request's own commits, because the cost is seconds and a leak on someone else's branch, or in a pull request whose branch a squash merge deleted, is published too: it turns every pull request red until it is revoked and acknowledged, which is the point, and a rule added later reaches all of it.
The canary is there because a scan that reports nothing is only evidence if it could have reported something: a regex edited into matching nothing or an allowlist widened over \texttt{.\allowbreak{}env} would otherwise turn the job silently green.
Its first version planted the tokens as files and scanned the directory, and an independent review showed what that cannot see: the real scan reads paths relative to the repository, so an allowlist anchored at the root, \texttt{\textasciicircum{}\textbackslash{}.\allowbreak{}env\$}, hid a committed \texttt{.\allowbreak{}env} token from it while the canary still passed.
The same review found that \texttt{git log -\allowbreak{}p}, which gitleaks runs, shows a merge commit no diff, so a token written into a conflict resolution and committed with \texttt{-\allowbreak{}-\allowbreak{}no-\allowbreak{}verify} was in the published history and in no scanned diff.
\texttt{-\allowbreak{}-\allowbreak{}remerge-\allowbreak{}diff} closes that by showing what a merge changed against an automatic remerge of its parents; \texttt{-\allowbreak{}-\allowbreak{}diff-\allowbreak{}merges=\allowbreak{}first-\allowbreak{}parent} caught it as well, but it reports again, for every merge of \texttt{main} into a branch, everything \texttt{main} ever carried, and each of those would need its own \texttt{.\allowbreak{}gitleaksignore} entry.
The canary now plants a history and scans it with the same command as the real scan, and it was checked against each way of blinding it: with \texttt{.\allowbreak{}gitleaks.\allowbreak{}toml} it finds all 28 tokens; with the default rules alone it misses 20; with the \texttt{\textasciicircum{}\textbackslash{}.\allowbreak{}env\$} allowlist, or a \texttt{.\allowbreak{}gitleaksignore} entry without its commit, it misses 7; and without \texttt{-\allowbreak{}-\allowbreak{}remerge-\allowbreak{}diff} it misses the two merge commits.
The hook is the \texttt{golang} variant because it needs nothing installed by hand, where \texttt{gitleaks-\allowbreak{}system} needs gitleaks on the path and \texttt{gitleaks-\allowbreak{}docker} needs Docker and an unpinned image.
CI skips it in the lint job, because it scans only staged changes, of which CI has none, and would report a scan of nothing as passed.
The job becomes one of the checks branch protection requires once this pull request is merged, which the owner approved; it is a repository setting, so it is not in the diff.
NFR-15 is \textbf{[manual]} in the specification because its evidence is this CI job and not a Pytest test.},
  ai = {Information seeking, Alternative generation, Alternative assessment, Recommendation, Solution generation},
  response = {Accepted with modifications},
  assessment = {AI laid out options A to D with their trade-offs and recommended A, and the owner chose it.
AI then tested the assumptions the choice rested on before implementing it, and one of them did not hold as stated: detect-secrets had been described as weaker on detection and noisy on the MD5s of \texttt{dvc.\allowbreak{}lock}, and measured, its defaults detect the DagsHub token in more places than gitleaks' and it skips \texttt{dvc.\allowbreak{}lock} entirely.
The decision survives on the argument that did hold, history scanning, and this entry records the measured reasons rather than the assumed ones.
The measurement also produced the parts of the design that were not in the recommendation: the \texttt{dagshub-\allowbreak{}token} rule, after the default rules turned out to miss the token in six places outright and in one case in twenty elsewhere, and the canary that keeps that rule from failing silently.
The modification came from an independent AI review of the pull request, which reproduced two holes in that first implementation, the directory-mode canary and the merge commits nothing scanned, and both were closed before merging.},
  aievidence = {Claude Code session on 2026-10-05 implementing issue \#60: the token's shape was read from the gitignored \texttt{.\allowbreak{}env} as length and character classes only, without printing the value; fake tokens of that shape were planted and scanned with both tools (\texttt{reports/\allowbreak{}analysis/\allowbreak{}secret\_\allowbreak{}scan\_\allowbreak{}dagshub\_\allowbreak{}token.\allowbreak{}py} and its results file); gitleaks-action's licence terms were read from its README and the repository's owner type from the GitHub API; the hook's install cost was measured with \texttt{/\allowbreak{}usr/\allowbreak{}bin/\allowbreak{}time} on a clean pre-commit cache, once with Go removed from the path and once with Go 1.27 on it; and a token committed with \texttt{-\allowbreak{}-\allowbreak{}no-\allowbreak{}verify} and deleted in the next commit was confirmed caught by the history scan.
The canary and each blinding variant were run in a container with git 2.45, because the laptop's git 2.34 predates \texttt{-\allowbreak{}-\allowbreak{}remerge-\allowbreak{}diff}; what an older Go does with gitleaks' \texttt{go.\allowbreak{}mod} was checked the same way with Go 1.22.},
  evidence = {\href{https://github.com/mlops-2627q1-mds-upc/MLOPS_Recommenditos/issues/60}{issue \#60}; \href{https://github.com/mlops-2627q1-mds-upc/MLOPS_Recommenditos/pull/70}{PR \#70}; \texttt{.\allowbreak{}gitleaks.\allowbreak{}toml}; \texttt{tools/\allowbreak{}secret\_\allowbreak{}scan.\allowbreak{}py}; the \texttt{secret-\allowbreak{}scan} job in \href{https://github.com/mlops-2627q1-mds-upc/MLOPS_Recommenditos/blob/main/.github/workflows/ci.yml}{\texttt{.\allowbreak{}github/\allowbreak{}workflows/\allowbreak{}ci.\allowbreak{}yml}} and the gitleaks hook in \href{https://github.com/mlops-2627q1-mds-upc/MLOPS_Recommenditos/blob/main/.pre-commit-config.yaml}{\texttt{.\allowbreak{}pre-\allowbreak{}commit-\allowbreak{}config.\allowbreak{}yaml}}; \href{https://github.com/mlops-2627q1-mds-upc/MLOPS_Recommenditos/blob/main/reports/analysis/secret_scan_dagshub_token.py}{\texttt{reports/\allowbreak{}analysis/\allowbreak{}secret\_\allowbreak{}scan\_\allowbreak{}dagshub\_\allowbreak{}token.\allowbreak{}py}} and \href{https://github.com/mlops-2627q1-mds-upc/MLOPS_Recommenditos/blob/main/reports/analysis/secret_scan_dagshub_token_results.txt}{its results}; \href{https://github.com/mlops-2627q1-mds-upc/MLOPS_Recommenditos/blob/main/docs/docs/specification.md}{specification} NFR-15; the Secrets section of \href{https://github.com/mlops-2627q1-mds-upc/MLOPS_Recommenditos/blob/main/CONTRIBUTING.md\#secrets}{CONTRIBUTING.md}; \hyperref[EDN-46]{EDN-46}, \hyperref[EDN-66]{EDN-66}.},
}

% Generated by tools/edn_latex.py from reports/edn.md.
% Do not edit: change reports/edn.md and run `make edn`.
\ednentry{
  id = {EDN-72},
  title = {The mileage range is a tolerant check on both frames, not a preprocessing filter},
  date = {2026-10-05},
  milestone = {M3: Quality Assurance},
  topic = {Data Validation, Data Preparation},
  participants = {@lukas2510, decided by the repository owner},
  decision = {The FR-03 mileage range, 0 to 1,000,000 km, is asserted on the raw and the interim frame alike as a tolerant check, \texttt{validate.\allowbreak{}mileage\_\allowbreak{}mostly =\allowbreak{} 0.\allowbreak{}999}, and preprocessing does not filter on it. The three listings above it stay in the cleaned frame and are reported by both suites.},
  rationale = {\emph{Alternatives considered:}
\begin{itemize}[leftmargin=*, topsep=2pt]
\item \textbf{Option A: drop them in preprocessing, as EDN-22 does for late registrations, and make the interim rule hard.}
\newline Pros: the cleaned frame becomes FR-03's domain exactly, so no training row is one the API refuses; the values are evidently entry errors -- a 2007 Volvo C30 at 1,560,000 km, a 2023 BYD Atto 3 at 2,424,696 km, a 2002 BMW 320 at 2,570,000 km -- and the interim suite would be hard on every bounded rule.
\newline Cons: it changes the training rows of every model and every published count (105,405 cleaned rows become 105,402), and it is a scope decision rather than part of putting \#25's rules into Great Expectations. Measured with the real downstream stages (\texttt{reports/\allowbreak{}analysis/\allowbreak{}mileage\_\allowbreak{}scope.\allowbreak{}py}): one of the three is a training row of a supported make, one a validation row and none a test row, and refitting without them moves \texttt{lgbm-\allowbreak{}basic}'s MdAPE from 6.83\,\% to 6.79\,\% and \texttt{b1}'s from 9.52\,\% to 9.60\,\%, leaves \texttt{lgbm-\allowbreak{}extended} and \texttt{b0} unchanged, and changes no gate verdict.
\item \textbf{Option B (chosen): a tolerant check on both frames.}
\newline Pros: no data or model changes; it matches how \texttt{params.\allowbreak{}yaml} already classed the range (``Bounds the Great Expectations suites assert'', while the price range ``lives under \texttt{preprocess}, because there it is a scope filter''); 0.999 tolerates more than thirty times today's count and still fails a scrape in another unit or with a shifted decimal.
\newline Cons: three rows the API would refuse stay in training, and the interim suite has one tolerant rule where the others are hard.
\item \textbf{Option C: a hard rule on the interim frame and no filter.}
\newline Pros: none beyond simplicity.
\newline Cons: it fails the stage on the current data, so the pipeline could not run at all.
\end{itemize}
\par
\emph{Why:} The measurement settled what was at stake: none of the three rows is a test row, and the model changes it causes are within a few hundredths of a percentage point with no verdict moving, so dropping them buys almost nothing that is measurable. What it would cost is a change to the training scope and to every count the documents and the report publish, a week before the first delivery, for an effect smaller than those documents round to. Keeping the range a check also keeps a systematic break visible as a failed rule instead of silently cleaning it away, the same reasoning EDN-22 gives for the raw suite's tolerance.},
  ai = {Information seeking, Alternative generation, Alternative assessment, Recommendation},
  response = {Accepted},
  assessment = {AI found the conflict by running the suites on the real snapshot, where a hard interim rule failed the stage on the three rows, and laid out the options instead of choosing one inside the Great Expectations work. It measured option A with the real stages rather than estimating it, which is where the MdAPE deltas come from, and recommended B for this PR. Lukas chose B.},
  aievidence = {Claude Code session on 2026-10-05 implementing issue \#25: the failing run, the measurement and the options were presented as a decision question in the session's report; Lukas answered ``keep the tolerant check''.},
  evidence = {\href{https://github.com/mlops-2627q1-mds-upc/MLOPS_Recommenditos/blob/main/reports/analysis/mileage_scope.py}{\texttt{reports/\allowbreak{}analysis/\allowbreak{}mileage\_\allowbreak{}scope.\allowbreak{}py}} and its results; \href{https://github.com/mlops-2627q1-mds-upc/MLOPS_Recommenditos/blob/main/reports/data-validation/summary.json}{\texttt{reports/\allowbreak{}data-\allowbreak{}validation/\allowbreak{}summary.\allowbreak{}json}}; \href{https://github.com/mlops-2627q1-mds-upc/MLOPS_Recommenditos/blob/main/recommenditos/data/gx_context_configuration.py}{\texttt{recommenditos/\allowbreak{}data/\allowbreak{}gx\_\allowbreak{}context\_\allowbreak{}configuration.\allowbreak{}py}}, \texttt{\_\allowbreak{}mileage\_\allowbreak{}in\_\allowbreak{}range}; \hyperref[EDN-22]{EDN-22}, \hyperref[EDN-68]{EDN-68}; \href{https://github.com/mlops-2627q1-mds-upc/MLOPS_Recommenditos/issues/25}{issue \#25}, \href{https://github.com/mlops-2627q1-mds-upc/MLOPS_Recommenditos/pull/76}{PR \#76}.},
}

% Generated by tools/edn_latex.py from reports/edn.md.
% Do not edit: change reports/edn.md and run `make edn`.
\ednentry{
  id = {EDN-73},
  title = {\texorpdfstring{\texttt{learning\_\allowbreak{}rate}}{learning\_rate} and \texorpdfstring{\texttt{num\_\allowbreak{}leaves}}{num\_leaves} are not tuned, and any later tuning follows a written protocol},
  note = {\textbf{Extended by EDN-78} (2026-10-06, issue \#88): the protocol below now covers both LightGBM variants and \texttt{min\_\allowbreak{}child\_\allowbreak{}samples}, its decision step is \texttt{recommenditos/\allowbreak{}modeling/\allowbreak{}tune.\allowbreak{}py}, and every fit of a search is measured and tracked.
This entry is kept as the record of the first sweep and of the protocol as it was written.},
  date = {2026-10-05},
  milestone = {M3: Quality Assurance},
  topic = {Hyperparameter Tuning, Evaluation Protocol, Experiment Tracking},
  participants = {@lukas2510, decided by the repository owner},
  decision = {\texttt{learning\_\allowbreak{}rate: 0.\allowbreak{}05} and \texttt{num\_\allowbreak{}leaves: 63} stay as they are for both LightGBM variants, because a bounded sweep around them found no point the validation split can tell apart from them.
The sweep is kept as the evidence that this region is flat, and any later tuning of a committed hyperparameter follows a protocol written down in \texttt{docs/\allowbreak{}docs/\allowbreak{}pipeline.\allowbreak{}md}, ``Tuning a hyperparameter'':
queued \texttt{dvc exp run -\allowbreak{}-\allowbreak{}queue -\allowbreak{}S .\allowbreak{}.\allowbreak{}.\allowbreak{} train@\allowbreak{}<variant>} experiments screen a sweep, so \texttt{evaluate} never runs and the test split takes no part in the choice, and they are tracked in their own MLflow experiment, \texttt{recommenditos-\allowbreak{}price-\allowbreak{}tuning}, so the ladder's \texttt{recommenditos-\allowbreak{}price} holds only pipeline runs;
\texttt{reports/\allowbreak{}analysis/\allowbreak{}tuning\_\allowbreak{}sweep.\allowbreak{}py} then re-fits the shortlist on the validation split and decides by \texttt{validation\_\allowbreak{}l1\_\allowbreak{}log\_\allowbreak{}price}, replacing a committed value only when its simultaneous 95\,\% interval lies below zero -- a paired bootstrap of the validation L1 difference, resampled by seller group, widened by the max-statistic bootstrap so that the intervals of all the sweep's points hold together.},
  rationale = {\emph{Alternatives considered:}
\begin{itemize}[leftmargin=*, topsep=2pt]
\item \textbf{Option A (chosen): do not tune now, and write the protocol down.}
\newline Pros: no number moves for a gain the data cannot resolve, so the report's numbers are final for this delivery.
It avoids the winner's curse: of eleven noisy comparisons against the committed point, the best-looking one is the one most likely to look better than it is, and a simultaneous interval is what keeps the protocol from adopting it.
It costs no compute, and the protocol makes the next tuning attempt cheap and honest instead of ad hoc.
\newline Cons: the model is untuned in the conventional sense, and the evidence covers a small grid of two knobs for one variant; \texttt{min\_\allowbreak{}child\_\allowbreak{}samples}, \texttt{feature\_\allowbreak{}fraction} and L2 regularisation are not parameters yet, and \texttt{lgbm-\allowbreak{}extended} was not swept.
\item \textbf{Option B: adopt the validation-L1 winner, \texttt{learning\_\allowbreak{}rate} 0.025 with \texttt{num\_\allowbreak{}leaves} 63, under the same protocol.}
\newline Pros: the lowest validation L1 of the grid, and a model that can be called tuned.
\newline Cons: it would fail the protocol it is adopted under.
The gain is inside the noise even point by point: L1 -0.00031 with a per-point 95\,\% interval of [-0.00086, +0.00025], better in 86\,\% of the draws, and MdAPE -0.009 pp with an interval of [-0.11, +0.09] pp; its simultaneous L1 interval is [-0.00109, +0.00047].
It costs 1.55x the trees and so the fit time and energy, a booster of about 12 MB and about 41 ms for one row's SHAP export, and it moves every number in the report for nothing measurable.
\item \textbf{Option C: a wider search, random or Optuna over more knobs, selected by seller-grouped k-fold cross-validation on train plus validation.}
\newline Pros: cross-validation is a less noisy selection than one validation split, and a larger space might hold a real gain.
\newline Cons: new code and new parameters, about five times the fits per point, more energy, and a departure from the problem specification's ``the validation set is used for early stopping and tuning'' -- all for a gain the flat grid suggests is small.
Left open for a later ticket if a criterion ever comes to depend on a fraction of a point.
\item \textbf{Option D: sweep with \texttt{dvc exp} over the whole pipeline and pick by the test metrics \texttt{dvc exp show} prints.}
\newline Pros: the least setup, and the metric is the one the gate reads.
\newline Cons: rejected, because it leaks the test split into the model the gate then judges, which makes SC-01 to SC-04 an optimistic estimate of a choice the test rows were part of.
\end{itemize}
\par
\emph{Why:} \texttt{reports/\allowbreak{}analysis/\allowbreak{}tuning\_\allowbreak{}sweep.\allowbreak{}py} fitted \texttt{lgbm-\allowbreak{}basic}, the candidate (EDN-62), at \texttt{learning\_\allowbreak{}rate} 0.025, 0.05 and 0.1 against \texttt{num\_\allowbreak{}leaves} 31, 63, 127 and 255, with the ceiling at 20,000 so early stopping decided every point (it did, at best rounds from 217 to 4,066), through the stage's own \texttt{fit\_\allowbreak{}variant} on the train and validation matrices only.
Each point's difference from the committed one was bootstrapped 2,000 times, paired and resampled by the 2,784 seller groups of the 8,859 validation rows, with the same draws for every point.
The region is flat: validation MdAPE spans 6.88\,\% to 7.04\,\% across all twelve points, every MdAPE interval spans zero and is about \ensuremath{\pm}0.13 pp wide, and no point improves validation L1 with an interval below zero.
Point by point, five are measurably \emph{worse} in L1: all four at \texttt{learning\_\allowbreak{}rate} 0.1, and 127 leaves at 0.05.
Simultaneously over the eleven, only 127 and 255 leaves at \texttt{learning\_\allowbreak{}rate} 0.1 remain worse, and no point is better.
The two metrics do not even agree on the best point -- \texttt{(0.\allowbreak{}025, 63)} by L1, \texttt{(0.\allowbreak{}05, 255)} by MdAPE -- which is the signature of a choice among noise.
\par
The protocol's parts each follow from the measurement.
Validation L1 is the selection metric because it is what early stopping reads and because its intervals are about 2.5 times narrower than MdAPE's relative to the value, so it resolves differences MdAPE cannot; MdAPE stays the gate's metric, read once on test after the choice.
The interval is simultaneous because a per-point one is not safe across a sweep: with eleven points no better than the committed one, the luckiest clears a per-point 95\,\% interval in about 24\,\% of sweeps if the comparisons were independent, and in at most 27.5\,\% by the union bound.
The max-statistic bootstrap reuses the paired, seller-grouped draws: each point's difference is standardised by its own bootstrap spread, the largest standardised deviation across the points is taken per draw, and its 95th percentile, 2.80 on this grid against 1.96 for one interval, sets every point's half-width.
\texttt{dvc exp run -\allowbreak{}-\allowbreak{}queue} targeted at \texttt{train@\allowbreak{}<variant>} was run once on 2026-10-05 to check the protocol is executable as written: only \texttt{train@\allowbreak{}lgbm-\allowbreak{}basic} ran, the stages before it were skipped as unchanged, \texttt{evaluate} was never reached, and the queue worker inherited the exported tracking variables, which matters because the queued workspace is built from the commit and does not see the gitignored \texttt{.\allowbreak{}env}.
It also showed one trap the protocol warns about: since \texttt{evaluate} does not run, every experiment carries its parent's \texttt{metrics.\allowbreak{}json}, so the test metrics \texttt{dvc exp show} prints belong to the committed model, not to the point.
So the queued runs are for screening, read in MLflow, and the decision is the analysis script's: it re-fits the shortlist on the validation split rather than reading the screening runs' bundles back, which MLflow (under \texttt{model/\allowbreak{}}) and the experiments' DVC outputs do keep, because a re-fit far past the ceiling gives every point its early-stopped model -- the same trees where the screening run early-stopped, and the uncut model where the committed ceiling bound it.
The script covers \texttt{learning\_\allowbreak{}rate} and \texttt{num\_\allowbreak{}leaves} only, and always fits the committed point, whether or not the grid names it.
\href{https://github.com/mlops-2627q1-mds-upc/MLOPS_Recommenditos/issues/88}{\#88} moved its statistics into \texttt{recommenditos/\allowbreak{}modeling/\allowbreak{}tune.\allowbreak{}py}, which runs the decision step for any grid and records every fit, and extended the protocol to \texttt{lgbm-\allowbreak{}extended} and \texttt{min\_\allowbreak{}child\_\allowbreak{}samples} (EDN-78).},
  ai = {Information seeking, Alternative generation, Alternative assessment, Recommendation, Solution generation},
  response = {Accepted with modifications},
  assessment = {AI designed and ran the bounded sweep, added the paired, seller-grouped bootstrap so the question became ``can the validation split tell these apart'' rather than ``which number is lowest'', and extended the grid to \texttt{learning\_\allowbreak{}rate} 0.025 after the first run showed the lower learning rate was the only direction not clearly worse.
It laid out options A to D, recommended A with the protocol, and stopped for the decision rather than committing a tuned value; Lukas chose A as proposed.
AI then ran one queued \texttt{dvc exp} experiment to confirm the documented commands do what the protocol claims before writing them into the pipeline documentation, which is how the \texttt{metrics.\allowbreak{}json} trap and the \texttt{.\allowbreak{}env} point were found, and wrote the decision step into the analysis script.
An independent review of the pull request, also by AI, corrected the first version of the protocol in two places.
Its rule used a per-point interval and called that safe against the winner's curse, which it is not across eleven comparisons.
And it left implicit that the decision step re-fits the shortlist rather than reading the queued runs back, while wrongly saying that neither DVC nor MLflow keeps what that step needs; MLflow keeps each run's bundle under \texttt{model/\allowbreak{}}, and each experiment's DVC outputs keep \texttt{models/\allowbreak{}}.
The rule became simultaneous, the sweep was re-run with the max-statistic intervals, and the conclusion did not change: a simultaneous interval is only wider, so ``no point distinguishable'' holds a fortiori.},
  aievidence = {Claude Code session on 2026-10-05 implementing issue \#64: the sweep's output is \texttt{reports/\allowbreak{}analysis/\allowbreak{}tuning\_\allowbreak{}sweep\_\allowbreak{}results.\allowbreak{}txt}; the options and the recommendation were reported to Lukas with those numbers, and he chose A.},
  evidence = {\href{https://github.com/mlops-2627q1-mds-upc/MLOPS_Recommenditos/issues/64}{issue \#64}; \href{https://github.com/mlops-2627q1-mds-upc/MLOPS_Recommenditos/pull/77}{PR \#77}; \href{https://github.com/mlops-2627q1-mds-upc/MLOPS_Recommenditos/blob/main/reports/analysis/tuning_sweep.py}{\texttt{reports/\allowbreak{}analysis/\allowbreak{}tuning\_\allowbreak{}sweep.\allowbreak{}py}} and its results file (2026-10-05); \href{https://github.com/mlops-2627q1-mds-upc/MLOPS_Recommenditos/blob/main/docs/docs/pipeline.md}{pipeline docs}, ``Tuning a hyperparameter''; \href{https://github.com/mlops-2627q1-mds-upc/MLOPS_Recommenditos/blob/main/docs/docs/model-card.md}{model card}, Hyperparameters as trained; \href{https://github.com/mlops-2627q1-mds-upc/MLOPS_Recommenditos/blob/main/params.yaml}{\texttt{params.\allowbreak{}yaml}}, \texttt{train.\allowbreak{}variants.\allowbreak{}lgbm-\allowbreak{}basic}; \href{https://github.com/mlops-2627q1-mds-upc/MLOPS_Recommenditos/blob/main/docs/docs/problem-spec.md}{problem specification} section 5; \hyperref[EDN-70]{EDN-70}, \hyperref[EDN-14]{EDN-14}.},
}

% Generated by tools/edn_latex.py from reports/edn.md.
% Do not edit: change reports/edn.md and run `make edn`.
\ednentry{
  id = {EDN-74},
  title = {A run is tagged with the hashes of its own stage's inputs, and \texorpdfstring{\texttt{git\_\allowbreak{}dirty}}{git\_dirty} ignores what DVC writes itself},
  date = {2026-10-05},
  milestone = {M3: Quality Assurance},
  topic = {Experiment Tracking, Reproducibility},
  participants = {@lukas2510, decided by the repository owner},
  decision = {Every MLflow run a pipeline stage opens or appends to is tagged with the hash of each dependency \texttt{dvc.\allowbreak{}yaml} declares for that stage, one \texttt{<stage>.\allowbreak{}deps.\allowbreak{}<path>} tag each, computed through DVC's own API when the stage starts; it is the value \texttt{dvc.\allowbreak{}lock} records under \texttt{stages.\allowbreak{}<stage>.\allowbreak{}deps} once the stage has finished.
\texttt{git\_\allowbreak{}dirty} is computed while ignoring the files \texttt{dvc repro} writes itself: each pipeline's lock file and every output declared \texttt{cache: false}, read off the pipeline rather than listed.
\texttt{dvc\_\allowbreak{}lock\_\allowbreak{}md5}, the MD5 of the whole lock, is dropped.
\texttt{evaluate}, which appends to the run \texttt{train} opened (EDN-55), records its own \texttt{evaluate.\allowbreak{}git\_\allowbreak{}commit}, \texttt{evaluate.\allowbreak{}git\_\allowbreak{}dirty} and \texttt{evaluate.\allowbreak{}deps.\allowbreak{}<path>} beside the fit's \texttt{git\_\allowbreak{}commit}, \texttt{git\_\allowbreak{}dirty} and \texttt{train.\allowbreak{}deps.\allowbreak{}<path>}.
Input hashes are recorded only when DVC runs the stage, which DVC signals by setting \texttt{DVC\_\allowbreak{}STAGE}.},
  rationale = {\emph{Alternatives considered:}
\begin{itemize}[leftmargin=*, topsep=2pt]
\item \textbf{Option A (chosen): the hashes of the stage's own inputs, taken when it starts, and \texttt{git\_\allowbreak{}dirty} without DVC's writes.}
\newline Pros: every tag names something a commit or the DVC remote reproduces: for the runs of a committed \texttt{dvc repro}, each input tag equals the committed lock line for line, and a clean commit is reported clean.
The tags exist from the moment the run does, so a stage that fails half-way still says what it read.
A data dependency's hash is its address in the DVC cache and on the remote, so \texttt{dvc pull} restores that input from it and it identifies the input whatever happens to the commits, a squash merge or a rebase included; a code file is restored by git at the commit, where a clean tree guarantees it has the recorded hash.
And one tag per input lets MLflow find every run that read a given matrix with an exact filter.
\newline Cons: a run carries 24 provenance tags instead of three, ten input hashes and a commit and a dirty flag for each of its two stages; the stage opens the DVC repository in-process, which costs about 0.3 s per run, measured at 0.28 to 0.41 s on the laptop against fits of 0.1 to 19 s; the code depends on DVC's Python API (\texttt{Index.\allowbreak{}stages}, \texttt{Output.\allowbreak{}get\_\allowbreak{}hash}), which DVC does not promise to keep stable, so a DVC upgrade can break it -- the test that runs a real \texttt{dvc repro} and compares with the lock is what would say so; and a hand edit of an uncached output, \texttt{metrics.\allowbreak{}json} say, no longer marks a run dirty, which costs nothing today, because no stage that opens or appends to a run depends on an uncached output; a stage that did would hash it among its inputs, so the edit would show there.
\item \textbf{Option B: keep the whole-lock MD5, and document what it does not identify.}
\newline Pros: no code change; the model card already warned that the four runs of one ladder carry four different digests.
\newline Cons: the tag identifies no committed state, so documenting it would only document that NFR-14's data-version clause is unmet.
Measured on the runs \texttt{main}'s models pointed at before this change: four \texttt{train} runs, four different digests, none equal to the committed lock; and the first stage's digest equals the lock from before the re-run.
It also leaves \texttt{git\_\allowbreak{}dirty} true for every stage after the first of a clean \texttt{dvc repro}, three of those four runs, only because \texttt{dvc.\allowbreak{}lock} had changed under them, so the report could not say its runs were made from a committed state either.
\item \textbf{Option C: tag after the run instead of before it.}
Either at the end of the stage's own process, or by a separate step after \texttt{dvc repro} that reads the committed lock and writes the tags onto the runs.
\newline Pros: the separate step reads the very lock it is compared with, so the tags would equal it by construction; and the stage-end variant needs no change to how the run is opened.
\newline Cons: at the end of its own process a stage still cannot read its lock entry, because DVC writes it only after the command has exited, so the stage-end variant would have to compute the same hashes A computes, later and with less of the run covered: a stage that fails half-way would carry none.
The separate step is a second command somebody has to remember after every \texttt{dvc repro}, and one forgotten leaves the runs without provenance and nothing failing; it would also describe the tree when the step ran, not when the stage did, and attribute anything changed in between to the run.
And it equals the lock by copying it, so it could not detect a run whose inputs differ from what the lock says, which is the one thing the check is for.
\end{itemize}
\par
\emph{Why:} NFR-14 promises that a run can be traced back to what produced it, and only a committed state can be traced back to.
A whole-lock digest taken during a \texttt{dvc repro} is never the lock committed with that run, so option B keeps a tag that cannot do its job, and C either recomputes A's hashes later or copies the lock instead of checking against it.
A is the only option whose tags are evidence: two programs, the stage at its start and DVC at its end, write the same values independently, and the check is that they agree.
\par
Three implementation choices are part of the decision.
The hashes are computed through DVC's API rather than read from upstream stages' outputs in \texttt{dvc.\allowbreak{}lock}, because only some dependencies are another stage's output: the code files are no stage's output, and \texttt{evaluate}'s \texttt{models} and \texttt{data/\allowbreak{}processed/\allowbreak{}features} are the parents of several outputs and the output of none, so their \texttt{.\allowbreak{}dir} hashes are in the lock only as \texttt{evaluate}'s own dependencies, written after it has run.
Reimplementing DVC's directory hash with \texttt{hashlib} would be one more thing that can disagree with the lock; calling the code that writes the lock cannot, and \texttt{tests/\allowbreak{}test\_\allowbreak{}provenance.\allowbreak{}py} proves the two equal on a three-stage pipeline DVC really runs, on a directory, a file, a \texttt{foreach} stage and an uncached output that a later stage depends on.
The tags are recorded only under DVC because every test calls the stages with paths under a temporary directory: the hashes of the files \texttt{dvc.\allowbreak{}yaml} declares would then describe files the run never opened.
And \texttt{evaluate} records its own commit beside the fit's because the two need not run at the same commit: a change to \texttt{evaluate.\allowbreak{}py} alone reruns only \texttt{evaluate}, against the runs of an earlier \texttt{train}, and overwriting \texttt{git\_\allowbreak{}commit} would then claim the fit came from code it did not come from.
\par
The implementation found more than the issue described.
DVC deletes a stage's outputs before running it, so a stage with an uncached output starts with that git-tracked file deleted: on a three-stage pipeline reproduced from a clean commit with the old code, even the first stage was tagged dirty, and only the first stage's digest matched any committed lock, the one from before the run.
And DVC adds a line to the \texttt{.\allowbreak{}gitignore} beside a cached output the first time it produces it, which for \texttt{gx/\allowbreak{}} is the repository root's hand-written \texttt{.\allowbreak{}gitignore}; ignoring that file would hide every edit to it, so the first run of a new cached output is reported dirty until its line is committed, which \texttt{pipeline.\allowbreak{}md} says.
And DVC addresses a stage relative to the directory \texttt{dvc repro} was started in, \texttt{.\allowbreak{}.\allowbreak{}/\allowbreak{}dvc.\allowbreak{}yaml:train@\allowbreak{}b0} from a subdirectory, so the stage is looked up by its name alone; an independent review of the pull request found that the first version resolved the path against the repository root, and that from a subdirectory every run would then have lost its tracking, under a warning that blamed the credentials.
\par
Checked on the real snapshot rather than only in a test: the four \texttt{train} stages and \texttt{evaluate} were re-run from a clean commit with \texttt{RECOMMENDITOS\_\allowbreak{}REQUIRE\_\allowbreak{}TRACKING=\allowbreak{}1}, and \texttt{reports/\allowbreak{}analysis/\allowbreak{}run\_\allowbreak{}provenance\_\allowbreak{}check.\allowbreak{}py} read the four runs back from DagsHub.
Before the re-run, the four runs \texttt{main}'s committed models pointed at carried none of the input tags, three of them were tagged dirty, and each \texttt{dvc\_\allowbreak{}lock\_\allowbreak{}md5} differed from the digest of the committed lock.
After it, each of the four carries ten \texttt{train.\allowbreak{}deps} and ten \texttt{evaluate.\allowbreak{}deps} tags, every one equal to the committed lock, and \texttt{git\_\allowbreak{}dirty} and \texttt{evaluate.\allowbreak{}git\_\allowbreak{}dirty} are \texttt{false} for all four, naming the clean commit the re-run started from (\texttt{reports/\allowbreak{}analysis/\allowbreak{}run\_\allowbreak{}provenance\_\allowbreak{}check\_\allowbreak{}results.\allowbreak{}txt}).
The re-run, 1 min 48 s at a peak RSS of 693 MB, moved no number: \texttt{metrics.\allowbreak{}json}, the segment table and the masking sweep are byte-identical, each per-variant record changed only in its MLflow run id, and inside each model bundle only \texttt{model.\allowbreak{}json} changed, by the provenance EDN-56 puts there, while every payload file kept its hash.},
  ai = {Information seeking, Alternative generation, Alternative assessment, Recommendation, Solution generation},
  response = {Accepted with modifications},
  assessment = {AI found the defect in an independent review of PR \#70, by reading the runs back from DagsHub instead of trusting the tests, which checked that the tags exist and that they change when \texttt{dvc.\allowbreak{}lock} changes, not that they identify a committed state.
It measured the four digests against the committed lock and the dirty flags against the files that had changed, generated options A to C and assessed them, and recommended A, fixed now in a pull request of its own; Lukas chose that, before the first delivery, so that the report's provenance claim is true and the CodeCarbon run of \#38 retrains only once afterwards.
The implementation details were AI's, as were the two further findings, the deleted uncached output and the \texttt{.\allowbreak{}gitignore} line, both reproduced before they were acted on.
An independent AI review of the pull request reproduced the evidence and found the subdirectory regression, two overstated claims -- that every hash retrieves its input from the remote, and that a hand edit of the validation summary would show in a run -- and three smaller gaps: a dependency path MLflow refuses as a tag key, a failed \texttt{git status} read as a clean tree, and a test pipeline without a parameter; each was fixed, test first where there was code, before merging, which is the modification.},
  aievidence = {Claude Code sessions on 2026-10-05: the independent review of PR \#70, whose measurement is the body of issue \#79, and the implementation of issue \#79, in which the issue's scenario was reproduced on a three-stage pipeline with \texttt{main}'s \texttt{run\_\allowbreak{}provenance} (every stage tagged dirty, the digests of the lock before the run and of two mid-run locks), the new tests were confirmed to fail on \texttt{main}, where the module they test does not exist, and the check script was run against DagsHub before and after the re-run.},
  evidence = {\href{https://github.com/mlops-2627q1-mds-upc/MLOPS_Recommenditos/issues/79}{issue \#79}; \href{https://github.com/mlops-2627q1-mds-upc/MLOPS_Recommenditos/pull/80}{PR \#80}; \href{https://github.com/mlops-2627q1-mds-upc/MLOPS_Recommenditos/blob/main/recommenditos/provenance.py}{\texttt{recommenditos/\allowbreak{}provenance.\allowbreak{}py}} and \texttt{recommenditos/\allowbreak{}tracking.\allowbreak{}py}; \href{https://github.com/mlops-2627q1-mds-upc/MLOPS_Recommenditos/blob/main/tests/test_provenance.py}{\texttt{tests/\allowbreak{}test\_\allowbreak{}provenance.\allowbreak{}py}} and the NFR-14 tests of \texttt{tests/\allowbreak{}test\_\allowbreak{}tracking.\allowbreak{}py}; \href{https://github.com/mlops-2627q1-mds-upc/MLOPS_Recommenditos/blob/main/reports/analysis/run_provenance_check.py}{\texttt{reports/\allowbreak{}analysis/\allowbreak{}run\_\allowbreak{}provenance\_\allowbreak{}check.\allowbreak{}py}} and \href{https://github.com/mlops-2627q1-mds-upc/MLOPS_Recommenditos/blob/main/reports/analysis/run_provenance_check_results.txt}{its results}; \href{https://github.com/mlops-2627q1-mds-upc/MLOPS_Recommenditos/blob/main/docs/docs/pipeline.md}{The DVC pipeline}, From a run to its inputs; \href{https://github.com/mlops-2627q1-mds-upc/MLOPS_Recommenditos/blob/main/docs/docs/specification.md}{specification} NFR-14; \hyperref[EDN-55]{EDN-55}, \hyperref[EDN-66]{EDN-66}.},
}



\end{document}
